\documentclass[11pt,a4paper,oneside]{report}
\usepackage[top=2.5cm,bottom=2.5cm,left=2.3cm,right=2.3cm,headsep=16pt,footskip=28pt]{geometry}

\usepackage{amsmath,amssymb}
\usepackage{fontspec}
\setmainfont{Liberation Serif}
\setsansfont{Liberation Sans}
\setmonofont{DejaVu Sans Mono}[Scale=0.84]

\usepackage{microtype}

\usepackage{xcolor}
\definecolor{brandblue}{RGB}{20, 55, 105}
\definecolor{accentblue}{RGB}{35, 95, 165}
\definecolor{codebg}{RGB}{248, 249, 251}
\definecolor{codeframe}{RGB}{215, 222, 230}
\definecolor{javakeyword}{RGB}{0, 45, 170}
\definecolor{javastring}{RGB}{5, 120, 20}
\definecolor{javacomment}{RGB}{120, 120, 120}
\definecolor{javanumber}{RGB}{20, 75, 220}

\definecolor{noteborder}{RGB}{30, 110, 65}
\definecolor{notebg}{RGB}{242, 249, 245}
\definecolor{warnborder}{RGB}{185, 75, 10}
\definecolor{warnbg}{RGB}{254, 248, 240}
\definecolor{tipborder}{RGB}{30, 95, 160}
\definecolor{tipbg}{RGB}{242, 247, 254}

\usepackage{fancyhdr}
\usepackage{lastpage}
\pagestyle{fancy}
\renewcommand{\chaptermark}[1]{\markboth{#1}{}}
\fancyhf{}
\lhead{\parbox[t]{0.58\textwidth}{\raggedright\small\sffamily\color{brandblue}\textbf{\nouppercase{\leftmark}}}}
\rhead{\small\sffamily\color{gray}Programmazione II -- UniVR (2025/2026)}
\cfoot{\small\sffamily Pagina \thepage\ di \pageref{LastPage}}
\renewcommand{\headrulewidth}{0.5pt}
\renewcommand{\headrule}{\hbox to\headwidth{\color{brandblue!70!black}\leaders\hrule height \headrulewidth\hfill}}

\fancypagestyle{plain}{
  \fancyhf{}
  \cfoot{\small\sffamily Pagina \thepage\ di \pageref{LastPage}}
  \renewcommand{\headrulewidth}{0pt}
}

\usepackage{titlesec}
\titleformat{\chapter}[display]
  {\normalfont\sffamily}
  {\flushright\color{accentblue}\Large\bfseries CAPITOLO \thechapter}
  {0.4ex}
  {\flushleft\Huge\bfseries\color{brandblue}}
  [\vspace{0.4ex}{\color{brandblue}\hrule height 1.2pt}\vspace{0.8ex}]
\titlespacing*{\chapter}{0pt}{-15pt}{22pt}

\titleformat{\section}{\Large\bfseries\sffamily\color{brandblue}}{\thesection}{0.8em}{}[\titlerule]
\titleformat{\subsection}{\large\bfseries\sffamily\color{accentblue}}{\thesubsection}{0.8em}{}
\titleformat{\subsubsection}{\normalsize\bfseries\sffamily\color{brandblue!85!black}}{\thesubsubsection}{0.8em}{}

\usepackage{tocloft}
\renewcommand{\cfttoctitlefont}{\Huge\bfseries\sffamily\color{brandblue}}
\renewcommand{\cftchapfont}{\large\bfseries\sffamily\color{brandblue}}
\renewcommand{\cftsecfont}{\normalsize\sffamily}
\renewcommand{\cftsubsecfont}{\small\sffamily}
\renewcommand{\cftchapleader}{\cftdotfill{\cftdotsep}}
\setlength{\cftbeforechapskip}{1.1ex}
\setlength{\cftbeforesecskip}{0.35ex}

\usepackage{needspace}
\usepackage{fancyvrb}
\DefineVerbatimEnvironment{asciibox}{Verbatim}{
  fontsize=\fontsize{7.5pt}{9.2pt}\selectfont,
  frame=single,
  rulecolor=\color{codeframe},
  fillcolor=\color{codebg},
  framesep=2.5mm
}
\usepackage{listings}
\lstset{
  backgroundcolor=\color{codebg},
  frame=single,
  rulecolor=\color{codeframe},
  basicstyle=\ttfamily\footnotesize,
  keywordstyle=\color{javakeyword}\bfseries,
  stringstyle=\color{javastring},
  commentstyle=\color{javacomment}\itshape,
  numberstyle=\tiny\color{gray},
  numbers=left,
  numbersep=7pt,
  stepnumber=1,
  breaklines=true,
  breakatwhitespace=true,
  showstringspaces=false,
  tabsize=4,
  xleftmargin=12pt,
  framexleftmargin=12pt
}

\newcommand{\passthrough}[1]{#1}
\providecommand{\tightlist}{\setlength{\itemsep}{0pt}\setlength{\parskip}{0pt}}

\usepackage{tcolorbox}
\tcbuselibrary{skins,breakable}
\newtcolorbox{studynote}[1][Nota]{
  enhanced,
  breakable,
  colback=notebg,
  colframe=noteborder,
  fonttitle=\bfseries\sffamily\small,
  coltitle=white,
  title={#1},
  arc=2mm,
  boxrule=0.8pt,
  left=9pt, right=9pt, top=7pt, bottom=7pt
}

\newtcolorbox{studywarning}[1][Attenzione]{
  enhanced,
  breakable,
  colback=warnbg,
  colframe=warnborder,
  fonttitle=\bfseries\sffamily\small,
  coltitle=white,
  title={#1},
  arc=2mm,
  boxrule=0.8pt,
  left=9pt, right=9pt, top=7pt, bottom=7pt
}

\newtcolorbox{studytip}[1][Suggerimento]{
  enhanced,
  breakable,
  colback=tipbg,
  colframe=tipborder,
  fonttitle=\bfseries\sffamily\small,
  coltitle=white,
  title={#1},
  arc=2mm,
  boxrule=0.8pt,
  left=9pt, right=9pt, top=7pt, bottom=7pt
}

\usepackage{booktabs}
\usepackage{longtable}
\usepackage{array}
\usepackage{calc}

\usepackage{hyperref}
\hypersetup{
  colorlinks=true,
  linkcolor=accentblue,
  urlcolor=accentblue,
  citecolor=brandblue,
  pdfborder={0 0 0},
  pdftitle={Programmazione II -- Compendio Generale Teorico e Pratico},
  pdfauthor={Wally -- Docente: Prof. Michele Pasqua},
  bookmarksnumbered=true
}

\begin{document}

% ==================== FRONTESPIZIO ====================
\begin{titlepage}
\centering
\vspace*{0.5cm}

{\Large\sffamily\scshape Università degli Studi di Verona}\\[0.3em]
{\large\sffamily Dipartimento di Informatica}\\[0.2em]
{\normalsize\sffamily Corso di Laurea in Informatica / Bioinformatica}\\[1.8cm]

\noindent{\color{accentblue}\rule{\textwidth}{2pt}}\\[0.8cm]

{\fontsize{28pt}{34pt}\selectfont\bfseries\sffamily\color{brandblue} PROGRAMMAZIONE II}\\[0.5cm]
{\LARGE\bfseries\sffamily\color{accentblue} COMPENDIO GENERALE TEORICO E PRATICO}\\[0.8cm]
{\large\sffamily Guida Completa alle Lezioni (Slide T01 -- T20), Analisi Minuziosa del Codice,\\ Trappole d'Esame e Best Practice di Progettazione ad Oggetti}\\[0.8cm]

\noindent{\color{accentblue}\rule{\textwidth}{0.8pt}}\\[2.2cm]

\begin{minipage}[t]{0.48\textwidth}
  {\small\sffamily\color{gray}DOCENTE DEL CORSO}\\[0.3em]
  {\large\bfseries\sffamily Prof. Michele Pasqua}\\[0.2em]
  {\small\sffamily Dipartimento di Informatica}\\
  {\small\ttfamily michele.pasqua@univr.it}
\end{minipage}
\hfill
\begin{minipage}[t]{0.48\textwidth}
  \begin{flushright}
    {\small\sffamily\color{gray}REVISIONE E APPUNTI DI STUDIO}\\[0.3em]
    {\large\bfseries\sffamily Wally}\\[0.2em]
    {\small\sffamily Corso di Programmazione II}\\
    {\small\sffamily Anno Accademico 2025/2026}
  \end{flushright}
\end{minipage}

\vfill

{\small\sffamily\color{gray} Settembre 2026 -- Edizione Integrale Unificata}\\[0.2em]
{\footnotesize\sffamily Compilato con XeLaTeX -- 19 Moduli Didattici -- Inclusione Totale di Codice, Teoria ed Esercizi}

\end{titlepage}

% ==================== PREFAZIONE ====================
\pagenumbering{roman}
\setcounter{page}{1}

\chapter*{Prefazione e Metodologia di Studio}
\addcontentsline{toc}{chapter}{Prefazione e Metodologia di Studio}

Questo volume nasce dall'esigenza di condensare, armonizzare e sistematizzare in un'unica opera organica l'intero patrimonio didattico dell'insegnamento di \textbf{Programmazione II} (A.A. 2025/2026, tenuto dal \textbf{Prof. Michele Pasqua} presso l'Università degli Studi di Verona).

Il corso si articola lungo tre direttrici formative sinergiche:
\begin{enumerate}
  \item \textbf{I Fondamenti e la Teoria dei Linguaggi ad Oggetti}: dalla scomposizione in classi e incapsulamento, alla semantica formale di memoria (\emph{Call Stack}, \emph{Heap}, garbage collection), fino al polimorfismo dinamico (\emph{dynamic binding}), all'astrazione mediante interfacce, ai tipi generici con polimorfismo parametrico e alle funzionalità dichiarative introdotte da Java 8+ (\emph{Functional Interfaces}, lambda e \emph{Streams API}).
  \item \textbf{L'Ingegneria del Software Incrementale}: il corso non si limita ad illustrare frammenti di codice isolati, ma costruisce un percorso evolutivo concreto. Partendo da una prima classe rudimentale (\texttt{SimpleDate}) con campi pubblici e logica procedurale, il codice viene via via rifattorizzato attraverso 19 passaggi chiave fino alla complessa architettura modulare di \texttt{MavenDate}, gestita con Apache Maven e conforme ai principi SOLID e ai classici Design Pattern (\emph{Singleton}, \emph{Builder}, \emph{Decorator}).
  \item \textbf{Robustezza, Testing e Qualità del Codice}: la gestione professionale delle eccezioni (\emph{checked} vs \emph{unchecked}, clausole di salvaguardia, \emph{try-with-resources}), la documentazione formale secondo lo standard \emph{Javadoc} e la verifica rigorosa della correttezza mediante unit test automatizzati con \emph{JUnit 5}.
\end{enumerate}

\section*{Convenzioni Tipografiche e Struttura dei Capitoli}
Ogni capitolo del compendio corrisponde a un blocco didattico unitario (\textbf{da Slide T01 a T20}) ed è suddiviso in tre macro-sezioni:
\begin{itemize}
  \item \textbf{Concetti Teorici}: sintesi completa e minuziosa di tutte le nozioni esposte nelle slide del docente, con schemi UML e puntualizzazioni semantiche.
  \item \textbf{Analisi del Codice}: disamina riga per riga dei sorgenti ufficiali del docente, evidenziando il razionale di ciascuna scelta architetturale e il confronto \emph{diff} rispetto alla versione precedente.
  \item \textbf{Verifica Pratica \& Output}: riscontro operativo da riga di comando tramite \texttt{mvn test} e \texttt{mvn exec:java}, con decodifica del comportamento della JVM.
\end{itemize}

Nel testo sono stati inseriti tre ambienti grafici dedicati per supportare lo studio mirato:
\begin{studynote}[Nota di Studio]
Approfondimento teorico, formalizzazione del modello di esecuzione Java o regola architetturale inderogabile.
\end{studynote}

\begin{studytip}[Suggerimento / Best Practice]
Consiglio pratico di scrittura del codice, convenzione di stile o pattern per ottimizzare leggibilità e manutenibilità.
\end{studytip}

\begin{studywarning}[Attenzione / Trappola d'Esame]
Caso limite insidioso, tipico errore di compilazione o comportamento runtime inatteso frequentemente oggetto di verifica d'esame.
\end{studywarning}

\newpage
% ==================== INDICE GENERALE ====================

\renewcommand{\contentsname}{Indice Generale}
\tableofcontents
\newpage

% ==================== CORPO DEL TESTO ====================
\pagenumbering{arabic}
\setcounter{page}{1}

\chapter[T01 --- Course Introduction]{T01 --- Course Introduction \& Metodologia del Corso}
\label{chap:t01}

{\small\sffamily\color{accentblue!80!black}\textit{Obiettivi Formativi, Modalità d'Esame al Calcolatore, Toolchain e Syllabus}}

\vspace{0.8em}

Benvenuto in questa revisione completa e metodica di
\textbf{Programmazione II} (Prof.~Michele Pasqua, UniVR 2025/2026).
Seguiremo fedelmente le regole stabilite: \textbf{nessuna omissione},
analisi minuziosa di ogni concetto teorico, confronto riga per riga del
codice, evidenziazione delle finezze d'esame e stop interattivo a ogni
blocco.

\section{\texorpdfstring{Teoria: Concetti Teorici dalle Slide
(\href{file:///home/wally/Documenti/GitHub/Programmazione-II/25-26/Teoria-e-Esercitazioni/Slides-20260902/T01\%20-\%20Course\%20Introduction.pdf}{T01
- Course
Introduction.pdf})}{Teoria: Concetti Teorici dalle Slide (T01 - Course Introduction.pdf)}}\label{chap01_teoria-concetti-teorici-dalle-slide-t01---course-introduction.pdf}

La slide introduttiva definisce il perimetro accademico, gli obiettivi
formativi e le modalità d'esame:

\begin{enumerate}
\def\labelenumi{\arabic{enumi}.}
\tightlist
\item
  \textbf{Docente e Background Scientifico}:

  \begin{itemize}
  \tightlist
  \item
    \textbf{Prof.~Michele Pasqua} (Dipartimento di Informatica, Cà
    Vignal 2, 2° piano, Stanza 24 ---
    \passthrough{\lstinline!michele.pasqua@univr.it!}).
  \item
    Aree di ricerca: \emph{Program verification and validation},
    \emph{Programming language design and implementation}, \emph{Code
    protection and software security}.
  \item
    Laboratori e progetti: co-fondatore del laboratorio
    \textbf{UNIVERSE} (\emph{University of Verona Software Engineering
    Lab}, diretto dal Prof.~Mariano Ceccato) e co-fondatore dello
    spin-off universitario \textbf{KonTrust Srl} (tracciamento IoT su
    blockchain e infrastrutture anti-manomissione); referente per Verona
    del programma nazionale di formazione e gare su cybersecurity
    (\emph{CyberChallenge.IT}).
  \end{itemize}
\item
  \textbf{Organizzazione della Didattica}:

  \begin{itemize}
  \tightlist
  \item
    \textbf{Modulo di Teoria (40+ ore)}: Lunedì (16:30) e Martedì
    (13:30) in Aula A (2 ore ciascuna).
  \item
    \textbf{Modulo di Laboratorio (12+ ore)}: Martedì (10:30) in Aula
    Delta (2 ore).
  \item
    \textbf{Materiale didattico di riferimento}:

    \begin{itemize}
    \tightlist
    \item
      Paul J. Deitel \& Harvey M. Deitel, \emph{Java -- How to Program},
      Pearson (11/e -- 2018).
    \item
      Claudio De Sio Cesari, \emph{Il nuovo Java -- Guida alla
      programmazione moderna}, Hoepli (1/e -- 2020).
    \item
      Moodle UniVR: slide, esercizi di laboratorio e soluzioni,
      simulazioni.
    \end{itemize}
  \end{itemize}
\item
  \textbf{Modalità d'Esame e Valutazione}:

  \begin{itemize}
  \tightlist
  \item
    \textbf{Formato}: Prova pratica al calcolatore in \textbf{Aula
    Delta} (ambiente Linux/IDE).
  \item
    \textbf{Contenuto}: Risoluzione di una serie di esercizi di
    programmazione in linguaggio Java (sviluppo classi, gerarchie a
    oggetti, implementazione interfacce, algoritmi su collezioni,
    gestione eccezioni, test con asserzioni/JUnit).
  \item
    \textbf{Voto finale}: Somma algebrica dei punti ottenuti nei singoli
    esercizi.
  \item
    È prevista una \textbf{simulazione d'esame} ufficiale a fine corso.
  \end{itemize}
\item
  \textbf{Syllabus del Corso (Mappa degli argomenti da T01 a T20)}:

  \begin{itemize}
  \tightlist
  \item
    \emph{Object-Orientation Principles} (T02)
  \item
    \emph{Building, Running, and Deployment} (T03)
  \item
    \emph{Types, Classes, and Objects} (T04)
  \item
    \emph{Java Syntax and Semantics} (T05)
  \item
    \emph{Java Classes and Objects} (T06)
  \item
    \emph{Libraries and String Class} (T07)
  \item
    \emph{Class Constructors and Encapsulation} (T08)
  \item
    \emph{Arrays and Enumerative Types} (T09)
  \item
    \emph{Java Variables and Call Stack} (T10)
  \item
    \emph{Class and Package Visibility} (T11)
  \item
    \emph{Inheritance and Polymorphism} (T12)
  \item
    \emph{Abstract Classes and Interfaces} (T13)
  \item
    \emph{Nested and Anonymous Classes} (T14)
  \item
    \emph{Generic Types} (T15)
  \item
    \emph{Collections} (T16)
  \item
    \emph{Error Handling with Exceptions} (T17)
  \item
    \emph{Functional Interfaces and Streams} (T18)
  \item
    \emph{Documentation and Unit Testing} (T19)
  \item
    \emph{Java IO} (T20)
  \end{itemize}
\end{enumerate}

\begin{center}\rule{0.5\linewidth}{0.5pt}\end{center}

\section{Codice: Analisi del Codice della Lezione
1}\label{chap01_codice-analisi-del-codice-della-lezione-1}

\begin{itemize}
\tightlist
\item
  Nella cartella ufficiale delle lezioni
  \href{file:///home/wally/Documenti/GitHub/Programmazione-II/25-26/Teoria-e-Esercitazioni/Codice-20260902}{\passthrough{\lstinline!Codice-20260902!}}:

  \begin{itemize}
  \tightlist
  \item
    \textbf{Non sono presenti file di codice per le Lezioni 01, 02 e
    03}.
  \item
    Il primo codice sorgente ufficiale del docente compare a partire
    dalla \textbf{Lezione 04} con
    \href{file:///home/wally/Documenti/GitHub/Programmazione-II/25-26/Teoria-e-Esercitazioni/Codice-20260902/Lezione04/Example.java}{\passthrough{\lstinline!Example.java!}}
    (dedicato a tipi primitivi, reference, casting e operatori), mentre
    il progetto a oggetti incrementale vero e proprio parte dal blocco
    \textbf{Lezioni 05-06} con il progetto
    \passthrough{\lstinline!SimpleDate!}
    (\passthrough{\lstinline!Date.java!} e
    \passthrough{\lstinline!MainDate.java!}).
  \end{itemize}
\item
  Nel tuo workspace
  \href{file:///home/wally/Documenti/GitHub/Programmazione-II/25-26/esercizi-wally-25-26}{\passthrough{\lstinline!esercizi-wally-25-26!}}:

  \begin{itemize}
  \tightlist
  \item
    La cartella
    \href{file:///home/wally/Documenti/GitHub/Programmazione-II/25-26/esercizi-wally-25-26/es01}{\passthrough{\lstinline!es01!}}
    è già predisposta come modulo Maven con la prima incarnazione di
    \passthrough{\lstinline!Date!} e \passthrough{\lstinline!MainDate!},
    sincronizzata con la struttura multi-modulo orchestrata dal
    \href{file:///home/wally/Documenti/GitHub/Programmazione-II/25-26/esercizi-wally-25-26/pom.xml}{\passthrough{\lstinline!pom.xml!}}
    di root e dallo script
    \href{file:///home/wally/Documenti/GitHub/Programmazione-II/25-26/esercizi-wally-25-26/esercizi.sh}{\passthrough{\lstinline!esercizi.sh!}}.
  \end{itemize}
\end{itemize}

\begin{center}\rule{0.5\linewidth}{0.5pt}\end{center}

\section{Setup: Configurazione Operativa
dell'Ambiente}\label{chap01_setup-configurazione-operativa-dellambiente}

Per garantire che tutti i passaggi da qui in avanti siano riproducibili
su IntelliJ IDEA e da CLI: 1. Verifica che la JDK attiva sia Java 17 o
superiore: \passthrough{\lstinline!bash    java -version    mvn -v!} 2.
Il file di root
\href{file:///home/wally/Documenti/GitHub/Programmazione-II/25-26/esercizi-wally-25-26/pom.xml}{\passthrough{\lstinline!pom.xml!}}
gestisce l'aggregazione di tutti i sottomoduli. Per ogni blocco
successivo, quando andremo ad analizzare le classi introdotte a lezione,
verificheremo che compilino ed eseguano pulitamente sia tramite:
\passthrough{\lstinline!bash    ./esercizi.sh run <numero\_modulo>!} sia
all'interno di IntelliJ IDEA.

\begin{center}\rule{0.5\linewidth}{0.5pt}\end{center}

\begin{studynote}[Nota di Studio] La \textbf{Lezione 1 / T01} esaurisce qui la parte
introduttiva/organizzativa. Il prossimo blocco è la \textbf{Lezione 2 /
Slide T02: \emph{Object-Orientation and UML}}, dove entreremo nel vivo
con: - Storia e tassonomia dei paradigmi: \emph{Imperativo/Procedurale},
\emph{Dichiarativo/Funzionale}, \emph{Object-Oriented}. - I 4 pilastri
dell'OOP: \emph{Astrazione, Incapsulamento, Ereditarietà, Polimorfismo}.
- Concetto di \emph{Classe} vs \emph{Oggetto} (stato, comportamento,
identità). - Notazione UML per classi e associazioni (\emph{Visibility
markers}, cardinalità, tipi di relazione).
\end{studynote}

\textbf{Dammi conferma quando sei pronto per passare a T02 (procederemo
analizzando il blocco concettuale un po' alla volta)!}

\newpage
\chapter[T02 --- Paradigma ad Oggetti e Modellazione]{T02 --- Paradigma ad Oggetti e Modellazione UML}
\label{chap:t02}

{\small\sffamily\color{accentblue!80!black}\textit{Astrazione, Incapsulamento, Modularità, Gerarchia e Sintassi dei Class Diagram}}

\vspace{0.8em}

In questo blocco analizziamo la transizione teorico-metodologica dal
paradigma procedurale a quello ad oggetti, i concetti cardine dell'OOP e
la modellazione grafica tramite \textbf{UML}
(\href{file:///home/wally/Documenti/GitHub/Programmazione-II/25-26/Teoria-e-Esercitazioni/Slides-20260902/T02\%20-\%20Object-Orientations\%20and\%20UML.pdf}{T02
- Object-Orientations and UML.pdf}).

\section{Teoria: Concetti Teorici dalle
Slide}\label{chap02_teoria-concetti-teorici-dalle-slide}

\subsection{La Tassonomia dei Linguaggi e il Bisogno di Astrazione
(Slide
1--13)}\label{chap02_la-tassonomia-dei-linguaggi-e-il-bisogno-di-astrazione-slide-113}

\begin{itemize}
\tightlist
\item
  \textbf{Nessun linguaggio universale}: citando Edsger W. Dijkstra
  (\emph{«Computer Science is no more about computers than Astronomy is
  about telescopes»}), il docente introduce il principio del
  \textbf{Golden Hammer} (\emph{«If all you have is a hammer, everything
  looks like a nail»}): non esiste un linguaggio perfetto per tutto.
\item
  \textbf{Evoluzione storica degli obiettivi}:

  \begin{itemize}
  \tightlist
  \item
    \emph{Anni '50/'60}: focus sull'\textbf{efficienza del codice
    macchina}. L'hardware costava milioni e i programmatori poco:
    l'obiettivo era ``tenere occupata la macchina'' con linguaggi vicini
    al silicio.
  \item
    \emph{Oggi}: focus sull'\textbf{efficienza dello sviluppo del
    software}. L'hardware costa poco e i programmatori costano molto:
    l'obiettivo è ``tenere produttivo il programmatore'', abbattendo
    complessità, bug e costi di manutenzione tramite
    l'\textbf{astrazione}.
  \end{itemize}
\item
  \textbf{I Tre Paradigmi Fondamentali}:

  \begin{enumerate}
  \def\labelenumi{\arabic{enumi}.}
  \tightlist
  \item
    \textbf{Imperativo / Procedurale} (C, Pascal, Fortran, COBOL):

    \begin{itemize}
    \tightlist
    \item
      Basato sul concetto di \textbf{stato} (celle di memoria che
      memorizzano valori) e \textbf{comandi sequenziali} che mutano tale
      stato (assegnamenti).
    \item
      Riuso del codice tramite procedure/funzioni che operano su dati
      esterni.
    \end{itemize}
  \item
    \textbf{Dichiarativo}:

    \begin{itemize}
    \tightlist
    \item
      \emph{Funzionale} (Haskell, Lisp): il programma è la valutazione
      di un'espressione matematica pura, basata su applicazione di
      funzioni e ricorsione (es.
      \passthrough{\lstinline!add x y = x + y!}). Assenza o
      minimizzazione degli effetti collaterali (\emph{side effects}).
    \item
      \emph{Logico} (Prolog): il programma definisce relazioni, fatti e
      regole di deduzione logica; l'esecuzione è una
      risoluzione/unificazione guidata da query (es.
      \passthrough{\lstinline!playsAirGuitar(sam) :- listens2Music(sam).!}).
    \end{itemize}
  \item
    \textbf{Object-Oriented (OOP)} (Smalltalk, Java, C++):

    \begin{itemize}
    \tightlist
    \item
      Unisce il nucleo imperativo (gli oggetti contengono uno stato
      interno) a tratti dichiarativi (gli oggetti espongono
      comportamenti invocabili senza conoscerne l'algoritmo interno).
    \item
      Basato su \textbf{istanze} che comunicano tramite
      \textbf{passaggio di messaggi}.
    \end{itemize}
  \end{enumerate}
\item
  \textbf{Lo Spettro dell'Astrazione e il Trade-off di Overhead}:
  \[\text{Machine Code} \longrightarrow \text{Assembly} \longrightarrow \text{Procedurale} \longrightarrow \text{Object-Oriented} \longrightarrow \text{Funzionale}\]

  \begin{itemize}
  \tightlist
  \item
    Più si sale nell'astrazione:

    \begin{itemize}
    \tightlist
    \item
      \textbf{Programmer Overhead \(\searrow\) (minimo)}: meno tempo
      perso a gestire registri, puntatori e memoria manuale.
    \item
      \textbf{Compiler/Runtime Overhead \(\nearrow\) (massimo)}: il
      compilatore/JVM deve effettuare controlli di tipo a runtime,
      garbage collection e risoluzione dinamica delle chiamate
      (\emph{dynamic dispatch}).
    \end{itemize}
  \end{itemize}
\item
  \textbf{Linea temporale}: dal 1949 (Assembly), Fortran (1957), C
  (1972), C++ (1983) fino all'avvento di \textbf{Java (1995)}, C\#
  (2000) e linguaggi moderni multiparadigma (Swift 2014).
\end{itemize}

\begin{center}\rule{0.5\linewidth}{0.5pt}\end{center}

\subsection{La Crisi del Procedurale e la Nascita degli Oggetti (Slide
14--23)}\label{chap02_la-crisi-del-procedurale-e-la-nascita-degli-oggetti-slide-1423}

Il prof. Pasqua presenta il caso di studio di un programma in C che
manipola un vettore globale:

\begin{lstlisting}[language=C]
int vect[20];
void sort() { /* ordina vect */ }
int search(int n) { /* cerca in vect */ }
void init() { /* inizializza vect */ }
int i;

void main() {
    init();
    sort();
    search(13);
}
\end{lstlisting}

\textbf{I 4 Difetti Letali dell'Approccio Procedurale evidenziati a
lezione}: 1. \textbf{Relazione Implicita Dati-Funzioni}:
\passthrough{\lstinline!vect!} e le funzioni
\passthrough{\lstinline!sort()!}, \passthrough{\lstinline!search()!}
sono entità slegate a livello di linguaggio. Non esiste un tipo autonomo
``Vettore'' che protegga se stesso. 2. \textbf{Nessun Controllo dei
Confini (Out-of-Bounds)}: un ciclo come
\passthrough{\lstinline!for (i=0; i<=21; i++) vect[i]=0;!} scrive in
memoria oltre i 20 elementi, corrompendo altre variabili o lo stack
senza alcun errore del compilatore (\emph{buffer overflow}). 3.
\textbf{Nessuna Garanzia sull'Inizializzazione}: se il programmatore
dimentica di invocare \passthrough{\lstinline!init()!}, le funzioni
\passthrough{\lstinline!sort()!} e \passthrough{\lstinline!search()!}
lavorano su memoria sporca non inizializzata. 4. \textbf{Perdita di
Confidenzialità e Integrità}: \passthrough{\lstinline!vect!} è
accessibile in lettura/scrittura da \emph{qualsiasi} funzione del
programma. Su sistemi reali questo genera il cosiddetto
\textbf{Spaghetti Code}. - \textbf{La Regola delle Linee di Codice
(LoC)}: - Per script \textless{} 30 LoC: lo stile procedurale è rapido e
facile da mantenere. - Per progetti \textgreater{} 1000 LoC con team
concorrenti: lo stile procedurale collassa; solo l'OOP permette di
modularizzare garantendo la manutenibilità. - \textbf{La Soluzione a
Oggetti}: - Unire nello stesso modulo sia i dati
(\textbf{attributi/campi}) sia le operazioni su di essi
(\textbf{metodi}). - L'oggetto diventa l'unico responsabile della
consistenza del proprio stato.

\begin{center}\rule{0.5\linewidth}{0.5pt}\end{center}

\subsection{Modellazione UML e Concetti Fondamentali di OOP (Slide
24--36)}\label{chap02_modellazione-uml-e-concetti-fondamentali-di-oop-slide-2436}

\begin{itemize}
\tightlist
\item
  \textbf{UML (\emph{Unified Modeling Language})}: linguaggio grafico
  standard per specificare e documentare sistemi ad oggetti.
\item
  \textbf{Differenza Ontologica tra Classe e Oggetto}:

  \begin{itemize}
  \tightlist
  \item
    \textbf{Classe}: è il \emph{tipo} (blueprint / modello astratto).
    Definisce gli attributi e i metodi. Non occupa memoria nello heap
    per i dati dell'istanza finché non viene istanziata.
  \item
    \textbf{Oggetto}: è l'\textbf{istanza concreta} nello \textbf{Heap}.
    Possiede:

    \begin{enumerate}
    \def\labelenumi{\arabic{enumi}.}
    \tightlist
    \item
      \textbf{Stato}: i valori correnti associati ai suoi attributi.
    \item
      \textbf{Comportamento}: i metodi che può eseguire.
    \item
      \textbf{Identità}: indirizzo di memoria univoco che lo distingue
      da tutti gli altri oggetti, anche se aventi lo stesso identico
      stato!
    \end{enumerate}
  \end{itemize}
\item
  \textbf{Rappresentazione grafica UML}:

  \begin{itemize}
  \tightlist
  \item
    \textbf{Classe (Rettangolo a 3 sezioni)}:

    \begin{enumerate}
    \def\labelenumi{\arabic{enumi}.}
    \tightlist
    \item
      Nome della classe (es. \passthrough{\lstinline!Car!}).
    \item
      Attributi con tipo (es.
      \passthrough{\lstinline!manufacturer : string!},
      \passthrough{\lstinline!model : string!}).
    \item
      Metodi con segnatura (es.
      \passthrough{\lstinline!canFly() : boolean!},
      \passthrough{\lstinline!setFrameNumber(num : uint)!}).
    \end{enumerate}
  \item
    \textbf{Oggetto / Istanza (Rettangolo con nome sottolineato)}:

    \begin{itemize}
    \tightlist
    \item
      Es. \passthrough{\lstinline!Car1: Car!} con
      \passthrough{\lstinline!manufacturer = DeLorean!},
      \passthrough{\lstinline!model = DMC-12!}.
    \end{itemize}
  \end{itemize}
\item
  \textbf{Interazione a Passaggio di Messaggi (\emph{Message Passing})}:

  \begin{itemize}
  \tightlist
  \item
    Gli oggetti non leggono né modificano la memoria altrui
    direttamente: inviano richieste di servizio (\emph{messaggi}).
  \item
    L'\textbf{interfaccia} dell'oggetto è l'insieme di messaggi che esso
    accetta. Se invio un messaggio non previsto, si genera un errore di
    compilazione.
  \end{itemize}
\item
  \textbf{Visibilità}:

  \begin{itemize}
  \tightlist
  \item
    \passthrough{\lstinline!+!} (\passthrough{\lstinline!public!}):
    accessibile da qualsiasi classe.
  \item
    \passthrough{\lstinline!-!} (\passthrough{\lstinline!private!}):
    accessibile esclusivamente dall'interno della classe stessa.
  \end{itemize}
\item
  \textbf{Incapsulamento vs Information Hiding}:

  \begin{itemize}
  \tightlist
  \item
    \textbf{Incapsulamento}: impacchettare dati e codice in una singola
    unità logica (la classe).
  \item
    \textbf{Information Hiding}: nascondere i dettagli implementativi
    (rendendo i campi \passthrough{\lstinline!private!}) ed esporre solo
    un'interfaccia controllata (metodi
    \passthrough{\lstinline!public!}).
  \item
    \emph{Vantaggio chiave}: se cambiamo la struttura dati interna (es.
    da array a tabella hash), il codice esterno che usa i metodi
    pubblici non subisce alcun impatto (\emph{zero refactoring
    esterno}).
  \end{itemize}
\item
  \textbf{Associazioni UML}:

  \begin{itemize}
  \tightlist
  \item
    \textbf{Associazione standard}: linea continua con cardinalità (es.
    \passthrough{\lstinline!1!} a \passthrough{\lstinline!0..n!} tra
    \passthrough{\lstinline!Insurance!} e
    \passthrough{\lstinline!Vehicle!}).
  \item
    \textbf{Ereditarietà (Generalizzazione)}: freccia con punta
    triangolare vuota orientata verso la superclasse
    (\passthrough{\lstinline!Car!} estende
    \passthrough{\lstinline!Vehicle!}, \passthrough{\lstinline!Truck!}
    estende \passthrough{\lstinline!Vehicle!}).
  \end{itemize}
\item
  \textbf{I 4 Pilastri dell'OOP}:

  \begin{enumerate}
  \def\labelenumi{\arabic{enumi}.}
  \tightlist
  \item
    \textbf{Incapsulamento} (\emph{Encapsulation}): bundling e
    protezione dello stato interno.
  \item
    \textbf{Astrazione} (\emph{Abstraction}): separazione netta tra cosa
    un componente fa (interfaccia pubblica) e come lo fa
    (implementazione privata).
  \item
    \textbf{Ereditarietà} (\emph{Inheritance}): meccanismo di riuso e
    specializzazione gerarchica di proprietà e comportamenti.
  \item
    \textbf{Polimorfismo} (\emph{Polymorphism}): capacità di trattare
    istanze di classi diverse in modo uniforme attraverso un tipo o
    interfaccia comune.
  \end{enumerate}
\end{itemize}

\begin{center}\rule{0.5\linewidth}{0.5pt}\end{center}

\section{Codice: Esempio Comparativo: Dal Vettore C all'Oggetto
Java}\label{chap02_codice-esempio-comparativo-dal-vettore-c-alloggetto-java}

Per apprezzare concretamente i concetti di T02, confrontiamo il
frammento C difettoso mostrato dal docente con la sua trasposizione
nativa ad oggetti in Java.

\subsection{Il Problema in C (Codice procedurale
insicuro)}\label{chap02_il-problema-in-c-codice-procedurale-insicuro}

\begin{lstlisting}[language=C]
// C procedurale: il dato è nudo e vulnerabile
int vect[20];
void init() { /* ... */ }

int main() {
    vect[25] = 999; // Corruzione di memoria silenziosa! Nessun errore a runtime.
    return 0;
}
\end{lstlisting}

\subsection{La Soluzione ad Oggetti in Java (Incapsulamento, Information
Hiding,
Costruttore)}\label{chap02_la-soluzione-ad-oggetti-in-java-incapsulamento-information-hiding-costruttore}

In Java, l'oggetto controlla se stesso. Non esiste memoria non
inizializzata e l'accesso fuori limite lancia un'eccezione esplicita a
runtime:

\begin{lstlisting}[language=Java]
public class SafeVector {
    // 1. Information Hiding: il dato interno è inaccessibile all'esterno
    private final int[] elements;

    // 2. Garanzia di Inizializzazione: il costruttore viene eseguito all'atto del 'new'
    public SafeVector(int size) {
        if (size <= 0) {
            throw new IllegalArgumentException("La dimensione deve essere positiva");
        }
        this.elements = new int[size]; // Nello Heap, azzerato di default a 0
    }

    // 3. Controllo dei confini e integrità: interfaccia controllata
    public void set(int index, int value) {
        // La JVM effettua l'ArrayBoundsCheck automatico, impedendo corruzioni di memoria
        this.elements[index] = value;
    }

    public int get(int index) {
        return this.elements[index];
    }

    public int size() {
        return this.elements.length;
    }
}
\end{lstlisting}

\begin{center}\rule{0.5\linewidth}{0.5pt}\end{center}

\begin{studynote}[Punto Importante per l'Esame] \textbf{Takeaway per l'esame}: - Una classe non alloca
memoria per lo stato fino alla chiamata a \passthrough{\lstinline!new!}
(che alloca l'oggetto nello \textbf{Heap}). - L'Information Hiding si
realizza marcando i campi di istanza come
\passthrough{\lstinline!private!}. - I 4 pilastri dell'OOP
(\emph{Encapsulation}, \emph{Abstraction}, \emph{Inheritance},
\emph{Polymorphism}) sono il filo conduttore dell'intero corso.
\end{studynote}

\begin{center}\rule{0.5\linewidth}{0.5pt}\end{center}

\begin{studynote}[Nota di Studio] Il prossimo blocco è la \textbf{Lezione 3 / Slide T03:
\emph{Building, Running and Deployment in Java}}, in cui studieremo: -
Il modello di esecuzione Java: \textbf{Sorgente (.java) \(\rightarrow\)
Bytecode (.class) \(\rightarrow\) JVM / JIT}. - I comandi CLI
fondamentali: \passthrough{\lstinline!javac!},
\passthrough{\lstinline!java!}, \passthrough{\lstinline!jar!},
\passthrough{\lstinline!javadoc!}. - I concetti di \textbf{Classpath}
(\passthrough{\lstinline!-cp!}) e gestione delle dipendenze. -
\textbf{Apache Maven}: cicli di vita (\emph{compile, test, package,
install}), struttura directory standard
(\passthrough{\lstinline!src/main/java!},
\passthrough{\lstinline!src/test/java!}) e il
\passthrough{\lstinline!pom.xml!}.
\end{studynote}

\textbf{Dimmi quando sei pronto per passare a T03!}

\newpage
\chapter[T03 --- Compilazione, Esecuzione e Deployme]{T03 --- Compilazione, Esecuzione e Deployment}
\label{chap:t03}

{\small\sffamily\color{accentblue!80!black}\textit{JVM, Bytecode, JDK vs JRE, Strumenti CLI (javac, java, jar) e Apache Maven}}

\vspace{0.8em}

In questo blocco analizziamo l'ecosistema tecnologico di Java,
l'architettura della \textbf{Java Virtual Machine (JVM)}, i meccanismi
di compilazione, il \textbf{Dynamic Class Loading}, la gestione del
\textbf{Classpath} e le strategie di impacchettamento con file
\textbf{JAR} ed eseguibili
(\href{file:///home/wally/Documenti/GitHub/Programmazione-II/25-26/Teoria-e-Esercitazioni/Slides-20260902/T03\%20-\%20Building,\%20Running\%20and\%20Deployment\%20in\%20Java.pdf}{T03
- Building, Running and Deployment in Java.pdf}).

\section{Teoria: Concetti Teorici dalle
Slide}\label{chap03_teoria-concetti-teorici-dalle-slide}

\subsection{Il Linguaggio Java e le sue Proprietà Fondamentali (Slide
1--11)}\label{chap03_il-linguaggio-java-e-le-sue-proprietuxe0-fondamentali-slide-111}

\begin{itemize}
\tightlist
\item
  \textbf{Cenni Storici}:

  \begin{itemize}
  \tightlist
  \item
    Creato da \textbf{James Gosling} presso Sun Microsystems (progetto
    \emph{Oak} nel 1991, rilasciato ufficialmente come Java nel
    \textbf{1995}, poi acquisito da Oracle).
  \item
    Motto: \emph{``Write Once, Run Everywhere''} (WORA).
  \end{itemize}
\item
  \textbf{Le Caratteristiche Chiave del Linguaggio}:

  \begin{enumerate}
  \def\labelenumi{\arabic{enumi}.}
  \tightlist
  \item
    \textbf{Robustezza}:

    \begin{itemize}
    \tightlist
    \item
      \emph{Tipizzazione forte e statica (Strongly Typed)}: ogni
      variabile ed espressione ha un tipo noto a tempo di compilazione;
      i vincoli di tipo sono verificati dal compilatore prima
      dell'esecuzione.
    \item
      \emph{Nessuna manipolazione esplicita di puntatori}: in Java non
      esistono l'operatore di indirizzo \passthrough{\lstinline!\&!},
      l'aritmetica dei puntatori né la deallocazione manuale
      (\passthrough{\lstinline!free()!}). Questo elimina alla radice
      \emph{dangling pointers}, \emph{segmentation fault} e \emph{buffer
      overflow}.
    \item
      \emph{Controlli a runtime}: la JVM esegue verifiche automatiche su
      ogni accesso ad array
      (\passthrough{\lstinline!ArrayIndexOutOfBoundsException!}) e sui
      riferimenti (\passthrough{\lstinline!NullPointerException!}).
    \item
      \emph{Garbage Collection (GC) automatica}: un thread demone a
      bassa priorità della JVM individua ed elimina gli oggetti non più
      raggiungibili nello Heap, riducendo drasticamente i \emph{memory
      leak}.
    \item
      \emph{Gestione strutturata degli errori}: gestione degli stati
      anomali tramite eccezioni controllate e non controllate
      (\passthrough{\lstinline!try-catch-finally!}).
    \end{itemize}
  \item
    \textbf{Dinamicità}:

    \begin{itemize}
    \tightlist
    \item
      \emph{Caricamento e linking dinamico (Dynamic Linking)}: le classi
      non vengono collegate in un unico binario monolitico a
      compile-time, ma sono caricate in memoria dalla JVM
      \emph{on-demand}, solo quando effettivamente referenziate dal
      codice in esecuzione.
    \item
      \emph{Allocazione dinamica}: la dimensione delle strutture dati e
      degli array può essere determinata a runtime
      (\passthrough{\lstinline!new int[size]!}).
    \end{itemize}
  \item
    \textbf{Portabilità e Architettura a Bytecode}:

    \begin{itemize}
    \tightlist
    \item
      Nei linguaggi puramente compilati (C/C++), il compilatore genera
      direttamente codice macchina specifico per l'architettura target
      (x86, x64, ARM) e per il sistema operativo (Linux, Windows,
      macOS). Il programmatore deve ricompilare o mantenere build
      cross-platform separate.
    \item
      In Java, il compilatore (\passthrough{\lstinline!javac!}) traduce
      il sorgente \passthrough{\lstinline!.java!} in un formato
      intermedio indipendente dall'hardware: il \textbf{Bytecode}
      (\passthrough{\lstinline!.class!}).
    \item
      È la \textbf{JVM} a farsi carico della traduzione del bytecode in
      istruzioni macchina della CPU ospitante. La portabilità cessa di
      essere un onere del programmatore e diventa un servizio fornito
      dal runtime.
    \end{itemize}
  \item
    \textbf{Prestazioni e Compilatore JIT (\emph{Just-In-Time})}:

    \begin{itemize}
    \tightlist
    \item
      Java non è puramente interpretato. La JVM include un compilatore
      \textbf{JIT} che monitora a runtime le porzioni di bytecode
      eseguite più frequentemente (\emph{hot spots}) e le compila al
      volo in codice macchina nativo ottimizzato, memorizzandole nella
      cache del codice.
    \end{itemize}
  \end{enumerate}
\item
  \textbf{Java vs JavaScript}:

  \begin{itemize}
  \tightlist
  \item
    Citazione di Christian Heilmann: \emph{«Java is to JavaScript what
    Car is to Carpet»} (hanno in comune solo le prime quattro lettere).
    JavaScript nacque da Netscape con quel nome per motivi puramente
    commerciali; i due linguaggi non condividono né modello dei tipi, né
    macchina virtuale, né filosofia di design.
  \end{itemize}
\end{itemize}

\begin{center}\rule{0.5\linewidth}{0.5pt}\end{center}

\subsection{Ecosistema Java: JDK vs JRE (Slide
11)}\label{chap03_ecosistema-java-jdk-vs-jre-slide-11}

La distinzione tra ambiente di sviluppo e ambiente di esecuzione: -
\textbf{JDK (\emph{Java Development Kit})}: ambiente per programmatori:
- \passthrough{\lstinline!javac!}: il compilatore che traduce sorgenti
\passthrough{\lstinline!.java!} in bytecode
\passthrough{\lstinline!.class!}. - \passthrough{\lstinline!java!}: il
launcher applicativo che avvia la JVM. -
\passthrough{\lstinline!javadoc!}: generatore di documentazione HTML a
partire dai commenti speciali \passthrough{\lstinline!/** ... */!}. -
\passthrough{\lstinline!javap!}: disassemblatore di bytecode (permette
di ispezionare il bytecode generato). - \passthrough{\lstinline!jdb!}:
debugger da riga di comando. - \textbf{JRE (\emph{Java Runtime
Environment})}: ambiente minimale per l'utente finale: - \textbf{JVM
(\emph{Java Virtual Machine})}: il motore di esecuzione. - Librerie
standard delle API Java (il modulo base
\passthrough{\lstinline!java.base!} contenente
\passthrough{\lstinline!java.lang!},
\passthrough{\lstinline!java.util!}, \passthrough{\lstinline!java.io!},
ecc.). - Compilatore JIT integrato.

\begin{center}\rule{0.5\linewidth}{0.5pt}\end{center}

\subsection{Compilazione, Esecuzione e Dynamic Class Loading (Slide
12--16)}\label{chap03_compilazione-esecuzione-e-dynamic-class-loading-slide-1216}

\begin{enumerate}
\def\labelenumi{\arabic{enumi}.}
\item
  \textbf{La Struttura del Programma Minimo}:

\begin{lstlisting}[language=Java]
// Salvato obbligatoriamente in MyClass.java (stesso nome della classe pubblica)
public class MyClass {
    public static void main(String[] args) {
        System.out.println("Hello World!");
    }
}
\end{lstlisting}

  \begin{itemize}
  \tightlist
  \item
    \emph{Regola aurea}: In un file \passthrough{\lstinline!.java!} può
    essere presente al massimo \textbf{una sola classe pubblica}, e il
    nome del file deve coincidere esattamente (incluso il
    maiuscolo/minuscolo) con il nome di tale classe.
  \item
    \emph{Firma del metodo \passthrough{\lstinline!main!}}:

    \begin{itemize}
    \tightlist
    \item
      \passthrough{\lstinline!public!}: deve essere accessibile
      dall'esterno da parte del launcher della JVM.
    \item
      \passthrough{\lstinline!static!}: la JVM lo invoca direttamente
      sulla classe, senza dover prima istanziare un oggetto
      (\passthrough{\lstinline!MyClass obj = new MyClass()!}).
    \item
      \passthrough{\lstinline!void!}: non restituisce un codice numerico
      di ritorno (l'uscita anomala si gestisce con
      \passthrough{\lstinline!System.exit(code)!}).
    \item
      \passthrough{\lstinline!String[] args!}: array di stringhe
      contenente i parametri passati da riga di comando.
    \end{itemize}
  \end{itemize}
\item
  \textbf{Pipeline di Esecuzione della JVM}:
  \[\text{Sorgente } (.java) \xrightarrow{\texttt{javac}} \text{Bytecode } (.class) \xrightarrow{\text{Loader}} \text{Bytecode Verifier} \xrightarrow{\text{Interpreter / JIT}} \text{OS / Hardware}\]

  \begin{itemize}
  \tightlist
  \item
    \textbf{Bytecode Verifier}: componente di sicurezza fondamentale
    della JVM. Prima di eseguire qualsiasi file
    \passthrough{\lstinline!.class!}, scansiona le istruzioni per
    verificare che non violino i limiti dello stack, non convertano
    puntatori illegalmente e rispettino le regole di visibilità.
  \end{itemize}
\item
  \textbf{Dynamic Class Loading e Classpath
  (\passthrough{\lstinline!-cp!})}:

  \begin{itemize}
  \item
    La JVM non cerca i file nel filesystem a caso: utilizza il
    \textbf{Classpath} (una lista ordinata di directory e archivi
    \passthrough{\lstinline!.jar!}).
  \item
    Quando durante l'esecuzione il codice fa riferimento per la prima
    volta a una classe \passthrough{\lstinline!X!}:

    \begin{enumerate}
    \def\labelenumii{\arabic{enumii}.}
    \tightlist
    \item
      Il \passthrough{\lstinline!ClassLoader!} cerca
      \passthrough{\lstinline!X.class!} nella prima cartella specificata
      nel Classpath.
    \item
      Se non la trova, passa alla successiva.
    \item
      Se la trova, la carica in memoria e ne inizializza le strutture
      statiche.
    \item
      Se non la trova in nessuna cartella del Classpath, solleva
      \passthrough{\lstinline!ClassNotFoundException!} o
      \passthrough{\lstinline!NoClassDefFoundError!}.
    \end{enumerate}
  \item
    Esempio di esecuzione esplicita:

\begin{lstlisting}[language=bash]
java -cp . MyClass
\end{lstlisting}

    Il parametro \passthrough{\lstinline!-cp .!} istruisce la JVM a
    cercare le classi compilate nella directory corrente
    (\passthrough{\lstinline!.!}). Le classi standard di sistema
    (\passthrough{\lstinline!System!}, \passthrough{\lstinline!String!})
    vengono invece caricate automaticamente dal bootstrap classloader
    della JDK.
  \end{itemize}
\end{enumerate}

\begin{center}\rule{0.5\linewidth}{0.5pt}\end{center}

\subsection{Deployment e Gestione dei File JAR (Slide
17--18)}\label{chap03_deployment-e-gestione-dei-file-jar-slide-1718}

\begin{itemize}
\item
  \textbf{Anatomia di un file JAR (\emph{Java Archive})}:

  \begin{itemize}
  \tightlist
  \item
    Un file \passthrough{\lstinline!.jar!} è fisicamente un
    \textbf{archivio compresso in formato standard ZIP}.
  \item
    Raggruppa decine o centinaia di file
    \passthrough{\lstinline!.class!}, risorse (immagini, file di
    configurazione) e metadati.
  \end{itemize}
\item
  \textbf{Creazione manuale da CLI}:

\begin{lstlisting}[language=bash]
jar cvf MyJar.jar MyClass.class
\end{lstlisting}

  \begin{itemize}
  \tightlist
  \item
    \passthrough{\lstinline!c!}: \emph{create} (crea nuovo archivio).
  \item
    \passthrough{\lstinline!v!}: \emph{verbose} (mostra i file
    aggiunti).
  \item
    \passthrough{\lstinline!f!}: \emph{file} (il parametro successivo è
    il nome del file \passthrough{\lstinline!.jar!} da produrre).
  \end{itemize}
\item
  \textbf{Esecuzione tramite Classpath}:

\begin{lstlisting}[language=bash]
java -cp MyJar.jar MyClass
\end{lstlisting}
\item
  \textbf{JAR Auto-Eseguibile (\passthrough{\lstinline!java -jar!})}:

  \begin{itemize}
  \item
    Per consentire all'utente di lanciare l'archivio direttamente con
    \passthrough{\lstinline!java -jar MyJar.jar!}, il JAR deve
    specificare quale classe ospita il metodo
    \passthrough{\lstinline!main!}.
  \item
    Questa informazione risiede nel file
    \passthrough{\lstinline!META-INF/MANIFEST.MF!} sotto la direttiva:

\begin{lstlisting}
Main-Class: MyClass
\end{lstlisting}
  \item
    Creazione con file manifest esterno:

\begin{lstlisting}[language=bash]
jar cvfm MyClass.jar manifest.txt MyClass.class
\end{lstlisting}
  \item
    Creazione con sintassi moderna senza file di testo intermedio (flag
    \passthrough{\lstinline!e!} per entrypoint):

\begin{lstlisting}[language=bash]
jar cfe MyClass.jar MyClass MyClass.class
\end{lstlisting}
  \end{itemize}
\end{itemize}

\begin{center}\rule{0.5\linewidth}{0.5pt}\end{center}

\subsection{Strumenti di Sviluppo: Da Editor a IntelliJ IDEA (Slide
19--26)}\label{chap03_strumenti-di-sviluppo-da-editor-a-intellij-idea-slide-1926}

\begin{itemize}
\tightlist
\item
  \textbf{CLI vs IDE}: per progetti complessi, l'uso del terminale e di
  un editor di testo semplice diventa poco scalabile.
\item
  \textbf{Vantaggi dell'IDE (\emph{Integrated Development
  Environment})}:

  \begin{itemize}
  \tightlist
  \item
    \emph{Syntax highlighting} e \emph{live code inspection/linting}
    (segnalazione immediata degli errori di tipo e code smell prima
    ancora di compilare).
  \item
    \emph{Debugger integrato} con breakpoint, ispezione dello stack dei
    frame e monitoraggio variabili.
  \item
    \emph{Integrazione con Version Control (Git)} e \emph{Build
    Automation Tools (Maven / Gradle)}.
  \item
    \emph{Monito del docente}: l'auto-completamento o la generazione
    automatica di codice da parte dell'IDE è utile in ambito
    professionale, ma va limitata durante la fase di apprendimento per
    comprendere a fondo le regole sintattiche e semantiche del
    linguaggio.
  \end{itemize}
\item
  \textbf{Scelta del corso}: \textbf{IntelliJ IDEA} (gratuito in
  versione Community o con licenza Education per studenti universitari).
\end{itemize}

\begin{center}\rule{0.5\linewidth}{0.5pt}\end{center}

\section{\texorpdfstring{Codice: Collegamento Pratico con il Tuo
Progetto: Maven e
\texttt{guida-jar.md}}{Codice: Collegamento Pratico con il Tuo Progetto: Maven e guida-jar.md}}\label{chap03_codice-collegamento-pratico-con-il-tuo-progetto-maven-e-guida-jar.md}

Nel tuo workspace
\href{file:///home/wally/Documenti/GitHub/Programmazione-II/25-26/esercizi-wally-25-26}{\passthrough{\lstinline!esercizi-wally-25-26!}},
hai formalizzato questi esatti concetti in modo eccellente in
\href{file:///home/wally/Documenti/GitHub/Programmazione-II/25-26/esercizi-wally-25-26/guida-jar.md}{\passthrough{\lstinline!guida-jar.md!}}
e nei file di configurazione Maven:

\begin{enumerate}
\def\labelenumi{\arabic{enumi}.}
\tightlist
\item
  \textbf{Il Fully Qualified Name e il Filesystem}:

  \begin{itemize}
  \tightlist
  \item
    Come visto a lezione, se una classe appartiene al package
    \passthrough{\lstinline!es01!}, il bytecode generato deve risiedere
    fisicamente nel percorso \passthrough{\lstinline!es01/Date.class!}.
  \item
    Quando si invoca la JVM, il nome della classe da passare deve essere
    il FQN (\passthrough{\lstinline!es01.MainDate!}).
  \end{itemize}
\item
  \textbf{Automazione del Packaging con Maven}:

  \begin{itemize}
  \tightlist
  \item
    Invece di lanciare manualmente \passthrough{\lstinline!jar cfe!},
    nei singoli moduli (es.
    \href{file:///home/wally/Documenti/GitHub/Programmazione-II/25-26/esercizi-wally-25-26/es01/pom.xml}{\passthrough{\lstinline!es01/pom.xml!}})
    usi il plugin
    \textbf{\passthrough{\lstinline!maven-assembly-plugin!}}:
  \end{itemize}

\begin{lstlisting}[language=XML]
<plugin>
    <groupId>org.apache.maven.plugins</groupId>
    <artifactId>maven-assembly-plugin</artifactId>
    <version>3.7.1</version>
    <configuration>
        <archive>
            <manifest>
                <mainClass>es01.MainDate</mainClass>
            </manifest>
        </archive>
        <descriptorRefs>
            <descriptorRef>jar-with-dependencies</descriptorRef>
        </descriptorRefs>
    </configuration>
    <executions>
        <execution>
            <id>make-assembly</id>
            <phase>package</phase>
            <goals>
                <goal>single</goal>
            </goals>
        </execution>
    </executions>
</plugin>
\end{lstlisting}

  Questo genera automaticamente nella cartella
  \passthrough{\lstinline!target/!} il cosiddetto \emph{Fat JAR}
  contenente sia il bytecode dell'applicazione sia il file
  \passthrough{\lstinline!META-INF/MANIFEST.MF!} precompilato.
\end{enumerate}

\begin{center}\rule{0.5\linewidth}{0.5pt}\end{center}

\section{{[}Test{]} Verifica dei Comandi da
Terminale}\label{chap03_test-verifica-dei-comandi-da-terminale}

Puoi verificare l'intero ciclo di build ed esecuzione da riga di comando
testando il modulo \passthrough{\lstinline!es01!} tramite lo script:

\begin{lstlisting}[language=bash]
./esercizi.sh package es01
\end{lstlisting}

Lo script esegue internamente: 1.
\passthrough{\lstinline!mvn clean package -pl es01!} 2. Individua il
file JAR autoprodotto in
\passthrough{\lstinline!es01/target/es01-*-jar-with-dependencies.jar!}
3. Esegue \passthrough{\lstinline!java -jar ...!} verificando che
l'entrypoint \passthrough{\lstinline!es01.MainDate!} risponda
correttamente.

\begin{center}\rule{0.5\linewidth}{0.5pt}\end{center}

\begin{studynote}[Nota di Studio] Con T03 abbiamo concluso l'intera panoramica propedeutica
(Introduzione, Paradigmi/UML e Tooling/JVM).

A partire dal prossimo blocco entriamo nel codice sorgente vero e
proprio con la \textbf{Lezione 4 / Slide T04: \emph{Types and Objects}}
e il file ufficiale del docente
\href{file:///home/wally/Documenti/GitHub/Programmazione-II/25-26/Teoria-e-Esercitazioni/Codice-20260902/Lezione04/Example.java}{\passthrough{\lstinline!Example.java!}}:
- Tipi primitivi (dimensioni, range, valori di default) vs Reference
Types. - Rappresentazione dei caratteri (\passthrough{\lstinline!char!}
come intero senza segno a 16 bit, codifica ASCII e Unicode
\passthrough{\lstinline!\\u0056!}). - Regole di inizializzazione:
variabili di istanza/statiche (default zero/null) vs variabili locali
dello stack (nessun default \(\rightarrow\) errore a compile-time!). -
Operatori logici standard (\passthrough{\lstinline!\&!},
\passthrough{\lstinline!|!}) vs operatori a corto circuito
(\emph{short-circuit} \passthrough{\lstinline!\&\&!},
\passthrough{\lstinline!||!}) e relativi effetti collaterali
(\passthrough{\lstinline!k++!}). - Scope dei blocchi di codice annidati
(\passthrough{\lstinline!\{ ... \}!}) e shadowing/visibilità delle
variabili.
\end{studynote}

\textbf{Dammi conferma per aprire la Lezione 4 / T04 e analizzare teoria
e codice di \passthrough{\lstinline!Example.java!}!}

\newpage
\chapter[T04 --- Sistema di Tipi, Primitive e Refere]{T04 --- Sistema di Tipi, Primitive e Reference}
\label{chap:t04}

{\small\sffamily\color{accentblue!80!black}\textit{Tipi Primitivi vs Reference, Casting, Wrapper, Operatori e Analisi di Example.java}}

\vspace{0.8em}

Con questo blocco entriamo nel vivo della semantica del linguaggio:
analizziamo come la JVM gestisce i tipi, la memoria (Stack vs Heap), i
puntatori/reference, gli operatori logici e lo scope, confrontando la
teoria di
\href{file:///home/wally/Documenti/GitHub/Programmazione-II/25-26/Teoria-e-Esercitazioni/Slides-20260902/T04\%20-\%20Types\%20and\%20Objects.pdf}{T04
- Types and Objects.pdf} con il primo file di codice ufficiale del
docente:
\href{file:///home/wally/Documenti/GitHub/Programmazione-II/25-26/Teoria-e-Esercitazioni/Codice-20260902/Lezione04/Example.java}{\passthrough{\lstinline!Example.java!}}.

\section{Teoria: Concetti Teorici dalle
Slide}\label{chap04_teoria-concetti-teorici-dalle-slide}

\subsection{Tipizzazione dei Linguaggi e il ``Contratto'' di Tipo (Slide
1--7)}\label{chap04_tipizzazione-dei-linguaggi-e-il-contratto-di-tipo-slide-17}

\begin{itemize}
\tightlist
\item
  \textbf{Definizione formale informale di Tipo}: Un tipo è una coppia
  formata da un \textbf{insieme di valori ammissibili} e da un
  \textbf{insieme di operazioni consentite} su di essi.

  \begin{itemize}
  \tightlist
  \item
    Es. \passthrough{\lstinline!int!}: valori da \(-2^{31}\) a
    \(2^{31}-1\), operazioni \passthrough{\lstinline!+!},
    \passthrough{\lstinline!-!}, \passthrough{\lstinline!*!},
    \passthrough{\lstinline!/!}, \passthrough{\lstinline!\%!}.
  \item
    Il tipo definisce un \textbf{contratto}: vincola cosa una variabile
    può contenere e in quali espressioni può comparire.
  \end{itemize}
\item
  \textbf{Weakly Typed (es. Python, JavaScript) vs Strongly Typed
  (Java)}:

  \begin{itemize}
  \tightlist
  \item
    Nei linguaggi \emph{weakly typed}, le variabili sono contenitori
    generici slegati dal tipo: una variabile \passthrough{\lstinline!x!}
    può contenere un intero, poi una stringa, poi una tupla.
  \item
    Nei linguaggi \emph{strongly typed} come Java, ogni variabile
    richiede una dichiarazione esplicita e immutabile di tipo a tempo di
    compilazione (\passthrough{\lstinline!int x;!}). Assegnare un tipo
    incompatibile genera un \textbf{Compile-Time Error}.
  \item
    \emph{Meme di JavaScript a lezione}:

    \begin{itemize}
    \tightlist
    \item
      \passthrough{\lstinline!"11" + 1!} \(\implies\)
      \passthrough{\lstinline!"111"!} (coercizione implicita a stringa e
      concatenazione).
    \item
      \passthrough{\lstinline!"11" - 1!} \(\implies\)
      \passthrough{\lstinline!10!} (coercizione implicita a numero e
      sottrazione numerica).
    \item
      In Java queste ambiguità ed errori silenziosi sono impossibili: il
      compilatore blocca ogni operazione non conforme prima
      dell'esecuzione.
    \end{itemize}
  \end{itemize}
\end{itemize}

\begin{center}\rule{0.5\linewidth}{0.5pt}\end{center}

\subsection{I Tipi Primitivi in Java e le Conversioni (Slide
8--13)}\label{chap04_i-tipi-primitivi-in-java-e-le-conversioni-slide-813}

Java dispone di \textbf{8 tipi primitivi} (allocati direttamente per
valore nello Stack o inline negli oggetti, non sono oggetti):

\begin{longtable}[]{@{}
  >{\raggedright\arraybackslash}p{(\linewidth - 6\tabcolsep) * \real{0.2500}}
  >{\raggedright\arraybackslash}p{(\linewidth - 6\tabcolsep) * \real{0.2500}}
  >{\raggedright\arraybackslash}p{(\linewidth - 6\tabcolsep) * \real{0.2500}}
  >{\raggedright\arraybackslash}p{(\linewidth - 6\tabcolsep) * \real{0.2500}}@{}}
\toprule\noalign{}
\begin{minipage}[b]{\linewidth}\raggedright
Tipo Primitivo
\end{minipage} & \begin{minipage}[b]{\linewidth}\raggedright
Dimensione
\end{minipage} & \begin{minipage}[b]{\linewidth}\raggedright
Range di Valori
\end{minipage} & \begin{minipage}[b]{\linewidth}\raggedright
Default (campi)
\end{minipage} \\
\midrule\noalign{}
\endhead
\bottomrule\noalign{}
\endlastfoot
\passthrough{\lstinline!byte!} & 8 bit (1 byte) & da \(-128\) a \(+127\)
(\(-2^7 \dots 2^7-1\)) & \passthrough{\lstinline!0!} \\
\passthrough{\lstinline!short!} & 16 bit (2 byte) & da \(-32.768\) a
\(+32.767\) (\(-2^{15} \dots 2^{15}-1\)) &
\passthrough{\lstinline!0!} \\
\textbf{\passthrough{\lstinline!int!}} & 32 bit (4 byte) & da
\(-2.147.483.648\) a \(+2.147.483.647\) (\(-2^{31} \dots 2^{31}-1\)) &
\passthrough{\lstinline!0!} \\
\passthrough{\lstinline!long!} & 64 bit (8 byte) & da \(-2^{63}\) a
\(+2^{63}-1\) (suffisso \passthrough{\lstinline!L!} o
\passthrough{\lstinline!l!}) & \passthrough{\lstinline!0L!} \\
\passthrough{\lstinline!float!} & 32 bit IEEE 754 & Precisione singola
(suffisso obbligatorio \passthrough{\lstinline!f!} o
\passthrough{\lstinline!F!}) & \passthrough{\lstinline!0.0f!} \\
\textbf{\passthrough{\lstinline!double!}} & 64 bit IEEE 754 & Precisione
doppia (default per numeri con la virgola) &
\passthrough{\lstinline!0.0d!} \\
\textbf{\passthrough{\lstinline!char!}} & 16 bit (2 byte) & da
\passthrough{\lstinline!0!} a \passthrough{\lstinline!65.535!} (Unicode
UTF-16, \textbf{unsigned}) & \passthrough{\lstinline!'\\u0000'!}
(NUL) \\
\textbf{\passthrough{\lstinline!boolean!}} & 1 bit logico &
\passthrough{\lstinline!true!} oppure \passthrough{\lstinline!false!} &
\passthrough{\lstinline!false!} \\
\end{longtable}

\begin{itemize}
\tightlist
\item
  \textbf{Aritmetica Modulare e Overflow degli
  \passthrough{\lstinline!int!}}:

  \begin{itemize}
  \tightlist
  \item
    In Java non viene lanciata alcuna eccezione se un calcolo intero
    supera il limite massimo: avviene il \textbf{wrap-around} secondo
    l'aritmetica del complemento a due a 32 bit.
  \item
    Se \passthrough{\lstinline!a = 2147483647!} (\(2^{31}-1\)),
    l'operazione \passthrough{\lstinline!a + 1!} restituisce
    \(-2147483648\) (\(-2^{31}\)).
  \end{itemize}
\item
  \textbf{Conversioni di Tipo (Casting)}:

  \begin{itemize}
  \tightlist
  \item
    \textbf{Widening (Allargamento / Implicito)}: da un tipo più stretto
    a uno più ampio (nessuna perdita d'ordine di grandezza):
    \[\text{byte} \longrightarrow \text{short} \longrightarrow \text{int} \longrightarrow \text{long} \longrightarrow \text{float} \longrightarrow \text{double}\]
    Es: \passthrough{\lstinline!float f = 3;!} converte automaticamente
    \passthrough{\lstinline!int 3!} in \passthrough{\lstinline!3.0f!}.
  \item
    \textbf{Narrowing (Restringimento / Esplicito obbligatorio)}: da un
    tipo più ampio a uno più stretto, richiede il cast sintattico
    \passthrough{\lstinline!(tipo)!} perché comporta potenziale
    troncamento o perdita di informazione: Es:
    \passthrough{\lstinline!int i = (int) 3.2;!} scarta la parte
    decimale e memorizza \passthrough{\lstinline!3!}.
  \end{itemize}
\item
  \textbf{I Caratteri sono Numeri Interi Unsigned a 16 bit}:

  \begin{itemize}
  \tightlist
  \item
    Un \passthrough{\lstinline!char!} memorizza il codice numerico
    Unicode (UTF-16 code unit).
  \item
    \passthrough{\lstinline!'V'!}, \passthrough{\lstinline!(char) 86!} e
    \passthrough{\lstinline!'\\u0056'!} sono tre rappresentazioni
    letterali dello \textbf{stesso identico dato binario} (valore
    decimale 86).
  \end{itemize}
\item
  \textbf{Tipi Non-Primitivi (Reference Types) e la Classe
  \passthrough{\lstinline!String!}}:

  \begin{itemize}
  \tightlist
  \item
    Tutto ciò che non è un tipo primitivo è un \textbf{oggetto}.
  \item
    \passthrough{\lstinline!String!} è una classe immutabile.
    L'operatore \passthrough{\lstinline!+!} tra una stringa e qualsiasi
    altro tipo (primitivo o oggetto) applica automaticamente
    l'overloading di concatenazione, invocando internamente la
    conversione a stringa.
  \end{itemize}
\end{itemize}

\begin{center}\rule{0.5\linewidth}{0.5pt}\end{center}

\subsection{Modello di Memoria, Puntatori e Garbage Collection (Slide
14--21)}\label{chap04_modello-di-memoria-puntatori-e-garbage-collection-slide-1421}

\begin{itemize}
\tightlist
\item
  \textbf{Assenza di Variabili Globali}: In Java non esiste lo scope
  globale in stile C. Ogni dato o funzione appartiene a una classe o a
  un'istanza.
\item
  \textbf{La Verità sui Puntatori in Java}:

  \begin{itemize}
  \tightlist
  \item
    \emph{«Java does have pointers, but there is no pointer
    arithmetic»}: ogni variabile di tipo non primitivo è tecnicamente un
    \textbf{puntatore/riferimento} (\emph{reference}) a un blocco di
    memoria nello \textbf{Heap}.
  \item
    La JVM protegge la memoria: il programmatore non può vedere
    l'indirizzo fisico, né fare \passthrough{\lstinline!ptr++!} o
    dereferenziare offset arbitrari.
  \item
    Il passaggio dei parametri ai metodi in Java è \textbf{sempre
    rigorosamente per valore} (\emph{pass-by-value}): per gli oggetti,
    viene copiato per valore il \textbf{riferimento} (l'indirizzo
    logico).
  \end{itemize}
\item
  \textbf{Rappresentazione di un Reference}:

  \begin{itemize}
  \tightlist
  \item
    Una variabile d'istanza non inizializzata contiene
    \passthrough{\lstinline!null!} (es.
    \passthrough{\lstinline!Vehicle@null!}).
  \item
    Quando viene invocato \passthrough{\lstinline!new Vehicle()!}, la
    JVM alloca l'oggetto nello Heap e restituisce l'indirizzo:
    \passthrough{\lstinline!myCar!} conterrà un valore del tipo
    \passthrough{\lstinline!Vehicle@25c7f37d!}.
  \end{itemize}
\item
  \textbf{Errori di Memoria e
  \passthrough{\lstinline!NullPointerException!} (NPE)}:

  \begin{itemize}
  \tightlist
  \item
    Se si tenta di invocare un metodo o accedere a un campo su una
    reference che vale \passthrough{\lstinline!null!} (es.
    \passthrough{\lstinline!myCar2.setFrameNumber(5)!}), la JVM solleva
    una \textbf{\passthrough{\lstinline!NullPointerException!}} a
    runtime, interrompendo l'esecuzione.
  \end{itemize}
\item
  \textbf{Perché i programmatori scrivono
  \passthrough{\lstinline!myCar = null;!}?}:

  \begin{itemize}
  \tightlist
  \item
    L'assegnamento a \passthrough{\lstinline!null!} rimuove il puntatore
    verso l'oggetto nello Heap. Se nessun'altra reference punta a
    quell'area di memoria, l'oggetto diventa \textbf{irraggiungibile}
    (\emph{unreachable}) ed eleggibile per la \textbf{Garbage Collection
    (GC)}, consentendo alla JVM di bonificare la memoria.
  \end{itemize}
\end{itemize}

\begin{center}\rule{0.5\linewidth}{0.5pt}\end{center}

\section{\texorpdfstring{Codice: Analisi Dettagliata del Codice:
\href{file:///home/wally/Documenti/GitHub/Programmazione-II/25-26/Teoria-e-Esercitazioni/Codice-20260902/Lezione04/Example.java}{\texttt{Example.java}}}{Codice: Analisi Dettagliata del Codice: Example.java}}\label{chap04_codice-analisi-dettagliata-del-codice-example.java}

Analizziamo riga per riga il sorgente della Lezione 04, evidenziando
tutte le trappole e i dettagli d'esame:

\begin{lstlisting}[language=Java]
public class Example {
    static float f1;
    static float f2 = 3.4f;
    static String s1;
    static final double PI = 3.14; 
    
    public static void main(String[] args) {
        int i1 = 6;
        int i2;
        char c1 = 'V';
        char c2 = 86;
        char c3 = '\u0056';
        String s2;
            
        System.out.println("i1: " + i1);
        // System.out.println("i2: " + i2); // compile-time error
        System.out.println("f1: " + f1);
        System.out.println("f2: " + f2);
        System.out.println("c1: " + c1 + " c2: " + c2 + " c3: " + c3);
        System.out.println("s1: " + s1);
        /*
         System.out.println("s2: " + s2); // compile-time error
         System.out.println(s1.toString()); // run-time error
         PI = 3.0; // compile-time error
        */
        int k = 0;
        System.out.println(true | k++ == 0);
        System.out.println("k (standard eval): " + k);
        k = 0;
        System.out.println(true || k++ == 0);
        System.out.println("k (short-circuit eval): " + k);
        //
        int n = 4;
        n ++;
        // int n = 8; // compile-time error
        {
            int m = 7;
            m ++;
        }
        {
            int m = 8;
            m --;
        }
        System.out.println("n: " + n);
        // System.out.println("m: " + m); // compile-time error
    }
}
\end{lstlisting}

\subsection{{[}Approfondimento{]} Le Finezze Implementative Spiegate
Riga per
Riga:}\label{chap04_approfondimento-le-finezze-implementative-spiegate-riga-per-riga}

\begin{enumerate}
\def\labelenumi{\arabic{enumi}.}
\tightlist
\item
  \textbf{Variabili Statiche/Campi vs Variabili Locali (\emph{Definite
  Assignment Analysis})}:

  \begin{itemize}
  \tightlist
  \item
    \passthrough{\lstinline!static float f1;!} e
    \passthrough{\lstinline!static String s1;!}: sono campi a livello di
    classe. \textbf{La JVM li inizializza SEMPRE automaticamente al loro
    default}: \passthrough{\lstinline!f1!} diventa
    \passthrough{\lstinline!0.0f!} e \passthrough{\lstinline!s1!}
    diventa \passthrough{\lstinline!null!}.
  \item
    \passthrough{\lstinline!int i2;!} e
    \passthrough{\lstinline!String s2;!}: sono variabili locali allocate
    nel frame dello \textbf{Stack}. \textbf{Le variabili locali NON
    hanno valori di default}.
  \item
    Se tentiamo di leggere \passthrough{\lstinline!i2!} o
    \passthrough{\lstinline!s2!} prima di aver assegnato loro un valore
    (\passthrough{\lstinline!System.out.println(i2)!}), il compilatore
    Java blocca la compilazione con l'errore:
    \passthrough{\lstinline!variable i2 might not have been initialized!}.
  \end{itemize}
\item
  \textbf{Suffisso dei Floating Point (\passthrough{\lstinline!3.4f!})}:

  \begin{itemize}
  \tightlist
  \item
    \passthrough{\lstinline!static float f2 = 3.4f;!}: il letterale
    \passthrough{\lstinline!3.4!} senza suffissi è un
    \passthrough{\lstinline!double!} a 64 bit. Scrivere
    \passthrough{\lstinline!float f = 3.4;!} genera errore di
    compilazione (\emph{loss of precision}). Il suffisso
    \passthrough{\lstinline!f!} forza il letterale a 32 bit.
  \end{itemize}
\item
  \textbf{Costanti con \passthrough{\lstinline!final!}}:

  \begin{itemize}
  \tightlist
  \item
    \passthrough{\lstinline!static final double PI = 3.14;!}: la keyword
    \passthrough{\lstinline!final!} rende la variabile a sola lettura
    dopo l'inizializzazione. Il tentativo di riassegnamento
    \passthrough{\lstinline!PI = 3.0;!} viene respinto dal compilatore
    (\passthrough{\lstinline!cannot assign a value to final variable PI!}).
    Per convenzione Java, le costanti
    \passthrough{\lstinline!static final!} si scrivono in
    \passthrough{\lstinline!MAIUSCOLO\_SNAKE\_CASE!}.
  \end{itemize}
\item
  \textbf{I 3 modi di rappresentare un \passthrough{\lstinline!char!}}:

  \begin{itemize}
  \tightlist
  \item
    \passthrough{\lstinline!c1 = 'V'!}: letterale carattere.
  \item
    \passthrough{\lstinline!c2 = 86!}: valore intero decimale
    ASCII/Unicode per la `V'.
  \item
    \passthrough{\lstinline!c3 = '\\u0056'!}: sequenza di escape Unicode
    a 16 bit in esadecimale (\(5 \times 16 + 6 = 86\)).
  \item
    Stampati a video, producono tutti e tre il carattere
    \passthrough{\lstinline!V!}.
  \end{itemize}
\item
  \textbf{Concatenazione di \passthrough{\lstinline!null!} vs
  \passthrough{\lstinline!NullPointerException!}}:

  \begin{itemize}
  \tightlist
  \item
    \passthrough{\lstinline!System.out.println("s1: " + s1);!} stampa
    \passthrough{\lstinline!"s1: null"!}. L'operatore
    \passthrough{\lstinline!+!} converte in sicurezza il riferimento
    nullo nella stringa letterale \passthrough{\lstinline!"null"!}.
  \item
    \passthrough{\lstinline!System.out.println(s1.toString());!} genera
    invece un \textbf{Run-Time Error}
    (\passthrough{\lstinline!NullPointerException!}): non è possibile
    dereferenziare un puntatore nullo per invocarvi un metodo d'istanza.
  \end{itemize}
\item
  \textbf{Operatori Logici Standard (\passthrough{\lstinline!|!},
  \passthrough{\lstinline!\&!}) vs Corto-Circuito (\emph{Short-Circuit}
  \passthrough{\lstinline!||!}, \passthrough{\lstinline!\&\&!}) e Side
  Effects}:

  \begin{itemize}
  \item
    Caso \passthrough{\lstinline!true | k++ == 0!}:

    \begin{itemize}
    \tightlist
    \item
      L'operatore singolo \passthrough{\lstinline!|!} (non a corto
      circuito) \textbf{valuta SEMPRE entrambi gli operandi}.
    \item
      Anche se la parte sinistra è già \passthrough{\lstinline!true!},
      la parte destra \passthrough{\lstinline!k++ == 0!} viene valutata:
      \passthrough{\lstinline!k!} viene confrontato con
      \passthrough{\lstinline!0!} (dando \passthrough{\lstinline!true!})
      e poi incrementato di \(1\).
    \item
      Stampa a video: \passthrough{\lstinline!true!}, seguito da
      \passthrough{\lstinline!k (standard eval): 1!}.
    \end{itemize}
  \item
    Caso \passthrough{\lstinline!true || k++ == 0!}:

    \begin{itemize}
    \tightlist
    \item
      L'operatore doppio \passthrough{\lstinline!||!} (a corto circuito)
      verifica la parte sinistra: poiché è
      \passthrough{\lstinline!true!}, il risultato globale dell'or
      logico è già matematicamente certo
      (\passthrough{\lstinline!true!}).
    \item
      La parte destra viene \textbf{completamente ignorata e mai
      valutata}!
    \item
      Di conseguenza, \passthrough{\lstinline!k++!} \textbf{non viene
      eseguito}.
    \item
      Stampa a video: \passthrough{\lstinline!true!}, seguito da
      \passthrough{\lstinline!k (short-circuit eval): 0!}.
    \end{itemize}
  \item
    \emph{Importanza pratica}: lo short-circuit
    \passthrough{\lstinline!\&\&!} è l'idioma fondamentale per evitare
    crash su \passthrough{\lstinline!null!}:

\begin{lstlisting}[language=Java]
if (s != null && s.length() > 0) { ... } // Se s è null, non chiama mai s.length()!
\end{lstlisting}
  \end{itemize}
\item
  \textbf{Scope delle Variabili e Divieto di Shadowing Locale}:

  \begin{itemize}
  \item
    In Java, il ciclo di vita e la visibilità delle variabili sono
    delimitati dalle graffe \passthrough{\lstinline!\{ ... \}!}.
  \item
    \passthrough{\lstinline!int n = 4; n++;!} dichiara
    \passthrough{\lstinline!n!} nello scope del metodo. Scrivere
    \passthrough{\lstinline!int n = 8;!} nello stesso scope o in un
    sottoblocco interno genera l'errore:
    \passthrough{\lstinline!Variable 'n' is already defined in the scope!}.
    In Java le variabili locali non possono oscurare altre variabili
    locali omonime.
  \item
    Blocchi anonimi indipendenti:

\begin{lstlisting}[language=Java]
{ int m = 7; m++; }
{ int m = 8; m--; }
\end{lstlisting}

    La variabile \passthrough{\lstinline!m!} del primo blocco viene
    distrutta alla chiusura della prima graffa. La seconda
    \passthrough{\lstinline!m!} è una nuova variabile in uno scope
    disgiunto. All'esterno dei blocchi, \passthrough{\lstinline!m!} non
    esiste (\passthrough{\lstinline!cannot find symbol variable m!}).
  \end{itemize}
\end{enumerate}

\begin{center}\rule{0.5\linewidth}{0.5pt}\end{center}

\section{{[}Test{]} Output Esatto di
Esecuzione}\label{chap04_test-output-esatto-di-esecuzione}

Eseguendo \passthrough{\lstinline!Example.java!} da terminale o
IntelliJ:

\begin{lstlisting}
i1: 6
f1: 0.0
f2: 3.4
c1: V c2: V c3: V
s1: null
true
k (standard eval): 1
true
k (short-circuit eval): 0
n: 5
\end{lstlisting}

\begin{center}\rule{0.5\linewidth}{0.5pt}\end{center}

\begin{studynote}[Nota di Studio] Con la \textbf{Lezione 4 / T04} abbiamo chiarito le basi di
tipo, stack/heap e operatori.

Il prossimo blocco raggruppa le \textbf{Lezioni 05-06 / Slide T05
(\emph{Java Basic Syntax}) e T06 (\emph{Java Classes and Objects})},
dove nasce il progetto incrementale centrale del corso: - Struttura
della prima classe ad oggetti:
\href{file:///home/wally/Documenti/GitHub/Programmazione-II/25-26/Teoria-e-Esercitazioni/Codice-20260902/Lezioni05-06/SimpleDate}{\passthrough{\lstinline!SimpleDate!}}
(\passthrough{\lstinline!Date.java!} e
\passthrough{\lstinline!MainDate.java!}). - La classe sperimentale
\href{file:///home/wally/Documenti/GitHub/Programmazione-II/25-26/Teoria-e-Esercitazioni/Codice-20260902/Lezioni05-06/StringPlayground.java}{\passthrough{\lstinline!StringPlayground.java!}}.
- La prima validazione delle date: algoritmo dei giorni del mese e
calcolo degli anni bisestili secondo il calendario Gregoriano. -
Sintassi Java: costrutti di controllo
(\passthrough{\lstinline!if-else!}, \passthrough{\lstinline!switch!},
\passthrough{\lstinline!for!}, \passthrough{\lstinline!while!},
\passthrough{\lstinline!do-while!}), etichette (\emph{labeled
break/continue}). - Metodi, parametri, ritorno e la keyword
\passthrough{\lstinline!this!}. - La classe
\passthrough{\lstinline!MainDate!}: test delle istanze, instanziazione
con \passthrough{\lstinline!new!} e stampa dello stato.
\end{studynote}

\textbf{Dammi conferma per aprire le Lezioni 05-06 / T05-T06 e
analizzare la prima versione di \passthrough{\lstinline!SimpleDate!} e
\passthrough{\lstinline!StringPlayground!}!}

\newpage
\chapter[T05-T06 --- Sintassi Java, Strutture di Control]{T05–T06 --- Sintassi Java, Strutture di Controllo e Classi Base}
\label{chap:t05-t06}

{\small\sffamily\color{accentblue!80!black}\textit{Costrutti Condizionali e Iterativi, Definizione di Classe e Nascita del Progetto SimpleDate}}

\vspace{0.8em}

In questo blocco analizziamo in dettaglio la sintassi fondamentale di
Java
(\href{file:///home/wally/Documenti/GitHub/Programmazione-II/25-26/Teoria-e-Esercitazioni/Slides-20260902/T05\%20-\%20Java\%20Basic\%20Syntax.pdf}{T05
- Java Basic Syntax.pdf}), la struttura delle classi e degli oggetti
(\href{file:///home/wally/Documenti/GitHub/Programmazione-II/25-26/Teoria-e-Esercitazioni/Slides-20260902/T06\%20-\%20Java\%20Classes\%20and\%20Objects.pdf}{T06
- Java Classes and Objects.pdf}) e il codice sorgente di debutto del
corso: il progetto
\href{file:///home/wally/Documenti/GitHub/Programmazione-II/25-26/Teoria-e-Esercitazioni/Codice-20260902/Lezioni05-06/SimpleDate}{\passthrough{\lstinline!SimpleDate!}}
e il file di laboratorio
\href{file:///home/wally/Documenti/GitHub/Programmazione-II/25-26/Teoria-e-Esercitazioni/Codice-20260902/Lezioni05-06/StringPlayground.java}{\passthrough{\lstinline!StringPlayground.java!}}.

\section{Teoria: Concetti Teorici dalle
Slide}\label{chap05_teoria-concetti-teorici-dalle-slide}

\subsection{Sintassi di Base e Costrutti Semplici (Slide T05,
1--19)}\label{chap05_sintassi-di-base-e-costrutti-semplici-slide-t05-119}

\begin{itemize}
\tightlist
\item
  \textbf{Variabili e Inizializzazione Automatica}:

  \begin{itemize}
  \tightlist
  \item
    Una variabile associa un nome e un tipo a una locazione di memoria.
  \item
    \textbf{Regola cruciale}: Solo i campi di classe/istanza
    (\emph{fields}) vengono inizializzati automaticamente al default
    (\passthrough{\lstinline!0!}, \passthrough{\lstinline!0.0!},
    \passthrough{\lstinline!false!},
    \passthrough{\lstinline!'\\u0000'!},
    \passthrough{\lstinline!null!}). Le variabili locali nello Stack
    richiedono l'assegnamento esplicito prima dell'uso.
  \end{itemize}
\item
  \textbf{Costanti (\passthrough{\lstinline!static final!})}:

  \begin{itemize}
  \tightlist
  \item
    \passthrough{\lstinline!final static int I = 5;!}: la variabile è a
    sola lettura; deve essere inizializzata contestualmente o nel
    costruttore e non può più essere riassegnata.
  \end{itemize}
\item
  \textbf{Identificatori e Convenzioni di Naming}:

  \begin{itemize}
  \tightlist
  \item
    Regola sintattica: sequenza arbitraria di caratteri
    \passthrough{\lstinline![a-z]!}, \passthrough{\lstinline![A-Z]!},
    \passthrough{\lstinline![0-9]!}, \passthrough{\lstinline!\_!} e
    \passthrough{\lstinline!$!}, che non inizi con una cifra e non
    coincida con una delle \textbf{50 parole chiave riservate}
    (\passthrough{\lstinline!class!}, \passthrough{\lstinline!static!},
    \passthrough{\lstinline!new!}, \passthrough{\lstinline!this!},
    \passthrough{\lstinline!super!}, \passthrough{\lstinline!switch!},
    \passthrough{\lstinline!goto!}, \passthrough{\lstinline!const!},
    ecc.).
  \item
    Convenzioni ufficiali Java:

    \begin{itemize}
    \tightlist
    \item
      \textbf{Classi}: \passthrough{\lstinline!UpperCamelCase!} (es.
      \passthrough{\lstinline!SimpleDate!},
      \passthrough{\lstinline!MainDate!}).
    \item
      \textbf{Metodi, variabili locali, campi e parametri}:
      \passthrough{\lstinline!lowerCamelCase!} (es.
      \passthrough{\lstinline!daysPerMonth!},
      \passthrough{\lstinline!myCar!}).
    \item
      \textbf{Costanti}: \passthrough{\lstinline!UPPER\_SNAKE\_CASE!}
      (es. \passthrough{\lstinline!FORMAT!},
      \passthrough{\lstinline!PI\_CONST!}).
    \end{itemize}
  \end{itemize}
\item
  \textbf{Caratteri di Escape Speciali}:

  \begin{itemize}
  \tightlist
  \item
    \passthrough{\lstinline!\\n!} (LF linefeed,
    \passthrough{\lstinline!\\u000a!}), \passthrough{\lstinline!\\r!}
    (CR carriage return, \passthrough{\lstinline!\\u000d!}),
    \passthrough{\lstinline!\\t!} (tab orizzontale,
    \passthrough{\lstinline!\\u0009!}), \passthrough{\lstinline!\\b!}
    (backspace, \passthrough{\lstinline!\\u0008!}),
    \passthrough{\lstinline!\\f!} (form feed,
    \passthrough{\lstinline!\\u000c!}), \passthrough{\lstinline!\\"!}
    (doppio apice), \passthrough{\lstinline!\\'!} (singolo apice),
    \passthrough{\lstinline!\\\\!} (backslash).
  \end{itemize}
\item
  \textbf{Commenti e Documentazione}:

  \begin{itemize}
  \tightlist
  \item
    \passthrough{\lstinline!// ...!} (singola riga).
  \item
    \passthrough{\lstinline!/* ... */!} (multi-riga).
  \item
    \passthrough{\lstinline!/** ... */!} (\textbf{Javadoc}): elaborato
    dal tool \passthrough{\lstinline!javadoc!} per generare
    automaticamente la documentazione HTML delle API.
  \end{itemize}
\item
  \textbf{Espressioni, Valutazione Sinistra-Destra e Side Effects}:

  \begin{itemize}
  \tightlist
  \item
    In Java le sotto-espressioni sono valutate \textbf{rigorosamente da
    sinistra a destra}:

    \begin{itemize}
    \tightlist
    \item
      \passthrough{\lstinline!int i = 1; int j = (i = 2) * i;!}
      \(\implies\) \passthrough{\lstinline!i!} diventa
      \passthrough{\lstinline!2!}, poi moltiplicato per
      \passthrough{\lstinline!i!} (\passthrough{\lstinline!2!})
      \(\implies\) \passthrough{\lstinline!j = 4!}.
    \item
      \passthrough{\lstinline!int t = 3; eval(t + (t = 2))!}
      \(\implies\) prima si valuta \passthrough{\lstinline!t!}
      (\passthrough{\lstinline!3!}), poi
      \passthrough{\lstinline!(t = 2)!} (\passthrough{\lstinline!2!})
      \(\implies 3 + 2 = 5\).
    \end{itemize}
  \end{itemize}
\item
  \textbf{Operatori Aritmetici Unari: Pre-incremento vs
  Post-incremento}:

  \begin{itemize}
  \item
    \passthrough{\lstinline!x++!} (post-incremento): restituisce il
    valore attuale di \passthrough{\lstinline!x!} e
    \emph{successivamente} incrementa \passthrough{\lstinline!x!} di 1.
  \item
    \passthrough{\lstinline!++x!} (pre-incremento): incrementa
    \emph{subito} \passthrough{\lstinline!x!} di 1 e restituisce il
    nuovo valore.
  \item
    \textbf{L'esercizio tranello delle slide}:

\begin{lstlisting}[language=Java]
int t = 3;
t = -t++ - ++t; // Quanto vale t?
\end{lstlisting}

    \begin{itemize}
    \tightlist
    \item
      Valutazione da sinistra a destra:

      \begin{enumerate}
      \def\labelenumi{\arabic{enumi}.}
      \tightlist
      \item
        \passthrough{\lstinline!-t++!}: valuta \(-3\), e incrementa
        \(t\) a \(4\).
      \item
        \passthrough{\lstinline!++t!}: incrementa subito \(t\) da \(4\)
        a \(5\), e valuta \(5\).
      \item
        Operazione: \((-3) - 5 = -8\). Assegnato a \(t\), il valore
        finale è \textbf{\passthrough{\lstinline!-8!}}.
      \end{enumerate}
    \end{itemize}
  \end{itemize}
\item
  \textbf{Operatore Ternario}:
  \passthrough{\lstinline!condizione ? expr1 : expr2!}.
\item
  \textbf{Operatori Logici e Corto-Circuito}:

  \begin{itemize}
  \tightlist
  \item
    \passthrough{\lstinline!\&!}, \passthrough{\lstinline!|!}: valutano
    sempre entrambi i rami.
  \item
    \passthrough{\lstinline!\&\&!}, \passthrough{\lstinline!||!}:
    valutazione a corto circuito (\emph{short-circuit}). Se il primo
    operando determina già il risultato, il secondo non viene nemmeno
    eseguito (essenziale per prevenire crash su puntatori nulli).
  \end{itemize}
\item
  \textbf{Assegnamenti Composti}: \passthrough{\lstinline!x += expr!}
  equivale a \passthrough{\lstinline!x = x + (expr)!}. Se \(x=4\),
  l'espressione \(x += x - 7\) assegna \(x = 4 + (4 - 7) = 1\).
\end{itemize}

\begin{center}\rule{0.5\linewidth}{0.5pt}\end{center}

\subsection{Costrutti di Controllo di Flusso (Slide T05,
20--29)}\label{chap05_costrutti-di-controllo-di-flusso-slide-t05-2029}

\begin{itemize}
\tightlist
\item
  \textbf{Blocchi \passthrough{\lstinline!\{ ... \}!}}: delimitano lo
  scope di vita delle variabili locali.
\item
  \textbf{Condizionali}:
  \passthrough{\lstinline!if (Bexpr) ... else ...!}.
\item
  \textbf{Costrutto \passthrough{\lstinline!switch-case!} e il
  Fall-Through}:

  \begin{itemize}
  \tightlist
  \item
    Seleziona il ramo da eseguire in base al valore di un'espressione
    (\passthrough{\lstinline!byte!}, \passthrough{\lstinline!short!},
    \passthrough{\lstinline!char!}, \passthrough{\lstinline!int!},
    \passthrough{\lstinline!String!}, \passthrough{\lstinline!enum!}).
  \item
    Se non si inserisce l'istruzione \passthrough{\lstinline!break!},
    l'esecuzione \textbf{continua nei rami successivi}
    (\emph{fall-through}). Questo meccanismo, se usato volontariamente,
    permette di raggruppare casi che condividono la stessa logica (come
    vedremo in \passthrough{\lstinline!daysPerMonth!}).
  \end{itemize}
\item
  \textbf{Cicli}:

  \begin{itemize}
  \item
    \passthrough{\lstinline!while (cond) \{ ... \}!}: test preliminare
    (può eseguire 0 volte).
  \item
    \passthrough{\lstinline!do \{ ... \} while (cond);!}: test
    posticipato (esegue almeno 1 volta).
  \item
    \passthrough{\lstinline!for (init; cond; step) \{ ... \}!}: ciclo
    controllato. Supporta inizializzazioni e passi multipli separati da
    virgola:

\begin{lstlisting}[language=Java]
for (int i = 0, j = 9; i <= 9 && j >= 0; i++, j--) { ... }
\end{lstlisting}
  \end{itemize}
\item
  \textbf{Salto di Controllo}:

  \begin{itemize}
  \tightlist
  \item
    \passthrough{\lstinline!break!}: interrompe il ciclo o lo switch ed
    esce dal blocco.
  \item
    \passthrough{\lstinline!continue!}: salta immediatamente alla
    successiva iterazione del ciclo.
  \item
    \passthrough{\lstinline!return!}: esce dal metodo corrente
    restituendo eventualmente un valore.
  \end{itemize}
\end{itemize}

\begin{center}\rule{0.5\linewidth}{0.5pt}\end{center}

\subsection{\texorpdfstring{Classi, Oggetti, \texttt{this} e
\texttt{toString} (Slide T06,
1--28)}{Classi, Oggetti, this e toString (Slide T06, 1--28)}}\label{chap05_classi-oggetti-this-e-tostring-slide-t06-128}

\begin{itemize}
\tightlist
\item
  \textbf{Perché modellare una data? (La vignetta a lezione, Slide 7)}:

  \begin{itemize}
  \tightlist
  \item
    Lo slide show mostra il fumetto geek: \emph{«Imagine an object much
    like a thing in the real world, like\ldots{} A DATE?»}. È da qui che
    nasce la classe \passthrough{\lstinline!Date!} come filo conduttore
    del corso.
  \end{itemize}
\item
  \textbf{Anatomia di una Classe}:

  \begin{itemize}
  \tightlist
  \item
    \textbf{Campi d'istanza (\emph{Fields})}: variabili che memorizzano
    lo stato di ciascun oggetto.
  \item
    \textbf{Metodi d'istanza (\emph{Methods})}: funzioni che definiscono
    il comportamento e i messaggi accettati dall'oggetto.
  \item
    \textbf{Firma del metodo}:
    \passthrough{\lstinline!tipoRitorno nomeMetodo(tipoParametro nomeParametro, ...)!}.
    Se non restituisce nulla si usa \passthrough{\lstinline!void!}.
  \item
    \textbf{Overloading}: possibilità di definire più metodi con lo
    stesso nome purché differiscano per numero, tipo o ordine dei
    parametri.
  \end{itemize}
\item
  \textbf{Entrypoint \passthrough{\lstinline!main!} e Metodi Statici
  (\passthrough{\lstinline!static!})}:

  \begin{itemize}
  \tightlist
  \item
    \passthrough{\lstinline!public static void main(String[] args)!}: il
    launcher della JVM non passa il nome del programma in
    \passthrough{\lstinline!args[0]!} (a differenza del C).
    \passthrough{\lstinline!args[0]!} è già il primo parametro utente.
  \item
    \textbf{Metodi Statici}: appartengono alla classe e non alle singole
    istanze. Si invocano con
    \passthrough{\lstinline!NomeClasse.metodo()!}. Non possono accedere
    allo stato d'istanza né usare \passthrough{\lstinline!this!}.
  \end{itemize}
\item
  \textbf{Creazione dell'Oggetto (\passthrough{\lstinline!new!}) e
  Costruttore}:

  \begin{itemize}
  \tightlist
  \item
    L'istruzione \passthrough{\lstinline!new Date(...)!}:

    \begin{enumerate}
    \def\labelenumi{\arabic{enumi}.}
    \tightlist
    \item
      Alloca memoria nello Heap per l'oggetto.
    \item
      Inizializza i campi ai valori di default.
    \item
      Invoca il \textbf{costruttore} (metodo speciale con lo stesso nome
      della classe e senza tipo di ritorno) per valorizzare lo stato.
    \item
      Restituisce il riferimento all'oggetto appena creato.
    \end{enumerate}
  \end{itemize}
\item
  \textbf{Aliasing e Garbage Collection}:

  \begin{itemize}
  \tightlist
  \item
    Se scriviamo \passthrough{\lstinline!Date d2 = d1;!}, le due
    variabili puntano allo stesso identico oggetto nello Heap
    (\emph{aliasing}).
  \item
    Quando una reference viene azzerata
    (\passthrough{\lstinline!d1 = null; d2 = null;!}), l'oggetto diventa
    orfano e il Garbage Collector può bonificarlo.
  \end{itemize}
\item
  \textbf{La Parola Chiave \passthrough{\lstinline!this!}}:

  \begin{itemize}
  \item
    È il riferimento implicito all'oggetto corrente su cui è stato
    invocato il metodo.
  \item
    È \textbf{obbligatoria} quando i parametri del costruttore/metodo
    hanno lo stesso identico nome dei campi d'istanza (\emph{shadowing
    dei parametri}):

\begin{lstlisting}[language=Java]
public Date(int day, int month, int year) {
    this.day = day;     // this.day è il campo dell'oggetto nello Heap
    this.month = month; // day è il parametro locale nello Stack
    this.year = year;
}
\end{lstlisting}
  \end{itemize}
\item
  \textbf{Rappresentazione a Stringa e
  \passthrough{\lstinline!toString()!}}:

  \begin{itemize}
  \tightlist
  \item
    Tutte le classi Java ereditano dalla superclasse universale
    \passthrough{\lstinline!java.lang.Object!} il metodo
    \passthrough{\lstinline!public String toString()!}.
  \item
    Il default di \passthrough{\lstinline!Object.toString()!} stampa
    \passthrough{\lstinline!NomeClasse@IndirizzoHashcode!} (es.
    \passthrough{\lstinline!Date@5acf9800!}).
  \item
    La JVM invoca \textbf{implicitamente}
    \passthrough{\lstinline!toString()!} ogni volta che un oggetto viene
    concatenato a una stringa
    (\passthrough{\lstinline!"Oggi: " + today!}).
  \item
    Sovrascrivendo (\emph{overriding})
    \passthrough{\lstinline!public String toString()!} nella classe,
    possiamo personalizzare la stampa dell'oggetto.
  \end{itemize}
\end{itemize}

\begin{center}\rule{0.5\linewidth}{0.5pt}\end{center}

\section{Codice: Analisi del Codice Ufficiale del
Docente}\label{chap05_codice-analisi-del-codice-ufficiale-del-docente}

\subsection{\texorpdfstring{\href{file:///home/wally/Documenti/GitHub/Programmazione-II/25-26/Teoria-e-Esercitazioni/Codice-20260902/Lezioni05-06/SimpleDate/src/Date.java}{\texttt{SimpleDate/src/Date.java}}}{SimpleDate/src/Date.java}}\label{chap05_simpledatesrcdate.java}

\begin{lstlisting}[language=Java]
public class Date {
    // 1. Stato dell'oggetto: campi di istanza (visibilità default di package per ora)
    int day;
    int month;
    int year;

    // 2. Costante di classe condivisa da tutte le istanze
    static final String FORMAT = "dd/mm/yyyy";

    // 3. Costruttore: usa 'this' per distinguere campi da parametri
    public Date(int day, int month, int year) {
        this.day = day;
        this.month = month;
        this.year = year;
        verify(); // Chiamata al metodo di validazione interna
    }

    // 4. Metodo di validazione dello stato
    void verify() {
        if (year < 0 || month < 1 || month > 12)
            System.out.println("Illegal date!");
        else
            if (day < 1 || day > daysPerMonth(month))
                System.out.println("Illegal date!");
    }

    String print() {
        return day + "/" + month + "/" + year + " [" + FORMAT + "]";
    }

    // 5. Override di toString(): delega al metodo print()
    public String toString() {
        return print();
    }
    
    // 6. Metodo statico: calcolo puro senza bisogno di istanziare Date
    static int daysPerMonth(int month) {
        int days;
        switch(month) {
            case 4:
            case 6:
            case 9:
            case 11:
                days = 30; // Fall-through controllato per i mesi da 30 giorni
                break;
            case 2:
                days = 28; // Febbraio fisso a 28 (il bisestile arriverà a Lezione 07!)
                break;
            default:
                days = 31;
                break;
        }
        return days;
    }
}
\end{lstlisting}

\subsection{{[}Approfondimento{]} Dettagli Tecnici e Scelte del
Docente:}\label{chap05_approfondimento-dettagli-tecnici-e-scelte-del-docente}

\begin{enumerate}
\def\labelenumi{\arabic{enumi}.}
\tightlist
\item
  \textbf{Perché \passthrough{\lstinline!daysPerMonth(month)!} è
  \passthrough{\lstinline!static!}?} Non ha bisogno di leggere
  \passthrough{\lstinline!this.day!} o
  \passthrough{\lstinline!this.year!}: riceve un mese intero e
  restituisce i giorni. Essendo una funzione pura di utilità, appartiene
  alla classe e può essere chiamata anche prima di creare oggetti (come
  fa il \passthrough{\lstinline!main!}).
\item
  \textbf{Lo \passthrough{\lstinline!switch!} con fall-through}: I
  \passthrough{\lstinline!case 4: case 6: case 9: case 11:!} sono privi
  di istruzione \passthrough{\lstinline!break!} intermedia; l'esecuzione
  scorre intenzionalmente fino a
  \passthrough{\lstinline!days = 30; break;!}.
\item
  \textbf{Validazione con \passthrough{\lstinline!verify()!}}: In questa
  prima versione didattica, se la data è errata viene stampato a video
  \passthrough{\lstinline'"Illegal date!"'}. Notare che l'oggetto non
  valido viene comunque creato in memoria! Solo più avanti (in Lezione
  17 con le eccezioni) il docente mostrerà come impedire l'istanziazione
  lanciando un'eccezione.
\end{enumerate}

\begin{center}\rule{0.5\linewidth}{0.5pt}\end{center}

\subsection{\texorpdfstring{\href{file:///home/wally/Documenti/GitHub/Programmazione-II/25-26/Teoria-e-Esercitazioni/Codice-20260902/Lezioni05-06/SimpleDate/src/MainDate.java}{\texttt{SimpleDate/src/MainDate.java}}}{SimpleDate/src/MainDate.java}}\label{chap05_simpledatesrcmaindate.java}

\begin{lstlisting}[language=Java]
import java.util.Scanner;

public class MainDate {
    public static void main(String[] args) {
        // Chiamata a metodo statico direttamente da classe
        System.out.println(Date.daysPerMonth(4));

        // Allocazione istanze nello Heap
        Date today = new Date(14, 10, 2025);
        Date tomorrow = new Date(15, 10, 2025);

        System.out.println(today.print());
        System.out.println(tomorrow.print());
        System.out.println(today.toString());

        String str = new String("ciao");
        System.out.println(str.toString());

        // Acquisizione input utente con java.util.Scanner
        Scanner sc = new Scanner(System.in);
        System.out.println("Enter the day: ");
        int day = sc.nextInt();
        System.out.println("Enter the month: ");
        int month = sc.nextInt();
        System.out.println("Enter the year: ");
        int year = sc.nextInt();
        sc.close(); // Chiusura dello stream di input

        Date date = new Date(day, month, year);
        System.out.println(date.toString());
    }
}
\end{lstlisting}

\begin{center}\rule{0.5\linewidth}{0.5pt}\end{center}

\subsection{\texorpdfstring{\href{file:///home/wally/Documenti/GitHub/Programmazione-II/25-26/Teoria-e-Esercitazioni/Codice-20260902/Lezioni05-06/StringPlayground.java}{\texttt{StringPlayground.java}}}{StringPlayground.java}}\label{chap05_stringplayground.java}

Questo file illustra 3 concetti cardine sulle stringhe e
l'ottimizzazione della memoria:

\begin{lstlisting}[language=Java]
public class StringPlayground {
    public static void main(String[] args) {
        String empty = new String();
        String str = "str";
        System.out.println("len(empty): " + empty.length());
        System.out.println("len(\"str\"): " + str.length());
        System.out.println(str.charAt(1)); // Restituisce 't'

        // str.charAt(1) = 'c'; // COMPILE-TIME ERROR: Le String in Java sono IMMUTABILI!

        String same = new String("str");
        System.out.println("str: " + str + " same: " + same);

        // Confronto di puntatori (identità di memoria) vs contenuto semantico
        System.out.println("str ?= same: " + (str == same));           // false! Oggetti diversi nello Heap
        String verySame = str;
        System.out.println("str ?= verySame: " + (str == verySame));   // true! Puntano alla stessa cella
        System.out.println("str ?= same: " + str.equals(same));        // true! Stessi caratteri

        // StringBuilder e Fluent Interface (Chaining)
        int n = 5;
        int sum = n * (n + 1) / 2;
        StringBuilder sb = new StringBuilder();
        sb.append("sum(").append(sum).append(")["); // Notazione a catena (fluent)
        for (int i = 0; i < n; i++)
            sb.append(" ").append(i);
        sb.append(" ]");
        System.out.println(sb.toString()); // sum(15)[ 0 1 2 3 4 ]
    }
}
\end{lstlisting}

\subsection{\texorpdfstring{{[}Approfondimento{]} Le Finezze di
\texttt{StringPlayground}:}{{[}Approfondimento{]} Le Finezze di StringPlayground:}}\label{chap05_approfondimento-le-finezze-di-stringplayground}

\begin{enumerate}
\def\labelenumi{\arabic{enumi}.}
\tightlist
\item
  \textbf{Immutabilità}: in Java l'oggetto
  \passthrough{\lstinline!String!} non può essere modificato dopo la
  creazione. Non esiste un metodo \passthrough{\lstinline!setCharAt()!}.
\item
  \textbf{\passthrough{\lstinline!==!} vs
  \passthrough{\lstinline!.equals()!}}:

  \begin{itemize}
  \tightlist
  \item
    \passthrough{\lstinline!str == same!} confronta se le due variabili
    contengono lo \textbf{stesso indirizzo di memoria}. Restituisce
    \passthrough{\lstinline!false!} perché
    \passthrough{\lstinline!same!} è stato creato con
    \passthrough{\lstinline!new String(...)!} forzando una nuova
    allocazione nello Heap.
  \item
    \passthrough{\lstinline!str.equals(same)!} confronta i caratteri uno
    a uno e restituisce \passthrough{\lstinline!true!}.
  \item
    \textbf{Regola d'esame assoluta}: per confrontare gli oggetti e le
    stringhe si usa \textbf{sempre \passthrough{\lstinline!.equals()!}},
    mai \passthrough{\lstinline!==!}!
  \end{itemize}
\item
  \textbf{\passthrough{\lstinline!StringBuilder!} e Fluent Notation}:

  \begin{itemize}
  \tightlist
  \item
    Ogni concatenazione con \passthrough{\lstinline!+!} su stringhe
    ordinarie crea un nuovo oggetto \passthrough{\lstinline!String!}
    temporaneo. Nei cicli questo genera un enorme sovraccarico di
    memoria per il Garbage Collector.
  \item
    \passthrough{\lstinline!StringBuilder!} alloca un buffer mutabile
    ridimensionabile internamente.
  \item
    Il metodo \passthrough{\lstinline!.append()!} restituisce un
    riferimento a \passthrough{\lstinline!this!} (l'istanza corrente del
    builder stesso), permettendo di concatenare le chiamate in cascata:
    \passthrough{\lstinline!sb.append(...).append(...)!} (\textbf{Fluent
    API / Method Chaining}).
  \end{itemize}
\end{enumerate}

\begin{center}\rule{0.5\linewidth}{0.5pt}\end{center}

\section{\texorpdfstring{🔄 Confronto con il tuo modulo
\texttt{es01}}{🔄 Confronto con il tuo modulo es01}}\label{chap05_confronto-con-il-tuo-modulo-es01}

Nel tuo file
\href{file:///home/wally/Documenti/GitHub/Programmazione-II/25-26/esercizi-wally-25-26/es01/src/main/java/es01/Date.java}{\passthrough{\lstinline!es01/src/main/java/es01/Date.java!}}:
1. \textbf{Switch Moderno vs Tradizionale}: Avevi già applicato la
sintassi Java 14+ con la \emph{Switch Expression}:
\passthrough{\lstinline!java    static int daysPerMonth(int month) \{        return switch (month) \{            case 4, 6, 9, 11 -> 30;            case 2 -> 28;            default -> 31;        \};    \}!}
Questa forma è sintatticamente pulitissima e idiomatica in Java moderno.
2. \textbf{Allineamento con la Lezione 05-06}: Per essere al 100\%
fedeli alla versione del docente: - Aggiungere la chiamata a
\passthrough{\lstinline!verify()!} nel costruttore per il controllo di
validità. - Correggere il refuso nella costante di formato:
\passthrough{\lstinline!"dd/mm/yyyy"!} invece di
\passthrough{\lstinline!"dd/hh/yyyy"!}.

\begin{center}\rule{0.5\linewidth}{0.5pt}\end{center}

\begin{studynote}[Nota di Studio] Con le \textbf{Lezioni 05-06} abbiamo stabilito
l'architettura iniziale di \passthrough{\lstinline!Date!}, lo scope, i
metodi statici, le stringhe e \passthrough{\lstinline!StringBuilder!}.

Il prossimo blocco è la \textbf{Lezione 07 / Slide T07: \emph{Libraries
and String Class}}: - Evoluzione di
\passthrough{\lstinline!SimpleDate!}: introduzione dell'algoritmo
completo del \textbf{calendario Gregoriano per gli anni bisestili}
(\emph{leap year}). - Analisi della classe statica ausiliaria di calcolo
matematico \passthrough{\lstinline!java.lang.Math!}. - Metodi avanzati
della classe \passthrough{\lstinline!String!}
(\passthrough{\lstinline!substring!}, \passthrough{\lstinline!indexOf!},
\passthrough{\lstinline!split!}, \passthrough{\lstinline!trim!},
\passthrough{\lstinline!replace!}). - Wrapper classes per i primitivi
(\passthrough{\lstinline!Integer!}, \passthrough{\lstinline!Double!},
\passthrough{\lstinline!Character!}), \emph{Autoboxing} e
\emph{Unboxing}.
\end{studynote}

\textbf{Dammi conferma per procedere alla Lezione 07 / T07!}

\newpage
\chapter[T07 --- Librerie Standard e Manipolazione d]{T07 --- Librerie Standard e Manipolazione delle Stringhe}
\label{chap:t07}

{\small\sffamily\color{accentblue!80!black}\textit{Java Class Library, Immutabilità della Classe String, Metodi Fondamentali e SimpleDate v2}}

\vspace{0.8em}

In questo blocco analizziamo la gestione delle librerie standard e la
classe \passthrough{\lstinline!String!}
(\href{file:///home/wally/Documenti/GitHub/Programmazione-II/25-26/Teoria-e-Esercitazioni/Slides-20260902/T07\%20-\%20Libraries\%20and\%20String\%20Class.pdf}{T07
- Libraries and String Class.pdf}), insieme all'evoluzione del codice
nella cartella
\href{file:///home/wally/Documenti/GitHub/Programmazione-II/25-26/Teoria-e-Esercitazioni/Codice-20260902/Lezione07/SimpleDate}{\passthrough{\lstinline!Lezione07!}},
dove il docente introduce \textbf{incapsulamento
(\passthrough{\lstinline!private!})}, \textbf{costruttori sovraccaricati
con delega \passthrough{\lstinline!this(...)!}}, \textbf{copy
constructor}, formattazione internazionalizzata e il primo test con
\textbf{asserzioni (\passthrough{\lstinline!assert!} e flag
\passthrough{\lstinline!-ea!})}.

\section{Teoria: Concetti Teorici dalle Slide (Slide
T07)}\label{chap06_teoria-concetti-teorici-dalle-slide-slide-t07}

\subsection{Librerie Standard e Meccanismo di Import (Slide
1--9)}\label{chap06_librerie-standard-e-meccanismo-di-import-slide-19}

\begin{itemize}
\tightlist
\item
  \textbf{Librerie Standard della JDK}:

  \begin{itemize}
  \tightlist
  \item
    \passthrough{\lstinline!java.lang!}: classi fondamentali del runtime
    (oggetti base, tipi primitivi wrapper, stringhe, thread, sistema).
    \textbf{È l'unico package importato automaticamente e implicitamente
    in ogni file sorgente Java}.
  \item
    \passthrough{\lstinline!java.util!}: strutture dati (collezioni),
    utilità temporali, parser e scanner.
  \item
    \passthrough{\lstinline!java.io!}: gestione di stream di
    input/output e filesystem.
  \item
    \passthrough{\lstinline!java.math!}: calcolo a precisione arbitraria
    (\passthrough{\lstinline!BigInteger!},
    \passthrough{\lstinline!BigDecimal!}).
  \end{itemize}
\item
  \textbf{Sintassi di Importazione}:

  \begin{itemize}
  \tightlist
  \item
    Selettiva: \passthrough{\lstinline!import java.util.Scanner;!}
    (consigliata per chiarezza e pulizia).
  \item
    Wildcard su package: \passthrough{\lstinline!import java.util.*;!}
    (rende visibili tutte le classi pubbliche del package).
  \item
    \emph{Regola del Classpath}: le classi definite dall'utente che
    risiedono nello stesso package e nella stessa cartella del Classpath
    non necessitano di alcuna direttiva
    \passthrough{\lstinline!import!}.
  \end{itemize}
\item
  \textbf{Classi di Sistema Notevoli}:

  \begin{itemize}
  \tightlist
  \item
    \passthrough{\lstinline!java.lang.System!}:

    \begin{itemize}
    \tightlist
    \item
      Tre stream predefiniti: \passthrough{\lstinline!System.in!}
      (standard input, byte stream da tastiera),
      \passthrough{\lstinline!System.out!} (standard output a video),
      \passthrough{\lstinline!System.err!} (standard error).
    \item
      \passthrough{\lstinline!System.currentTimeMillis()!}: restituisce
      un \passthrough{\lstinline!long!} con il timestamp \textbf{Unix
      Epoch} (millisecondi trascorsi dalle ore 00:00:00 UTC del 1°
      Gennaio 1970).
    \end{itemize}
  \item
    \passthrough{\lstinline!java.util.Scanner!}:

    \begin{itemize}
    \tightlist
    \item
      Parser per estrarre tipi primitivi e stringhe da flussi di testo:

      \begin{itemize}
      \tightlist
      \item
        \passthrough{\lstinline!sc.nextLine()!}: legge un'intera riga
        fino al terminatore \passthrough{\lstinline!\\n!} e restituisce
        una \passthrough{\lstinline!String!}.
      \item
        \passthrough{\lstinline!sc.nextInt()!},
        \passthrough{\lstinline!sc.nextFloat()!},
        \passthrough{\lstinline!sc.nextDouble()!}: leggono il token
        successivo effettuando il parsing del tipo. Se l'utente
        inserisce caratteri incompatibili, lancia un'eccezione
        (\passthrough{\lstinline!InputMismatchException!}).
      \end{itemize}
    \item
      Buona norma: invocare \passthrough{\lstinline!sc.close()!} per
      rilasciare il descrittore del file o dello stream ed evitare
      \emph{resource leak}.
    \end{itemize}
  \end{itemize}
\end{itemize}

\begin{center}\rule{0.5\linewidth}{0.5pt}\end{center}

\subsection{\texorpdfstring{La Classe \texttt{String} e le sue
Operazioni (Slide
10--26)}{La Classe String e le sue Operazioni (Slide 10--26)}}\label{chap06_la-classe-string-e-le-sue-operazioni-slide-1026}

\begin{itemize}
\tightlist
\item
  \textbf{Immutabilità}:

  \begin{itemize}
  \tightlist
  \item
    Le istanze di \passthrough{\lstinline!String!} sono immutabili.
    Nessun metodo della classe altera i caratteri dell'oggetto
    esistente: qualunque operazione (es.
    \passthrough{\lstinline!replace!},
    \passthrough{\lstinline!substring!}) restituisce \textbf{una nuova
    istanza di \passthrough{\lstinline!String!}} nello Heap.
  \item
    Tentare \passthrough{\lstinline!str.charAt(1) = 'c';!} produce un
    \textbf{Compile-Time Error}.
  \end{itemize}
\item
  \textbf{Operazioni Fondamentali}:

  \begin{itemize}
  \tightlist
  \item
    \passthrough{\lstinline!str.length()!}: numero di caratteri (nota: è
    un metodo con parentesi tonde, a differenza di
    \passthrough{\lstinline!arr.length!} per gli array).
  \item
    \passthrough{\lstinline!str.charAt(int index)!}: carattere in
    posizione \passthrough{\lstinline!index!} (0-indexed). Indici
    negativi o \(\ge\) \passthrough{\lstinline!length()!} lanciano
    \passthrough{\lstinline!StringIndexOutOfBoundsException!}.
  \item
    \passthrough{\lstinline!str.indexOf(char c)!}: prima occorrenza del
    carattere (restituisce \passthrough{\lstinline!-1!} se non
    presente).
  \item
    \passthrough{\lstinline!str.substring(int start, int end)!}: estrae
    la sottostringa nell'intervallo semi-aperto \([start, end)\), ovvero
    il carattere all'indice \passthrough{\lstinline!end!} è
    \textbf{escluso}. Se \passthrough{\lstinline!end!} è omesso
    (\passthrough{\lstinline!str.substring(start)!}), estrae fino al
    termine della stringa.
  \item
    \passthrough{\lstinline!str.replace(CharSequence target, CharSequence replacement)!}:
    sostituisce tutte le occorrenze da sinistra a destra.
  \item
    \passthrough{\lstinline!String.format("Format: \%s, \%d", str, num)!}:
    formattazione con sintassi analoga alla
    \passthrough{\lstinline!printf!} del C.
  \end{itemize}
\item
  \textbf{Confronto di Stringhe: \passthrough{\lstinline!==!} vs
  \passthrough{\lstinline!.equals()!}}:

  \begin{itemize}
  \item
    \passthrough{\lstinline!s == t!}: confronta l'\textbf{identità dei
    puntatori} (indirizzi nello Heap). È \passthrough{\lstinline!true!}
    solo se \passthrough{\lstinline!s!} e \passthrough{\lstinline!t!}
    puntano allo stesso identico oggetto.
  \item
    \passthrough{\lstinline!s.equals(t)!}: confronta
    l'\textbf{equivalenza semantica} dei contenuti carattere per
    carattere.
  \item
    \textbf{Lo String Constant Pool e l'ottimizzazione del compilatore}:

\begin{lstlisting}[language=Java]
String s = "hello";
String t = "hello";
System.out.println(s == t); // Stampa TRUE!
\end{lstlisting}

    \emph{Perché?} Le costanti letterali vengono inserite nel
    \emph{String Constant Pool}. Se due variabili dichiarano lo stesso
    letterale, la JVM riusa lo stesso identico oggetto in memoria per
    risparmiare RAM (\emph{interning}). Tuttavia, se scriviamo
    \passthrough{\lstinline!String r = new String("hello");!},
    costringiamo la JVM a creare un nuovo oggetto nello Heap:
    \passthrough{\lstinline!s == r!} sarà
    \passthrough{\lstinline!false!}! \textbf{Regola d'oro: usare SEMPRE
    \passthrough{\lstinline!.equals()!}}.
  \end{itemize}
\end{itemize}

\begin{center}\rule{0.5\linewidth}{0.5pt}\end{center}

\subsection{\texorpdfstring{Stringhe Mutabili con \texttt{StringBuilder}
(Slide
27--29)}{Stringhe Mutabili con StringBuilder (Slide 27--29)}}\label{chap06_stringhe-mutabili-con-stringbuilder-slide-2729}

\begin{itemize}
\tightlist
\item
  \textbf{Il problema delle prestazioni con
  \passthrough{\lstinline!String!}}:

  \begin{itemize}
  \tightlist
  \item
    L'operatore \passthrough{\lstinline!+!} tra stringhe invoca
    internamente \passthrough{\lstinline!String.valueOf()!}. Se usato in
    un ciclo di \(N\) iterazioni, alloca \(O(N)\) oggetti temporanei che
    intasano la memoria del Garbage Collector.
  \end{itemize}
\item
  \textbf{La soluzione: \passthrough{\lstinline!StringBuilder!}}:

  \begin{itemize}
  \tightlist
  \item
    Mantiene un buffer interno ridimensionabile mutabile.
  \item
    Metodi principali: \passthrough{\lstinline!.append(...)!},
    \passthrough{\lstinline!.delete(start, end)!},
    \passthrough{\lstinline!.insert(index, str)!},
    \passthrough{\lstinline!.reverse()!}.
  \item
    Supporta la \textbf{Fluent Notation (Method Chaining)} perché
    restituisce \passthrough{\lstinline!this!}.
  \end{itemize}
\end{itemize}

\begin{center}\rule{0.5\linewidth}{0.5pt}\end{center}

\section{\texorpdfstring{Codice: Evoluzione del Codice:
\texttt{SimpleDate} (Lezione
07)}{Codice: Evoluzione del Codice: SimpleDate (Lezione 07)}}\label{chap06_codice-evoluzione-del-codice-simpledate-lezione-07}

Nel passaggio da \passthrough{\lstinline!Lezioni05-06!} a
\passthrough{\lstinline!Lezione07!}, il docente trasforma radicalmente
la classe
\href{file:///home/wally/Documenti/GitHub/Programmazione-II/25-26/Teoria-e-Esercitazioni/Codice-20260902/Lezione07/SimpleDate/src/Date.java}{\passthrough{\lstinline!Date.java!}}
applicando i principi di Information Hiding e aggiunge per la prima
volta una classe di test dedicata:
\href{file:///home/wally/Documenti/GitHub/Programmazione-II/25-26/Teoria-e-Esercitazioni/Codice-20260902/Lezione07/SimpleDate/src/TestDate.java}{\passthrough{\lstinline!TestDate.java!}}.

\subsection{\texorpdfstring{{[}Approfondimento{]} Il Diff Concettuale di
\texttt{Date.java}:}{{[}Approfondimento{]} Il Diff Concettuale di Date.java:}}\label{chap06_approfondimento-il-diff-concettuale-di-date.java}

\begin{lstlisting}
 public class Date {
-    int day;
-    int month;
-    int year;
-    static final String FORMAT = "dd/mm/yyyy";
+    private int day;
+    private int month;
+    private int year;
+    static final String FORMAT_IT = "dd/mm/yyyy";
+    static final String FORMAT_US = "mm/dd/yyyy";
+    private byte lang; // 0: IT, 1: US

     // Costruttore principale
     public Date(int day, int month, int year) {
         this.day = day;
         this.month = month;
         this.year = year;
         verify();
+        lang = 0;
     }

+    // Costruttore sovraccaricato: delega tramite this(...)
+    public Date(int day, int month) {
+        this(day, month, 2025); // default year
+    }

+    // Copy Constructor: clona lo stato da un'altra istanza
+    public Date(Date other) {
+        this.day = other.day;
+        this.month = other.month;
+        this.year = other.year;
+        verify();
+        lang = other.lang;
+    }

+    // Getters
+    public int getDay() { return day; }
+    public int getMonth() { return day; } // [Attenzione]  Attenzione al refuso del prof!
+    public int getYear() { return day; }  // [Attenzione]  Ritorna 'day' anziché 'year'!
+    public byte getlang() { return lang; }

+    // Setters con validazione dell'invariante
+    public void setDay(int day) {
+        this.day = day;
+        verify();
+    }
+    public void setMonth(int month) {
+        this.month = month;
+        verify();
+    }
+    public void setYear(int year) {
+        this.year = year;
+        verify();
+    }
+    public void setLang(byte lang) {
+        if (lang == 1) this.lang = 1;
+        else this.lang = 0;
+    }
+
+    public String printFormat() {
+        if (lang == 0) return FORMAT_IT;
+        else return FORMAT_US;
+    }

     public String toString() {
-        return print();
+        if (lang == 0) return day + "/" + month + "/" + year;
+        else return month + "/" + day + "/" + year;
     }
\end{lstlisting}

\subsection{{[}Idea{]} Le Novità e le Finezze Implementative da
Notare:}\label{chap06_idea-le-novituxe0-e-le-finezze-implementative-da-notare}

\begin{enumerate}
\def\labelenumi{\arabic{enumi}.}
\item
  \textbf{Incapsulamento e Protezione dello Stato
  (\passthrough{\lstinline!private!})}:

  \begin{itemize}
  \tightlist
  \item
    I campi non sono più accessibili direttamente dall'esterno
    (\passthrough{\lstinline!d.day = -7;!} genera ora un
    \textbf{Compile-Time Error} in \passthrough{\lstinline!MainDate!}).
  \end{itemize}
\item
  \textbf{Delega tra Costruttori con
  \passthrough{\lstinline!this(...)!}}:

  \begin{itemize}
  \tightlist
  \item
    Il costruttore \passthrough{\lstinline!Date(int day, int month)!}
    invoca \passthrough{\lstinline!this(day, month, 2025);!}.
  \item
    \textbf{Regola tassativa per l'esame}: la chiamata
    \passthrough{\lstinline!this(...)!} verso un altro costruttore deve
    essere \textbf{la prima istruzione assoluta} del costruttore, pena
    errore di compilazione (\emph{«Call to `this()' must be first
    statement in constructor body»}).
  \end{itemize}
\item
  \textbf{Copy Constructor}:

  \begin{itemize}
  \tightlist
  \item
    \passthrough{\lstinline!public Date(Date other)!} permette di creare
    una nuova istanza duplicando i valori di un oggetto esistente,
    evitando l'aliasing accidentale dei puntatori.
  \end{itemize}
\item
  \textbf{Validazione Continua dell'Invariante nei Setter}:

  \begin{itemize}
  \tightlist
  \item
    Ogni volta che un setter muta lo stato
    (\passthrough{\lstinline!setDay!},
    \passthrough{\lstinline!setMonth!},
    \passthrough{\lstinline!setYear!}), invoca subito
    \passthrough{\lstinline!verify()!}. Lo stato dell'oggetto viene così
    monitorato costantemente.
  \end{itemize}
\item
  \textbf{Il Refuso di Copia-Incolla del Docente nei Getter}:

  \begin{itemize}
  \item
    Nei getter di \passthrough{\lstinline!Date.java!} (righe 38--39 del
    codice ufficiale), il prof ha inavvertitamente scritto:

\begin{lstlisting}[language=Java]
public int getMonth() { return day; }
public int getYear() { return day; }
\end{lstlisting}

    Sia \passthrough{\lstinline!getMonth()!} che
    \passthrough{\lstinline!getYear()!} restituiscono il campo
    \passthrough{\lstinline!day!} anziché
    \passthrough{\lstinline!month!} e \passthrough{\lstinline!year!}!
    Nel tuo codice assicurati di restituire i campi corretti
    (\passthrough{\lstinline!return month;!} e
    \passthrough{\lstinline!return year;!}).
  \end{itemize}
\item
  \textbf{Testing con Asserzioni e Flag \passthrough{\lstinline!-ea!}
  (\href{file:///home/wally/Documenti/GitHub/Programmazione-II/25-26/Teoria-e-Esercitazioni/Codice-20260902/Lezione07/SimpleDate/src/TestDate.java}{\passthrough{\lstinline!TestDate.java!}})}:

\begin{lstlisting}[language=Java]
public class TestDate {
    private static Date date = new Date(1, 1, 1970);

    private static void printFormatTest() {
        date.setLang((byte) 1);
        assert date.printFormat().equals(Date.FORMAT_US);
    }

    public static void main(String[] args) {
        printFormatTest();
    }
}
\end{lstlisting}

  \begin{itemize}
  \item
    \textbf{Concetto Chiave d'Esame}: La keyword
    \passthrough{\lstinline!assert!} verifica che una condizione
    booleana sia vera; se è falsa, solleva un
    \passthrough{\lstinline!AssertionError!}.
  \item
    \textbf{Attenzione}: Nella JVM le asserzioni sono
    \textbf{disabilitate di default}. Se lanci semplicemente
    \passthrough{\lstinline!java TestDate!}, l'istruzione
    \passthrough{\lstinline!assert!} viene ignorata. Per attivarle, è
    obbligatorio passare il flag \passthrough{\lstinline!-ea!}
    (\emph{Enable Assertions}):

\begin{lstlisting}[language=bash]
java -ea TestDate
\end{lstlisting}
  \item
    Nel tuo script
    \href{file:///home/wally/Documenti/GitHub/Programmazione-II/25-26/esercizi-wally-25-26/esercizi.sh}{\passthrough{\lstinline!esercizi.sh!}},
    infatti, il comando \passthrough{\lstinline!./esercizi.sh test!} usa
    proprio \passthrough{\lstinline!-ea!}!
  \end{itemize}
\end{enumerate}

\begin{center}\rule{0.5\linewidth}{0.5pt}\end{center}

\section{🔄 Corrispondenza con la Struttura del tuo
Workspace}\label{chap06_corrispondenza-con-la-struttura-del-tuo-workspace}

Nel tuo repository
\href{file:///home/wally/Documenti/GitHub/Programmazione-II/25-26/esercizi-wally-25-26}{\passthrough{\lstinline!esercizi-wally-25-26!}},
hai modularizzato perfettamente questi argomenti: -
\textbf{\passthrough{\lstinline!es02!}}:
\href{file:///home/wally/Documenti/GitHub/Programmazione-II/25-26/esercizi-wally-25-26/es02/src/main/java/es02/MainScanner.java}{\passthrough{\lstinline!MainScanner.java!}}
copre l'uso di \passthrough{\lstinline!java.util.Scanner!} introdotto in
T07. - \textbf{\passthrough{\lstinline!es04!}}:
\href{file:///home/wally/Documenti/GitHub/Programmazione-II/25-26/esercizi-wally-25-26/es04/src/main/java/es04/MainString.java}{\passthrough{\lstinline!MainString.java!}}
testa tutti i metodi di \passthrough{\lstinline!String!}
(\passthrough{\lstinline!length!}, \passthrough{\lstinline!charAt!},
\passthrough{\lstinline!indexOf!}, \passthrough{\lstinline!substring!},
\passthrough{\lstinline!replace!}, constant pool
\passthrough{\lstinline!==!} vs \passthrough{\lstinline!.equals()!}) e
le mutazioni con \passthrough{\lstinline!StringBuilder!}
(\passthrough{\lstinline!append!}, \passthrough{\lstinline!delete!},
\passthrough{\lstinline!insert!}, \passthrough{\lstinline!reverse!}). -
\textbf{\passthrough{\lstinline!es05!}}:
\href{file:///home/wally/Documenti/GitHub/Programmazione-II/25-26/esercizi-wally-25-26/es05/src/main/java/es05/Date.java}{\passthrough{\lstinline!Date.java!}}
implementa l'overloading dei costruttori, il constructor chaining con
\passthrough{\lstinline!this(...)!} e il copy constructor.

\begin{center}\rule{0.5\linewidth}{0.5pt}\end{center}

\begin{studynote}[Nota di Studio] Con la \textbf{Lezione 07} abbiamo completato le librerie
standard, le stringhe, i costruttori multipli e le asserzioni.

Il prossimo blocco è la \textbf{Lezione 08 / Slide T08: \emph{Class
Constructors and Encapsulation}}: - Studio approfondito della teoria su
costruttori (default constructor, costruttore falso, copy constructor,
overloading). - Analisi dettagliata dell'Information Hiding e dei
\textbf{modificatori di accesso} (\passthrough{\lstinline!public!},
\passthrough{\lstinline!protected!}, package-private di default,
\passthrough{\lstinline!private!}). - Evoluzione del codice in
\passthrough{\lstinline!Lezione08!}: sostituzione del
\passthrough{\lstinline!byte lang!} con un tipo enumerato ad hoc:
\textbf{\passthrough{\lstinline!Language.java!}}
(\passthrough{\lstinline!enum Language \{ IT, US \}!})! - Array costanti
per i nomi dei mesi (\passthrough{\lstinline!MONTHS\_IT!},
\passthrough{\lstinline!MONTHS\_US!}) e nuovo metodo
\passthrough{\lstinline!getMonthAsString()!}.
\end{studynote}

\textbf{Dammi conferma quando vuoi procedere alla Lezione 08 / T08!}

\newpage
\chapter[T08 --- Costruttori, Inizializzazione ed In]{T08 --- Costruttori, Inizializzazione ed Incapsulamento}
\label{chap:t08}

{\small\sffamily\color{accentblue!80!black}\textit{Ciclo di Inizializzazione, Overloading dei Costruttori, Modificatori e SimpleDate v3}}

\vspace{0.8em}

In questo blocco analizziamo in dettaglio la teoria dei costruttori,
dell'incapsulamento e dell'information hiding
(\href{file:///home/wally/Documenti/GitHub/Programmazione-II/25-26/Teoria-e-Esercitazioni/Slides-20260902/T08\%20-\%20Class\%20Constructors\%20and\%20Encapsulation.pdf}{T08
- Class Constructors and Encapsulation.pdf}), confrontandola con il
codice ufficiale in
\href{file:///home/wally/Documenti/GitHub/Programmazione-II/25-26/Teoria-e-Esercitazioni/Codice-20260902/Lezione08/SimpleDate/src}{\passthrough{\lstinline!Lezione08!}},
dove nasce l'enum \textbf{\passthrough{\lstinline!Language.java!}},
vengono aggiunti gli array dei mesi per la stampa estesa
(\passthrough{\lstinline!prettyPrint()!}) e si introduce
l'\textbf{Enhanced For (for-each)}.

\section{Teoria: Concetti Teorici dalle Slide (Slide
T08)}\label{chap07_teoria-concetti-teorici-dalle-slide-slide-t08}

\subsection{Il Ciclo di Vita e i Meccanismi dei Costruttori (Slide
1--12)}\label{chap07_il-ciclo-di-vita-e-i-meccanismi-dei-costruttori-slide-112}

\begin{itemize}
\tightlist
\item
  \textbf{Cosa avviene esattamente dietro le quinte con
  \passthrough{\lstinline!new Clazz()!}}:

  \begin{enumerate}
  \def\labelenumi{\arabic{enumi}.}
  \tightlist
  \item
    La JVM calcola l'ingombro in byte di tutti i campi e alloca un
    blocco contiguo di memoria nello \textbf{Heap}.
  \item
    Genera un riferimento (\emph{reference}) univoco che punta a
    quell'area di memoria.
  \item
    Inizializza tutti i campi ai rispettivi valori di default
    (\passthrough{\lstinline!0!}, \passthrough{\lstinline!false!},
    \passthrough{\lstinline!null!}).
  \item
    Invoca il \textbf{costruttore} designato passando implicitamente il
    puntatore \passthrough{\lstinline!this!}.
  \item
    Restituisce l'indirizzo dell'oggetto alla variabile a sinistra
    dell'assegnamento.
  \end{enumerate}
\item
  \textbf{Caratteristiche Tassative di un Costruttore}:

  \begin{itemize}
  \tightlist
  \item
    Deve avere lo \textbf{stesso identico nome della classe}.
  \item
    \textbf{Non ha alcun tipo di ritorno} (nemmeno
    \passthrough{\lstinline!void!}!).
  \item
    Viene eseguito una sola volta, all'atto della creazione
    dell'istanza.
  \end{itemize}
\item
  \textbf{Default Constructor (Costruttore di Default)}:

  \begin{itemize}
  \tightlist
  \item
    Se il programmatore \textbf{non dichiara alcun costruttore}, il
    compilatore Java inserisce automaticamente un costruttore vuoto
    senza argomenti (\passthrough{\lstinline!public Clazz() \{ \}!}).
  \item
    {[}Attenzione{]} \textbf{Regola aurea d'esame}: se il programmatore
    dichiara \emph{anche un solo} costruttore con parametri (es.
    \passthrough{\lstinline!public Car(String color)!}), il compilatore
    \textbf{non genera più il costruttore di default}! Invocare
    \passthrough{\lstinline!new Car()!} senza argomenti provocherà un
    \textbf{Compile-Time Error}.
  \end{itemize}
\item
  \textbf{False Constructor (Il Falso Costruttore --- Trabocchetto
  d'Esame)}:

  \begin{itemize}
  \item
    Se si inserisce un tipo di ritorno (ad esempio
    \passthrough{\lstinline!void Car(String color)!}):

\begin{lstlisting}[language=Java]
public class Car {
    void Car(String color) { ... } // [Attenzione]  NON È UN COSTRUTTORE!
}
\end{lstlisting}
  \item
    Per il compilatore questo è un \textbf{normale metodo d'istanza}
    (avente incidentalmente lo stesso nome della classe). Il compilatore
    aggiungerà comunque il costruttore di default
    \passthrough{\lstinline!Car()!}, e il metodo
    \passthrough{\lstinline!void Car(...)!} non verrà mai invocato al
    momento del \passthrough{\lstinline!new!}!
  \end{itemize}
\item
  \textbf{Copy Constructor}:

  \begin{itemize}
  \tightlist
  \item
    Costruttore che accetta un'altra istanza della medesima classe per
    clonarne lo stato (\passthrough{\lstinline!Car(Car other)!}). In
    Java non esiste la copia automatica bit-a-bit del C++: deve essere
    scritta esplicitamente dal programmatore.
  \end{itemize}
\item
  \textbf{Constructor Overloading e Chaining con
  \passthrough{\lstinline!this(...)!}}:

  \begin{itemize}
  \tightlist
  \item
    Più costruttori con firme distinte.
  \item
    La chiamata \passthrough{\lstinline!this(...)!} verso un costruttore
    fratello serve a evitare duplicazione di logica di validazione e
    \textbf{deve essere tassativamente la prima riga di codice nel corpo
    del costruttore}.
  \item
    \emph{Meme delle slide (p.~11)}: Il matematico in lacrime che grida
    \emph{``Abuse of notation''} vs il Chad Programmer che usa
    l'overloading per semplificare il codice.
  \end{itemize}
\item
  \textbf{Distruttori e il Metodo \passthrough{\lstinline!finalize()!}}:

  \begin{itemize}
  \tightlist
  \item
    In Java non esiste la deallocazione manuale: non esistono
    distruttori (\passthrough{\lstinline!\~Car()!}).
  \item
    Esisteva il metodo protetto
    \passthrough{\lstinline!protected void finalize()!} invocato prima
    della rimozione da parte del Garbage Collector.
  \item
    {[}Attenzione{]} \textbf{Attenzione}:
    \passthrough{\lstinline!finalize()!} è ufficialmente
    \textbf{deprecato da Java 9} in poi (e rimosso/reso no-op nelle
    versioni moderne) per l'imprevedibilità temporale dell'esecuzione
    del GC e rischi di deadlock. Non va mai utilizzato.
  \end{itemize}
\end{itemize}

\begin{center}\rule{0.5\linewidth}{0.5pt}\end{center}

\subsection{Incapsulamento, Information Hiding e Modificatori di Accesso
(Slide
13--24)}\label{chap07_incapsulamento-information-hiding-e-modificatori-di-accesso-slide-1324}

\begin{itemize}
\tightlist
\item
  \textbf{Incapsulamento}: aggregazione fisica nello stesso file
  sorgente di dati (campi) e procedure (metodi).
\item
  \textbf{Information Hiding}: occultamento della rappresentazione
  interna dello stato.

  \begin{itemize}
  \tightlist
  \item
    \emph{Meme della Cipolla (p.~18)}: \emph{«Most software is written
    like an onion: The more layers you peel back, the more you want to
    cry»}. L'Information Hiding crea confini netti per proteggere il
    codice dal disfacimento.
  \item
    \emph{Meme di Gandalf (p.~24)}: \emph{«YOU SHALL NOT ACCESS MY
    MEMBERS»}.
  \end{itemize}
\item
  \textbf{I 4 Livelli di Visibilità in Java}:

  \begin{enumerate}
  \def\labelenumi{\arabic{enumi}.}
  \tightlist
  \item
    \passthrough{\lstinline!private!}: visibile esclusivamente
    all'interno della stessa classe.
  \item
    \emph{Default (Package-Private)}: nessun modificatore; visibile
    all'interno della classe e a tutte le classi residenti nello stesso
    package.
  \item
    \passthrough{\lstinline!protected!}: visibile nel package e a tutte
    le sottoclassi derivate (anche in package diversi).
  \item
    \passthrough{\lstinline!public!}: visibile ovunque, da qualsiasi
    package del Classpath.
  \end{enumerate}
\item
  \textbf{Vantaggi Architetturali di Getter e Setter}:

  \begin{itemize}
  \tightlist
  \item
    \emph{Nei Setter}: possibilità di intercettare valori anomali,
    applicare regole di business e loggare le modifiche prima di mutare
    lo stato.
  \item
    \emph{Nei Getter}: possibilità di alterare la rappresentazione
    interna della memoria senza rompere il codice dei client
    (\textbf{Retro-compatibilità / Loose Coupling}). Ad esempio,
    cambiare il campo interno da \passthrough{\lstinline!int age!} a
    \passthrough{\lstinline!short age!} per dimezzare l'occupazione di
    RAM, mantenendo il getter
    \passthrough{\lstinline!public int getAge() \{ return age; \}!}
    (promozione implicita di tipo) senza che il mondo esterno debba
    ricompilare o cambiare una sola riga di codice.
  \end{itemize}
\end{itemize}

\begin{center}\rule{0.5\linewidth}{0.5pt}\end{center}

\subsection{Asserzioni e Testing (Slide
25--26)}\label{chap07_asserzioni-e-testing-slide-2526}

\begin{itemize}
\item
  La parola chiave \passthrough{\lstinline!assert!}:

\begin{lstlisting}[language=Java]
assert person.getAge() != 0 : "Age cannot be 0";
\end{lstlisting}

  \begin{itemize}
  \item
    Se l'espressione a sinistra valuta \passthrough{\lstinline!false!},
    la JVM interrompe l'esecuzione e lancia un
    \passthrough{\lstinline!AssertionError!} mostrando il messaggio a
    destra dei due punti.
  \item
    Richiede obbligatoriamente l'attivazione a runtime con il flag JVM
    \textbf{\passthrough{\lstinline!-ea!}} (\emph{Enable Assertions}):

\begin{lstlisting}[language=bash]
java -ea TestPerson
\end{lstlisting}
  \end{itemize}
\end{itemize}

\begin{center}\rule{0.5\linewidth}{0.5pt}\end{center}

\section{\texorpdfstring{Codice: Evoluzione del Codice:
\href{file:///home/wally/Documenti/GitHub/Programmazione-II/25-26/Teoria-e-Esercitazioni/Codice-20260902/Lezione08/SimpleDate/src}{\texttt{Lezione08}}}{Codice: Evoluzione del Codice: Lezione08}}\label{chap07_codice-evoluzione-del-codice-lezione08}

In \passthrough{\lstinline!Lezione08!} il docente applica un refactoring
strutturale a \passthrough{\lstinline!SimpleDate!}:

\subsection{\texorpdfstring{Introduzione del Tipo Enumerato
\href{file:///home/wally/Documenti/GitHub/Programmazione-II/25-26/Teoria-e-Esercitazioni/Codice-20260902/Lezione08/SimpleDate/src/Language.java}{\texttt{Language.java}}}{Introduzione del Tipo Enumerato Language.java}}\label{chap07_introduzione-del-tipo-enumerato-language.java}

Viene abbandonato il primitivo \passthrough{\lstinline!byte lang!} (dove
\passthrough{\lstinline!0!} era IT e \passthrough{\lstinline!1!} era US,
fragile e poco leggibile) a favore di un tipo enum fortemente tipizzato:

\begin{lstlisting}[language=Java]
public enum Language {
    IT, US;

    public String toString() {
        switch (this) {
            case IT: return "italiano";
            case US: return "american";
            default: return null;
        }
    }
}
\end{lstlisting}

\emph{Finezza}: Anche gli enum in Java sono tipi reference completi e
possono sovrascrivere \passthrough{\lstinline!toString()!} per
restituire una descrizione testuale appropriata.

\begin{center}\rule{0.5\linewidth}{0.5pt}\end{center}

\subsection{\texorpdfstring{Refactoring di
\href{file:///home/wally/Documenti/GitHub/Programmazione-II/25-26/Teoria-e-Esercitazioni/Codice-20260902/Lezione08/SimpleDate/src/Date.java}{\texttt{Date.java}}}{Refactoring di Date.java}}\label{chap07_refactoring-di-date.java}

\begin{lstlisting}
 public class Date {
-    private byte lang; // 0: IT, 1: US
+    private Language lang;
+    private static final String[] MONTHS_IT = { "gennaio", "febbraio", "marzo", "aprile", "maggio", "giugno", "luglio", "agosto", "setembre", "ottobre", "novembre", "dicembre" };
+    private static final String[] MONTHS_US = { "January", "February", "March", "April", "May", "June", "July", "August", "September", "October", "November", "December"};

     public Date(int day, int month, int year) {
         this.day = day;
         this.month = month;
         this.year = year;
         verify();
-        lang = 0;
+        lang = Language.IT;
     }

-    public void setLang(byte lang) { ... }
+    public void setAmerican() { lang = Language.US; }
+    public void setItalian() { lang = Language.IT; }

+    public String getMonthAsString() {
+        if (lang == Language.IT) return MONTHS_IT[month-1];
+        else return MONTHS_US[month-1];
+    }

+    public String prettyPrint() {
+        if (lang == Language.IT) return day + " " + MONTHS_IT[month-1] + " " + year;
+        else return MONTHS_US[month-1] + " " + day + ", " + year;
+    }
\end{lstlisting}

\subsection{\texorpdfstring{{[}Approfondimento{]} Dettagli Tecnici e
Trabocchetti nel Codice di
\texttt{Date}:}{{[}Approfondimento{]} Dettagli Tecnici e Trabocchetti nel Codice di Date:}}\label{chap07_approfondimento-dettagli-tecnici-e-trabocchetti-nel-codice-di-date}

\begin{enumerate}
\def\labelenumi{\arabic{enumi}.}
\item
  \textbf{Separazione Semantica dei Metodi di Configurazione}: Invece di
  un setter generico che accetta valori arbitrari, il prof definisce due
  metodi espliciti auto-esplicativi:
  \passthrough{\lstinline!setItalian()!} e
  \passthrough{\lstinline!setAmerican()!}.
\item
  \textbf{Array di Costanti di Classe
  (\passthrough{\lstinline!private static final String[]!})}: Gli
  elenchi dei nomi dei mesi sono \passthrough{\lstinline!static final!}
  (condivisi da tutte le istanze senza duplicare memoria nello Heap) e
  \passthrough{\lstinline!private!} (inaccessibili all'esterno).
  \emph{Curiosità nel codice del prof}: nell'array
  \passthrough{\lstinline!MONTHS\_IT!} c'è un piccolo refuso:
  \passthrough{\lstinline!"setembre"!} con una sola `t'.
\item
  \textbf{Persistenza del Refuso del Prof nei Getter}: Anche in questa
  lezione, alle righe 43--44 di \passthrough{\lstinline!Date.java!}:

\begin{lstlisting}[language=Java]
public int getDay() { return day; }
public int getMonth() { return day; } // Ritorna day anziché month!
public int getYear() { return day; }  // Ritorna day anziché year!
\end{lstlisting}

  Un chiaro refuso di copia-incolla che non è stato corretto dal docente
  nel rilascio ufficiale. Nel proprio codice personale è doveroso
  restituire i rispettivi campi \passthrough{\lstinline!month!} e
  \passthrough{\lstinline!year!}.
\item
  \textbf{\passthrough{\lstinline!prettyPrint()!} e Formattazione
  Culturale}:

  \begin{itemize}
  \tightlist
  \item
    In italiano: giorno, mese per esteso, anno (es.
    \passthrough{\lstinline!7 gennaio 2025!}).
  \item
    In inglese/americano: mese per esteso, giorno, virgola, anno (es.
    \passthrough{\lstinline!February 20, 2025!}).
  \end{itemize}
\end{enumerate}

\begin{center}\rule{0.5\linewidth}{0.5pt}\end{center}

\subsection{\texorpdfstring{Novità in
\href{file:///home/wally/Documenti/GitHub/Programmazione-II/25-26/Teoria-e-Esercitazioni/Codice-20260902/Lezione08/SimpleDate/src/MainDate.java}{\texttt{MainDate.java}}:
Array di Oggetti ed Enhanced
For}{Novità in MainDate.java: Array di Oggetti ed Enhanced For}}\label{chap07_novituxe0-in-maindate.java-array-di-oggetti-ed-enhanced-for}

\begin{lstlisting}[language=Java]
// Inizializzazione rapida di un array di oggetti nello Heap
Date[] dates = { d1, d2, d3, new Date(14, 2, 2024) };

// 1. Scansione all'indietro classica tramite indice
for (int i = dates.length - 1; i >= 0; i--)
    System.out.println(dates[i].toString() + ": " + dates[i].getMonthAsString());

// 2. Scansione in avanti con Enhanced For (For-Each)
for (Date date : dates)
    System.out.println(date.toString() + ": " + date.getMonthAsString());
\end{lstlisting}

\begin{itemize}
\tightlist
\item
  Il costrutto
  \textbf{\passthrough{\lstinline!for (Tipo elemento : collezione)!}}
  evita di manipolare manualmente indici di scorrimento, rendendo il
  codice più leggibile ed eliminando il rischio di errori di off-by-one
  o out-of-bounds.
\end{itemize}

\begin{center}\rule{0.5\linewidth}{0.5pt}\end{center}

\section{🔄 Corrispondenza con i Tuoi Moduli
Workspace}\label{chap07_corrispondenza-con-i-tuoi-moduli-workspace}

Nel tuo repository
\href{file:///home/wally/Documenti/GitHub/Programmazione-II/25-26/esercizi-wally-25-26}{\passthrough{\lstinline!esercizi-wally-25-26!}}:
- \textbf{\passthrough{\lstinline!es06!}}:
\href{file:///home/wally/Documenti/GitHub/Programmazione-II/25-26/esercizi-wally-25-26/es06/src/main/java/es06/Person.java}{\passthrough{\lstinline!Person.java!}}
implementa l'esempio accademico della persona con i metodi di
validazione \passthrough{\lstinline!verifyAge!} e
\passthrough{\lstinline!verifyName!}, incapsulamento rigoroso e gestione
dei setter/getter. - \textbf{\passthrough{\lstinline!es07!}}:
\href{file:///home/wally/Documenti/GitHub/Programmazione-II/25-26/esercizi-wally-25-26/es07/src/main/java/es07/Language.java}{\passthrough{\lstinline!Language.java!}}
definisce l'enum \passthrough{\lstinline!Language!} sfruttando
modernamente la \emph{Switch Expression} di Java 14+
(\passthrough{\lstinline!case IT -> "italiano";!}).

\begin{center}\rule{0.5\linewidth}{0.5pt}\end{center}

\begin{studynote}[Nota di Studio] Con la \textbf{Lezione 08} abbiamo completato l'intero ciclo
di vita dei costruttori, l'incapsulamento dei dati, gli enum base e la
gestione dei formati estesi.

Il prossimo blocco è la \textbf{Lezione 09 / Slide T09: \emph{Arrays and
Enumerative Types}}: - Approfondimento teorico sistematico sugli
\textbf{Array in Java} (dichiarazione, allocazione nello Heap, proprietà
immutabile \passthrough{\lstinline!.length!}, array multidimensionali,
copie superficiali vs profonde). - Analisi rigorosa degli
\textbf{Enumerative Types (\passthrough{\lstinline!enum!})}: costruttori
di enum, campi interni, metodi \passthrough{\lstinline!values()!} e
\passthrough{\lstinline!ordinal()!}. - Evoluzione del codice in
\passthrough{\lstinline!Lezione09!}: - Potenziamento di
\passthrough{\lstinline!Language.java!} con campi interni
(\passthrough{\lstinline!format!},
\passthrough{\lstinline!description!}) e costruttore privato. -
Riscrittura ed espansione di \passthrough{\lstinline!Date.java!}.
\end{studynote}

\textbf{Dammi conferma per aprire la Lezione 09 / T09 e continuare!}

\newpage
\chapter[T09 --- Array, Matrici e Tipi Enumerativi (]{T09 --- Array, Matrici e Tipi Enumerativi (enum)}
\label{chap:t09}

{\small\sffamily\color{accentblue!80!black}\textit{Array Mono e Multidimensionali, Enum Tipizzati, Immutabilità e Singleton in SimpleDate v4}}

\vspace{0.8em}

In questo blocco approfondiamo la gestione a basso livello degli array
nello Heap, i tipi enumerati (\passthrough{\lstinline!enum!})
(\href{file:///home/wally/Documenti/GitHub/Programmazione-II/25-26/Teoria-e-Esercitazioni/Slides-20260902/T09\%20-\%20Arrays\%20and\%20Enumerative\%20Types.pdf}{T09
- Arrays and Enumerative Types.pdf}) e analizziamo il codice di
\href{file:///home/wally/Documenti/GitHub/Programmazione-II/25-26/Teoria-e-Esercitazioni/Codice-20260902/Lezione09/SimpleDate/src}{\passthrough{\lstinline!Lezione09!}},
dove \passthrough{\lstinline!Date!} diventa una \textbf{classe
immutabile} (campi \passthrough{\lstinline!final!} e rimozione dei
setter) e compare per la prima volta un design pattern fondamentale: il
\textbf{Singleton Pattern} implementato in
\href{file:///home/wally/Documenti/GitHub/Programmazione-II/25-26/Teoria-e-Esercitazioni/Codice-20260902/Lezione09/SimpleDate/src/BirthDay.java}{\passthrough{\lstinline!BirthDay.java!}}.

\section{Teoria: Concetti Teorici dalle Slide (Slide
T09)}\label{chap08_teoria-concetti-teorici-dalle-slide-slide-t09}

\subsection{Gli Array in Java (Slide
1--13)}\label{chap08_gli-array-in-java-slide-113}

\begin{itemize}
\item
  \textbf{Natura degli Array}:

  \begin{itemize}
  \tightlist
  \item
    Un array è una sequenza contigua e ordinata di variabili omogenee
    (dello stesso tipo), accessibili tramite indice numerico intero che
    parte da \passthrough{\lstinline!0!}.
  \item
    In Java gli array \textbf{non sono tipi primitivi}: sono
    \textbf{oggetti speciali allocati nello Heap}.
  \item
    Possono memorizzare valori primitivi oppure
    \textbf{riferimenti/puntatori a oggetti}, mai oggetti per valore
    diretto.
  \item
    La dimensione di un array può essere stabilita dinamicamente a
    runtime all'atto dell'allocazione
    (\passthrough{\lstinline!new int[size]!}), ma \textbf{è immutabile
    una volta creata}.
  \end{itemize}
\item
  \textbf{Dichiarazione vs Creazione}:

  \begin{itemize}
  \tightlist
  \item
    Sintassi raccomandata: \passthrough{\lstinline!int[] arr;!} (la
    notazione \passthrough{\lstinline!int arr[];!} è ammessa per
    retrocompatibilità col C, ma sconsigliata).
  \item
    La dichiarazione \passthrough{\lstinline!int[] arr;!} alloca solo
    una variabile reference nello Stack (inizializzata a
    \passthrough{\lstinline!null!}), \textbf{senza riservare spazio per
    gli elementi nello Heap}.
  \item
    Allocazione:

    \begin{itemize}
    \tightlist
    \item
      Tramite \passthrough{\lstinline!new!}:
      \passthrough{\lstinline!float[] arr = new float[10];!} (gli
      elementi assumono il valore di default del tipo, es.
      \passthrough{\lstinline!0.0f!}).
    \item
      Tramite inizializzazione statica:
      \passthrough{\lstinline!int[] primes = \{2, 3, 5, 7\};!} (la
      dimensione viene inferita automaticamente dal compilatore).
    \end{itemize}
  \end{itemize}
\item
  \textbf{Il Pseudo-Campo \passthrough{\lstinline!.length!}}:

  \begin{itemize}
  \tightlist
  \item
    Ogni array possiede la proprietà
    \passthrough{\lstinline!public final int length!}.
  \item
    {[}Attenzione{]} \textbf{Distinzione d'esame}: per gli array si
    scrive \passthrough{\lstinline!arr.length!} (\textbf{senza
    parentesi}, è una variabile \passthrough{\lstinline!final!}), mentre
    per le stringhe si invoca il metodo
    \passthrough{\lstinline!str.length()!} (\textbf{con le parentesi}).
  \item
    Se si tenta di leggere \passthrough{\lstinline!arr.length!} su un
    array nullo (\passthrough{\lstinline!arr = null;!}), la JVM lancia
    una \passthrough{\lstinline!NullPointerException!}.
  \end{itemize}
\item
  \textbf{Confronto e Stampa degli Array (I Trabocchetti Classici)}:

  \begin{enumerate}
  \def\labelenumi{\arabic{enumi}.}
  \tightlist
  \item
    \passthrough{\lstinline!arr1 == arr2!}: confronta solo gli indirizzi
    nello Heap. Se due array contengono gli stessi identici elementi ma
    risiedono in aree diverse, restituisce
    \passthrough{\lstinline!false!}.
  \item
    \passthrough{\lstinline!arr1.equals(arr2)!}: \textbf{NON funziona}.
    Gli array non sovrascrivono il metodo
    \passthrough{\lstinline!.equals()!} di
    \passthrough{\lstinline!Object!}, quindi
    \passthrough{\lstinline!equals!} si comporta esattamente come
    \passthrough{\lstinline!==!}.
  \item
    \textbf{Soluzione standard}: per confrontare il contenuto si usa il
    metodo statico
    \textbf{\passthrough{\lstinline!java.util.Arrays.equals(arr1, arr2)!}}.
  \item
    \passthrough{\lstinline!arr.toString()!}: stampa la firma bytecode
    della reference (es. \passthrough{\lstinline![I@6ce253f1!}, dove
    \passthrough{\lstinline![!} indica un array e
    \passthrough{\lstinline!I!} indica il tipo
    \passthrough{\lstinline!int!}).
  \item
    Per visualizzare gli elementi a video in formato leggibile:
    \textbf{\passthrough{\lstinline!java.util.Arrays.toString(arr)!}}.
  \end{enumerate}
\item
  \textbf{Enhanced For Loop (For-Each, Java 5+)}:

\begin{lstlisting}[language=Java]
for (String arg : args) {
    System.out.println(arg);
}
\end{lstlisting}

  Sostituisce il for contatore eliminando il rischio di errori di
  \emph{off-by-one} e \emph{ArrayIndexOutOfBoundsException}.
\item
  \textbf{Array Multidimensionali e Matrici Frastagliate (\emph{Jagged
  Arrays})}:

  \begin{itemize}
  \tightlist
  \item
    In Java non esistono matrici bidimensionali a blocco contiguo in
    stile Fortran/C: \textbf{un array multidimensionale è un array di
    riferimenti ad altri array}.
  \item
    \emph{Conseguenze pratiche}:

    \begin{enumerate}
    \def\labelenumi{\arabic{enumi}.}
    \item
      \textbf{Scambio di righe in tempo \(O(1)\)}: per invertire la
      prima e l'ultima riga di una matrice non serve copiare gli
      elementi uno per uno, basta scambiare i loro puntatori:

\begin{lstlisting}[language=Java]
int[] temp = matrix[0];
matrix[0] = matrix[matrix.length - 1];
matrix[matrix.length - 1] = temp;
\end{lstlisting}
    \item
      \textbf{Righe a lunghezza eterogenea (\emph{Jagged/Pyramid
      Arrays})}:

\begin{lstlisting}[language=Java]
int[][] pyramid = new int[3][];
pyramid[0] = new int[1]; // Riga 0 ha 1 elemento
pyramid[1] = new int[2]; // Riga 1 ha 2 elementi
pyramid[2] = new int[3]; // Riga 2 ha 3 elementi
\end{lstlisting}
    \item
      Per stampare matrici annidate,
      \passthrough{\lstinline!Arrays.toString(matrix)!} stampa gli
      indirizzi delle singole righe. Si deve iterare sulle righe o usare
      \textbf{\passthrough{\lstinline!java.util.Arrays.deepToString(matrix)!}}.
    \end{enumerate}
  \end{itemize}
\end{itemize}

\begin{center}\rule{0.5\linewidth}{0.5pt}\end{center}

\subsection{\texorpdfstring{Tipi Enumerati (\texttt{enum}) (Slide
14--19)}{Tipi Enumerati (enum) (Slide 14--19)}}\label{chap08_tipi-enumerati-enum-slide-1419}

\begin{itemize}
\tightlist
\item
  Introdotti in Java 5 con la keyword \passthrough{\lstinline!enum!} per
  gestire insiemi chiusi e prefissati di valori costanti (es. giorni
  della settimana, punti cardinali, lingue).
\item
  \textbf{Regole e Caratteristiche}:

  \begin{itemize}
  \tightlist
  \item
    Possono essere dichiarati in un file \passthrough{\lstinline!.java!}
    dedicato (come classe autonoma) o all'interno di una classe come
    membro statico, ma \textbf{mai all'interno di un metodo}.
  \item
    Non sono semplici interi o stringhe: sono classi speciali che
    estendono implicitamente \passthrough{\lstinline!java.lang.Enum!}.
    Il loro costruttore è privato e non può essere invocato con
    \passthrough{\lstinline!new!}.
  \item
    \textbf{Metodi Predefiniti Fondamentali}:

    \begin{itemize}
    \tightlist
    \item
      \passthrough{\lstinline!e.name()!}: restituisce il nome letterale
      della costante come \passthrough{\lstinline!String!} (es.
      \passthrough{\lstinline!"WEST"!}).
    \item
      \passthrough{\lstinline!e.ordinal()!}: restituisce la posizione
      ordinale intera partendo da 0 (es. \passthrough{\lstinline!3!}).
    \item
      \passthrough{\lstinline!EnumClass.values()!}: restituisce un array
      contenente tutte le costanti dichiarate nell'enum, comodo per
      cicli for-each.
    \end{itemize}
  \end{itemize}
\item
  \textbf{Confronto tra Enum: Perché usare \passthrough{\lstinline!==!}
  invece di \passthrough{\lstinline!.equals()!}}:

  \begin{itemize}
  \tightlist
  \item
    La JVM garantisce che in memoria esista \textbf{una e una sola
    istanza} per ciascuna costante di un enum (\emph{garanzia di unicità
    singleton}).
  \item
    Pertanto, il confronto con \textbf{\passthrough{\lstinline!==!}} è
    perfettamente sicuro ed è \textbf{più robusto di
    \passthrough{\lstinline!.equals()!}}: se una variabile vale
    \passthrough{\lstinline!null!}, l'espressione
    \passthrough{\lstinline!move == Direction.WEST!} restituisce
    \passthrough{\lstinline!false!} senza errori, mentre
    \passthrough{\lstinline!move.equals(...)!} causerebbe un crash con
    \passthrough{\lstinline!NullPointerException!}!
  \end{itemize}
\item
  \textbf{Uso di Enum negli \passthrough{\lstinline!switch!}}:

  \begin{itemize}
  \tightlist
  \item
    All'interno dei blocchi \passthrough{\lstinline!case!}, \textbf{non
    si qualifica} il tipo: si scrive \passthrough{\lstinline!case IT:!}
    e non \passthrough{\lstinline!case Language.IT:!}.
  \end{itemize}
\end{itemize}

\begin{center}\rule{0.5\linewidth}{0.5pt}\end{center}

\section{\texorpdfstring{Codice: Evoluzione del Codice:
\texttt{SimpleDate} (Lezione
09)}{Codice: Evoluzione del Codice: SimpleDate (Lezione 09)}}\label{chap08_codice-evoluzione-del-codice-simpledate-lezione-09}

In
\href{file:///home/wally/Documenti/GitHub/Programmazione-II/25-26/Teoria-e-Esercitazioni/Codice-20260902/Lezione09/SimpleDate/src}{\passthrough{\lstinline!Lezione09!}},
il docente introduce due cambiamenti strutturali fondamentali:

\subsection{\texorpdfstring{Immutabilità di
\href{file:///home/wally/Documenti/GitHub/Programmazione-II/25-26/Teoria-e-Esercitazioni/Codice-20260902/Lezione09/SimpleDate/src/Date.java}{\texttt{Date.java}}}{Immutabilità di Date.java}}\label{chap08_immutabilituxe0-di-date.java}

\begin{lstlisting}
 public class Date {
-    private int day;
-    private int month;
-    private int year;
+    private final int day;
+    private final int month;
+    private final int year;

     // Costruttori...

-    public void setDay(int day) { ... }
-    public void setMonth(int month) { ... }
-    public void setYear(int year) { ... }
\end{lstlisting}

\begin{itemize}
\tightlist
\item
  \textbf{Campi \passthrough{\lstinline!final!}}: Una volta
  inizializzati nel costruttore, \passthrough{\lstinline!day!},
  \passthrough{\lstinline!month!} e \passthrough{\lstinline!year!} non
  possono più essere riassegnati.
\item
  \textbf{Rimozione dei Setter}: I metodi mutatori
  \passthrough{\lstinline!setDay!}, \passthrough{\lstinline!setMonth!} e
  \passthrough{\lstinline!setYear!} vengono eliminati.
\item
  \textbf{Pattern Immutabile}: Un oggetto \passthrough{\lstinline!Date!}
  è ora a prova di manomissione. Se un programma vuole calcolare ``il
  giorno successivo'', non può mutare l'oggetto esistente, ma
  \textbf{deve instanziare un nuovo oggetto
  \passthrough{\lstinline!Date!}} con le coordinate aggiornate
  (esattamente come avviene per le \passthrough{\lstinline!String!}).
\end{itemize}

\begin{center}\rule{0.5\linewidth}{0.5pt}\end{center}

\subsection{\texorpdfstring{Il Design Pattern Singleton:
\href{file:///home/wally/Documenti/GitHub/Programmazione-II/25-26/Teoria-e-Esercitazioni/Codice-20260902/Lezione09/SimpleDate/src/BirthDay.java}{\texttt{BirthDay.java}}}{Il Design Pattern Singleton: BirthDay.java}}\label{chap08_il-design-pattern-singleton-birthday.java}

\begin{lstlisting}[language=Java]
public class BirthDay {
    private static Date date = new Date(1, 1, 1970);
    private static BirthDay instance; // Riferimento all'unica istanza

    // 1. Costruttore privato: impedisce istanziazioni esterne con 'new BirthDay()'
    private BirthDay() { }

    // 2. Metodo statico di fabbrica con Lazy Initialization
    public static BirthDay getInstance() {
        if (instance == null) 
            instance = new BirthDay();
        return instance;
    }

    public Date getDate() { 
        return date; 
    }
}
\end{lstlisting}

\subsection{{[}Approfondimento{]} Analisi Architetturale del
Singleton:}\label{chap08_approfondimento-analisi-architetturale-del-singleton}

\begin{itemize}
\tightlist
\item
  \textbf{Scopo del Pattern}: Garantire che per tutta la durata
  dell'applicazione esista \textbf{una e una sola istanza} di una classe
  nello Heap e fornire un punto di accesso globale ad essa.
\item
  \textbf{I Tre Pilastri del Singleton in Java}:

  \begin{enumerate}
  \def\labelenumi{\arabic{enumi}.}
  \tightlist
  \item
    \textbf{Costruttore \passthrough{\lstinline!private!}}: blocca
    l'accesso a \passthrough{\lstinline!new BirthDay()!} da qualsiasi
    altra classe (tentare di chiamarlo genera un errore a compile-time).
  \item
    \textbf{Campo statico privato (\passthrough{\lstinline!instance!})}:
    contiene l'unico riferimento all'oggetto.
  \item
    \textbf{Metodo factory pubblico statico
    (\passthrough{\lstinline!getInstance()!})}: implementa la \emph{Lazy
    Initialization} (crea l'oggetto solo alla prima richiesta,
    riusandolo per tutte le successive).
  \end{enumerate}
\end{itemize}

\begin{center}\rule{0.5\linewidth}{0.5pt}\end{center}

\subsection{\texorpdfstring{Verifica in
\href{file:///home/wally/Documenti/GitHub/Programmazione-II/25-26/Teoria-e-Esercitazioni/Codice-20260902/Lezione09/SimpleDate/src/MainDate.java}{\texttt{MainDate.java}}
e il Bug dei
Getter}{Verifica in MainDate.java e il Bug dei Getter}}\label{chap08_verifica-in-maindate.java-e-il-bug-dei-getter}

\begin{lstlisting}[language=Java]
BirthDay bDay = BirthDay.getInstance();
// BirthDay day1 = new BirthDay(); // COMPILE-TIME ERROR!
BirthDay day1 = BirthDay.getInstance();
System.out.println("bDay ?= day1: " + (bDay == day1)); // STAMPA TRUE!

Date today = new Date(27, 10, 2025);
// today.setDay(today.getDay() + 1); // COMPILE-TIME ERROR: l'oggetto è immutabile!
Date tomorrow = new Date(today.getDay() + 1, today.getMonth(), today.getYear());
\end{lstlisting}

\begin{studywarning}[Attenzione / Trappola d'Esame] \textbf{Il Bug Persistente dei Getter e l'Effetto a
Runtime}: Anche in \passthrough{\lstinline!Lezione09!},
\passthrough{\lstinline!Date.java!} ha ancora il refuso di copia-incolla
alle righe 43--44:

\begin{lstlisting}[language=Java]
public int getMonth() { return day; } // Restituisce il campo day!
public int getYear() { return day; }  // Restituisce il campo day!
\end{lstlisting}

Quando \passthrough{\lstinline!MainDate!} crea
\passthrough{\lstinline!tomorrow!} invocando
\passthrough{\lstinline!today.getMonth()!} su una data con
\passthrough{\lstinline!day = 27!}, ottiene \passthrough{\lstinline!27!}
come mese! Di conseguenza, il costruttore riceve
\passthrough{\lstinline!(28, 27, 27)!} e
\passthrough{\lstinline!verify()!} stampa:
\passthrough{\lstinline"Illegal date!"}
\passthrough{\lstinline!tomorrow: 28/27/27!} Nel tuo codice è opportuno
correggere questi getter (\passthrough{\lstinline!return month;!} e
\passthrough{\lstinline!return year;!}).
\end{studywarning}

\begin{center}\rule{0.5\linewidth}{0.5pt}\end{center}

\section{🔄 Corrispondenza con i Tuoi Moduli
Workspace}\label{chap08_corrispondenza-con-i-tuoi-moduli-workspace}

Nel tuo repository
\href{file:///home/wally/Documenti/GitHub/Programmazione-II/25-26/esercizi-wally-25-26}{\passthrough{\lstinline!esercizi-wally-25-26!}}:
- Nel modulo \textbf{\passthrough{\lstinline!es09!}}: hai formalizzato
questo esatto pattern creando
\href{file:///home/wally/Documenti/GitHub/Programmazione-II/25-26/esercizi-wally-25-26/es09/src/main/java/es09/SingletonPattern.java}{\passthrough{\lstinline!SingletonPattern.java!}}
con costruttore privato e metodo factory
\passthrough{\lstinline!getInstance()!}.

\begin{center}\rule{0.5\linewidth}{0.5pt}\end{center}

\begin{studynote}[Nota di Studio] Con la \textbf{Lezione 09} abbiamo consolidato la
manipolazione avanzata degli array, le proprietà degli enum,
l'immutabilità e il pattern Singleton.

Il prossimo blocco è la \textbf{Lezione 10 / Slide T10: \emph{Java
Variables and Call Stack}}: - Modello formale della memoria JVM:
\textbf{Call Stack (Stack Frames)} vs \textbf{Heap (Objects \& Class
Data)}. - Variabili d'istanza, variabili di classe
(\passthrough{\lstinline!static!}), variabili locali e parametri di
metodo. - Il passaggio dei parametri per valore (\emph{pass-by-value})
per tipi primitivi vs reference types. - Il blocco di inizializzazione
statica (\passthrough{\lstinline!static \{ ... \}!}). - \textbf{Il
grande spartiacque del corso}: il debutto del progetto
\textbf{\passthrough{\lstinline!MavenDate!}}, con l'adozione ufficiale
di Apache Maven, \passthrough{\lstinline!pom.xml!}, layout directory
standard e primo packaging!
\end{studynote}

\textbf{Dammi conferma per aprire la Lezione 10 / T10 e analizzare Call
Stack e MavenDate!}

\newpage
\chapter[T10 --- Variabili, Scope, Call Stack e Memo]{T10 --- Variabili, Scope, Call Stack e Memoria Heap}
\label{chap:t10}

{\small\sffamily\color{accentblue!80!black}\textit{Stack Frame, Passaggio dei Parametri per Valore, Garbage Collection e Debutto di MavenDate}}

\vspace{0.8em}

In questo blocco analizziamo in profondità il modello di memoria della
JVM (Stack vs Heap), il passaggio dei parametri per valore, gli oggetti
immutabili e i tipi wrapper
(\href{file:///home/wally/Documenti/GitHub/Programmazione-II/25-26/Teoria-e-Esercitazioni/Slides-20260902/T10\%20-\%20Java\%20Variables\%20and\%20Call\%20Stack.pdf}{T10
- Java Variables and Call Stack.pdf}). Sul fronte pratico analizziamo il
\textbf{grande spartiacque architetturale del corso}: il passaggio dal
progetto manuale a pacchetti
\href{file:///home/wally/Documenti/GitHub/Programmazione-II/25-26/Teoria-e-Esercitazioni/Codice-20260902/Lezione10/SimpleDate}{\passthrough{\lstinline!SimpleDate!}}
al progetto ufficiale
\textbf{\href{file:///home/wally/Documenti/GitHub/Programmazione-II/25-26/Teoria-e-Esercitazioni/Codice-20260902/Lezione10/MavenDate}{\passthrough{\lstinline!MavenDate!}}}
con Apache Maven, file
\href{file:///home/wally/Documenti/GitHub/Programmazione-II/25-26/Teoria-e-Esercitazioni/Codice-20260902/Lezione10/MavenDate/pom.xml}{\passthrough{\lstinline!pom.xml!}},
layout \passthrough{\lstinline!src/main/java!} e collisione di nomi tra
package diversi.

\section{Teoria: Concetti Teorici dalle Slide (Slide
T10)}\label{chap09_teoria-concetti-teorici-dalle-slide-slide-t10}

\subsection{Il Modello di Memoria della JVM: Stack vs Heap (Slide
1--11)}\label{chap09_il-modello-di-memoria-della-jvm-stack-vs-heap-slide-111}

La memoria a disposizione di un processo Java è partizionata in tre aree
principali: 1. \textbf{Call Stack (Stack dei Thread)}: - Memoria
estremamente veloce, gestita a politica LIFO (Push/Pop). - È suddivisa
in \textbf{Stack Frame} (o record di attivazione): ogni invocazione di
metodo crea un nuovo frame nello stack contenente: - La tabella delle
variabili locali. - I parametri formali passati al metodo. - Lo stack
degli operandi (per i calcoli intermedi della CPU/bytecode). - I
riferimenti alla costante del metodo e l'indirizzo di ritorno. - Quando
il metodo termina (con \passthrough{\lstinline!return!} o eccezione), il
frame viene rimosso istantaneamente dallo Stack e tutte le sue variabili
locali vengono distrutte. 2. \textbf{Heap (Memoria Dinamica degli
Oggetti)}: - Memoria globale condivisa da tutti i thread
dell'applicazione. - Ospita tutti gli \textbf{oggetti} istanziati con
\passthrough{\lstinline!new!}, inclusi gli \textbf{array}. - I dati
nello Heap non si deallocano all'uscita dal metodo: sopravvivono finché
esiste almeno un riferimento attivo nello Stack o in altri oggetti che
punta ad essi. Quando diventano irraggiungibili (\emph{unreachable}),
vengono eliminati dal Garbage Collector. 3. \textbf{Method Area /
Metaspace (Memoria di Classe)}: - Ospita il bytecode compilato delle
classi caricate dal ClassLoader, i metadati di tipo e le
\textbf{variabili statiche} (\passthrough{\lstinline!static!}). -
\textbf{Blocco di Inizializzazione Statica
(\passthrough{\lstinline!static \{ ... \}!})}: - Blocco speciale
eseguito \textbf{una sola volta} nel ciclo di vita dell'applicazione, al
momento in cui la classe viene caricata in memoria dalla JVM:
\passthrough{\lstinline!java     static int[] arr;     static \{         arr = new int[3];         arr[0] = 1; arr[1] = 2; arr[2] = 0;     \}!}
- Si usa per inizializzare strutture statiche complesse che richiedono
cicli o controlli prima dell'uso.

\begin{center}\rule{0.5\linewidth}{0.5pt}\end{center}

\subsection{Immutabilità e Passaggio Parametri: Pass-by-Value (Slide
15--25)}\label{chap09_immutabilituxe0-e-passaggio-parametri-pass-by-value-slide-1525}

\begin{itemize}
\tightlist
\item
  \textbf{Oggetti Immutabili (\emph{Immutable Objects})}:

  \begin{itemize}
  \tightlist
  \item
    Un oggetto il cui stato interno non può essere modificato dopo la
    costruzione (es. \passthrough{\lstinline!String!},
    \passthrough{\lstinline!Integer!}, \passthrough{\lstinline!Date!} da
    Lezione 09).
  \item
    \emph{Regole per creare una classe immutabile}:

    \begin{enumerate}
    \def\labelenumi{\arabic{enumi}.}
    \tightlist
    \item
      Rendere tutti i campi \passthrough{\lstinline!private final!}.
    \item
      Non fornire metodi mutatori (nessun setter).
    \item
      Dichiarare la classe \passthrough{\lstinline!final!} (per impedire
      che le sottoclassi violino l'immutabilità con override malevoli).
    \item
      Se l'oggetto contiene riferimenti a strutture mutabili (es. un
      array o un'altra data mutabile), eseguire copie difensive
      (\emph{defensive copy}) nel costruttore e nei getter.
    \end{enumerate}
  \item
    \emph{Vantaggi}: Immuni da side-effects da aliasing, intrinsecamente
    \emph{thread-safe} (possono essere condivisi tra thread concorrenti
    senza sincronizzazione o lock).
  \end{itemize}
\item
  \textbf{Come Java passa i parametri ai metodi: SEMPRE Strictly
  Pass-by-Value}:

  \begin{itemize}
  \tightlist
  \item
    In Java \textbf{non esiste il passaggio per riferimento del C++
    (\passthrough{\lstinline!int\& x!})}.
  \item
    \textbf{Per i tipi primitivi (\passthrough{\lstinline!int!},
    \passthrough{\lstinline!double!})}: il metodo riceve una copia del
    valore binario. Qualsiasi modifica locale sul parametro non ha alcun
    effetto sulla variabile del chiamante.
  \item
    \textbf{Per i tipi reference (oggetti/array)}: il metodo riceve
    \textbf{una copia per valore del riferimento} (dell'indirizzo
    puntato):

    \begin{itemize}
    \item
      \emph{Caso 1 (Mutazione dell'oggetto)}: se usiamo il riferimento
      ricevuto per chiamare un metodo mutatore
      (\passthrough{\lstinline!p.setAge(25)!}), l'oggetto puntato nello
      Heap viene effettivamente modificato per il chiamante!
    \item
      \emph{Caso 2 (Riassegnamento del puntatore --- Il tranello
      d'esame)}:

\begin{lstlisting}[language=Java]
void reset(Person p) {
    p = new Person("Bob"); // Assegna un nuovo oggetto al puntatore locale!
}
\end{lstlisting}

      All'esterno, la variabile del chiamante \textbf{non cambia
      affatto}! Ha semplicemente mutato la copia locale del puntatore
      nel proprio stack frame.
    \end{itemize}
  \end{itemize}
\end{itemize}

\begin{center}\rule{0.5\linewidth}{0.5pt}\end{center}

\subsection{Classi Wrapper, Autoboxing e i loro Trabocchetti (Slide
26--35)}\label{chap09_classi-wrapper-autoboxing-e-i-loro-trabocchetti-slide-2635}

\begin{itemize}
\tightlist
\item
  Per ciascuno degli 8 tipi primitivi, Java mette a disposizione una
  corrispondente \textbf{Wrapper Class} nel package
  \passthrough{\lstinline!java.lang!}:

  \begin{itemize}
  \tightlist
  \item
    \passthrough{\lstinline!byte!} \(\rightarrow\)
    \passthrough{\lstinline!Byte!}
  \item
    \passthrough{\lstinline!short!} \(\rightarrow\)
    \passthrough{\lstinline!Short!}
  \item
    \passthrough{\lstinline!int!} \(\rightarrow\)
    \passthrough{\lstinline!Integer!}
  \item
    \passthrough{\lstinline!long!} \(\rightarrow\)
    \passthrough{\lstinline!Long!}
  \item
    \passthrough{\lstinline!float!} \(\rightarrow\)
    \passthrough{\lstinline!Float!}
  \item
    \passthrough{\lstinline!double!} \(\rightarrow\)
    \passthrough{\lstinline!Double!}
  \item
    \passthrough{\lstinline!char!} \(\rightarrow\)
    \passthrough{\lstinline!Character!}
  \item
    \passthrough{\lstinline!boolean!} \(\rightarrow\)
    \passthrough{\lstinline!Boolean!}
  \end{itemize}
\item
  Tutte le classi wrapper sono \textbf{immutabili}.
\item
  \textbf{Autoboxing e Unboxing (Java 5+)}:

  \begin{itemize}
  \tightlist
  \item
    \emph{Autoboxing}: conversione automatica da primitivo a wrapper
    (\passthrough{\lstinline!Integer x = 5;!} \(\implies\) compilato
    come \passthrough{\lstinline!Integer.valueOf(5)!}).
  \item
    \emph{Unboxing}: estrazione automatica del valore primitivo dal
    wrapper (\passthrough{\lstinline!int y = x;!} \(\implies\) compilato
    come \passthrough{\lstinline!x.intValue()!}).
  \end{itemize}
\item
  {[}Attenzione{]} \textbf{I Due Grandi Pericoli delle Wrapper Classes}:

  \begin{enumerate}
  \def\labelenumi{\arabic{enumi}.}
  \item
    \textbf{Il Trabocchetto dell'Integer Cache (\(-128 \dots +127\))}:

\begin{lstlisting}[language=Java]
Integer a = 100, b = 100;
System.out.println(a == b); // TRUE! (La JVM ricicla lo stesso oggetto cacheato)

Integer c = 200, d = 200;
System.out.println(c == d); // FALSE! (Fuori range [-128, 127], alloca 2 oggetti distinti nello Heap!)
\end{lstlisting}

    Conclusione: \textbf{Mai usare \passthrough{\lstinline!==!} per
    confrontare i wrapper}, usare sempre
    \passthrough{\lstinline!.equals()!}!
  \item
    \textbf{Crash da Unboxing su \passthrough{\lstinline!null!}}:

\begin{lstlisting}[language=Java]
Integer obj = null;
int val = obj; // RUNTIME ERROR: NullPointerException!
\end{lstlisting}

    La JVM tenta di invocare \passthrough{\lstinline!obj.intValue()!},
    fallendo rovinosamente se \passthrough{\lstinline!obj!} è nullo.
  \end{enumerate}
\end{itemize}

\begin{center}\rule{0.5\linewidth}{0.5pt}\end{center}

\section{\texorpdfstring{Codice: Evoluzione del Codice: Il Passaggio a
\texttt{MavenDate}}{Codice: Evoluzione del Codice: Il Passaggio a MavenDate}}\label{chap09_codice-evoluzione-del-codice-il-passaggio-a-mavendate}

In
\href{file:///home/wally/Documenti/GitHub/Programmazione-II/25-26/Teoria-e-Esercitazioni/Codice-20260902/Lezione10}{\passthrough{\lstinline!Lezione10!}},
il docente archivia il vecchio progetto non strutturato e adotta lo
standard industriale \textbf{Apache Maven}.

\subsection{Riorganizzazione dei
Package}\label{chap09_riorganizzazione-dei-package}

Il progetto viene suddiviso in due package distinti: -
\textbf{\passthrough{\lstinline!it.oop.core!}}: il dominio applicativo
(le classi di modello:
\href{file:///home/wally/Documenti/GitHub/Programmazione-II/25-26/Teoria-e-Esercitazioni/Codice-20260902/Lezione10/MavenDate/src/main/java/it/oop/core/Date.java}{\passthrough{\lstinline!Date.java!}},
\href{file:///home/wally/Documenti/GitHub/Programmazione-II/25-26/Teoria-e-Esercitazioni/Codice-20260902/Lezione10/MavenDate/src/main/java/it/oop/core/Language.java}{\passthrough{\lstinline!Language.java!}},
\href{file:///home/wally/Documenti/GitHub/Programmazione-II/25-26/Teoria-e-Esercitazioni/Codice-20260902/Lezione10/MavenDate/src/main/java/it/oop/core/BirthDay.java}{\passthrough{\lstinline!BirthDay.java!}}).
- \textbf{\passthrough{\lstinline!it.oop.ui!}}: l'interfaccia utente
contenente il \passthrough{\lstinline!main!}
(\href{file:///home/wally/Documenti/GitHub/Programmazione-II/25-26/Teoria-e-Esercitazioni/Codice-20260902/Lezione10/MavenDate/src/main/java/it/oop/ui/MainDate.java}{\passthrough{\lstinline!MainDate.java!}}).

\subsection{Il Test Didattico sulla Collisione dei Nomi di
Package}\label{chap09_il-test-didattico-sulla-collisione-dei-nomi-di-package}

Nel package \passthrough{\lstinline!it.oop.ui!}, il docente aggiunge
appositamente una seconda classe:
\href{file:///home/wally/Documenti/GitHub/Programmazione-II/25-26/Teoria-e-Esercitazioni/Codice-20260902/Lezione10/MavenDate/src/main/java/it/oop/ui/Date.java}{\passthrough{\lstinline!it/oop/ui/Date.java!}}:

\begin{lstlisting}[language=Java]
package it.oop.ui;

class Date { // Visibilità package-private
    public int time;
    private int start = 0;

    public Date(int time) { this.time = time; }
}
\end{lstlisting}

E in
\href{file:///home/wally/Documenti/GitHub/Programmazione-II/25-26/Teoria-e-Esercitazioni/Codice-20260902/Lezione10/MavenDate/src/main/java/it/oop/ui/MainDate.java}{\passthrough{\lstinline!MainDate.java!}}
mostra come la JVM risolve l'ambiguità:

\begin{lstlisting}[language=Java]
package it.oop.ui;

import it.oop.core.Date; // 1. L'import esplicito ha la precedenza!

public class MainDate {
    public static void main(String[] args) {
        Date d1 = new Date(7, 1, 2025); // Istanzia it.oop.core.Date!

        // 2. Per istanziare la classe Date del package it.oop.ui,
        // è obbligatorio usare il Fully Qualified Name (FQN):
        it.oop.ui.Date d = new it.oop.ui.Date(12345);
        System.out.println(d.time);      // OK: 'time' è public
        // System.out.println(d.start);  // COMPILE-TIME ERROR: 'start' è private!
    }
}
\end{lstlisting}

\subsection{\texorpdfstring{Anatomia del
\href{file:///home/wally/Documenti/GitHub/Programmazione-II/25-26/Teoria-e-Esercitazioni/Codice-20260902/Lezione10/MavenDate/pom.xml}{\texttt{pom.xml}}
di
\texttt{MavenDate}}{Anatomia del pom.xml di MavenDate}}\label{chap09_anatomia-del-pom.xml-di-mavendate}

\begin{lstlisting}[language=XML]
<project ...>
    <modelVersion>4.0.0</modelVersion>
    <groupId>it.oop</groupId>
    <artifactId>MavenDate</artifactId>
    <version>1.0-SNAPSHOT</version>

    <properties>
        <maven.compiler.source>17</maven.compiler.source>
        <maven.compiler.target>17</maven.compiler.target>
        <project.build.sourceEncoding>UTF-8</project.build.sourceEncoding>
    </properties>

    <build>
        <plugins>
            <!-- Plugin per creare il JAR eseguibile con dipendenze -->
            <plugin>
                <groupId>org.apache.maven.plugins</groupId>
                <artifactId>maven-assembly-plugin</artifactId>
                <version>3.7.1</version>
                <configuration>
                    <descriptorRefs>
                        <descriptorRef>jar-with-dependencies</descriptorRef>
                    </descriptorRefs>
                    <archive>
                        <manifest>
                            <addClasspath>true</addClasspath>
                            <!-- Entrypoint del JAR -->
                            <mainClass>it.oop.ui.MainDate</mainClass>
                        </manifest>
                    </archive>
                </configuration>
                <executions>
                    <execution>
                        <id>assemble-all</id>
                        <phase>package</phase>
                        <goals>
                            <goal>single</goal>
                        </goals>
                    </execution>
                </executions>
            </plugin>
        </plugins>
    </build>
</project>
\end{lstlisting}

Questo file è il prototipo esatto che hai esteso in tutti i moduli del
tuo workspace \passthrough{\lstinline!esercizi-wally-25-26!}: -
\passthrough{\lstinline!compile!}: compila da
\passthrough{\lstinline!src/main/java!} verso
\passthrough{\lstinline!target/classes!}. -
\passthrough{\lstinline!test!}: compila ed esegue i test da
\passthrough{\lstinline!src/test/java!} verso
\passthrough{\lstinline!target/test-classes!}. -
\passthrough{\lstinline!package!}: genera il JAR auto-eseguibile con
\passthrough{\lstinline!META-INF/MANIFEST.MF!} configurato con
\passthrough{\lstinline!Main-Class: it.oop.ui.MainDate!}.

\begin{center}\rule{0.5\linewidth}{0.5pt}\end{center}

\section{🔄 Corrispondenza con i Tuoi Moduli
Workspace}\label{chap09_corrispondenza-con-i-tuoi-moduli-workspace}

Nel tuo repository
\href{file:///home/wally/Documenti/GitHub/Programmazione-II/25-26/esercizi-wally-25-26}{\passthrough{\lstinline!esercizi-wally-25-26!}}:
- \textbf{\passthrough{\lstinline!es08!}}:
\href{file:///home/wally/Documenti/GitHub/Programmazione-II/25-26/esercizi-wally-25-26/es08/src/main/java/es08/Memory.java}{\passthrough{\lstinline!Memory.java!}}
implementa l'allocazione Stack vs Heap e il blocco
\passthrough{\lstinline!static \{ ... \}!}. -
\textbf{\passthrough{\lstinline!es10!}}:
\href{file:///home/wally/Documenti/GitHub/Programmazione-II/25-26/esercizi-wally-25-26/es10/src/main/java/es10/ImmutablePerson.java}{\passthrough{\lstinline!ImmutablePerson.java!}}
applica formalmente tutte le regole di immutabilità di T10
(\passthrough{\lstinline!final class!}, campi
\passthrough{\lstinline!final!}, validazione e assenza di setter).

\begin{center}\rule{0.5\linewidth}{0.5pt}\end{center}

\begin{studynote}[Nota di Studio] Con la \textbf{Lezione 10} abbiamo completato l'intero
quadro su memoria, call stack, tipi wrapper e la transizione a Maven.

Il prossimo blocco è la \textbf{Lezione 11 / Slide T11: \emph{Packages
and Class Visibility}}: - Tassonomia gerarchica dei package e
convenzioni DNS invertite (\passthrough{\lstinline!it.univr...!}). -
Regole di visibilità e accessibilità inter-package: membri
\passthrough{\lstinline!public!}, \passthrough{\lstinline!protected!},
package-private e \passthrough{\lstinline!private!}. - Evoluzione di
\passthrough{\lstinline!MavenDate!}: - Refactoring dell'algoritmo
Gregoriano: introduzione definitiva del \textbf{calcolo dell'anno
bisestile (\passthrough{\lstinline!isLeapYear!})} in
\passthrough{\lstinline!Date.java!}! - Integrazione di
\passthrough{\lstinline!isLeapYear(year)!} in
\passthrough{\lstinline!daysPerMonth(month, year)!} e fine dei mesi
errati per Febbraio!
\end{studynote}

\textbf{Dammi conferma per procedere alla Lezione 11 / T11!}

\newpage
\chapter[T11 --- Package, Moduli e Regole di Visibil]{T11 --- Package, Moduli e Regole di Visibilità}
\label{chap:t11}

{\small\sffamily\color{accentblue!80!black}\textit{Struttura Cartelle, Modificatori public/protected/default/private e MavenDate v2 con equals()}}

\vspace{0.8em}

In questo blocco analizziamo la modularità avanzata offerta dai package,
le regole di risoluzione dei nomi e i livelli di visibilità dei membri
(\href{file:///home/wally/Documenti/GitHub/Programmazione-II/25-26/Teoria-e-Esercitazioni/Slides-20260902/T11\%20-\%20Packages\%20and\%20Class\%20Visibility.pdf}{T11
- Packages and Class Visibility.pdf}). Sul piano pratico, in
\href{file:///home/wally/Documenti/GitHub/Programmazione-II/25-26/Teoria-e-Esercitazioni/Codice-20260902/Lezione11/MavenDate}{\passthrough{\lstinline!Lezione11/MavenDate!}}
assistiamo a una grande svolta architetturale: 1. Nascono le sottoclassi
specializzate
\textbf{\href{file:///home/wally/Documenti/GitHub/Programmazione-II/25-26/Teoria-e-Esercitazioni/Codice-20260902/Lezione11/MavenDate/src/main/java/it/oop/core/ItalianDate.java}{\passthrough{\lstinline!ItalianDate!}}}
e
\textbf{\href{file:///home/wally/Documenti/GitHub/Programmazione-II/25-26/Teoria-e-Esercitazioni/Codice-20260902/Lezione11/MavenDate/src/main/java/it/oop/core/AmericanDate.java}{\passthrough{\lstinline!AmericanDate!}}}
tramite estensione (\passthrough{\lstinline!extends Date!}). 2. Viene
introdotto il modificatore \textbf{\passthrough{\lstinline!protected!}}.
3. Viene implementato l'algoritmo canonico per l'overriding del metodo
\textbf{\passthrough{\lstinline!equals(Object)!}}. 4. Viene risolto il
bug storico dei getter \passthrough{\lstinline!getMonth()!} e
\passthrough{\lstinline!getYear()!}.

\section{Teoria: Concetti Teorici dalle Slide (Slide
T11)}\label{chap10_teoria-concetti-teorici-dalle-slide-slide-t11}

\subsection{Modularità, Gerarchia e Convenzioni dei Package (Slide
1--11)}\label{chap10_modularituxe0-gerarchia-e-convenzioni-dei-package-slide-111}

\begin{itemize}
\tightlist
\item
  \textbf{Definizione di Package}:

  \begin{itemize}
  \tightlist
  \item
    Un package è un raggruppamento logico di classi correlate che scala
    il concetto di incapsulamento a livello di libreria.
  \item
    Definisce un \textbf{namespace isolato}: classi con lo stesso nome
    possono coesistere senza conflitti se collocate in package
    differenti (es. \passthrough{\lstinline!it.oop.ui.Date!} vs
    \passthrough{\lstinline!it.oop.core.Date!}).
  \end{itemize}
\item
  \textbf{Convenzione di Naming Standard}:

  \begin{itemize}
  \tightlist
  \item
    Si usa il nome di dominio Internet invertito in caratteri minuscoli:
    \passthrough{\lstinline!it.univr.mypackage!},
    \passthrough{\lstinline!it.oop.core!}.
  \end{itemize}
\item
  \textbf{Dichiarazione e Corrispondenza col Filesystem}:

  \begin{itemize}
  \tightlist
  \item
    La direttiva \passthrough{\lstinline!package nomepkg;!} deve essere
    \textbf{la prima istruzione non di commento} all'inizio del file
    \passthrough{\lstinline!.java!}.
  \item
    Java impone una corrispondenza 1-a-1 tra la gerarchia dei package e
    la struttura delle cartelle su disco: la classe
    \passthrough{\lstinline!it.oop.core.Date!} deve trovarsi nel
    percorso \passthrough{\lstinline!it/oop/core/Date.java!}.
  \end{itemize}
\item
  \textbf{Sotto-Package e Wildcard}:

  \begin{itemize}
  \tightlist
  \item
    I package possono essere annidati (es.
    \passthrough{\lstinline!java.awt.event!} è un sottopackage di
    \passthrough{\lstinline!java.awt!}).
  \item
    {[}Attenzione{]} \textbf{Regola d'esame}: la direttiva
    \passthrough{\lstinline!import pkg.*;!} importa solo le classi
    direttamente contenute in \passthrough{\lstinline!pkg!}. \textbf{Non
    importa le classi contenute nei suoi sotto-package} (per importare
    \passthrough{\lstinline!java.awt.event!} è obbligatorio scrivere
    \passthrough{\lstinline!import java.awt.event.*;!}).
  \end{itemize}
\item
  \textbf{Il Default Package (Unnamed Package)}:

  \begin{itemize}
  \tightlist
  \item
    Se un file \passthrough{\lstinline!.java!} omette la dichiarazione
    \passthrough{\lstinline!package!}, la classe appartiene al package
    predefinito senza nome.
  \item
    \textbf{Limitazione critica}: Le classi nel default package
    \textbf{non possono essere importate né utilizzate} da classi
    situate in package con nome. Nel software professionale e agli esami
    il suo uso è fortemente scoraggiato.
  \end{itemize}
\item
  \textbf{Compilazione e Packaging con Strumenti Standard e Maven}:

  \begin{itemize}
  \tightlist
  \item
    \passthrough{\lstinline!javac -d bin/ src/pkg/MyClass.java!}: il
    flag \passthrough{\lstinline!-d!} crea automaticamente l'albero
    delle sottocartelle nella directory di destinazione
    \passthrough{\lstinline!bin/!}.
  \item
    Con Maven: il layout \passthrough{\lstinline!src/main/java!} e
    \passthrough{\lstinline!src/test/java!} unito al
    \passthrough{\lstinline!pom.xml!} automatizza compilazione, test
    surefire e packaging in Fat JAR.
  \end{itemize}
\end{itemize}

\begin{center}\rule{0.5\linewidth}{0.5pt}\end{center}

\subsection{Visibilità di Classe e Matrice di Accesso dei Membri (Slide
12--19)}\label{chap10_visibilituxe0-di-classe-e-matrice-di-accesso-dei-membri-slide-1219}

\begin{itemize}
\tightlist
\item
  \textbf{Visibilità a Livello di Classe (Top-Level Classes)}:

  \begin{itemize}
  \tightlist
  \item
    \passthrough{\lstinline!public class Clazz!}: accessibile e
    istanziabile da qualsiasi package del Classpath.
  \item
    \passthrough{\lstinline!class Clazz!} (\emph{Package-Private} /
    default): visibile e utilizzabile \textbf{esclusivamente all'interno
    dello stesso package}. Non può essere usata all'esterno, anche se i
    suoi metodi interni sono dichiarati
    \passthrough{\lstinline!public!}.
  \item
    \emph{Nota}: le classi di primo livello non possono essere
    dichiarate \passthrough{\lstinline!private!} né
    \passthrough{\lstinline!protected!} (questi modificatori sono
    ammessi solo per classi annidate/interne).
  \end{itemize}
\item
  \textbf{La Matrice dei Modificatori di Accesso (Fields e Methods)}:
\end{itemize}

\begin{longtable}[]{@{}
  >{\raggedright\arraybackslash}p{(\linewidth - 8\tabcolsep) * \real{0.1667}}
  >{\centering\arraybackslash}p{(\linewidth - 8\tabcolsep) * \real{0.2083}}
  >{\centering\arraybackslash}p{(\linewidth - 8\tabcolsep) * \real{0.2083}}
  >{\centering\arraybackslash}p{(\linewidth - 8\tabcolsep) * \real{0.2083}}
  >{\centering\arraybackslash}p{(\linewidth - 8\tabcolsep) * \real{0.2083}}@{}}
\toprule\noalign{}
\begin{minipage}[b]{\linewidth}\raggedright
Modificatore
\end{minipage} & \begin{minipage}[b]{\linewidth}\centering
Stessa Classe
\end{minipage} & \begin{minipage}[b]{\linewidth}\centering
Stesso Package
\end{minipage} & \begin{minipage}[b]{\linewidth}\centering
Sottoclasse (Altro Package)
\end{minipage} & \begin{minipage}[b]{\linewidth}\centering
Mondo Esterno
\end{minipage} \\
\midrule\noalign{}
\endhead
\bottomrule\noalign{}
\endlastfoot
\textbf{\passthrough{\lstinline!private!}} & {[}OK{]} & ❌ & ❌ & ❌ \\
\textbf{Default} \emph{(package-private)} & {[}OK{]} & {[}OK{]} & ❌ &
❌ \\
\textbf{\passthrough{\lstinline!protected!}} & {[}OK{]} & {[}OK{]} &
{[}OK{]} & ❌ \\
\textbf{\passthrough{\lstinline!public!}} & {[}OK{]} & {[}OK{]} &
{[}OK{]} & {[}OK{]} \\
\end{longtable}

\begin{itemize}
\tightlist
\item
  \textbf{Risoluzione dei Nomi}:

  \begin{itemize}
  \tightlist
  \item
    Se una classe importa due tipi omonimi o usa una classe locale
    avente lo stesso nome di una classe importata, per disambiguare è
    obbligatorio usare il \textbf{Fully Qualified Name (FQN)}:
    \passthrough{\lstinline!it.oop.ui.Date d = new it.oop.ui.Date(12345);!}
  \end{itemize}
\end{itemize}

\begin{center}\rule{0.5\linewidth}{0.5pt}\end{center}

\section{\texorpdfstring{Codice: Evoluzione del Codice:
\href{file:///home/wally/Documenti/GitHub/Programmazione-II/25-26/Teoria-e-Esercitazioni/Codice-20260902/Lezione11/MavenDate}{\texttt{Lezione11/MavenDate}}}{Codice: Evoluzione del Codice: Lezione11/MavenDate}}\label{chap10_codice-evoluzione-del-codice-lezione11mavendate}

In \passthrough{\lstinline!Lezione11!}, il codice di
\passthrough{\lstinline!MavenDate!} compie un balzo in avanti con
l'introduzione di gerarchie di classi e l'overriding:

\subsection{\texorpdfstring{Generalizzazione in
\href{file:///home/wally/Documenti/GitHub/Programmazione-II/25-26/Teoria-e-Esercitazioni/Codice-20260902/Lezione11/MavenDate/src/main/java/it/oop/core/Date.java}{\texttt{Date.java}}}{Generalizzazione in Date.java}}\label{chap10_generalizzazione-in-date.java}

\begin{lstlisting}[language=Java]
package it.oop.core;

public class Date {
    // 1. Campo 'protected': accessibile direttamente dalle sottoclassi
    protected final int day;
    // 2. Campi 'private': incapsulati e accessibili solo tramite getter
    private final int month;
    private final int year;

    public Date(int day, int month, int year) {
        this.day = day;
        this.month = month;
        this.year = year;
        verify();
    }

    // 3. I GETTER SONO STATI CORRETTI!
    public int getDay() { return day; }
    public int getMonth() { return month; } // Non restituisce più day!
    public int getYear() { return year; }   // Non restituisce più day!

    // 4. Implementazione canonica di equals(Object)
    @Override
    public boolean equals(Object other) {
        if (other == null) return false;              // Controllo null-safety
        if (this == other) return true;               // Controllo identità (stesso indirizzo)
        if (!(other instanceof Date)) return false;   // Controllo di tipo/ereditarietà
        Date otherAsDate = (Date) other;              // Downcast sicuro
        return day == otherAsDate.getDay()            // Confronto dei campi
                && month == otherAsDate.getMonth()
                && year == otherAsDate.getYear();
    }
}
\end{lstlisting}

\subsection{\texorpdfstring{{[}Approfondimento{]} I 5 Passi Fondamentali
dell'Algoritmo
\texttt{equals(Object)}:}{{[}Approfondimento{]} I 5 Passi Fondamentali dell'Algoritmo equals(Object):}}\label{chap10_approfondimento-i-5-passi-fondamentali-dellalgoritmo-equalsobject}

\begin{enumerate}
\def\labelenumi{\arabic{enumi}.}
\tightlist
\item
  \passthrough{\lstinline!if (other == null) return false;!}: garantisce
  che \passthrough{\lstinline!d.equals(null)!} ritorni
  \passthrough{\lstinline!false!} senza lanciare
  \passthrough{\lstinline!NullPointerException!}.
\item
  \passthrough{\lstinline!if (this == other) return true;!}:
  ottimizzazione immediata (riflessività). Se sono lo stesso oggetto
  nello Heap, sono identici.
\item
  \passthrough{\lstinline"if (!(other instanceof Date)) return false;"}:
  verifica se l'oggetto passato è compatibile con la gerarchia
  \passthrough{\lstinline!Date!}.
\item
  \passthrough{\lstinline!Date otherAsDate = (Date) other;!}: downcast
  necessario perché la firma di \passthrough{\lstinline!equals!} accetta
  un generico \passthrough{\lstinline!Object!}.
\item
  Confronto logico dei campi che definiscono lo stato:
  \passthrough{\lstinline!day!}, \passthrough{\lstinline!month!},
  \passthrough{\lstinline!year!}.
\end{enumerate}

\begin{center}\rule{0.5\linewidth}{0.5pt}\end{center}

\subsection{\texorpdfstring{Le Nuove Sottoclassi:
\href{file:///home/wally/Documenti/GitHub/Programmazione-II/25-26/Teoria-e-Esercitazioni/Codice-20260902/Lezione11/MavenDate/src/main/java/it/oop/core/ItalianDate.java}{\texttt{ItalianDate.java}}
e
\href{file:///home/wally/Documenti/GitHub/Programmazione-II/25-26/Teoria-e-Esercitazioni/Codice-20260902/Lezione11/MavenDate/src/main/java/it/oop/core/AmericanDate.java}{\texttt{AmericanDate.java}}}{Le Nuove Sottoclassi: ItalianDate.java e AmericanDate.java}}\label{chap10_le-nuove-sottoclassi-italiandate.java-e-americandate.java}

\begin{lstlisting}[language=Java]
package it.oop.core;

public class ItalianDate extends Date {
    private static final String FORMAT = "dd/mm/yyyy";
    private static final String[] MONTHS = { "gennaio", "febbraio", ... };

    // Invocazione del costruttore genitore con 'super'
    public ItalianDate(int day, int month, int year) {
        super(day, month, year);
    }

    public String printFormat() { return FORMAT; }

    public String prettyPrint() {
        // 'day' è accessibile direttamente perché 'protected'!
        // 'getYear()' è invocato tramite metodo perché 'year' è 'private' in Date!
        return day + " " + getMonthAsString() + " " + getYear();
    }

    @Override
    public String toString() {
        return day + "/" + getMonth() + "/" + getYear();
    }
}
\end{lstlisting}

E simmetricamente per \passthrough{\lstinline!AmericanDate!}:

\begin{lstlisting}[language=Java]
public class AmericanDate extends Date {
    private static final String FORMAT = "mm/dd/yyyy";

    public AmericanDate(int day, int month, int year) {
        super(day, month, year);
    }

    @Override
    public String toString() {
        // Mese prima del giorno!
        return getMonth() + "/" + getDay() + "/" + getYear();
    }
}
\end{lstlisting}

\begin{center}\rule{0.5\linewidth}{0.5pt}\end{center}

\subsection{\texorpdfstring{Polimorfismo, Upcasting e Downcasting in
\href{file:///home/wally/Documenti/GitHub/Programmazione-II/25-26/Teoria-e-Esercitazioni/Codice-20260902/Lezione11/MavenDate/src/main/java/it/oop/ui/MainDate.java}{\texttt{MainDate.java}}}{Polimorfismo, Upcasting e Downcasting in MainDate.java}}\label{chap10_polimorfismo-upcasting-e-downcasting-in-maindate.java}

Analizziamo i comportamenti a runtime testati dal docente:

\begin{lstlisting}[language=Java]
Date date = new Date(3, 11, 2025);
ItalianDate itDate = new ItalianDate(3, 11, 2025);
AmericanDate usDate = new AmericanDate(3, 11, 2025);

// 1. Confronti con equals():
System.out.println("date ?= itDate: " + date.equals(itDate));     // Stampa TRUE!
System.out.println("itDate ?= usDate: " + itDate.equals(usDate)); // Stampa TRUE!
\end{lstlisting}

\emph{Perché stampano \passthrough{\lstinline!true!}?} Perché sia
\passthrough{\lstinline!itDate!} che \passthrough{\lstinline!usDate!}
estendono \passthrough{\lstinline!Date!}: l'operatore
\passthrough{\lstinline!instanceof Date!} è soddisfatto, e i campi
(giorno 3, mese 11, anno 2025) coincidono perfettamente!

\begin{lstlisting}[language=Java]
// 2. Upcasting (Polimorfismo di Inclusione):
Date d = new ItalianDate(1, 1, 1970); // LECITO: ItalianDate È UN Date!

// 3. Tipi Incompatibili tra rami fratelli:
// AmericanDate ad = new ItalianDate(1, 1, 1970); // COMPILE-TIME ERROR!

// 4. Downcasting esplicito lecito:
ItalianDate id = (ItalianDate) d; // LECITO a runtime: l'oggetto nello Heap è davvero ItalianDate!
System.out.println(id.printFormat()); // Stampa dd/mm/yyyy

// 5. Downcasting illegale a runtime:
// ItalianDate idFalso = (ItalianDate) date; // CRASH RUNTIME: ClassCastException!
// (perché 'date' è un Date generico, non contiene la specializzazione ItalianDate)
\end{lstlisting}

\begin{center}\rule{0.5\linewidth}{0.5pt}\end{center}

\section{🔄 Corrispondenza con i Tuoi Moduli
Workspace}\label{chap10_corrispondenza-con-i-tuoi-moduli-workspace}

Nel tuo repository
\href{file:///home/wally/Documenti/GitHub/Programmazione-II/25-26/esercizi-wally-25-26}{\passthrough{\lstinline!esercizi-wally-25-26!}}:
- \textbf{\passthrough{\lstinline!es13!}}:
\href{file:///home/wally/Documenti/GitHub/Programmazione-II/25-26/esercizi-wally-25-26/es13/src/main/java/es13/model/Calculator.java}{\passthrough{\lstinline!Calculator.java!}}
e
\href{file:///home/wally/Documenti/GitHub/Programmazione-II/25-26/esercizi-wally-25-26/es13/src/test/java/es13/TestCalculator.java}{\passthrough{\lstinline!TestCalculator.java!}}
applicano la separazione \passthrough{\lstinline!es13.model!} e
\passthrough{\lstinline!es13!} illustrata a lezione (Slide T11,
pp.~3--14). - \textbf{\passthrough{\lstinline!es14!}}: implementa
esattamente questa architettura a oggetti multi-package
(\passthrough{\lstinline!it.oop.core!} e
\passthrough{\lstinline!it.oop.ui!}) con
\href{file:///home/wally/Documenti/GitHub/Programmazione-II/25-26/esercizi-wally-25-26/es14/src/main/java/it/oop/core/ItalianDate.java}{\passthrough{\lstinline!ItalianDate!}},
\href{file:///home/wally/Documenti/GitHub/Programmazione-II/25-26/esercizi-wally-25-26/es14/src/main/java/it/oop/core/AmericanDate.java}{\passthrough{\lstinline!AmericanDate!}},
\href{file:///home/wally/Documenti/GitHub/Programmazione-II/25-26/esercizi-wally-25-26/es14/src/main/java/it/oop/core/Date.java}{\passthrough{\lstinline!Date!}}
e i relativi test in
\href{file:///home/wally/Documenti/GitHub/Programmazione-II/25-26/esercizi-wally-25-26/es14/src/test/java/it/oop/core/TestItalianDate.java}{\passthrough{\lstinline!src/test/java/it/oop/core/TestItalianDate.java!}}.

\begin{center}\rule{0.5\linewidth}{0.5pt}\end{center}

\begin{studynote}[Nota di Studio] Con la \textbf{Lezione 11} abbiamo visto l'introduzione
dell'ereditarietà (\passthrough{\lstinline!extends!}), la visibilità
\passthrough{\lstinline!protected!}, l'overriding di
\passthrough{\lstinline!equals()!} e le regole di cast tra classi.

Il prossimo blocco è la \textbf{Lezione 12 / Slide T12:
\emph{Inheritance and Polymorphism}}: - Teoria approfondita
sull'ereditarietà: riuso, estensione e specializzazione. - La classe
radice universale \textbf{\passthrough{\lstinline!java.lang.Object!}} e
i suoi metodi (\passthrough{\lstinline!toString!},
\passthrough{\lstinline!equals!}, \passthrough{\lstinline!hashCode!},
\passthrough{\lstinline!getClass!}, \passthrough{\lstinline!clone!}). -
Il costrutto \passthrough{\lstinline!super!}: invocazione di costruttori
(\passthrough{\lstinline!super(...)!}) e invocazione di metodi della
superclasse (\passthrough{\lstinline!super.metodo()!}). -
\textbf{Polimorfismo e Dynamic Method Dispatch (Late Binding)}: come la
JVM decide a runtime quale metodo eseguire tramite la \emph{vtable}. -
Regole di overriding vs overloading e l'annotazione
\passthrough{\lstinline!@Override!}. - Evoluzione del codice in
\passthrough{\lstinline!Lezione12!}: - La classe
\passthrough{\lstinline!Date!} diventa la superclasse polimorfica di
riferimento. - Aggiunta del metodo
\passthrough{\lstinline!isLeapYear(int year)!} per la validazione
completa dei bisestili gregoriani!
\end{studynote}

\textbf{Dammi conferma per aprire la Lezione 12 / T12 e proseguire!}

\newpage
\chapter[T12 --- Ereditarietà, Dynamic Binding e Pol]{T12 --- Ereditarietà, Dynamic Binding e Polimorfismo}
\label{chap:t12}

{\small\sffamily\color{accentblue!80!black}\textit{Estensione di Classi, super, Overriding vs Overloading, Polimorfismo e MavenDate v3}}

\vspace{0.8em}

In questo blocco analizziamo in modo esaustivo i due pilastri più
importanti dell'OOP in Java: l'\textbf{Ereditarietà} e il
\textbf{Polimorfismo}
(\href{file:///home/wally/Documenti/GitHub/Programmazione-II/25-26/Teoria-e-Esercitazioni/Slides-20260902/T12\%20-\%20Inheritance\%20and\%20Polymorphism.pdf}{T12
- Inheritance and Polymorphism.pdf}). Nel codice di
\href{file:///home/wally/Documenti/GitHub/Programmazione-II/25-26/Teoria-e-Esercitazioni/Codice-20260902/Lezione12/MavenDate}{\passthrough{\lstinline!Lezione12/MavenDate!}},
la gerarchia evolve introducendo la classe astratta intermedia
\textbf{\href{file:///home/wally/Documenti/GitHub/Programmazione-II/25-26/Teoria-e-Esercitazioni/Codice-20260902/Lezione12/MavenDate/src/main/java/it/oop/core/FormattedDate.java}{\passthrough{\lstinline!FormattedDate!}}},
l'interfaccia
\textbf{\href{file:///home/wally/Documenti/GitHub/Programmazione-II/25-26/Teoria-e-Esercitazioni/Codice-20260902/Lezione12/MavenDate/src/main/java/it/oop/core/Time.java}{\passthrough{\lstinline!Time!}}},
la classe estesa
\textbf{\href{file:///home/wally/Documenti/GitHub/Programmazione-II/25-26/Teoria-e-Esercitazioni/Codice-20260902/Lezione12/MavenDate/src/main/java/it/oop/core/TimeStamp.java}{\passthrough{\lstinline!TimeStamp!}}}
e l'interfaccia standard \textbf{\passthrough{\lstinline!Comparable!}}
per l'ordinamento naturale degli array.

\section{Teoria: Concetti Teorici dalle Slide (Slide
T12)}\label{chap11_teoria-concetti-teorici-dalle-slide-slide-t12}

\subsection{\texorpdfstring{Ereditarietà di Classe (\texttt{extends}) e
Riuso (Slide
1--17)}{Ereditarietà di Classe (extends) e Riuso (Slide 1--17)}}\label{chap11_ereditarietuxe0-di-classe-extends-e-riuso-slide-117}

\begin{itemize}
\tightlist
\item
  \textbf{Relazione ``IS-A'' vs ``HAS-A''}:

  \begin{itemize}
  \tightlist
  \item
    \emph{Ereditarietà (``IS-A'')}: un'auto \emph{è un} veicolo
    (\passthrough{\lstinline!Car extends Vehicle!}); una data italiana
    \emph{è una} data
    (\passthrough{\lstinline!ItalianDate extends Date!}).
  \item
    \emph{Composizione (``HAS-A'')}: un'auto \emph{ha un} motore
    (\passthrough{\lstinline!Car!} ha un campo
    \passthrough{\lstinline!Engine!}); una persona \emph{ha una} data di
    nascita (\passthrough{\lstinline!Person!} ha un campo
    \passthrough{\lstinline!Date!}).
  \end{itemize}
\item
  \textbf{Ereditarietà Singola in Java}:

  \begin{itemize}
  \tightlist
  \item
    In Java \textbf{non esiste l'ereditarietà multipla tra classi} (a
    differenza del C++). Una classe può estendere al massimo \textbf{una
    sola classe genitore diretta}
    (\passthrough{\lstinline!extends SuperClass!}). Questo previene il
    cosiddetto \emph{Diamond Problem}.
  \end{itemize}
\item
  \textbf{Overriding vs Overloading}:

  \begin{itemize}
  \tightlist
  \item
    \emph{Method Overloading}: metodi nella stessa classe con lo stesso
    nome ma parametri diversi (risolto staticamente a compile-time).
  \item
    \emph{Method Overriding}: una sottoclasse riscrive l'implementazione
    di un metodo ereditato mantenendo \textbf{la stessa identica
    segnatura} (nome, tipi dei parametri, ordine) e un tipo di ritorno
    identico o covariante.
  \item
    \textbf{L'Annotazione \passthrough{\lstinline!@Override!}}: Non è
    sintatticamente obbligatoria per far funzionare l'overriding, ma è
    una best practice fondamentale: dice al compilatore di verificare
    che quel metodo esista davvero nella superclasse, intercettando
    errori di battitura (es. scrivere per sbaglio
    \passthrough{\lstinline!tostring()!} con la `s' minuscola verrebbe
    segnalato come errore).
  \end{itemize}
\item
  \textbf{Costruttori ed Ereditarietà: La Parola Chiave
  \passthrough{\lstinline!super!}}:

  \begin{itemize}
  \tightlist
  \item
    {[}Attenzione{]} \textbf{I costruttori NON vengono mai ereditati
    dalle sottoclassi}.
  \item
    All'interno del costruttore di una sottoclasse, la prima istruzione
    deve essere una chiamata a \passthrough{\lstinline!super(...)!}
    oppure a un altro costruttore con
    \passthrough{\lstinline!this(...)!}.
  \item
    Se il programmatore non scrive nulla, il compilatore inserisce
    automaticamente \textbf{\passthrough{\lstinline!super();!}}
    (invocazione del costruttore vuoto della superclasse).
  \item
    {[}Attenzione{]} \textbf{Trappola d'esame}: Se la superclasse non
    possiede un costruttore senza parametri (perché ne ha definito uno
    con parametri e non ha aggiunto quello di default), omettere
    \passthrough{\lstinline!super(...)!} genera un \textbf{Compile-Time
    Error}!
  \item
    \passthrough{\lstinline!super.metodo()!} permette inoltre di
    invocare l'implementazione originaria della superclasse, bypassando
    l'override locale.
  \end{itemize}
\item
  \textbf{Classi e Metodi \passthrough{\lstinline!final!}}:

  \begin{itemize}
  \tightlist
  \item
    \passthrough{\lstinline!final class!}: non può essere estesa da
    nessun'altra classe (es. \passthrough{\lstinline!String!},
    \passthrough{\lstinline!Integer!}).
  \item
    \passthrough{\lstinline!final method!}: non può essere sovrascritto
    (\emph{overridden}) da nessuna sottoclasse.
  \end{itemize}
\end{itemize}

\begin{center}\rule{0.5\linewidth}{0.5pt}\end{center}

\subsection{\texorpdfstring{La Classe Radice Universale:
\texttt{java.lang.Object} (Slide
18--31)}{La Classe Radice Universale: java.lang.Object (Slide 18--31)}}\label{chap11_la-classe-radice-universale-java.lang.object-slide-1831}

\begin{itemize}
\tightlist
\item
  \emph{«The One Ring»}: in Java qualsiasi classe estende implicitamente
  \passthrough{\lstinline!java.lang.Object!}. Se una classe non
  specifica \passthrough{\lstinline!extends!}, estende direttamente
  \passthrough{\lstinline!Object!}.
\item
  \textbf{I Metodi Universali di \passthrough{\lstinline!Object!}}:

  \begin{itemize}
  \tightlist
  \item
    \passthrough{\lstinline!public String toString()!}: rappresentazione
    testuale (\passthrough{\lstinline!NomeClasse@hashcode!}).
  \item
    \passthrough{\lstinline!public boolean equals(Object obj)!}:
    confronto semantico (di default confronta l'identità con
    \passthrough{\lstinline!==!}).
  \item
    \passthrough{\lstinline!public int hashCode()!}: valore hash
    associato all'oggetto.
  \item
    \passthrough{\lstinline!public final Class<?> getClass()!}:
    restituisce il descrittore a runtime del tipo della classe.
  \end{itemize}
\item
  \textbf{Il Contratto Formale di \passthrough{\lstinline!equals!}}:
  Ogni implementazione di \passthrough{\lstinline!equals!} deve
  rispettare 5 proprietà matematiche:

  \begin{enumerate}
  \def\labelenumi{\arabic{enumi}.}
  \tightlist
  \item
    \emph{Riflessiva}: \passthrough{\lstinline!x.equals(x)!} è sempre
    \passthrough{\lstinline!true!}.
  \item
    \emph{Simmetrica}: \passthrough{\lstinline!x.equals(y)!} ritorna
    \passthrough{\lstinline!true!} se e solo se
    \passthrough{\lstinline!y.equals(x)!} ritorna
    \passthrough{\lstinline!true!}.
  \item
    \emph{Transitiva}: se \passthrough{\lstinline!x.equals(y)!} è
    \passthrough{\lstinline!true!} e
    \passthrough{\lstinline!y.equals(z)!} è
    \passthrough{\lstinline!true!}, allora
    \passthrough{\lstinline!x.equals(z)!} è
    \passthrough{\lstinline!true!}.
  \item
    \emph{Consistente}: invocazioni multiple producono lo stesso
    risultato finché lo stato non muta.
  \item
    \emph{Non-Nullità}: per qualsiasi
    \passthrough{\lstinline"x != null"}, la chiamata
    \passthrough{\lstinline!x.equals(null)!} \textbf{deve restituire
    \passthrough{\lstinline!false!}} senza lanciare eccezioni.
  \end{enumerate}
\end{itemize}

\begin{center}\rule{0.5\linewidth}{0.5pt}\end{center}

\subsection{Polimorfismo e Dynamic Method Dispatch (Late Binding) (Slide
32--50)}\label{chap11_polimorfismo-e-dynamic-method-dispatch-late-binding-slide-3250}

\begin{itemize}
\item
  \textbf{Principio di Sostituzione di Liskov (LSP)}: Se \(S\) è
  sottotipo di \(T\), qualsiasi porzione di codice scritta per
  interagire con \(T\) deve poter operare in modo trasparente e corretto
  con un'istanza di \(S\).
\item
  \textbf{Tipo Statico vs Tipo Dinamico}:

\begin{lstlisting}[language=Java]
Date d = new ItalianDate(1, 1, 1970);
\end{lstlisting}

  \begin{itemize}
  \tightlist
  \item
    \textbf{Tipo Statico (a Compile-Time)}: è il tipo dichiarato della
    variabile reference (\passthrough{\lstinline!Date!}). Il compilatore
    accetta solo invocazioni a metodi e accessi a campi definiti nella
    classe \passthrough{\lstinline!Date!}.
  \item
    \textbf{Tipo Dinamico (a Runtime)}: è il tipo effettivo dell'oggetto
    allocato nello Heap (\passthrough{\lstinline!ItalianDate!}).
  \end{itemize}
\item
  \textbf{Dynamic Method Dispatch (Late Binding)}: Quando a runtime
  viene eseguita una chiamata a un metodo d'istanza (es.
  \passthrough{\lstinline!d.toString()!}), la JVM consulta la
  \textbf{Virtual Method Table (vtable)} dell'oggetto concreto nello
  Heap e seleziona \textbf{la versione del metodo definita dal Tipo
  Dinamico}! Non importa che \passthrough{\lstinline!d!} sia dichiarata
  come \passthrough{\lstinline!Date!}: verrà eseguito il
  \passthrough{\lstinline!toString()!} specializzato di
  \passthrough{\lstinline!ItalianDate!}.
\item
  \textbf{Regole di Conversione di Tipo (Casting tra Oggetti)}:

  \begin{itemize}
  \item
    \textbf{Upcasting (verso la superclasse)}:
    \passthrough{\lstinline!Date d = new ItalianDate(...);!}
    \(\implies\) \textbf{Implicito, automatico e sempre sicuro al
    100\%}.
  \item
    \textbf{Downcasting (verso la sottoclasse)}:
    \passthrough{\lstinline!ItalianDate id = (ItalianDate) d;!}
    \(\implies\) \textbf{Esplicito obbligatorio}.

    \begin{itemize}
    \tightlist
    \item
      Se l'oggetto concreto nello Heap è effettivamente un'istanza
      compatibile, il cast ha successo.
    \item
      Se l'oggetto concreto nello Heap è di tipo diverso (es. un
      \passthrough{\lstinline!Date!} generico o un
      \passthrough{\lstinline!AmericanDate!}), la JVM interrompe
      l'esecuzione e lancia una
      \textbf{\passthrough{\lstinline!ClassCastException!}}!
    \end{itemize}
  \item
    \textbf{Downcasting Sicuro con
    \passthrough{\lstinline!instanceof!}}: Per evitare crash a runtime
    si verifica preventivamente la compatibilità:

\begin{lstlisting}[language=Java]
if (d instanceof ItalianDate) {
    ItalianDate id = (ItalianDate) d;
}
\end{lstlisting}
  \end{itemize}
\end{itemize}

\begin{center}\rule{0.5\linewidth}{0.5pt}\end{center}

\section{\texorpdfstring{Codice: Evoluzione del Codice:
\href{file:///home/wally/Documenti/GitHub/Programmazione-II/25-26/Teoria-e-Esercitazioni/Codice-20260902/Lezione12/MavenDate}{\texttt{Lezione12/MavenDate}}}{Codice: Evoluzione del Codice: Lezione12/MavenDate}}\label{chap11_codice-evoluzione-del-codice-lezione12mavendate}

In \passthrough{\lstinline!Lezione12!} la gerarchia del progetto diventa
ricca e completa:

\begin{lstlisting}
classDiagram
    class Object {
        +toString()
        +equals(Object)
    }
    class Comparable {
        <<interface>>
        +compareTo(Object)
    }
    class Time {
        <<interface>>
        +getSeconds()
        +getMinutes()
        +getHours()
    }
    class Date {
        #day: int
        -month: int
        -year: int
        +toString()
        +equals(Object)
        +compareTo(Object)
    }
    class FormattedDate {
        <<abstract>>
        #format: String
        #months: String[]
        +printFormat() final
        +getMonthAsString() final
        +prettyPrint()* String
    }
    class ItalianDate {
        +prettyPrint()
        +toString()
    }
    class AmericanDate {
        +prettyPrint()
        +toString()
    }
    class TimeStamp {
        #seconds: int
        #minutes: int
        #hours: int
        +toString()
        +equals(Object)
    }

    Object <|-- Date
    Comparable <|.. Date
    Date <|-- FormattedDate
    FormattedDate <|-- ItalianDate
    FormattedDate <|-- AmericanDate
    Date <|-- TimeStamp
    Time <|.. TimeStamp
\end{lstlisting}

\subsection{\texorpdfstring{L'Interfaccia \texttt{Comparable} in
\href{file:///home/wally/Documenti/GitHub/Programmazione-II/25-26/Teoria-e-Esercitazioni/Codice-20260902/Lezione12/MavenDate/src/main/java/it/oop/core/Date.java}{\texttt{Date.java}}}{L'Interfaccia Comparable in Date.java}}\label{chap11_linterfaccia-comparable-in-date.java}

\passthrough{\lstinline!Date!} implementa l'interfaccia standard
\passthrough{\lstinline!java.lang.Comparable!} per consentire
l'ordinamento naturale:

\begin{lstlisting}[language=Java]
public class Date implements Comparable {
    ...
    @Override
    public int compareTo(Object other) {
        Date otherAsDate = (Date) other;
        int diff = this.year - otherAsDate.getYear();
        if (diff != 0) return diff;
        diff = this.month - otherAsDate.getMonth();
        if (diff != 0) return diff;
        return this.day - otherAsDate.getDay();
    }
}
\end{lstlisting}

\emph{Logica di confronto}: confronta prima gli anni, se uguali
confronta i mesi, se uguali confronta i giorni. Restituisce un intero
negativo se \passthrough{\lstinline!this < other!}, zero se uguali,
positivo se \passthrough{\lstinline!this > other!}.

\begin{center}\rule{0.5\linewidth}{0.5pt}\end{center}

\subsection{\texorpdfstring{L'Astrazione con Classe Astratta:
\href{file:///home/wally/Documenti/GitHub/Programmazione-II/25-26/Teoria-e-Esercitazioni/Codice-20260902/Lezione12/MavenDate/src/main/java/it/oop/core/FormattedDate.java}{\texttt{FormattedDate.java}}}{L'Astrazione con Classe Astratta: FormattedDate.java}}\label{chap11_lastrazione-con-classe-astratta-formatteddate.java}

Centralizza i dati di formattazione evitando duplicazioni tra
\passthrough{\lstinline!ItalianDate!} e
\passthrough{\lstinline!AmericanDate!}:

\begin{lstlisting}[language=Java]
public abstract class FormattedDate extends Date {
    protected final String format;
    protected final String[] months;

    public FormattedDate(int day, int month, int year, String format, String[] months) {
        super(day, month, year);
        this.format = format;
        this.months = months;
    }

    // Metodi final: le sottoclassi non possono sovrascriverli (Template Method Pattern)
    public final String printFormat() { return format; }
    public final String getMonthAsString() { return months[getMonth() - 1]; }

    // Metodo astratto: DEVE essere implementato dalle sottoclassi concrete
    public abstract String prettyPrint();
}
\end{lstlisting}

\begin{center}\rule{0.5\linewidth}{0.5pt}\end{center}

\subsection{\texorpdfstring{Ereditarietà Multipla di Tipo:
\href{file:///home/wally/Documenti/GitHub/Programmazione-II/25-26/Teoria-e-Esercitazioni/Codice-20260902/Lezione12/MavenDate/src/main/java/it/oop/core/TimeStamp.java}{\texttt{TimeStamp.java}}}{Ereditarietà Multipla di Tipo: TimeStamp.java}}\label{chap11_ereditarietuxe0-multipla-di-tipo-timestamp.java}

Mostra come combinare l'estensione di una classe concreta con
l'implementazione di un'interfaccia:

\begin{lstlisting}[language=Java]
public class TimeStamp extends Date implements Time {
    protected final int seconds;
    protected final int minutes;
    protected final int hours;

    public TimeStamp(int seconds, int minutes, int hours, int day, int month, int year) {
        super(day, month, year);
        this.seconds = seconds;
        this.minutes = minutes;
        this.hours = hours;
    }

    @Override
    public String toString() {
        // Riusa il toString() della superclasse aggiungendo l'orario con zero-padding a due cifre:
        return String.format("%s[%02d:%02d:%02d]", super.toString(), hours, minutes, seconds);
    }

    @Override
    public boolean equals(Object other) {
        // 1. Verifica prima l'uguaglianza della data con super.equals():
        if (super.equals(other) == false) return false;
        if (!(other instanceof TimeStamp)) return false;
        TimeStamp otherAsTimeStamp = (TimeStamp) other;
        // 2. Poi verifica l'orario:
        return hours == otherAsTimeStamp.getHours() &&
               minutes == otherAsTimeStamp.getMinutes() &&
               seconds == otherAsTimeStamp.getSeconds();
    }
}
\end{lstlisting}

\begin{center}\rule{0.5\linewidth}{0.5pt}\end{center}

\subsection{\texorpdfstring{Dimostrazione Pratica in
\href{file:///home/wally/Documenti/GitHub/Programmazione-II/25-26/Teoria-e-Esercitazioni/Codice-20260902/Lezione12/MavenDate/src/main/java/it/oop/ui/MainDate.java}{\texttt{MainDate.java}}}{Dimostrazione Pratica in MainDate.java}}\label{chap11_dimostrazione-pratica-in-maindate.java}

\begin{lstlisting}[language=Java]
Time ts1 = new TimeStamp(0, 0, 10, 10, 11, 2025);
System.out.println(ts1.toString()); // Esegue TimeStamp.toString() tramite Late Binding!

// ts1.getDay(); // COMPILE-TIME ERROR: il tipo statico Time non possiede il metodo getDay()!
System.out.println(((Date) ts1).getDay()); // OK: cast esplicito a Date!

Date[] arr = { ts2, itDate, new Date(9, 9, 2025) };
Arrays.sort(arr); // Ordina l'array eterogeneo grazie a compareTo()!
System.out.println(Arrays.toString(arr));
\end{lstlisting}

\textbf{Output a video prodotto dall'ordinamento polimorfico}:

\begin{lstlisting}
[y2025m9d9, 3/11/2025, y2025m11d10[10:00:00]]
\end{lstlisting}

\begin{itemize}
\tightlist
\item
  L'ordinamento colloca correttamente prima il 9 settembre, poi il 3
  novembre (\passthrough{\lstinline!itDate!}) e infine il 10 novembre
  (\passthrough{\lstinline!TimeStamp!}).
\item
  In fase di stampa, ciascun elemento viene formattato dal
  \textbf{proprio} metodo \passthrough{\lstinline!toString()!} dinamico.
\end{itemize}

\begin{center}\rule{0.5\linewidth}{0.5pt}\end{center}

\begin{studynote}[Nota di Studio] Con la \textbf{Lezione 12} abbiamo analizzato la gerarchia
completa di ereditarietà, l'overriding con
\passthrough{\lstinline!super!}, il late binding, la classe astratta
\passthrough{\lstinline!FormattedDate!} e l'ordinamento con
\passthrough{\lstinline!Comparable!}.

Il prossimo blocco è la \textbf{Lezione 13 / Slide T13: \emph{Abstract
Classes and Interfaces}}: - Differenze ontologiche e strutturali tra
\textbf{Classi Astratte} e \textbf{Interfacce}. - Metodi
\passthrough{\lstinline!default!} e metodi
\passthrough{\lstinline!static!} nelle interfacce (Java 8+). -
Ereditarietà multipla tramite interfacce e risoluzione dei conflitti
(\emph{Diamond Problem} sui metodi default). - Evoluzione di
\passthrough{\lstinline!MavenDate!} in
\passthrough{\lstinline!Lezione13!}: - La classe
\passthrough{\lstinline!Date!} viene ripensata per gestire il calendario
completo. - Analisi delle nuove interfacce e classi di supporto.
\end{studynote}

\textbf{Dammi conferma per aprire la Lezione 13 / T13 e analizzare a
fondo Classi Astratte e Interfacce!}

\newpage
\chapter[T13 --- Classi Astratte ed Interfacce]{T13 --- Classi Astratte ed Interfacce}
\label{chap:t13}

{\small\sffamily\color{accentblue!80!black}\textit{Contratti Comportamentali, Metodi Astratti, Default e Static Methods, Builder Pattern e Lambda}}

\vspace{0.8em}

In questo blocco analizziamo a fondo i meccanismi di astrazione pura di
Java: le \textbf{Classi Astratte}, le \textbf{Interfacce}, il design
pattern \textbf{Template Method} e le interfacce cardine della libreria
standard (\passthrough{\lstinline!Comparable!},
\passthrough{\lstinline!Iterable!}, \passthrough{\lstinline!Iterator!})
(\href{file:///home/wally/Documenti/GitHub/Programmazione-II/25-26/Teoria-e-Esercitazioni/Slides-20260902/T13\%20-\%20Abstract\%20Classes\%20and\%20Interfaces.pdf}{T13
- Abstract Classes and Interfaces.pdf}). Nel codice di
\href{file:///home/wally/Documenti/GitHub/Programmazione-II/25-26/Teoria-e-Esercitazioni/Codice-20260902/Lezione13/MavenDate}{\passthrough{\lstinline!Lezione13/MavenDate!}},
il docente compie un'evoluzione straordinaria introducendo: 1. La classe
statica annidata
\textbf{\href{file:///home/wally/Documenti/GitHub/Programmazione-II/25-26/Teoria-e-Esercitazioni/Codice-20260902/Lezione13/MavenDate/src/main/java/it/oop/core/Date.java}{\passthrough{\lstinline!Date.Builder!}}}
per la costruzione controllata degli oggetti. 2. L'interfaccia
funzionale
\textbf{\href{file:///home/wally/Documenti/GitHub/Programmazione-II/25-26/Teoria-e-Esercitazioni/Codice-20260902/Lezione13/MavenDate/src/main/java/it/oop/core/FormattedDateConverter.java}{\passthrough{\lstinline!FormattedDateConverter!}}}.
3. L'uso congiunto in
\href{file:///home/wally/Documenti/GitHub/Programmazione-II/25-26/Teoria-e-Esercitazioni/Codice-20260902/Lezione13/MavenDate/src/main/java/it/oop/ui/MainDate.java}{\passthrough{\lstinline!MainDate.java!}}
di \textbf{Classi Anonime}
(\passthrough{\lstinline!new Time() \{ ... \}!}) e \textbf{Espressioni
Lambda} (\passthrough{\lstinline!d -> new AmericanDate(...)!}).

\section{Teoria: Concetti Teorici dalle Slide (Slide
T13)}\label{chap12_teoria-concetti-teorici-dalle-slide-slide-t13}

\subsection{Classi Astratte e Template Method Pattern (Slide
1--10)}\label{chap12_classi-astratte-e-template-method-pattern-slide-110}

\begin{itemize}
\tightlist
\item
  \textbf{Metodo Astratto}:

  \begin{itemize}
  \item
    Un metodo dichiarato con la parola chiave
    \passthrough{\lstinline!abstract!} e \textbf{privo di corpo}
    (termina con il punto e virgola \passthrough{\lstinline!;!} anziché
    con le parentesi graffe \passthrough{\lstinline!\{ ... \}!}):

\begin{lstlisting}[language=Java]
public abstract String prettyPrint();
\end{lstlisting}
  \item
    Definisce una ``firma obbligatoria'': impone alle sottoclassi
    concrete il compito di implementarne la logica.
  \end{itemize}
\item
  \textbf{Classe Astratta (\passthrough{\lstinline!abstract class!})}:

  \begin{itemize}
  \tightlist
  \item
    Se una classe contiene anche un solo metodo astratto, \textbf{deve
    obbligatoriamente essere dichiarata
    \passthrough{\lstinline!abstract!}}.
  \item
    {[}Attenzione{]} \textbf{Regola aurea d'esame}: Una classe astratta
    \textbf{non può MAI essere istanziata direttamente}
    (\passthrough{\lstinline!new FormattedDate(...)!} genera un
    \textbf{Compile-Time Error}!).
  \item
    Può contenere costruttori (invocabili dalle sottoclassi tramite
    \passthrough{\lstinline!super(...)!}), campi di istanza
    (\passthrough{\lstinline!protected!}/\passthrough{\lstinline!private!}),
    metodi concreti e metodi astratti.
  \item
    Una sottoclasse che estende una classe astratta deve:

    \begin{itemize}
    \tightlist
    \item
      O implementare \textbf{tutti} i metodi astratti ereditati.
    \item
      Oppure essere dichiarata essa stessa
      \passthrough{\lstinline!abstract!}.
    \end{itemize}
  \end{itemize}
\item
  \textbf{Il Template Method Pattern}:

  \begin{itemize}
  \tightlist
  \item
    Pattern comportamentale in cui la superclasse definisce lo scheletro
    immutabile di un algoritmo all'interno di un metodo concreto marcato
    \textbf{\passthrough{\lstinline!final!}} (il template), mentre
    delega i singoli passi variabili o dipendenti dal contesto a metodi
    \passthrough{\lstinline!abstract!} implementati dalle sottoclassi.
  \item
    Nel nostro progetto: \passthrough{\lstinline!printFormat()!} e
    \passthrough{\lstinline!getMonthAsString()!} sono
    \passthrough{\lstinline!final!} in
    \passthrough{\lstinline!FormattedDate!}, mentre
    \passthrough{\lstinline!prettyPrint()!} è
    \passthrough{\lstinline!abstract!}.
  \end{itemize}
\end{itemize}

\begin{center}\rule{0.5\linewidth}{0.5pt}\end{center}

\subsection{\texorpdfstring{Le Interfacce in Java (\texttt{interface})
(Slide
11--17)}{Le Interfacce in Java (interface) (Slide 11--17)}}\label{chap12_le-interfacce-in-java-interface-slide-1117}

\begin{itemize}
\tightlist
\item
  \textbf{Definizione e Filosofia}:

  \begin{itemize}
  \tightlist
  \item
    Un'interfaccia è un \textbf{contratto puro di comportamento}.
    Rappresenta ciò che una classe \emph{sa fare} (\emph{can-do}), non
    la sua identità ontologica.
  \item
    \textbf{Membri di un'interfaccia standard}:

    \begin{itemize}
    \tightlist
    \item
      \textbf{Campi}: sono implicitamente ed esclusivamente
      \textbf{\passthrough{\lstinline!public static final!}} (costanti
      di classe). Non possono esistere variabili d'istanza o stato
      mutabile.
    \item
      \textbf{Metodi}: sono implicitamente
      \textbf{\passthrough{\lstinline!public abstract!}} (non serve
      specificarlo).
    \end{itemize}
  \end{itemize}
\item
  \textbf{Ereditarietà Multipla di Tipo
  (\passthrough{\lstinline!implements!})}:

  \begin{itemize}
  \item
    Una classe può estendere al massimo una sola superclasse, ma può
    implementare \textbf{un numero arbitrario di interfacce separate da
    virgola}:

\begin{lstlisting}[language=Java]
public class TimeStamp extends Date implements Time, Serializable, Cloneable
\end{lstlisting}
  \item
    Questo realizza in Java l'ereditarietà multipla dei tipi senza
    incorrere nei problemi di ambiguità di memoria del C++.
  \end{itemize}
\item
  \textbf{Ereditarietà tra Interfacce}:

  \begin{itemize}
  \tightlist
  \item
    Un'interfaccia può estendere altre interfacce tramite
    \passthrough{\lstinline!extends!} (anche multiple
    contemporaneamente:
    \passthrough{\lstinline!interface C extends A, B!}).
  \end{itemize}
\end{itemize}

\begin{center}\rule{0.5\linewidth}{0.5pt}\end{center}

\subsection{Confronto Sistematico: Classe Astratta vs
Interfaccia}\label{chap12_confronto-sistematico-classe-astratta-vs-interfaccia}

\begin{longtable}[]{@{}
  >{\raggedright\arraybackslash}p{(\linewidth - 4\tabcolsep) * \real{0.3333}}
  >{\raggedright\arraybackslash}p{(\linewidth - 4\tabcolsep) * \real{0.3333}}
  >{\raggedright\arraybackslash}p{(\linewidth - 4\tabcolsep) * \real{0.3333}}@{}}
\toprule\noalign{}
\begin{minipage}[b]{\linewidth}\raggedright
Proprietà
\end{minipage} & \begin{minipage}[b]{\linewidth}\raggedright
Classe Astratta (\passthrough{\lstinline!abstract class!})
\end{minipage} & \begin{minipage}[b]{\linewidth}\raggedright
Interfaccia (\passthrough{\lstinline!interface!})
\end{minipage} \\
\midrule\noalign{}
\endhead
\bottomrule\noalign{}
\endlastfoot
\textbf{Istanziabilità} & ❌ No (\passthrough{\lstinline!new!} vietato)
& ❌ No (\passthrough{\lstinline!new!} vietato) \\
\textbf{Stato d'istanza (campi)} & {[}OK{]} Sì (qualsiasi visibilità) &
❌ No (solo costanti \passthrough{\lstinline!public static final!}) \\
\textbf{Costruttori} & {[}OK{]} Sì (invocabili con
\passthrough{\lstinline!super(...)!}) & ❌ No \\
\textbf{Ereditarietà} & Singola (\passthrough{\lstinline!extends!} una
sola classe) & Multipla (\passthrough{\lstinline!implements!} \(N\)
interfacce) \\
\textbf{Relazione concettuale} & Identità ontologica forte (\emph{IS-A})
& Capacità / Contratto comportamentale (\emph{CAN-DO}) \\
\end{longtable}

\begin{center}\rule{0.5\linewidth}{0.5pt}\end{center}

\subsection{Interfacce Standard Fondamentali della JDK (Slide
18--22)}\label{chap12_interfacce-standard-fondamentali-della-jdk-slide-1822}

\begin{enumerate}
\def\labelenumi{\arabic{enumi}.}
\tightlist
\item
  \textbf{\passthrough{\lstinline!java.lang.Comparable<T>!}}:

  \begin{itemize}
  \tightlist
  \item
    Definisce il metodo
    \passthrough{\lstinline!int compareTo(T other)!}.
  \item
    Permette agli algoritmi di ordinamento
    (\passthrough{\lstinline!Arrays.sort()!},
    \passthrough{\lstinline!Collections.sort()!}) e agli insiemi
    ordinati (\passthrough{\lstinline!TreeSet!},
    \passthrough{\lstinline!TreeMap!}) di ordinare gli oggetti secondo
    il loro ``ordine naturale''.
  \end{itemize}
\item
  \textbf{\passthrough{\lstinline!java.lang.Iterable<T>!} e
  \passthrough{\lstinline!java.util.Iterator<T>!}}:

  \begin{itemize}
  \tightlist
  \item
    \textbf{\passthrough{\lstinline!Iterable<T>!}}: espone il metodo
    \passthrough{\lstinline!Iterator<T> iterator()!}.
  \item
    \textbf{Regola cruciale}: Qualsiasi classe che implementi
    \passthrough{\lstinline!Iterable!} può essere utilizzata come
    sorgente nel \textbf{ciclo For-Each
    (\passthrough{\lstinline!for (T elem : collection)!})}!
  \item
    \textbf{\passthrough{\lstinline!Iterator<T>!}}: l'oggetto cursore
    che scorre la sequenza:

    \begin{itemize}
    \tightlist
    \item
      \passthrough{\lstinline!boolean hasNext()!}: verifica se vi sono
      ulteriori elementi.
    \item
      \passthrough{\lstinline!T next()!}: restituisce il prossimo
      elemento avanzando il cursore.
    \item
      \passthrough{\lstinline!default void remove()!}: rimuove l'ultimo
      elemento restituito.
    \end{itemize}
  \end{itemize}
\end{enumerate}

\begin{center}\rule{0.5\linewidth}{0.5pt}\end{center}

\section{\texorpdfstring{Codice: Evoluzione del Codice:
\href{file:///home/wally/Documenti/GitHub/Programmazione-II/25-26/Teoria-e-Esercitazioni/Codice-20260902/Lezione13/MavenDate}{\texttt{Lezione13/MavenDate}}}{Codice: Evoluzione del Codice: Lezione13/MavenDate}}\label{chap12_codice-evoluzione-del-codice-lezione13mavendate}

\subsection{\texorpdfstring{Il Design Pattern Builder con Classe Statica
Annidata
(\href{file:///home/wally/Documenti/GitHub/Programmazione-II/25-26/Teoria-e-Esercitazioni/Codice-20260902/Lezione13/MavenDate/src/main/java/it/oop/core/Date.java}{\texttt{Date.java}})}{Il Design Pattern Builder con Classe Statica Annidata (Date.java)}}\label{chap12_il-design-pattern-builder-con-classe-statica-annidata-date.java}

All'interno di \passthrough{\lstinline!Date.java!}, il docente definisce
una \textbf{Static Nested Class} per costruire date in modo protetto:

\begin{lstlisting}[language=Java]
public class Date implements Comparable {
    ...
    public static class Builder {
        private final int year;

        public Builder(int year) {
            if (year > 0)
                this.year = year;
            else
                this.year = 1970; // Anno di default sicuro
        }

        public Date build(int day, int month) {
            // Se i parametri sono illegali, ricade sulla data sicura 1/1/year!
            if (month < 1 || month > 12 || day < 1 || day > daysPerMonth(month))
                return new Date(1, 1, year);
            return new Date(day, month, year);
        }
    }
}
\end{lstlisting}

\emph{Vantaggio del Builder}: Incapsula e pre-configura l'anno; se
l'utente fornisce parametri non validi
(\passthrough{\lstinline!month = -1!}), il costruttore non fallisce a
video ma restituisce una data valida di fallback
(\passthrough{\lstinline!01/01/year!}).

\begin{center}\rule{0.5\linewidth}{0.5pt}\end{center}

\subsection{\texorpdfstring{L'Interfaccia Funzionale:
\href{file:///home/wally/Documenti/GitHub/Programmazione-II/25-26/Teoria-e-Esercitazioni/Codice-20260902/Lezione13/MavenDate/src/main/java/it/oop/core/FormattedDateConverter.java}{\texttt{FormattedDateConverter.java}}}{L'Interfaccia Funzionale: FormattedDateConverter.java}}\label{chap12_linterfaccia-funzionale-formatteddateconverter.java}

\begin{lstlisting}[language=Java]
package it.oop.core;

public interface FormattedDateConverter {
    FormattedDate convert(FormattedDate date);
}
\end{lstlisting}

Questa interfaccia dichiara \textbf{un solo metodo astratto} (SAM:
\emph{Single Abstract Method}). In Java è a tutti gli effetti
un'\textbf{Interfaccia Funzionale}, target ideale per le espressioni
Lambda!

\begin{center}\rule{0.5\linewidth}{0.5pt}\end{center}

\subsection{\texorpdfstring{Sintesi Pratica in
\href{file:///home/wally/Documenti/GitHub/Programmazione-II/25-26/Teoria-e-Esercitazioni/Codice-20260902/Lezione13/MavenDate/src/main/java/it/oop/ui/MainDate.java}{\texttt{MainDate.java}}}{Sintesi Pratica in MainDate.java}}\label{chap12_sintesi-pratica-in-maindate.java}

Questo file è una miniera di concetti d'esame avanzati:

\begin{lstlisting}[language=Java]
package it.oop.ui;

import it.oop.core.*;
import it.oop.core.Date;

public class MainDate {
    public static void main(String[] args) {
        // 1. Uso della Static Nested Class Builder:
        Date.Builder dateBuilder = new Date.Builder(2025);
        Date d1 = dateBuilder.build(10, 9);  // Giorno 10, Mese 9 -> Valida!
        Date d2 = dateBuilder.build(10, -1); // Mese -1 illegale -> Fallback su 1/1/2025!
        System.out.println(d1.toString());   // y2025m9d10
        System.out.println(d2.toString());   // y2025m1d1

        // 2. Classe Anonima (Anonymous Inner Class):
        // Implementa al volo l'interfaccia Time senza creare un file .java separato!
        Time init = new Time() {
            @Override
            public int getHours() { return 0; }
            @Override
            public int getMinutes() { return 0; }
            @Override
            public int getSeconds() { return 0; }
            @Override
            public String toString() {
                return String.format("%02d:%02d:%02d", getHours(), getMinutes(), getSeconds());
            }
        };
        System.out.println(init.toString()); // Stampa 00:00:00

        // 3. Espressione Lambda (Sintassi compatta per Interfaccia Funzionale):
        FormattedDateConverter toAmerican =
                d -> new AmericanDate(d.getDay(), d.getMonth(), d.getYear());

        // Test di conversione polimorfica:
        System.out.println(toAmerican.convert(new ItalianDate(11, 11, 2025)) instanceof AmericanDate);
        // Stampa TRUE!
    }
}
\end{lstlisting}

\begin{center}\rule{0.5\linewidth}{0.5pt}\end{center}

\section{{[}Test{]} Output di Esecuzione Verificato con
Maven}\label{chap12_test-output-di-esecuzione-verificato-con-maven}

Eseguendo
\passthrough{\lstinline!mvn exec:java -Dexec.mainClass="it.oop.ui.MainDate"!}:

\begin{lstlisting}
y2025m9d10
y2025m1d1
00:00:00
true
\end{lstlisting}

\begin{center}\rule{0.5\linewidth}{0.5pt}\end{center}

\begin{studynote}[Nota di Studio] Con la \textbf{Lezione 13} abbiamo chiuso il cerchio su
Classi Astratte, Interfacce, Template Method e abbiamo anticipato le
classi annidate e le lambda.

Il prossimo blocco è la \textbf{Lezione 14 / Slide T14: \emph{Nested and
Anonymous Classes}}: - Tassonomia completa delle classi interne:
\textbf{Static Nested Classes}, \textbf{Inner Member Classes} (non
statiche), \textbf{Local Classes} e \textbf{Anonymous Classes}. -
Accesso allo stato della classe contenitore (\emph{Enclosing instance})
e la sintassi speciale \passthrough{\lstinline!EnclosingClass.this!}. -
Cattura delle variabili locali nelle classi locali/anonime e il vincolo
di essere \textbf{\passthrough{\lstinline!effectively final!}}. -
Evoluzione del codice in \passthrough{\lstinline!Lezione14!}:
consolidamento delle classi interne e preparativi per i Generics.
\end{studynote}

\textbf{Dammi conferma per aprire la Lezione 14 / T14 e continuare la
revisione!}

\newpage
\chapter[T14 --- Classi Annidate, Interne, Locali ed]{T14 --- Classi Annidate, Interne, Locali ed Anonime}
\label{chap:t14}

{\small\sffamily\color{accentblue!80!black}\textit{Static Nested vs Inner Classes, Enclosing Pointer, Shadowing, Effectively Final e Classi Anonime}}

\vspace{0.8em}

\section{Analisi Teorica Approfondita (Slide
T14)}\label{chap13_analisi-teorica-approfondita-slide-t14}

La lezione 14 esplora la gerarchia delle classi definite all'interno di
altre classi (meccanismo introdotto fin da Java 1.1), le classi anonime,
l'evoluzione verso le \textbf{espressioni Lambda} (Java 8) e la feature
dei \textbf{Varargs}.

\needspace{7cm}
\begin{asciibox}
                 ┌───────────────────────────┐
                 │  Classi Annidate in Java  │
                 └─────────────┬─────────────┘
                               │
            ┌──────────────────┴──────────────────┐
            ▼                                     ▼
 ┌───────────────────────┐             ┌───────────────────────┐
 │ Static Nested Classes │             │     Inner Classes     │
 │   (static class ...)  │             │     (non-static)      │
 └───────────────────────┘             └───────────┬───────────┘
                                                   │
                  ┌────────────────────────────────┼─────────────────┐
                  ▼                                ▼                 ▼
       ┌──────────────────┐ ┌──────────────────┐ ┌──────────────────┐
       │Member Inner Class│ │   Local Class    │ │ Anonymous Class  │
       │(a livello classe)│ │(dentro un metodo)│ │(senza nome/corpo)│
       └──────────────────┘ └──────────────────┘ └────────┬─────────┘
                                                           │ (SAM Interface)
                                                           ▼
                                                 ┌──────────────────┐
                                                 │Lambda Expression │
                                                 │  (arg -> body)   │
                                                 └──────────────────┘
\end{asciibox}

\begin{center}\rule{0.5\linewidth}{0.5pt}\end{center}

\subsection{\texorpdfstring{Static Nested Classes
(\texttt{static\ class\ Nested})}{Static Nested Classes (static class Nested)}}\label{chap13_static-nested-classes-static-class-nested}

Una classe dichiarata con il modificatore
\passthrough{\lstinline!static!} all'interno di un'altra classe
contenitore (\emph{enclosing class}):

\begin{lstlisting}[language=Java]
package pkg;
class Outer {
    static class Nested {
        // ...
    }
}
\end{lstlisting}

\begin{itemize}
\item
  \textbf{Natura}: È a tutti gli effetti una classe di primo livello
  (\emph{top-level}), ma inserita nel \emph{namespace} della classe che
  la ospita.
\item
  \textbf{Nome completo (Fully Qualified Name)}:
  \passthrough{\lstinline!pkg.Outer.Nested!} (a livello di bytecode su
  disco viene generato il file
  \passthrough{\lstinline!Outer$Nested.class!}).
\item
  \textbf{Accesso ai membri}:

  \begin{itemize}
  \tightlist
  \item
    Trattandosi di un membro statico, \textbf{ha accesso a tutti i
    membri statici} della classe contenitore, \textbf{compresi quelli
    privati (\passthrough{\lstinline!private!})}!
  \item
    \textbf{NON ha accesso} ai membri di istanza (campi o metodi non
    statici) di \passthrough{\lstinline!Outer!}, poiché non esiste alcun
    riferimento implicito ad un'istanza di
    \passthrough{\lstinline!Outer!}.
  \end{itemize}
\item
  \textbf{Istanziazione}: Può essere istanziata in modo totalmente
  indipendente dall'esistenza di un oggetto
  \passthrough{\lstinline!Outer!}:

\begin{lstlisting}[language=Java]
Outer.Nested nestedObject = new Outer.Nested();
\end{lstlisting}
\item
  \textbf{Caso d'uso tipico}: Il pattern \textbf{Builder} (visto in
  \href{file:///home/wally/Documenti/GitHub/Programmazione-II/25-26/Teoria-e-Esercitazioni/Codice-20260902/Lezione14/MavenDate/src/main/java/it/oop/core/Date.java\#L85-L106}{Date.java}),
  classi di utilità, nodi interni di strutture dati (es.
  \passthrough{\lstinline!static class Node<E>!} in un albero o grafo).
\end{itemize}

\begin{center}\rule{0.5\linewidth}{0.5pt}\end{center}

\subsection{\texorpdfstring{Member Inner Classes (\texttt{class\ Inner}
non-statica)}{Member Inner Classes (class Inner non-statica)}}\label{chap13_member-inner-classes-class-inner-non-statica}

Una classe dichiarata all'interno di un'altra classe \textbf{senza} la
parola chiave \passthrough{\lstinline!static!}:

\begin{lstlisting}[language=Java]
package pkg;
class Outer {
    class Inner {
        // ...
    }
}
\end{lstlisting}

\begin{itemize}
\item
  \textbf{Natura}: Un'istanza di \passthrough{\lstinline!Inner!}
  \textbf{non può esistere senza un'istanza associata di
  \passthrough{\lstinline!Outer!}}.
\item
  \textbf{Riferimento implicito (\passthrough{\lstinline!Outer.this!} o
  \passthrough{\lstinline!this$0!})}:

  \begin{itemize}
  \tightlist
  \item
    Ogni oggetto \passthrough{\lstinline!Inner!} memorizza internamente
    un puntatore sintetico nascosto all'oggetto
    \passthrough{\lstinline!Outer!} che l'ha generato.
  \item
    Tramite questo puntatore, \passthrough{\lstinline!Inner!} ha accesso
    a \textbf{tutti} i membri di \passthrough{\lstinline!Outer!}, sia
    statici che di istanza, inclusi quelli
    \passthrough{\lstinline!private!}!
  \end{itemize}
\item
  \textbf{Sintassi di istanziazione da codice esterno}:

\begin{lstlisting}[language=Java]
Outer outerObject = new Outer();
Outer.Inner innerObject = outerObject.new Inner(); // Sintassi specifica con outerObject.new
\end{lstlisting}
\end{itemize}

\subsubsection{\texorpdfstring{Confronto: \texttt{Nested} vs
\texttt{Inner} (Slide
7)}{Confronto: Nested vs Inner (Slide 7)}}\label{chap13_confronto-nested-vs-inner-slide-7}

\begin{lstlisting}[language=Java]
class Outer {
    private static int I = 5;
    private int i;

    static class Nested {
        // NON ha accesso al campo di istanza i!
        static void print() { System.out.println(I); } // OK: accede a membro statico privato I
    }

    class Inner {
        // Ha accesso sia a I che al campo di istanza i della sua Outer associata
        void print() { System.out.println(i); }
    }

    public void outerMethod() {
        Nested.print();
        Inner innerObject = new Inner(); // Qui 'this.new Inner()' è implicito
        innerObject.print();
    }
}
\end{lstlisting}

\begin{center}\rule{0.5\linewidth}{0.5pt}\end{center}

\subsection{Local Classes (Classi Locali a un
Metodo)}\label{chap13_local-classes-classi-locali-a-un-metodo}

Una classe definita all'interno del corpo di un blocco di codice
(tipicamente un metodo):

\begin{lstlisting}[language=Java]
public void outerMethod() {
    final int j = 1;

    class Local { // NESSUN modificatore di visibilità ammesso (no public, private, protected)
        int addOne() { return j + 1; }
    }

    Local localObject = new Local();
    System.out.println(localObject.addOne()); // Stampa 2
}
\end{lstlisting}

\begin{itemize}
\tightlist
\item
  \textbf{Regole di visibilità}: La classe è visibile e istanziabile
  \textbf{esclusivamente all'interno del blocco/metodo} in cui è
  dichiarata.
\item
  \textbf{Cattura delle variabili locali (\emph{Variable Capture} \&
  \emph{Closures})}:

  \begin{itemize}
  \tightlist
  \item
    Può accedere alle variabili locali e ai parametri del metodo
    contenitore \textbf{solo se sono dichiarati
    \passthrough{\lstinline!final!} oppure sono \emph{effectively
    final}} (ovvero il loro valore non viene mai modificato dopo
    l'inizializzazione).
  \item
    \textbf{Perché questo vincolo severo? (Domanda classica d'esame)}:

    \begin{itemize}
    \tightlist
    \item
      Le variabili locali risiedono nello \textbf{Stack Frame} del
      metodo e vengono distrutte non appena il metodo termina
      (\passthrough{\lstinline!return!}).
    \item
      L'oggetto della classe locale
      (\passthrough{\lstinline!localObject!}) risiede nell'\textbf{Heap}
      e potrebbe sopravvivere ben oltre la fine del metodo (ad esempio
      se restituito come interfaccia o registrato come listener).
    \item
      Per permettere all'oggetto di accedere a
      \passthrough{\lstinline!j!}, il compilatore Java crea un campo
      sintetico nascosto dentro \passthrough{\lstinline!Local!} e vi
      \textbf{copia} il valore di \passthrough{\lstinline!j!} al momento
      della creazione.
    \item
      Se \passthrough{\lstinline!j!} potesse essere riassegnato
      (\passthrough{\lstinline!j++!}), il valore nello stack e la copia
      nell'heap diverrebbero disallineati, rompendo la coerenza del
      linguaggio. Rendendo la variabile immutabile
      (\passthrough{\lstinline!final!}), la copia rimane identica e
      coerente per sempre.
    \end{itemize}
  \end{itemize}
\end{itemize}

\begin{center}\rule{0.5\linewidth}{0.5pt}\end{center}

\subsection{\texorpdfstring{Il Meme della Slide 10: \emph{Objects inside
objects}}{Il Meme della Slide 10: Objects inside objects}}\label{chap13_il-meme-della-slide-10-objects-inside-objects}

Nella Slide 10, il prof. Pasqua inserisce il celebre meme di Xzibit
(\emph{Pimp My Ride}): \textgreater{} \emph{``YO DAWG, I HEAR YOU LIKE
OBJECT ORIENTED PROGRAMMING SO I NESTED A CLASS INSIDE YOUR CLASS SO YOU
CAN CREATE OBJECTS WHILE YOU CREATE OBJECTS''}

A sottolineare con ironia la tentazione (e il potere) di creare
gerarchie di classi dentro classi per incapsulare comportamenti
strettamente accoppiati.

\begin{center}\rule{0.5\linewidth}{0.5pt}\end{center}

\subsection{\texorpdfstring{Motivazioni Architetturali \& Il Caso di
Studio \emph{Iterator Pattern} (Slide
11-12)}{Motivazioni Architetturali \& Il Caso di Studio Iterator Pattern (Slide 11-12)}}\label{chap13_motivazioni-architetturali-il-caso-di-studio-iterator-pattern-slide-11-12}

Perché annidare le classi? 1. \textbf{Migliore organizzazione del
codice}: collocare classi di supporto logico all'interno del namespace
della sola classe che le utilizza. 2. \textbf{Politica di visibilità
privilegiata (\emph{Privileged Access})}: le classi annidate possono
leggere lo stato privato della classe ospitante senza costringere a
esporre metodi \passthrough{\lstinline!public!} o
\passthrough{\lstinline!package-private!} pericolosi verso l'esterno. 3.
\textbf{Meno complessità nel codebase}: evita la proliferazione di
decine di minuscoli file \passthrough{\lstinline!.java!} separati.

\subsubsection{\texorpdfstring{L'Esempio del pattern \texttt{Iterable} /
\texttt{Iterator} (Slide
12)}{L'Esempio del pattern Iterable / Iterator (Slide 12)}}\label{chap13_lesempio-del-pattern-iterable-iterator-slide-12}

\begin{itemize}
\item
  \textbf{Senza annidamento (Due classi separate)}:
  \passthrough{\lstinline!CRIterator!} deve essere dichiarata
  all'esterno, riceve l'array \passthrough{\lstinline!Student[] arr!}
  nel costruttore e necessita di visibilità speciale:

\begin{lstlisting}[language=Java]
class ClassRoom implements Iterable {
    private Student[] students;
    @Override
    public Iterator iterator() { return new CRIterator(students); }
}
class CRIterator implements Iterator {
    private Student[] arr;
    private int next = 0;
    CRIterator(Student[] arr) { this.arr = arr; }
    @Override public boolean hasNext() { return next < arr.length; }
    @Override public Object next() { return new Student(arr[next++]); }
}
\end{lstlisting}
\item
  \textbf{Con Inner Class (Soluzione idiomatica Java)}:
  \passthrough{\lstinline!CRIterator!} è una
  \passthrough{\lstinline!Inner Class!} dentro
  \passthrough{\lstinline!ClassRoom!}. Non serve memorizzare un array
  duplicato: accede direttamente al campo privato
  \passthrough{\lstinline!students!} di
  \passthrough{\lstinline!ClassRoom!}!

\begin{lstlisting}[language=Java]
class ClassRoom implements Iterable {
    private Student[] students;
    @Override
    public Iterator iterator() { return new CRIterator(); }

    class CRIterator implements Iterator {
        private int next = 0;
        @Override public boolean hasNext() { return next < students.length; }
        @Override public Object next() { return new Student(students[next++]); }
    }
}
\end{lstlisting}
\end{itemize}

\begin{center}\rule{0.5\linewidth}{0.5pt}\end{center}

\subsection{Anonymous Classes (Classi Anonime, Slide
13-16)}\label{chap13_anonymous-classes-classi-anonime-slide-13-16}

Una classe locale \textbf{priva di nome}, definita e istanziata
contestualmente all'interno di una singola espressione. *
\textbf{Vincolo}: Può essere creata solo estendendo una classe esistente
o implementando un'interfaccia: ```java // 1. Estendendo una classe
esistente: Person sam = new Person(``Sam'') \{ @Override public String
toString() \{ return ``It's Sam!''; \} \};

// 2. Implementando un'interfaccia: Time init = new Time() \{ @Override
public int getHours() \{ return 0; \} @Override public int getMinutes()
\{ return 0; \} @Override public int getSeconds() \{ return 0; \} \};
\passthrough{\lstinline"* **Regola aurea sui Costruttori**: **Una classe anonima NON può dichiarare alcun costruttore esplicito!** (Poiché i costruttori devono avere per definizione lo stesso nome della classe, e una classe anonima non ha nome). Eventuali inizializzazioni devono sfruttare gli *instance initializer blocks* `\{ ... \}` o i parametri passati al costruttore della superclasse (`new SuperClass(arg1, arg2) \{ ... \}`). * **Espressioni**: Essendo espressioni, devono terminare con un punto e virgola `;` (se assegnate a una variabile) o possono essere passate direttamente come parametri attuali di un metodo:"}java
tasks.add(new Runnable() \{ @Override public void run() \{
System.out.println(``Running!''); \} \}); ```

\begin{center}\rule{0.5\linewidth}{0.5pt}\end{center}

\subsection{Lambda Expressions (Slide
17-19)}\label{chap13_lambda-expressions-slide-17-19}

Introdotte in Java 8, sono una notazione sintattica ultra-compatta per
implementare \textbf{interfacce funzionali} (interfacce con un
\textbf{unico metodo astratto}, note come SAM --- \emph{Single Abstract
Method}).

\begin{lstlisting}
(parametri) -> corpo
\end{lstlisting}

\begin{itemize}
\tightlist
\item
  \textbf{Parametri}:

  \begin{itemize}
  \tightlist
  \item
    Nessun parametro: \passthrough{\lstinline!() -> 42!}
  \item
    Un parametro: \passthrough{\lstinline!x -> x + 1!} (le parentesi
    tonde sono opzionali se non si specifica il tipo esplicito)
  \item
    Due o più parametri: \passthrough{\lstinline!(x, y) -> x + y!}
  \end{itemize}
\item
  \textbf{Corpo}:

  \begin{itemize}
  \tightlist
  \item
    Espressione singola: \passthrough{\lstinline!param -> param - 1!}
    (restituzione implicita del valore calcolato, niente
    \passthrough{\lstinline!return!} né graffe)
  \item
    Blocco di codice:
    \passthrough{\lstinline!param -> \{ int res = param - 1; return res; \}!}
    (richiede graffe esplicite e \passthrough{\lstinline!return!})
  \end{itemize}
\item
  \textbf{Inferenza di tipo (\emph{Type Inference})}: Il compilatore
  Java deduce i tipi dei parametri e del valore di ritorno dal contesto
  d'uso (\emph{target typing} dell'interfaccia funzionale).
\item
  \textbf{Differenza tra Classi Anonime e Lambda}:

  \begin{itemize}
  \tightlist
  \item
    La classe anonima può implementare interfacce con più metodi
    astratti o estendere classi concrete/astratte; la lambda può
    implementare \textbf{solo interfacce funzionali SAM}.
  \item
    La classe anonima crea un nuovo scope per
    \passthrough{\lstinline!this!} (che punta all'istanza anonima
    stessa); nella lambda \passthrough{\lstinline!this!} è
    \textbf{lessicale} (punta all'istanza della classe circostante).
  \item
    La classe anonima può dichiarare campi di istanza che mantengono
    stato; la lambda è per natura \emph{stateless}.
  \end{itemize}
\end{itemize}

\begin{center}\rule{0.5\linewidth}{0.5pt}\end{center}

\subsection{\texorpdfstring{Varargs (\texttt{type...\ args}, Slide
20-21)}{Varargs (type... args, Slide 20-21)}}\label{chap13_varargs-type...-args-slide-20-21}

Consente a un metodo di accettare un numero variabile (da zero a molti)
di argomenti dello stesso tipo:

\begin{lstlisting}[language=Java]
static int min(int... values) {
    int res = Integer.MAX_VALUE;
    for (int v : values)
        if (v < res) res = v;
    return res;
}
\end{lstlisting}

\begin{itemize}
\tightlist
\item
  \textbf{Dietro le quinte del compilatore}: Il parametro
  \passthrough{\lstinline!int... values!} viene tradotto esattamente in
  un array \passthrough{\lstinline!int[] values!}. Al momento della
  chiamata:

  \begin{itemize}
  \tightlist
  \item
    \passthrough{\lstinline!min(1, 0, 5)!} \(\rightarrow\) il
    compilatore inietta il codice:
    \passthrough{\lstinline!min(new int[]\{ 1, 0, 5 \})!}.
  \item
    \passthrough{\lstinline!min()!} (senza argomenti) \(\rightarrow\)
    crea un array vuoto \passthrough{\lstinline!min(new int[]\{\})!}.
  \end{itemize}
\item
  \textbf{Regole tassative}:

  \begin{enumerate}
  \def\labelenumi{\arabic{enumi}.}
  \tightlist
  \item
    Ci può essere \textbf{al massimo un} parametro varargs per metodo.
  \item
    Il parametro varargs deve essere tassativamente l'\textbf{ultimo}
    nella lista dei parametri formali
    (\passthrough{\lstinline!void foo(String s, int... nums)!} è valido;
    \passthrough{\lstinline!void foo(int... nums, String s)!} genera
    errore di compilazione).
  \end{enumerate}
\end{itemize}

\begin{center}\rule{0.5\linewidth}{0.5pt}\end{center}

\section{\texorpdfstring{Analisi Dettagliata del Codice
(\texttt{Lezione14/MavenDate})}{Analisi Dettagliata del Codice (Lezione14/MavenDate)}}\label{chap13_analisi-dettagliata-del-codice-lezione14mavendate}

Nella Lezione 14 il progetto
\href{file:///home/wally/Documenti/GitHub/Programmazione-II/25-26/Teoria-e-Esercitazioni/Codice-20260902/Lezione14/MavenDate}{MavenDate}
introduce il primo ponte fondamentale verso i \textbf{Tipi Generici}
(anticipando la teoria formale di Slide T15) e sfrutta le classi
annidate, le classi anonime e le lambda viste nella Lezione 13.

\subsection{\texorpdfstring{Diff di
\href{file:///home/wally/Documenti/GitHub/Programmazione-II/25-26/Teoria-e-Esercitazioni/Codice-20260902/Lezione14/MavenDate/src/main/java/it/oop/core/Date.java\#L1-L80}{Date.java}
(Lezione 13 vs Lezione
14)}{Diff di Date.java (Lezione 13 vs Lezione 14)}}\label{chap13_diff-di-date.java-lezione-13-vs-lezione-14}

Fino alla Lezione 13, \passthrough{\lstinline!Date!} implementava
l'interfaccia non tipizzata (\emph{raw})
\passthrough{\lstinline!Comparable!}:

\begin{lstlisting}
--- Lezione13/Date.java
+++ Lezione14/Date.java
@@ -1,6 +1,6 @@
 package it.oop.core;
 
-public class Date implements Comparable {
+public class Date implements Comparable<Date> {
 
@@ -73,15 +73,14 @@
     @Override
-    public int compareTo(Object other) {
-        Date otherAsDate = (Date) other;
-        int diff = this.year - otherAsDate.getYear();
+    public int compareTo(Date other) {
+        int diff = this.year - other.getYear();
         if (diff != 0)
             return diff;
-        diff = this.month - otherAsDate.getMonth();
+        diff = this.month - other.getMonth();
         if (diff != 0)
             return diff;
-        return this.day - otherAsDate.getDay();
+        return this.day - other.getDay();
     }
\end{lstlisting}

\begin{itemize}
\tightlist
\item
  \textbf{Beneficio immediato del Generics}:

  \begin{enumerate}
  \def\labelenumi{\arabic{enumi}.}
  \tightlist
  \item
    La firma del metodo diventa fortemente tipizzata a compile-time:
    \passthrough{\lstinline!compareTo(Date other)!} anziché
    \passthrough{\lstinline!compareTo(Object other)!}.
  \item
    Viene eliminato il cast esplicito
    \passthrough{\lstinline!(Date) other!}, azzerando il rischio di
    incorrere a runtime in una
    \passthrough{\lstinline!ClassCastException!}.
  \end{enumerate}
\end{itemize}

\begin{center}\rule{0.5\linewidth}{0.5pt}\end{center}

\subsection{\texorpdfstring{La Nuova Classe Generica
\href{file:///home/wally/Documenti/GitHub/Programmazione-II/25-26/Teoria-e-Esercitazioni/Codice-20260902/Lezione14/MavenDate/src/main/java/it/oop/core/Pair.java}{Pair.java}}{La Nuova Classe Generica Pair.java}}\label{chap13_la-nuova-classe-generica-pair.java}

\begin{lstlisting}[language=Java]
package it.oop.core;

public class Pair<F, S> {
    private final F first;
    private final S second;

    public Pair(F first, S second) {
        this.first = first;
        this.second = second;
    }

    public F getFirst() { return first; }
    public S getSecond() { return second; }
}
\end{lstlisting}

\begin{itemize}
\tightlist
\item
  \passthrough{\lstinline!F!} (\emph{First}) e
  \passthrough{\lstinline!S!} (\emph{Second}) sono due \textbf{Type
  Parameters} indipendenti.
\item
  I campi sono \passthrough{\lstinline!private final!}, rendendo la
  coppia immutabile una volta istanziata.
\end{itemize}

\begin{center}\rule{0.5\linewidth}{0.5pt}\end{center}

\subsection{\texorpdfstring{La Classe Generica con Bounded Type
\href{file:///home/wally/Documenti/GitHub/Programmazione-II/25-26/Teoria-e-Esercitazioni/Codice-20260902/Lezione14/MavenDate/src/main/java/it/oop/core/OrderedPair.java}{OrderedPair.java}}{La Classe Generica con Bounded Type OrderedPair.java}}\label{chap13_la-classe-generica-con-bounded-type-orderedpair.java}

\begin{lstlisting}[language=Java]
package it.oop.core;

public class OrderedPair <T extends Comparable<T>> {
    private final T first;
    private final T second;

    public OrderedPair(T first, T second) {
        this.first = first;
        this.second = second;
        if (first.compareTo(second) > 0)
            System.out.println("Pair not ordered!");
    }

    public T getFirst() { return first; }
    public T getSecond() { return second; }
}
\end{lstlisting}

\begin{itemize}
\tightlist
\item
  \textbf{Bounded Type Parameter
  (\passthrough{\lstinline!<T extends Comparable<T>>!})}:

  \begin{itemize}
  \tightlist
  \item
    Il tipo \passthrough{\lstinline!T!} non può essere un qualsiasi tipo
    \passthrough{\lstinline!Object!}: deve tassativamente implementare
    l'interfaccia \passthrough{\lstinline!Comparable<T>!}.
  \item
    Grazie a questo vincolo, il compilatore consente all'interno del
    costruttore di invocare legalmente
    \passthrough{\lstinline!first.compareTo(second)!}.
  \item
    Se si tentasse di istanziare
    \passthrough{\lstinline!OrderedPair<Object>!}, il compilatore
    bloccherebbe l'operazione con un errore prima ancora di eseguire.
  \end{itemize}
\end{itemize}

\begin{center}\rule{0.5\linewidth}{0.5pt}\end{center}

\subsection{\texorpdfstring{Le Sottoclassi Specializzate
\href{file:///home/wally/Documenti/GitHub/Programmazione-II/25-26/Teoria-e-Esercitazioni/Codice-20260902/Lezione14/MavenDate/src/main/java/it/oop/core/DatePair.java}{DatePair.java}
e
\href{file:///home/wally/Documenti/GitHub/Programmazione-II/25-26/Teoria-e-Esercitazioni/Codice-20260902/Lezione14/MavenDate/src/main/java/it/oop/core/DateInterval.java}{DateInterval.java}}{Le Sottoclassi Specializzate DatePair.java e DateInterval.java}}\label{chap13_le-sottoclassi-specializzate-datepair.java-e-dateinterval.java}

\subsubsection{\texorpdfstring{\href{file:///home/wally/Documenti/GitHub/Programmazione-II/25-26/Teoria-e-Esercitazioni/Codice-20260902/Lezione14/MavenDate/src/main/java/it/oop/core/DatePair.java}{DatePair.java}}{DatePair.java}}\label{chap13_datepair.java}

\begin{lstlisting}[language=Java]
package it.oop.core;

public class DatePair extends Pair<Date, Date> {
    public DatePair(Date left, Date right) {
        super(left, right);
    }

    public Date getLeft() { return getFirst(); }
    public Date getRight() { return getSecond(); }

    @Override
    public String toString() {
        Date l = getLeft();
        Date r = getRight();
        return String.format("[%s .. %s]", l.toString(), r.toString());
    }
}
\end{lstlisting}

\begin{itemize}
\tightlist
\item
  Fissa sia \passthrough{\lstinline!F!} che \passthrough{\lstinline!S!}
  al tipo concreto \passthrough{\lstinline!Date!}
  (\passthrough{\lstinline!extends Pair<Date, Date>!}).
\item
  Offre metodi con semantica di dominio più chiara:
  \passthrough{\lstinline!getLeft()!} e
  \passthrough{\lstinline!getRight()!}.
\end{itemize}

\subsubsection{\texorpdfstring{\href{file:///home/wally/Documenti/GitHub/Programmazione-II/25-26/Teoria-e-Esercitazioni/Codice-20260902/Lezione14/MavenDate/src/main/java/it/oop/core/DateInterval.java}{DateInterval.java}}{DateInterval.java}}\label{chap13_dateinterval.java}

\begin{lstlisting}[language=Java]
package it.oop.core;

public class DateInterval extends OrderedPair<Date> {
    public DateInterval(Date left, Date right) {
        super(left, right);
    }

    public Date getLeft() { return getFirst(); }
    public Date getRight() { return getSecond(); }

    @Override
    public String toString() {
        Date l = getLeft();
        Date r = getRight();
        return String.format("[%s .. %s]", l.toString(), r.toString());
    }
}
\end{lstlisting}

\begin{itemize}
\tightlist
\item
  Eredita da \passthrough{\lstinline!OrderedPair<Date>!}: è legale
  perché \passthrough{\lstinline!Date implements Comparable<Date>!}.
\item
  Se le date passate non sono cronologicamente ordinate
  (\passthrough{\lstinline!left > right!}), il costruttore di
  \passthrough{\lstinline!OrderedPair!} emette il messaggio d'avviso
  \passthrough{\lstinline'"Pair not ordered!"'}.
\end{itemize}

\begin{center}\rule{0.5\linewidth}{0.5pt}\end{center}

\subsection{\texorpdfstring{\href{file:///home/wally/Documenti/GitHub/Programmazione-II/25-26/Teoria-e-Esercitazioni/Codice-20260902/Lezione14/MavenDate/src/main/java/it/oop/ui/MainDate.java}{MainDate.java}}{MainDate.java}}\label{chap13_maindate.java}

Raccoglie in un unico punto di ingresso tutte le feature viste nelle
lezioni 13 e 14:

\begin{lstlisting}[language=Java]
package it.oop.ui;

import it.oop.core.*;
import it.oop.core.Date;

public class MainDate {

    public static void main(String[] args) {
        // 1. Static Nested Class (Builder Pattern da Lezione 13)
        Date.Builder dateBuilder = new Date.Builder(2025);
        Date d1 = dateBuilder.build(10,9);
        Date d2 = dateBuilder.build(10,-1);
        System.out.println(d1.toString());
        System.out.println(d2.toString());

        // 2. Anonymous Class che implementa l'interfaccia Time (Lezione 13/14)
        Time init = new Time() {
            @Override public int getHours() { return 0; }
            @Override public int getMinutes() { return 0; }
            @Override public int getSeconds() { return 0; }
            @Override public String toString() {
                return String.format("%02d:%02d:%02d", getHours(), getMinutes(), getSeconds());
            }
        };
        System.out.println(init.toString());

        // 3. Lambda Expression che implementa FormattedDateConverter (Lezione 13/14)
        FormattedDateConverter toAmerican =
                d -> new AmericanDate(d.getDay(), d.getMonth(), d.getYear());
        System.out.println(toAmerican.convert(new ItalianDate(11, 11, 2025)) instanceof AmericanDate);

        // 4. Generics: Pair con Diamond Operator <>
        Pair<String, FormattedDate> event =
                new Pair<>("OOP exam", new ItalianDate(2, 2, 2026));
        System.out.println(event.getFirst() + " on " + event.getSecond().prettyPrint());

        // 5. DatePair e DateInterval
        DatePair dp = new DatePair(new Date(1,2,2025), new Date(1,1,2025));
        System.out.println("date pair: " + dp.toString());

        // Qui scatta l'avviso di ordinamento perché 1/2/2025 è successivo a 1/1/2025!
        DateInterval di = new DateInterval(new Date(1,2,2025), new Date(1,1,2025));

        // 6. Wildcard con Bounded Type (anticipazione Lezione 15)
        Pair<?, ? extends FormattedDate> unknown = event;
        System.out.println(unknown.getFirst().toString() + " on " + unknown.getSecond().prettyPrint());
    }
}
\end{lstlisting}

\begin{center}\rule{0.5\linewidth}{0.5pt}\end{center}

\section{Risultato di Compilazione ed Esecuzione
Reale}\label{chap13_risultato-di-compilazione-ed-esecuzione-reale}

Comando eseguito nel workspace del corso:

\begin{lstlisting}[language=bash]
mvn compile exec:java -Dexec.mainClass="it.oop.ui.MainDate"
\end{lstlisting}

Output ottenuto:

\begin{lstlisting}
y2025m9d10
y2025m1d1
00:00:00
true
OOP exam on 2 febbraio 2026
date pair: [y2025m2d1 .. y2025m1d1]
Pair not ordered!
OOP exam on 2 febbraio 2026
\end{lstlisting}

\subsection{Dettaglio riga per riga
dell'output:}\label{chap13_dettaglio-riga-per-riga-delloutput}

\begin{enumerate}
\def\labelenumi{\arabic{enumi}.}
\tightlist
\item
  \passthrough{\lstinline!y2025m9d10!} \(\rightarrow\) generato da
  \passthrough{\lstinline!d1 = dateBuilder.build(10, 9)!} (10 settembre
  2025).
\item
  \passthrough{\lstinline!y2025m1d1!} \(\rightarrow\) generato da
  \passthrough{\lstinline!d2 = dateBuilder.build(10, -1)!} (mese errato
  fallback su 1/1/2025).
\item
  \passthrough{\lstinline!00:00:00!} \(\rightarrow\) generato
  dall'override di \passthrough{\lstinline!toString()!} della
  \textbf{classe anonima} \passthrough{\lstinline!init!} che implementa
  \passthrough{\lstinline!Time!}.
\item
  \passthrough{\lstinline!true!} \(\rightarrow\) la \textbf{lambda}
  \passthrough{\lstinline!toAmerican!} converte l'istanza di
  \passthrough{\lstinline!ItalianDate!} in
  \passthrough{\lstinline!AmericanDate!}.
\item
  \passthrough{\lstinline!OOP exam on 2 febbraio 2026!} \(\rightarrow\)
  \passthrough{\lstinline!Pair<String, FormattedDate>!} con
  formattazione italiana.
\item
  \passthrough{\lstinline!date pair: [y2025m2d1 .. y2025m1d1]!}
  \(\rightarrow\) \passthrough{\lstinline!DatePair!} stampa le due date
  senza controllo di sequenzialità.
\item
  \passthrough{\lstinline"Pair not ordered!"} \(\rightarrow\) emesso dal
  costruttore di \passthrough{\lstinline!OrderedPair!} dentro
  \passthrough{\lstinline!new DateInterval(1/2/2025, 1/1/2025)!} perché
  \passthrough{\lstinline!1/2/2025 > 1/1/2025!}.
\item
  \passthrough{\lstinline!OOP exam on 2 febbraio 2026!} \(\rightarrow\)
  invocazione polimorfica tramite wildcard generica
  \passthrough{\lstinline!Pair<?, ? extends FormattedDate>!}.
\end{enumerate}

\begin{center}\rule{0.5\linewidth}{0.5pt}\end{center}

Dimmi \textbf{``vai''} per procedere con il \textbf{Blocco 14} (Lezione
15 / Slide T15 --- \emph{Generic Types}, Type Erasure, Wildcards, PECS e
diff in \passthrough{\lstinline!Lezione15/MavenDate!}).

\newpage
\chapter[T15 --- Tipi Generici, Wildcard e Type Eras]{T15 --- Tipi Generici, Wildcard e Type Erasure}
\label{chap:t15}

{\small\sffamily\color{accentblue!80!black}\textit{Polimorfismo Parametrico, Bounded Type Parameters, Invarianza, Wildcard (PECS) e Type Erasure}}

\vspace{0.8em}

\section{Analisi Teorica Approfondita (Slide
T15)}\label{chap14_analisi-teorica-approfondita-slide-t15}

La lezione 15 introduce formalmente i \textbf{Tipi Generici}
(\emph{Generics}), introdotti in Java 5 per abilitare il
\textbf{polimorfismo parametrico}, eliminare la fragilità dei cast a
\passthrough{\lstinline!Object!} e garantire la \textbf{Type Safety} a
tempo di compilazione.

\begin{center}\rule{0.5\linewidth}{0.5pt}\end{center}

\subsection{\texorpdfstring{L'Idea Guida \& La Tazza
\texttt{Cup\textless{}T\textgreater{}} (Slide
2-3)}{L'Idea Guida \& La Tazza Cup\textless T\textgreater{} (Slide 2-3)}}\label{chap14_lidea-guida-la-tazza-cupt-slide-2-3}

Nella Slide 3 campeggia l'immagine di una tazza con la scritta:
\[\mathbf{Cup\langle T\rangle}\] accompagnata dal motto: \textbf{``Same
container, different content''}. Un contenitore (una tazza, una coppia,
una lista, un albero) possiede una logica strutturale identica a
prescindere dal tipo di dato che custodisce (caffè, tè, stringhe,
interi, date). Prima di Java 5, per ottenere contenitori riusabili si
usava \passthrough{\lstinline!Object!}, con conseguenze disastrose
sull'affidabilità del software.

\begin{center}\rule{0.5\linewidth}{0.5pt}\end{center}

\subsection{\texorpdfstring{Polimorfismo Parametrico vs Approccio Basato
su \texttt{Object} (Slide
4-7)}{Polimorfismo Parametrico vs Approccio Basato su Object (Slide 4-7)}}\label{chap14_polimorfismo-parametrico-vs-approccio-basato-su-object-slide-4-7}

Consideriamo la funzione identità matematica \(\lambda x.\, x\): * Se
riceve un \passthrough{\lstinline!int!}, restituisce lo stesso
\passthrough{\lstinline!int!}. * Se riceve una
\passthrough{\lstinline!String!}, restituisce la stessa
\passthrough{\lstinline!String!}.

\subsubsection{\texorpdfstring{L'approccio ingenuo (pre-Java 5):
segnaposto
\texttt{Object}}{L'approccio ingenuo (pre-Java 5): segnaposto Object}}\label{chap14_lapproccio-ingenuo-pre-java-5-segnaposto-object}

\begin{lstlisting}[language=Java]
public class Pair {
    private Object first, second;
    public Pair(Object first, Object second) {
        this.first = first;
        this.second = second;
    }
    public Object getFirst() { return first; }
    public Object getSecond() { return second; }
}
\end{lstlisting}

\begin{itemize}
\item
  \textbf{Casting obbligatorio}:

\begin{lstlisting}[language=Java]
Pair p = new Pair("one", "two");
String first = (String) p.getFirst(); // Cast esplicito verboso e tedioso
\end{lstlisting}
\item
  \textbf{Il disastro della mancanza di Type Safety}:

\begin{lstlisting}[language=Java]
Pair q = new Pair(1, 2); // Autoboxing da int a Integer
String second = (String) q.getSecond(); // COMPILA SENZA ERRORI!
// Ma a RUNTIME crasha con: java.lang.ClassCastException: class java.lang.Integer cannot be cast to class java.lang.String
\end{lstlisting}
\end{itemize}

\begin{center}\rule{0.5\linewidth}{0.5pt}\end{center}

\subsection{Classi e Metodi Generici (Slide
8-10)}\label{chap14_classi-e-metodi-generici-slide-8-10}

Con i Generics, la classe introduce una o più \textbf{variabili di tipo}
(\emph{Type Parameters}):

\begin{lstlisting}[language=Java]
public class Pair<T> {
    private T first, second;
    public Pair(T first, T second) {
        this.first = first;
        this.second = second;
    }
    public T getFirst() { return first; }
    public T getSecond() { return second; }
}
\end{lstlisting}

\begin{itemize}
\item
  \textbf{Sintassi d'uso}:

\begin{lstlisting}[language=Java]
Pair<String> p = new Pair<String>("one", "two");
Pair<Integer> q = new Pair<Integer>(1, 2);
\end{lstlisting}
\item
  \textbf{Controllo a compile-time}:

\begin{lstlisting}[language=Java]
String first = p.getFirst(); // Nessun cast necessario: p restituisce String garantita!
String second = q.getSecond(); // COMPILE-TIME ERROR: Type mismatch (Integer non convertibile a String)
\end{lstlisting}

  L'errore viene intercettato prima ancora di mandare in esecuzione il
  programma!
\item
  \textbf{Regola fondamentale sui Tipi Primitivi}: \textbf{Un parametro
  di tipo \(T\) può essere sostituito solo da tipi riferimento
  (oggetti), MAI da tipi primitivi!} Non è consentito scrivere
  \passthrough{\lstinline!Pair<int>!}, ma si deve usare il wrapper
  \passthrough{\lstinline!Pair<Integer>!}.
\item
  \textbf{Convenzioni di Naming}: Si usano lettere maiuscole singole:

  \begin{itemize}
  \tightlist
  \item
    \passthrough{\lstinline!T!}: \emph{Type} generico.
  \item
    \passthrough{\lstinline!E!}: \emph{Element} (usato nelle collezioni:
    \passthrough{\lstinline!List<E>!},
    \passthrough{\lstinline!Set<E>!}).
  \item
    \passthrough{\lstinline!K, V!}: \emph{Key} e \emph{Value} (usato
    nelle mappe: \passthrough{\lstinline!Map<K, V>!}).
  \item
    \passthrough{\lstinline!N!}: \emph{Number}.
  \item
    \passthrough{\lstinline!S, U!}: Tipi secondari ausiliari.
  \end{itemize}
\end{itemize}

\begin{center}\rule{0.5\linewidth}{0.5pt}\end{center}

\subsection{Interfacce della Libreria Standard Riscritte con i Generics
(Slide
11-13)}\label{chap14_interfacce-della-libreria-standard-riscritte-con-i-generics-slide-11-13}

\begin{itemize}
\item
  \textbf{\passthrough{\lstinline!Comparable<T>!}}:

\begin{lstlisting}[language=Java]
public interface Comparable<T> {
    int compareTo(T obj);
}
\end{lstlisting}

  Consente alla classe che la implementa (es.
  \passthrough{\lstinline!Student implements Comparable<Student>!}) di
  ricevere direttamente un parametro di tipo
  \passthrough{\lstinline!Student!} nel metodo
  \passthrough{\lstinline!compareTo(Student other)!}, eliminando
  controlli \passthrough{\lstinline!instanceof!} e downcast forzati.
\item
  \textbf{\passthrough{\lstinline!Iterable<E>!} e
  \passthrough{\lstinline!Iterator<E>!}}:

\begin{lstlisting}[language=Java]
public interface Iterable<E> {
    Iterator<E> iterator();
}
public interface Iterator<E> {
    boolean hasNext();
    E next();
}
\end{lstlisting}

  Permette a costrutti come il ciclo \passthrough{\lstinline!for-each!}
  di iterare in modo fortemente tipizzato senza alcun cast:

\begin{lstlisting}[language=Java]
for (Student student : classRoom) {
    System.out.println(student.toString());
}
\end{lstlisting}
\end{itemize}

\begin{center}\rule{0.5\linewidth}{0.5pt}\end{center}

\subsection{\texorpdfstring{Il Diamond Operator
\texttt{\textless{}\textgreater{}} (Slide
14)}{Il Diamond Operator \textless\textgreater{} (Slide 14)}}\label{chap14_il-diamond-operator-slide-14}

Introdotto in Java 7, evita la ridondanza tra la dichiarazione del tipo
di riferimento e l'istanziazione:

\begin{lstlisting}[language=Java]
// Java 5 e 6:
Pair<Integer> intPair = new Pair<Integer>(2, 4);

// Java 7+:
Pair<Integer> intPair = new Pair<>(2, 4); // Il compilatore inferisce il tipo Integer dai membri a sinistra
\end{lstlisting}

\begin{center}\rule{0.5\linewidth}{0.5pt}\end{center}

\subsection{Bounded Type Parameters: Vincolare i Parametri di Tipo
(Slide
15-17)}\label{chap14_bounded-type-parameters-vincolare-i-parametri-di-tipo-slide-15-17}

\begin{itemize}
\item
  \textbf{Unbounded Type Parameter (Slide 15)}: Se dichiariamo
  \passthrough{\lstinline!public class Pair<T>!}, \(T\) è privo di
  vincoli. Il compilatore applica il meccanismo del \textbf{Type
  Erasure} assumendo che \(T\) sia un generico
  \passthrough{\lstinline!Object!}. Di conseguenza, su un oggetto di
  tipo \(T\) è possibile invocare \textbf{esclusivamente i metodi
  definiti in \passthrough{\lstinline!Object!}}
  (\passthrough{\lstinline!equals!}, \passthrough{\lstinline!hashCode!},
  \passthrough{\lstinline!toString!},
  \passthrough{\lstinline!getClass!}). Se tentiamo di chiamare
  \passthrough{\lstinline!first.getName()!}, il compilatore rigetta il
  codice con errore:
  \passthrough{\lstinline!cannot find symbol: method getName() in T!}.
\item
  \textbf{Upper Bound con \passthrough{\lstinline!extends!} (Slide
  16-17)}: Possiamo imporre che \(T\) sia un sottotipo di una
  determinata classe o implementi una o più interfacce:

\begin{lstlisting}[language=Java]
public class Pair<T extends Person> {
    private T first, second;
    // ...
    public String getFirstName() {
        return first.getName(); // PERFETTAMENTE LEGALE: T è garantito essere almeno una Person!
    }
}
\end{lstlisting}

  \begin{itemize}
  \tightlist
  \item
    \textbf{Vincoli multipli}:
    \passthrough{\lstinline!<T extends ClassA \& InterfaceB \& InterfaceC>!}.
    La classe (se presente) deve comparire per prima, seguita dalle
    interfacce separate da \passthrough{\lstinline!\&!}.
  \item
    \emph{Nota terminologica della Slide 16}: La slide cita
    \passthrough{\lstinline!class Clazz<T super S>!}, ma in Java a
    livello di dichiarazione di classe generica è ammesso solo
    \passthrough{\lstinline!extends!}. La parola chiave
    \passthrough{\lstinline!super!} è utilizzabile
    \textbf{esclusivamente nelle Wildcards} (che vediamo subito sotto).
  \end{itemize}
\end{itemize}

\begin{center}\rule{0.5\linewidth}{0.5pt}\end{center}

\subsection{Ereditarietà e Generics: La Trappola dell'Invarianza (Slide
18-20)}\label{chap14_ereditarietuxe0-e-generics-la-trappola-dellinvarianza-slide-18-20}

Questo è uno dei punti più critici dell'intero corso e delle domande
d'esame.

\needspace{7cm}
\begin{asciibox}
   Ereditarietà Normale             Ereditarietà con Generics
      ┌────────────┐                     ┌───────────────┐
      │   Person   │                     │ Pair<Person>  │
      └─────▲──────┘                     └───────▲───────┘
            │ extends                            │ (Nessun legame
            │                                    │  di parentela!)
      ┌─────┴──────┐                     ┌───────┴───────┐
      │  Student   │                     │ Pair<Student> │
      └────────────┘                     └───────────────┘
   Student È UN Person             Pair<Student> NON È Pair<Person>
\end{asciibox}

\begin{itemize}
\item
  \textbf{Invarianza dei Generics (Slide 19)}: Anche se
  \passthrough{\lstinline!Student!} è sottotipo di
  \passthrough{\lstinline!Person!},
  \textbf{\passthrough{\lstinline!Pair<Student>!} NON è un sottotipo di
  \passthrough{\lstinline!Pair<Person>!}}!

\begin{lstlisting}[language=Java]
Student s = new Student("Sam", 123);
Person p = s; // OK: principio di sostituzione di Liskov

Pair<Student> pairS = new Pair<>(s, new Student("Paul", 456));
Pair<Person> pairP = pairS; // COMPILE-TIME ERROR: Incompatible types!
\end{lstlisting}
\item
  \textbf{Perché Java vieta questo assegnamento?}: Se
  \passthrough{\lstinline!pairP = pairS!} fosse consentito, potremmo
  eseguire:

\begin{lstlisting}[language=Java]
pairP.setFirst(new Professor("Mario", 999)); // Legale su Pair<Person>, poiché un Professor è una Person!
\end{lstlisting}

  Ma poiché \passthrough{\lstinline!pairP!} e
  \passthrough{\lstinline!pairS!} punterebbero allo stesso identico
  oggetto nell'Heap, \passthrough{\lstinline!pairS!} (che pretende di
  contenere solo \passthrough{\lstinline!Student!}) si troverebbe a
  contenere un \passthrough{\lstinline!Professor!}, distruggendo la
  garanzia di tipo a runtime!
\item
  \textbf{Il Contrasto con gli Array (Array Covariance, Slide 20)}: In
  Java gli array sono storicamente \textbf{covarianti}:

\begin{lstlisting}[language=Java]
Student[] arrS = { new Student("Sam", 123) };
Person[] arrP = arrS; // COMPILA PERFETTAMENTE!
arrP[0] = new Professor("Mario", 999); // BOOM! A runtime genera: java.lang.ArrayStoreException
\end{lstlisting}

  La covarianza degli array è universalmente riconosciuta come una
  debolezza storica di Java (introdotta in Java 1.0 prima dell'avvento
  dei generics). I Generics sono stati resi \textbf{invarianti} proprio
  per risolvere questa falla a tempo di compilazione.
\end{itemize}

\begin{center}\rule{0.5\linewidth}{0.5pt}\end{center}

\subsection{\texorpdfstring{Wildcards:
\texttt{\textless{}?\textgreater{}},
\texttt{\textless{}?\ extends\ T\textgreater{}},
\texttt{\textless{}?\ super\ T\textgreater{}} (Slide
21-23)}{Wildcards: \textless?\textgreater, \textless? extends T\textgreater, \textless? super T\textgreater{} (Slide 21-23)}}\label{chap14_wildcards-extends-t-super-t-slide-21-23}

Per recuperare la flessibilità polimorfica senza rinunciare alla Type
Safety, Java ha introdotto il carattere jolly (\emph{Wildcard})
\passthrough{\lstinline!?!}.

\begin{enumerate}
\def\labelenumi{\arabic{enumi}.}
\tightlist
\item
  \textbf{Unbounded Wildcard (\passthrough{\lstinline!<?>!}, Slide 21)}:

  \begin{itemize}
  \item
    Rappresenta un tipo totalmente sconosciuto:

\begin{lstlisting}[language=Java]
Pair<?> pairUnknown = pairS; // Legale!
\end{lstlisting}
  \item
    Possiamo \textbf{leggere} elementi, ma il loro tipo a compile-time è
    unicamente \passthrough{\lstinline!Object!}.
  \item
    \textbf{Non possiamo scrivere o aggiungere nulla} (tranne il
    letterale \passthrough{\lstinline!null!}), perché non sappiamo quale
    sia il tipo effettivo accettabile.
  \end{itemize}
\item
  \textbf{Upper-Bounded Wildcard
  (\passthrough{\lstinline!<? extends T>!}, Slide 22-23)}:

  \begin{itemize}
  \item
    Accetta \passthrough{\lstinline!T!} o qualsiasi sottoclasse di
    \passthrough{\lstinline!T!}:

\begin{lstlisting}[language=Java]
Pair<? extends Person> pairP = pairS; // Legale! Student extends Person
\end{lstlisting}
  \item
    Ora possiamo invocare in sicurezza tutti i metodi di
    \passthrough{\lstinline!Person!}:

\begin{lstlisting}[language=Java]
System.out.println(pairP.getFirst().getName()); // OK!
\end{lstlisting}
  \item
    \textbf{Sintassi pulita nei metodi}: Invece di dichiarare un
    parametro formale di tipo esplicito a livello di metodo:

\begin{lstlisting}[language=Java]
<T extends Person> void printPair(Pair<T> p) { ... }
\end{lstlisting}

    possiamo scrivere in modo molto più compatto e leggibile:

\begin{lstlisting}[language=Java]
void printPair(Pair<? extends Person> p) { ... }
\end{lstlisting}
  \end{itemize}
\item
  \textbf{Il Principio Guida PECS (\emph{Producer Extends, Consumer
  Super})}:

  \begin{itemize}
  \tightlist
  \item
    Se la struttura generica funge da \textbf{produttore} di dati
    (vogliamo solo \emph{leggere} dati di tipo
    \passthrough{\lstinline!T!}), si usa
    \passthrough{\lstinline!<? extends T>!}.
  \item
    Se la struttura funge da \textbf{consumatore} di dati (vogliamo
    \emph{scrivere/aggiungere} elementi di tipo
    \passthrough{\lstinline!T!}), si usa
    \passthrough{\lstinline!<? super T>!}.
  \end{itemize}
\end{enumerate}

\begin{center}\rule{0.5\linewidth}{0.5pt}\end{center}

\section{\texorpdfstring{Analisi Dettagliata del Codice
(\texttt{Lezione15/MavenDate})}{Analisi Dettagliata del Codice (Lezione15/MavenDate)}}\label{chap14_analisi-dettagliata-del-codice-lezione15mavendate}

Tra la Lezione 14 e la Lezione 15, il prof. Pasqua compie due modifiche
sostanziali: 1. Implementa \passthrough{\lstinline!hashCode()!} in
\href{file:///home/wally/Documenti/GitHub/Programmazione-II/25-26/Teoria-e-Esercitazioni/Codice-20260902/Lezione15/MavenDate/src/main/java/it/oop/core/Date.java\#L102-L105}{Date.java}.
2. Riscrive ed espande
\href{file:///home/wally/Documenti/GitHub/Programmazione-II/25-26/Teoria-e-Esercitazioni/Codice-20260902/Lezione15/MavenDate/src/main/java/it/oop/ui/MainDate.java}{MainDate.java}
per integrare la libreria delle \textbf{Java Collections}
(\passthrough{\lstinline!List!}, \passthrough{\lstinline!Set!},
\passthrough{\lstinline!Map!}), il contratto
\passthrough{\lstinline!equals!}/\passthrough{\lstinline!hashCode!}, e i
meccanismi di ordinamento con \passthrough{\lstinline!Comparable!} e
\passthrough{\lstinline!Comparator!} (usando classi anonime e lambdas).

\begin{center}\rule{0.5\linewidth}{0.5pt}\end{center}

\subsection{\texorpdfstring{Il Contratto \texttt{equals} /
\texttt{hashCode} in
\href{file:///home/wally/Documenti/GitHub/Programmazione-II/25-26/Teoria-e-Esercitazioni/Codice-20260902/Lezione15/MavenDate/src/main/java/it/oop/core/Date.java\#L102-L105}{Date.java}}{Il Contratto equals / hashCode in Date.java}}\label{chap14_il-contratto-equals-hashcode-in-date.java}

Nel passaggio da Lezione 14 a Lezione 15 viene aggiunto a
\passthrough{\lstinline!Date!}:

\begin{lstlisting}[language=Java]
    @Override
    public int hashCode() {
        return (day + month) * year;
    }
\end{lstlisting}

\subsubsection{Perché questo metodo è vitale e non
opzionale?}\label{chap14_perchuxe9-questo-metodo-uxe8-vitale-e-non-opzionale}

In Java vige la regola ferrea del contratto di
\passthrough{\lstinline!Object!}: \textgreater{} \textbf{Se due oggetti
sono considerati uguali dal metodo
\passthrough{\lstinline!equals(Object)!}
(\passthrough{\lstinline!a.equals(b) == true!}), allora le chiamate a
\passthrough{\lstinline!hashCode()!} su entrambi gli oggetti DEVONO
restituire lo stesso valore intero
(\passthrough{\lstinline!a.hashCode() == b.hashCode()!}).}

Se si fa l'override di \passthrough{\lstinline!equals!} senza fare
l'override di \passthrough{\lstinline!hashCode!}, strutture dati basate
su hash come \passthrough{\lstinline!HashSet!} e
\passthrough{\lstinline!HashMap!} falliscono miseramente: considerano
due date identiche come memorizzate in \emph{bucket} differenti,
inserendo duplicati o non trovando chiavi presenti.

Nel codice di \passthrough{\lstinline!MainDate.java!}:

\begin{lstlisting}[language=Java]
AmericanDate usD1 = new AmericanDate(7, 11, 2025);
AmericanDate usD2 = new AmericanDate(7, 11, 2025);
System.out.println("usD1 ?= usD2: " + usD1.equals(usD2)); // true
System.out.println("hash(usD1): " + usD1.hashCode());      // 36450
System.out.println("hash(usD2): " + usD2.hashCode());      // 36450
\end{lstlisting}

Calcolo: \((7 + 11) \times 2025 = 18 \times 2025 = 36450\). Entrambi gli
oggetti hanno identico hash code!

\begin{center}\rule{0.5\linewidth}{0.5pt}\end{center}

\subsection{\texorpdfstring{\texttt{List} vs \texttt{Set}: Accettazione
vs Rifiuto dei Duplicati in
\href{file:///home/wally/Documenti/GitHub/Programmazione-II/25-26/Teoria-e-Esercitazioni/Codice-20260902/Lezione15/MavenDate/src/main/java/it/oop/ui/MainDate.java\#L43-L55}{MainDate.java}}{List vs Set: Accettazione vs Rifiuto dei Duplicati in MainDate.java}}\label{chap14_list-vs-set-accettazione-vs-rifiuto-dei-duplicati-in-maindate.java}

\begin{lstlisting}[language=Java]
// 1. Creazione di una List (ammette elementi duplicati e mantiene l'ordine di inserimento)
List<Date> dateList = new ArrayList<>();
dateList.add(new AmericanDate(7, 11, 2025));
dateList.add(new Date(9, 11, 2025));
dateList.add(new TimeStamp(0, 0, 0, 8, 11, 2025));
dateList.add(new AmericanDate(7, 11, 2025)); // Duplicato!
System.out.println("list(" + dateList.size() + "): " + dateList.toString());

// 2. Creazione di un Set (insieme matematico: nessun duplicato consentito)
Set<Date> dateSet = new HashSet<>();
dateSet.add(new AmericanDate(7, 11, 2025));
dateSet.add(new Date(9, 11, 2025));
dateSet.add(new TimeStamp(0, 0, 0, 8, 11, 2025));
dateSet.add(new AmericanDate(7, 11, 2025)); // Duplicato: viene scartato silenziosamente!
System.out.println("set(" + dateSet.size() + "): " + dateSet.toString());
\end{lstlisting}

\begin{itemize}
\tightlist
\item
  \textbf{Risultato a runtime}:

  \begin{itemize}
  \tightlist
  \item
    \passthrough{\lstinline!dateList.size()!} vale \textbf{4} (contiene
    entrambe le istanze di \passthrough{\lstinline!7/11/2025!}).
  \item
    \passthrough{\lstinline!dateSet.size()!} vale \textbf{3} (la seconda
    istanza è stata rifiutata perché \passthrough{\lstinline!equals!} ha
    restituito \passthrough{\lstinline!true!} e
    l'\passthrough{\lstinline!hashCode!} coincideva).
  \end{itemize}
\end{itemize}

\begin{center}\rule{0.5\linewidth}{0.5pt}\end{center}

\subsection{\texorpdfstring{\texttt{Map\textless{}K,\ V\textgreater{}} e
Lambdas Registrate Dinamicamente in
\href{file:///home/wally/Documenti/GitHub/Programmazione-II/25-26/Teoria-e-Esercitazioni/Codice-20260902/Lezione15/MavenDate/src/main/java/it/oop/ui/MainDate.java\#L61-L69}{MainDate.java}}{Map\textless K, V\textgreater{} e Lambdas Registrate Dinamicamente in MainDate.java}}\label{chap14_mapk-v-e-lambdas-registrate-dinamicamente-in-maindate.java}

\begin{lstlisting}[language=Java]
Map<String, FormattedDateConverter> converter = new HashMap<>();
// Registrazione di strategie di conversione tramite espressioni lambda
converter.put("usToIt", d -> new ItalianDate(d.getDay(), d.getMonth(), d.getYear()));
converter.put("itToUs", d -> new AmericanDate(d.getDay(), d.getMonth(), d.getYear()));

FormattedDate usDate = new AmericanDate(1, 1, 1970);
FormattedDate itDate = converter.get("usToIt").convert(usDate);
System.out.println(itDate.prettyPrint()); // Stampa "1 gennaio 1970"

for (String k : converter.keySet())
    System.out.println(converter.get(k).toString());
\end{lstlisting}

\begin{itemize}
\tightlist
\item
  \textbf{Finezza di esame sul \passthrough{\lstinline!toString()!}
  delle Lambdas}: Quando si invoca \passthrough{\lstinline!toString()!}
  su un oggetto lambda in Java, la JVM non stampa il codice sorgente,
  bensì il descrittore sintetico della classe generata a runtime
  (\passthrough{\lstinline!invokedynamic!}):
  \passthrough{\lstinline!it.oop.ui.MainDate$$Lambda/0x...!} (la
  notazione interna con i due dollari \passthrough{\lstinline!$$Lambda!}
  è la convenzione standard con cui la JVM battezza le classi sintetiche
  generate dinamicamente).
\end{itemize}

\begin{center}\rule{0.5\linewidth}{0.5pt}\end{center}

\subsection{\texorpdfstring{Ordinamento: \texttt{Comparable} vs
\texttt{Comparator} (Slide 11 \&
\href{file:///home/wally/Documenti/GitHub/Programmazione-II/25-26/Teoria-e-Esercitazioni/Codice-20260902/Lezione15/MavenDate/src/main/java/it/oop/ui/MainDate.java\#L70-L88}{MainDate.java})}{Ordinamento: Comparable vs Comparator (Slide 11 \& MainDate.java)}}\label{chap14_ordinamento-comparable-vs-comparator-slide-11-maindate.java}

In Java esistono due modi per ordinare oggetti: 1. \textbf{Ordinamento
Naturale con \passthrough{\lstinline!Comparable<T>!}}: * Definito
all'interno della classe stessa tramite
\passthrough{\lstinline!compareTo(T other)!}. *
\passthrough{\lstinline!Date!} implementa
\passthrough{\lstinline!Comparable<Date>!}:
\passthrough{\lstinline"java      Collections.sort(dateList); // Funziona direttamente!"}
La lista disordinata
\passthrough{\lstinline![11/7/2025, y2025m11d9, y2025m11d8, 11/7/2025]!}
viene riordinata cronologicamente:
\passthrough{\lstinline![11/7/2025, 11/7/2025, y2025m11d8[00:00:00], y2025m11d9]!}.

\begin{enumerate}
\def\labelenumi{\arabic{enumi}.}
\setcounter{enumi}{1}
\tightlist
\item
  \textbf{Ordinamento Personalizzato con
  \passthrough{\lstinline!Comparator<T>!}}:

  \begin{itemize}
  \item
    Consideriamo \passthrough{\lstinline!DateInterval!}: eredita da
    \passthrough{\lstinline!OrderedPair<Date>!}, ma \textbf{NON
    implementa \passthrough{\lstinline!Comparable<DateInterval>!}}!
  \item
    Se si prova a chiamare:

\begin{lstlisting}[language=Java]
// Collections.sort(intvList); // COMPILE-TIME ERROR:
// no suitable method found for sort(List<DateInterval>)
// reason: cannot infer type-variable(s) T (DateInterval is not Comparable)
\end{lstlisting}
  \item
    Per risolvere il problema, il prof. Pasqua applica due stili per
    fornire un \passthrough{\lstinline!Comparator<DateInterval>!}
    esterno:

    \begin{itemize}
    \item
      \textbf{Stile 1: Classe Anonima} (ordinamento per data di inizio
      \passthrough{\lstinline!left!} crescente):

\begin{lstlisting}[language=Java]
Collections.sort(intvList, new Comparator<DateInterval>() {
    @Override
    public int compare(DateInterval di1, DateInterval di2) {
        return di1.getLeft().compareTo(di2.getLeft());
    }
});
\end{lstlisting}
    \item
      \textbf{Stile 2: Espressione Lambda} (ordinamento per data di fine
      \passthrough{\lstinline!right!} decrescente):

\begin{lstlisting}[language=Java]
Collections.sort(intvList, (di1, di2) -> di2.getRight().compareTo(di1.getRight()));
\end{lstlisting}
    \end{itemize}
  \end{itemize}
\end{enumerate}

\begin{center}\rule{0.5\linewidth}{0.5pt}\end{center}

\section{Risultato di Compilazione ed Esecuzione
Reale}\label{chap14_risultato-di-compilazione-ed-esecuzione-reale}

Comando eseguito nel progetto
\href{file:///home/wally/Documenti/GitHub/Programmazione-II/25-26/Teoria-e-Esercitazioni/Codice-20260902/Lezione15/MavenDate}{Lezione15/MavenDate}:

\begin{lstlisting}[language=bash]
mvn compile exec:java -Dexec.mainClass="it.oop.ui.MainDate"
\end{lstlisting}

Output ottenuto:

\begin{lstlisting}
y2025m9d10
y2025m1d1
00:00:00
true
OOP exam on 2 febbraio 2026
date pair: [y2025m2d1 .. y2025m1d1]
Pair not ordered!
OOP exam on 2 febbraio 2026
list(4): [11/7/2025, y2025m11d9, y2025m11d8[00:00:00], 11/7/2025]
set(3): [11/7/2025, y2025m11d9, y2025m11d8[00:00:00]]
usD1 ?= usD2: true
hash(usD1): 36450
hash(usD2): 36450
1 gennaio 1970
it.oop.ui.MainDate$$Lambda/0x000000008f383b48@21043ee9
it.oop.ui.MainDate$$Lambda/0x000000008f383920@6dc5f73f
[11/7/2025, 11/7/2025, y2025m11d8[00:00:00], y2025m11d9]
[[y2025m1d1 .. y2025m1d20], [y2025m1d10 .. y2025m1d31]]
[[y2025m1d10 .. y2025m1d31], [y2025m1d1 .. y2025m1d20]]
\end{lstlisting}

\subsection{Decodifica puntuale
dell'output:}\label{chap14_decodifica-puntuale-delloutput}

\begin{enumerate}
\def\labelenumi{\arabic{enumi}.}
\tightlist
\item
  \passthrough{\lstinline!list(4): [...]!} vs
  \passthrough{\lstinline!set(3): [...]!} \(\rightarrow\) Dimostrazione
  pratica che \passthrough{\lstinline!HashSet!} elimina il duplicato
  \passthrough{\lstinline!11/7/2025!} grazie al corretto funzionamento
  congiunto di \passthrough{\lstinline!equals()!} e
  \passthrough{\lstinline!hashCode()!}.
\item
  \passthrough{\lstinline!hash(usD1): 36450!} e
  \passthrough{\lstinline!hash(usD2): 36450!} \(\rightarrow\) Rispetto
  formale del contratto di \passthrough{\lstinline!Object!}.
\item
  \passthrough{\lstinline!1 gennaio 1970!} \(\rightarrow\) La mappa
  \passthrough{\lstinline!converter!} recupera il converter
  \passthrough{\lstinline!"usToIt"!} memorizzato come lambda e trasforma
  \passthrough{\lstinline!1/1/1970!} in formato italiano.
\item
  \passthrough{\lstinline!it.oop.ui.MainDate$$Lambda/...!}
  \(\rightarrow\) Nome sintetico dell'oggetto funzionale generato dalla
  JVM a runtime (in cui il doppio dollaro
  \passthrough{\lstinline!$$Lambda!} rappresenta il discriminatore
  standard del bytecode).
\item
  \passthrough{\lstinline![11/7/2025, 11/7/2025, y2025m11d8[00:00:00], y2025m11d9]!}
  \(\rightarrow\) Lista ordinata secondo l'ordinamento naturale di
  \passthrough{\lstinline!Date!} (7 nov, poi 8 nov, poi 9 nov).
\item
  \passthrough{\lstinline![[y2025m1d1 .. y2025m1d20], [y2025m1d10 .. y2025m1d31]]!}
  \(\rightarrow\) Intervalli ordinati per data d'inizio crescente (1 gen
  prima del 10 gen) tramite \passthrough{\lstinline!Comparator!} in
  classe anonima.
\item
  \passthrough{\lstinline![[y2025m1d10 .. y2025m1d31], [y2025m1d1 .. y2025m1d20]]!}
  \(\rightarrow\) Intervalli ordinati per data di fine decrescente (31
  gen prima del 20 gen) tramite \passthrough{\lstinline!Comparator!} in
  espressione lambda.
\end{enumerate}

\begin{center}\rule{0.5\linewidth}{0.5pt}\end{center}

Dimmi \textbf{``vai''} per procedere con il \textbf{Blocco 15} (Lezione
16 / Slide T16 --- \emph{Collections Framework}, gerarchia completa di
\passthrough{\lstinline!List!}, \passthrough{\lstinline!Set!},
\passthrough{\lstinline!Queue!}, \passthrough{\lstinline!Map!},
complessità asintotiche, iteratori e diff in
\passthrough{\lstinline!Lezione16/MavenDate!}).

\newpage
\chapter[T16 --- Java Collections Framework (JCF)]{T16 --- Java Collections Framework (JCF)}
\label{chap:t16}

{\small\sffamily\color{accentblue!80!black}\textit{Gerarchia Collection, List, Set, Map, Queue; Prestazioni e Scelta delle Strutture Dati}}

\vspace{0.8em}

\section{Analisi Teorica Approfondita (Slide
T16)}\label{chap15_analisi-teorica-approfondita-slide-t16}

La lezione 16 è dedicata al \textbf{Java Collections Framework (JCF)}
(\passthrough{\lstinline!java.util!}), una delle architetture cardine
del linguaggio Java: un insieme unificato e standardizzato di
interfacce, implementazioni concrete e algoritmi polimorfici per la
gestione di collezioni di oggetti.

\begin{center}\rule{0.5\linewidth}{0.5pt}\end{center}

\subsection{Architettura e Gerarchia delle Collezioni (Slide
3-4)}\label{chap15_architettura-e-gerarchia-delle-collezioni-slide-3-4}

Tutte le collezioni di elementi singoli discendono dall'interfaccia
radice \passthrough{\lstinline!Iterable<E>!}:

\needspace{7cm}
\begin{asciibox}
┌───────────────────┐
│    Iterable<E>    │
└─────────▲─────────┘
          │
┌─────────┴─────────┐
│   Collection<E>   │
└─────────▲─────────┘
          │
┌─────────┴───────────────┬─────────────────────────┐
│                         │                         │
┌─────────┴─────────┐     ┌─────────┴─────────┐     ┌─────────┴─────────┐
│      List<E>      │     │      Set<E>       │     │     Queue<E>      │
└─────────▲─────────┘     └─────────▲─────────┘     └─────────▲─────────┘
          │                         │                         │
   ┌──────┴──────┐           ┌──────┴──────┐           ┌──────┴──────┐
   │             │           │             │           │             │
┌──┴────┐   ┌────┴────┐ ┌────┴────┐   ┌────┴────┐ ┌────┴────┐   ┌────┴────┐
│Array- │   │ Linked- │ │ HashSet │   │ TreeSet │ │Priority-│   │  Deque  │
│ List  │   │  List   │ │         │   │(Sorted) │ │  Queue  │   │(Double) │
└───────┘   └────┬────┘ └─────────┘   └─────────┘ └─────────┘   └────┬────┘
                 │                                                   │
                 └───────────────────┬───────────────────────────────┘
                                     │ (LinkedList implementa
                                     ▼  sia List che Deque/Queue)
\end{asciibox}

\begin{studynote}[Nota di Studio] L'interfaccia \passthrough{\lstinline!Map<K, V>!} (dizionari
chiave-valore) \textbf{NON estende
\passthrough{\lstinline!Collection<E>!}}, poiché memorizza coppie
\((K, V)\) anziché singoli elementi. Tuttavia, fa parte integrante a
pieno titolo del \emph{Java Collections Framework}.
\end{studynote}

\begin{center}\rule{0.5\linewidth}{0.5pt}\end{center}

\subsection{\texorpdfstring{L'Interfaccia Radice
\texttt{Collection\textless{}E\textgreater{}} (Slide
6-7)}{L'Interfaccia Radice Collection\textless E\textgreater{} (Slide 6-7)}}\label{chap15_linterfaccia-radice-collectione-slide-6-7}

Definisce il contratto comune a tutti i contenitori di elementi:

\begin{lstlisting}[language=Java]
public interface Collection<E> extends Iterable<E> {
    int size();
    boolean isEmpty();
    boolean contains(Object o);                    // Ricerca basata su equals()
    boolean containsAll(Collection<?> c);
    boolean add(E e);                              // Aggiunge un elemento (restituisce false se rifiutato)
    boolean addAll(Collection<? extends E> c);     // Wildcard covariante!
    boolean remove(Object o);
    boolean removeAll(Collection<?> c);
    void clear();
    Object[] toArray();
    <T> T[] toArray(T[] a);                        // es. set.toArray(new Person[0])
    Iterator<E> iterator();                        // Ereditato da Iterable<E>
}
\end{lstlisting}

\subsubsection{Interscambiabilità tra Collezioni (Slide
7)}\label{chap15_interscambiabilituxe0-tra-collezioni-slide-7}

Tutte le implementazioni forniscono per convenzione un costruttore che
accetta un'altra \passthrough{\lstinline!Collection<? extends E>!},
permettendo conversioni istantanee:

\begin{lstlisting}[language=Java]
Collection<Person> list = new LinkedList<>();
list.add(new Person("Joe"));
list.add(new Person("Sam"));

// Conversione da List a Set (elimina al volo eventuali duplicati):
Collection<Person> set = new HashSet<>(list);

// Esportazione in Array tipizzato:
Person[] arrayP = set.toArray(new Person[0]); // Idioma standard Java
\end{lstlisting}

\begin{center}\rule{0.5\linewidth}{0.5pt}\end{center}

\subsection{\texorpdfstring{L'Interfaccia
\texttt{List\textless{}E\textgreater{}} e le sue Implementazioni (Slide
8-9)}{L'Interfaccia List\textless E\textgreater{} e le sue Implementazioni (Slide 8-9)}}\label{chap15_linterfaccia-liste-e-le-sue-implementazioni-slide-8-9}

Una \textbf{Lista} è una sequenza ordinata di elementi che: 1.
\textbf{Preserva l'ordine di inserimento}. 2. \textbf{Accetta elementi
duplicati} (possono esistere più elementi \(e_1, e_2\) tali che
\passthrough{\lstinline!e1.equals(e2)!} sia vero). 3. \textbf{Fornisce
accesso posizionale indicizzato} tramite indice intero
\([0 \dots \text{size}-1]\): *
\passthrough{\lstinline!E get(int index)!}: legge l'elemento all'indice
specificato. * \passthrough{\lstinline!E set(int index, E element)!}:
rimpiazza l'elemento alla posizione data. *
\passthrough{\lstinline!void add(int index, E element)!}: inserisce
shiftando a destra gli elementi successivi. *
\passthrough{\lstinline!E remove(int index)!}: elimina shiftando a
sinistra. * \passthrough{\lstinline!int indexOf(Object o)!}: restituisce
l'indice della prima occorrenza (o \(-1\)). *
\passthrough{\lstinline!List<E> subList(int fromIndex, int toIndex)!}:
vista sulla porzione \([fromIndex, toIndex)\).

\subsubsection{\texorpdfstring{Confronto delle Implementazioni:
\texttt{ArrayList} vs \texttt{LinkedList} (Slide 5 e
9)}{Confronto delle Implementazioni: ArrayList vs LinkedList (Slide 5 e 9)}}\label{chap15_confronto-delle-implementazioni-arraylist-vs-linkedlist-slide-5-e-9}

\begin{longtable}[]{@{}
  >{\raggedright\arraybackslash}p{(\linewidth - 4\tabcolsep) * \real{0.3333}}
  >{\raggedright\arraybackslash}p{(\linewidth - 4\tabcolsep) * \real{0.3333}}
  >{\raggedright\arraybackslash}p{(\linewidth - 4\tabcolsep) * \real{0.3333}}@{}}
\toprule\noalign{}
\begin{minipage}[b]{\linewidth}\raggedright
Operazione
\end{minipage} & \begin{minipage}[b]{\linewidth}\raggedright
\passthrough{\lstinline!ArrayList<E>!} (Array ridimensionabile)
\end{minipage} & \begin{minipage}[b]{\linewidth}\raggedright
\passthrough{\lstinline!LinkedList<E>!} (Lista doppiamente concatenata)
\end{minipage} \\
\midrule\noalign{}
\endhead
\bottomrule\noalign{}
\endlastfoot
\textbf{Accesso per indice (\passthrough{\lstinline!get(i)!})} &
\textbf{\(O(1)\)} (accesso diretto in memoria contigua) &
\textbf{\(O(n)\)} (deve scorrere i nodi dal capo più vicino) \\
\textbf{Inserimento in coda (\passthrough{\lstinline!add(e)!})} &
\textbf{\(O(1)\)} ammortizzato & \textbf{\(O(1)\)} \\
\textbf{Inserimento in testa (\passthrough{\lstinline!add(0, e)!})} &
\textbf{\(O(n)\)} (deve copiare e shiftare l'array) & \textbf{\(O(1)\)}
(cambio puntatori dei nodi) \\
\textbf{Cancellazione (\passthrough{\lstinline!remove(i)!})} &
\textbf{\(O(n)\)} (shift dei successivi) & \textbf{\(O(1)\)} se si ha il
riferimento al nodo, \(O(n)\) per cercarlo \\
\textbf{Overhead di memoria} & Minimo (array contiguo con eventuale
\emph{capacity} vuota) & Alto (ogni nodo crea un oggetto con puntatori
\passthrough{\lstinline!prev!} e \passthrough{\lstinline!next!}) \\
\textbf{Funzionalità extra} &
\passthrough{\lstinline!ensureCapacity(int minCapacity)!} & Implementa
anche \passthrough{\lstinline!Queue!} e \passthrough{\lstinline!Deque!}:
\passthrough{\lstinline!addFirst()!},
\passthrough{\lstinline!addLast()!}, \passthrough{\lstinline!pop()!},
\passthrough{\lstinline!push()!} \\
\end{longtable}

\begin{center}\rule{0.5\linewidth}{0.5pt}\end{center}

\subsection{\texorpdfstring{L'Interfaccia
\texttt{Set\textless{}E\textgreater{}} e le sue Implementazioni (Slide
10-16)}{L'Interfaccia Set\textless E\textgreater{} e le sue Implementazioni (Slide 10-16)}}\label{chap15_linterfaccia-sete-e-le-sue-implementazioni-slide-10-16}

Un \textbf{Insieme (Set)} modella il concetto matematico di insieme: 1.
\textbf{Nessun duplicato consentito}: non possono mai coesistere due
elementi tali che \passthrough{\lstinline!e1.equals(e2)!} sia vero. 2.
\textbf{Nessun accesso posizionale}: non esistono indici, non si può
fare \passthrough{\lstinline!get(i)!}. 3. L'interfaccia
\passthrough{\lstinline!Set<E>!} non aggiunge nuovi metodi rispetto a
\passthrough{\lstinline!Collection<E>!}, ma ne restringe il contratto
semantico.

\subsubsection{\texorpdfstring{\texttt{HashSet\textless{}E\textgreater{}}
(Slide
11-14)}{HashSet\textless E\textgreater{} (Slide 11-14)}}\label{chap15_hashsete-slide-11-14}

\begin{itemize}
\tightlist
\item
  \textbf{Struttura interna}: Mantiene internamente una tabella hash
  (\passthrough{\lstinline!HashMap<E, Object>!}).
\item
  \textbf{Complessità}: Ricerca (\passthrough{\lstinline!contains!}),
  inserimento (\passthrough{\lstinline!add!}) e rimozione
  (\passthrough{\lstinline!remove!}) avvengono in \textbf{tempo costante
  medio \(O(1)\)}.
\item
  \textbf{Ordine}: \textbf{Nessun ordine garantito}; l'ordine di
  scansione dell'iteratore può cambiare nel tempo se la tabella viene
  riallocata (\emph{rehash}).
\item
  \textbf{Vincolo Fondamentale (Slide 12-14)}: L'univocità si basa su
  \textbf{\passthrough{\lstinline!equals()!} e
  \passthrough{\lstinline!hashCode()!}}.

  \begin{itemize}
  \tightlist
  \item
    \textbf{La Regola Aurea}: \emph{``Se fai l'override di
    \passthrough{\lstinline!equals()!}, DEVI fare l'override anche di
    \passthrough{\lstinline!hashCode()!}, e viceversa''}.
  \item
    Proprietà vincolante: \passthrough{\lstinline!a.equals(b)!}
    \(\implies\) \passthrough{\lstinline!a.hashCode() == b.hashCode()!}.
  \item
    Proprietà desiderabile (riduzione collisioni):
    \passthrough{\lstinline"!a.equals(b)"} \(\implies\) preferibilmente
    \passthrough{\lstinline"a.hashCode() != b.hashCode()"}.
  \item
    Da Java 7 si usa l'utility di sistema:
    \passthrough{\lstinline!Objects.hash(campo1, campo2, ...)!}.
  \end{itemize}
\end{itemize}

\subsubsection{\texorpdfstring{\texttt{TreeSet\textless{}E\textgreater{}}
(Slide
15-16)}{TreeSet\textless E\textgreater{} (Slide 15-16)}}\label{chap15_treesete-slide-15-16}

\begin{itemize}
\tightlist
\item
  \textbf{Struttura interna}: Albero binario di ricerca auto-bilanciante
  (\textbf{Red-Black Tree}).
\item
  \textbf{Complessità}: \passthrough{\lstinline!add!},
  \passthrough{\lstinline!remove!}, \passthrough{\lstinline!contains!}
  richiedono tempo logaritmico \textbf{\(O(\log n)\)}.
\item
  \textbf{Ordinamento garantito}: Gli elementi sono
  \textbf{costantemente mantenuti ordinati}:

  \begin{itemize}
  \tightlist
  \item
    Ordinamento naturale: richiede che gli elementi implementino
    \passthrough{\lstinline!Comparable<E>!}.
  \item
    Ordinamento esplicito: passando un
    \passthrough{\lstinline!Comparator<E>!} al costruttore:
    \passthrough{\lstinline!new TreeSet<>(comparator)!}.
  \end{itemize}
\item
  \textbf{Metodi speciali per insiemi ordinati}:

  \begin{itemize}
  \tightlist
  \item
    \passthrough{\lstinline!E first()!}: elemento minimo.
  \item
    \passthrough{\lstinline!E last()!}: elemento massimo.
  \item
    \passthrough{\lstinline!TreeSet<E> subSet(E fromElement, E toElement)!}:
    sottoinsieme nell'intervallo \([fromElement, toElement)\).
  \end{itemize}
\end{itemize}

\begin{center}\rule{0.5\linewidth}{0.5pt}\end{center}

\subsection{\texorpdfstring{Mappe Associative:
\texttt{Map\textless{}K,\ V\textgreater{}} (Slide
17-19)}{Mappe Associative: Map\textless K, V\textgreater{} (Slide 17-19)}}\label{chap15_mappe-associative-mapk-v-slide-17-19}

Una mappa modella una funzione matematica \(K \to V\) (tabella di
associazioni chiave-valore): * Le \textbf{chiavi (\(K\)) sono univoche}
(formano un \passthrough{\lstinline!Set<K>!}). * A ogni chiave è
associato esattamente un valore (\(V\)). I valori possono ripetersi. *
Metodi cardine: * \passthrough{\lstinline!V put(K key, V value)!}:
inserisce o rimpiazza l'associazione. *
\passthrough{\lstinline!V get(Object key)!}: restituisce il valore (o
\passthrough{\lstinline!null!} se la chiave non esiste). *
\passthrough{\lstinline!boolean containsKey(Object key)!} /
\passthrough{\lstinline!boolean containsValue(Object value)!} *
\passthrough{\lstinline!Set<K> keySet()!}: restituisce la vista
dell'insieme delle chiavi. *
\passthrough{\lstinline!Collection<V> values()!}: restituisce la
collezione dei valori. *
\passthrough{\lstinline!Set<Map.Entry<K, V>> entrySet()!}: restituisce
le coppie chiave-valore. *
\textbf{\passthrough{\lstinline!HashMap<K, V>!}}: implementazione
standard basata su hashing delle chiavi, con operazioni
\passthrough{\lstinline!put!} e \passthrough{\lstinline!get!} in
\textbf{\(O(1)\)}.

\begin{center}\rule{0.5\linewidth}{0.5pt}\end{center}

\subsection{\texorpdfstring{Iteratori e l'Errore
\texttt{ConcurrentModificationException} (Slide
21-23)}{Iteratori e l'Errore ConcurrentModificationException (Slide 21-23)}}\label{chap15_iteratori-e-lerrore-concurrentmodificationexception-slide-21-23}

Tutte le collezioni forniscono un \passthrough{\lstinline!Iterator<E>!}
(\passthrough{\lstinline!hasNext()!}, \passthrough{\lstinline!next()!},
\passthrough{\lstinline!remove()!}). Il ciclo
\passthrough{\lstinline!for-each!} di Java è uno zucchero sintattico
compilato tramite \passthrough{\lstinline!Iterator!}:

\begin{lstlisting}[language=Java]
for (Integer i : list) {
    System.out.println(i);
}
\end{lstlisting}

\subsubsection{\texorpdfstring{La Trappola della Modifica Concorrente
(\emph{Fail-Fast Iterator}, Slide
23)}{La Trappola della Modifica Concorrente (Fail-Fast Iterator, Slide 23)}}\label{chap15_la-trappola-della-modifica-concorrente-fail-fast-iterator-slide-23}

\begin{lstlisting}[language=Java]
// CODICE ERRATO:
int count = 0;
for (Integer i : list) {
    System.out.println(i);
    if (count == 2)
        list.remove(count); // RUNTIME ERROR: java.util.ConcurrentModificationException!
    count++;
}
\end{lstlisting}

\begin{itemize}
\item
  \textbf{Perché crasha?}: Le collezioni Java mantengono internamente un
  contatore delle modifiche strutturali
  (\passthrough{\lstinline!modCount!}). Quando si invoca
  \passthrough{\lstinline!list.remove()!},
  \passthrough{\lstinline!modCount!} viene incrementato. Al ciclo
  successivo, la chiamata implicita a
  \passthrough{\lstinline!it.next()!} scopre che la collezione è
  cambiata alle sue spalle (\emph{fail-fast}) e lancia immediatamente
  \passthrough{\lstinline!ConcurrentModificationException!} per
  prevenire stati incoerenti.
\item
  \textbf{La Soluzione Corretta (Slide 23)}: Usare esplicitamente
  l'iteratore e invocare il \textbf{suo} metodo
  \passthrough{\lstinline!remove()!}, che aggiorna coerentemente lo
  stato interno dell'iteratore:

\begin{lstlisting}[language=Java]
for (Iterator<Integer> it = list.iterator(); it.hasNext(); ) {
    Integer i = it.next();
    if (count == 2)
        it.remove(); // PERFETTAMENTE LEGALE E SICURO!
    count++;
}
\end{lstlisting}
\end{itemize}

\begin{center}\rule{0.5\linewidth}{0.5pt}\end{center}

\subsection{\texorpdfstring{Algoritmi di Ordinamento:
\texttt{Comparable} vs \texttt{Comparator} (Slide
24-29)}{Algoritmi di Ordinamento: Comparable vs Comparator (Slide 24-29)}}\label{chap15_algoritmi-di-ordinamento-comparable-vs-comparator-slide-24-29}

La classe di utilità \passthrough{\lstinline!Collections!}
(\passthrough{\lstinline!java.util.Collections!}) fornisce due varianti
sovraccaricate del metodo statico \passthrough{\lstinline!sort!}:

\subsubsection{Variante 1: Ordinamento
Naturale}\label{chap15_variante-1-ordinamento-naturale}

\begin{lstlisting}[language=Java]
public static <T extends Comparable<? super T>> void sort(List<T> list)
\end{lstlisting}

\begin{itemize}
\tightlist
\item
  \textbf{Perché \passthrough{\lstinline!<? super T>!} nella firma?
  (Slide 26 - Domanda Classica d'Esame)}: Supponiamo di avere
  \passthrough{\lstinline!class Student extends Person!}. Se solo
  \passthrough{\lstinline!Person!} implementa
  \passthrough{\lstinline!Comparable<Person>!}, allora
  \passthrough{\lstinline!Student!} eredita il confronto basato su
  \passthrough{\lstinline!Person!}. Se la firma fosse
  \passthrough{\lstinline!<T extends Comparable<T>>!}, allora
  \passthrough{\lstinline!Student!} dovrebbe implementare tassativamente
  \passthrough{\lstinline!Comparable<Student>!}. Con
  \passthrough{\lstinline!Comparable<? super T>!}, una lista di
  \passthrough{\lstinline!Student!} può essere ordinata legalmente
  usando il \passthrough{\lstinline!compareTo!} della superclasse
  \passthrough{\lstinline!Person!}
  (\passthrough{\lstinline!Comparable<Person>!}), garantendo il massimo
  riuso polimorfico!
\end{itemize}

\subsubsection{\texorpdfstring{Variante 2: Ordinamento On-Demand con
\texttt{Comparator}}{Variante 2: Ordinamento On-Demand con Comparator}}\label{chap15_variante-2-ordinamento-on-demand-con-comparator}

\begin{lstlisting}[language=Java]
public static <T> void sort(List<T> list, Comparator<? super T> cmp)
\end{lstlisting}

\begin{itemize}
\tightlist
\item
  Permette di definire criteri di ordinamento arbitrari senza toccare la
  classe sorgente, passando:

  \begin{enumerate}
  \def\labelenumi{\arabic{enumi}.}
  \tightlist
  \item
    Una classe dedicata
    (\passthrough{\lstinline!class StudentComparator implements Comparator<Student>!}).
  \item
    Una \textbf{classe anonima}
    (\passthrough{\lstinline!new Comparator<Student>() \{ ... \}!}).
  \item
    Un'\textbf{espressione Lambda}
    (\passthrough{\lstinline!(s1, s2) -> s2.getMatricola() - s1.getMatricola()!}).
  \end{enumerate}
\end{itemize}

\begin{center}\rule{0.5\linewidth}{0.5pt}\end{center}

\section{\texorpdfstring{Riscontro Pratico nel Codice
(\texttt{Lezione15}/\texttt{Lezione16})}{Riscontro Pratico nel Codice (Lezione15/Lezione16)}}\label{chap15_riscontro-pratico-nel-codice-lezione15lezione16}

Come osservato durante l'ispezione dei sorgenti: 1. Nella Lezione 15 il
docente ha introdotto tutte le strutture viste nella Slide T16
all'interno di
\href{file:///home/wally/Documenti/GitHub/Programmazione-II/25-26/Teoria-e-Esercitazioni/Codice-20260902/Lezione15/MavenDate/src/main/java/it/oop/ui/MainDate.java}{MainDate.java}:
* \passthrough{\lstinline!List<Date> dateList = new ArrayList<>();!}
(con duplicato \passthrough{\lstinline!AmericanDate(7,11,2025)!}
\(\implies\) dimensione 4). *
\passthrough{\lstinline!Set<Date> dateSet = new HashSet<>();!}
(duplicato scartato \(\implies\) dimensione 3). *
\passthrough{\lstinline!Map<String, FormattedDateConverter> converter = new HashMap<>();!}
* Ordinamento naturale:
\passthrough{\lstinline!Collections.sort(dateList);!} * Ordinamento con
\passthrough{\lstinline!Comparator!} via classe anonima e lambda su
\passthrough{\lstinline!DateInterval!}. 2. Nella Lezione 16 il docente
compie il passo architetturale successivo: separa la gestione degli
errori creando il package \passthrough{\lstinline!it.oop.exception!} con
eccezioni personalizzate (\passthrough{\lstinline!IllegalDateException!}
e \passthrough{\lstinline!OrderdPairException!}), preparando il terreno
per la teoria della Slide T17.

\begin{center}\rule{0.5\linewidth}{0.5pt}\end{center}

Dimmi \textbf{``vai''} per procedere con il \textbf{Blocco 16} (Lezione
17 / Slide T17 --- \emph{Error Handling with Exceptions}, gerarchia
\passthrough{\lstinline!Throwable!}/\passthrough{\lstinline!Exception!}/\passthrough{\lstinline!RuntimeException!},
eccezioni controllate vs non controllate, costrutto
\passthrough{\lstinline!try-catch-finally!},
\passthrough{\lstinline!try-with-resources!} e il codice di
\passthrough{\lstinline!Lezione16/MavenDate!}).

\newpage
\chapter[T17 --- Gestione degli Errori ed Architettu]{T17 --- Gestione degli Errori ed Architettura delle Eccezioni}
\label{chap:t17}

{\small\sffamily\color{accentblue!80!black}\textit{Gerarchia Throwable, Checked vs Unchecked, try-catch-finally, Try-with-resources e AutoCloseable}}

\vspace{0.8em}

\section{Analisi Teorica Approfondita (Slide
T17)}\label{chap16_analisi-teorica-approfondita-slide-t17}

La lezione 17 affronta la gestione degli errori e delle anomalie a
runtime attraverso il meccanismo delle \textbf{Eccezioni} in Java.

\begin{center}\rule{0.5\linewidth}{0.5pt}\end{center}

\subsection{Limiti della Gestione Tradizionale degli Errori (Slide
3-6)}\label{chap16_limiti-della-gestione-tradizionale-degli-errori-slide-3-6}

Nei linguaggi procedurali privi di eccezioni (come C) o nell'uso ingenuo
dei metodi, gli errori vengono segnalati tramite \textbf{valori
sentinella} (es. restituire \passthrough{\lstinline!-1!},
\passthrough{\lstinline!null!},
\passthrough{\lstinline!Float.MAX\_VALUE!}):

\begin{lstlisting}[language=Java]
float division(int num, int den) {
    if (den != 0)
        return (float) num / den;
    else
        return Float.MAX_VALUE; // Valore speciale di errore
}
\end{lstlisting}

\subsubsection{Perché questo approccio è
inadeguato?}\label{chap16_perchuxe9-questo-approccio-uxe8-inadeguato}

\begin{enumerate}
\def\labelenumi{\arabic{enumi}.}
\tightlist
\item
  \textbf{Inquinamento della logica applicativa}: Il chiamante deve
  ricordarsi di controllare continuamente i valori di ritorno con una
  ragnatela di \passthrough{\lstinline!if-else!}.
\item
  \textbf{Ambiguità semantica}: A volte il valore sentinella può essere
  un risultato matematico o di dominio assolutamente lecito.
\item
  \textbf{Inapplicabilità ai Costruttori}: Un costruttore non ha tipo di
  ritorno (\passthrough{\lstinline!void!} implicito); se i parametri
  sono non validi, non può restituire \passthrough{\lstinline!-1!} per
  segnalare il fallimento!
\item
  \textbf{Difficoltà di propagazione}: Se l'errore avviene a 10 livelli
  di profondità nello stack di chiamate, ogni singola funzione
  intermedia dovrebbe propagare manualmente il codice d'errore fino al
  punto in cui può essere gestito.
\end{enumerate}

\begin{center}\rule{0.5\linewidth}{0.5pt}\end{center}

\subsection{\texorpdfstring{La Filosofia delle Eccezioni in Java:
\texttt{try-catch-throw} (Slide
7-10)}{La Filosofia delle Eccezioni in Java: try-catch-throw (Slide 7-10)}}\label{chap16_la-filosofia-delle-eccezioni-in-java-try-catch-throw-slide-7-10}

Java disaccoppia la \textbf{logica di business} (\emph{cosa deve fare il
programma quando tutto va bene}) dalla \textbf{logica di gestione degli
errori} (\emph{cosa fare quando accade un'anomalia}): *
\textbf{\passthrough{\lstinline!try!}}: racchiude il blocco di codice a
rischio di eccezione. * \textbf{\passthrough{\lstinline!throw!}}:
solleva attivamente un'istanza di eccezione nel punto in cui l'anomalia
viene riscontrata. * \textbf{\passthrough{\lstinline!catch!}}: cattura
l'eccezione ed esegue il codice di ripristino/notifica. *
\textbf{Interruzione immediata}: quando viene eseguito
\passthrough{\lstinline!throw!}, il metodo \textbf{interrompe
immediatamente la propria esecuzione}, svuota lo stack frame corrente e
risale lungo la catena di chiamate (\emph{call stack}) fino al primo
blocco \passthrough{\lstinline!catch!} compatibile.

\begin{center}\rule{0.5\linewidth}{0.5pt}\end{center}

\subsection{La Gerarchia delle Eccezioni in Java (Slide
20-21)}\label{chap16_la-gerarchia-delle-eccezioni-in-java-slide-20-21}

Tutto l'albero discende da
\passthrough{\lstinline!java.lang.Throwable!}:

\needspace{7cm}
\begin{asciibox}
                 ┌─────────────────────┐
                 │ java.lang.Throwable │
                 └──────────▲──────────┘
                            │
        ┌───────────────────┴───────────────────┐
        │                                       │
┌───────┴─────────┐                   ┌─────────┴─────────┐
│      Error      │                   │     Exception     │
└───────▲─────────┘                   └─────────▲─────────┘
        │                                       │
┌───────┴───────────────┐               ┌───────┴───────────────┐
│  OutOfMemoryError     │               │  Checked Exceptions   │
│  StackOverflowError   │               │ (IOException, ecc.)   │
│  ThreadDeath          │               └───────────▲───────────┘
└───────────────────────┘                           │
                                        ┌───────────┴───────────┐
                                        │   RuntimeException    │
                                        │ (Unchecked Exceptions)│
                                        └───────────▲───────────┘
                                                    │
                                        ┌───────────┴───────────┐
                                        │ NullPointerException  │
                                        │ ClassCastException    │
                                        │ IndexOutOfBoundsExc.  │
                                        └───────────────────────┘
\end{asciibox}

\begin{enumerate}
\def\labelenumi{\arabic{enumi}.}
\tightlist
\item
  \textbf{\passthrough{\lstinline!Error!} (Unchecked)}:

  \begin{itemize}
  \tightlist
  \item
    Condizioni anomale e catastrofiche interne alla JVM o al sistema
    operativo (\passthrough{\lstinline!OutOfMemoryError!},
    \passthrough{\lstinline!StackOverflowError!}).
  \item
    Il codice utente non deve tentare di catturarle, poiché lo stato
    della memoria della JVM è compromesso.
  \end{itemize}
\item
  \textbf{\passthrough{\lstinline!Exception!}}:

  \begin{itemize}
  \tightlist
  \item
    Anomalie derivanti da condizioni esterne o logiche. Si ramificano in
    due grandi famiglie:

    \begin{itemize}
    \tightlist
    \item
      \textbf{Checked Exceptions} (sottoclassi dirette di
      \passthrough{\lstinline!Exception!}, escluse le
      \passthrough{\lstinline!RuntimeException!}).
    \item
      \textbf{Unchecked Exceptions} (sottoclassi di
      \passthrough{\lstinline!RuntimeException!}).
    \end{itemize}
  \end{itemize}
\end{enumerate}

\begin{center}\rule{0.5\linewidth}{0.5pt}\end{center}

\subsection{Checked vs Unchecked Exceptions (Slide
22-24)}\label{chap16_checked-vs-unchecked-exceptions-slide-22-24}

\begin{longtable}[]{@{}
  >{\raggedright\arraybackslash}p{(\linewidth - 4\tabcolsep) * \real{0.3333}}
  >{\raggedright\arraybackslash}p{(\linewidth - 4\tabcolsep) * \real{0.3333}}
  >{\raggedright\arraybackslash}p{(\linewidth - 4\tabcolsep) * \real{0.3333}}@{}}
\toprule\noalign{}
\begin{minipage}[b]{\linewidth}\raggedright
Caratteristica
\end{minipage} & \begin{minipage}[b]{\linewidth}\raggedright
Checked Exceptions (\passthrough{\lstinline!extends Exception!})
\end{minipage} & \begin{minipage}[b]{\linewidth}\raggedright
Unchecked Exceptions
(\passthrough{\lstinline!extends RuntimeException!})
\end{minipage} \\
\midrule\noalign{}
\endhead
\bottomrule\noalign{}
\endlastfoot
\textbf{Natura dell'errore} & Situazioni anomale \textbf{prevedibili} ma
esterne (es. \passthrough{\lstinline!FileNotFoundException!}, rete
disconnessa). & \textbf{Errori del programmatore} (bug logici,
precondizioni violate, puntatori nulli). \\
\textbf{Controllo del compilatore} & \textbf{Obbligatorio} (\emph{Catch
or Declare rule}). Il codice non compila se non gestite con
\passthrough{\lstinline!try-catch!} o dichiarate con
\passthrough{\lstinline!throws!}. & \textbf{Opzionale}. Il compilatore
non richiede alcuna dichiarazione o cattura. \\
\textbf{Firma del metodo} & Modificata:
\passthrough{\lstinline!public void foo() throws MyCheckedException!} &
Invariata: nessun \passthrough{\lstinline!throws!} richiesto. \\
\textbf{Diffusione (\emph{Propagazione})} & \emph{``Si propagano come un
virus''} (Slide 24): costringono tutti i chiamanti intermedi a
dichiarare \passthrough{\lstinline!throws!} o gestire. & Risalgono
naturalmente lo stack fino al \passthrough{\lstinline!main!} o al
gestore globale. \\
\textbf{Raccomandazione moderna (Slide 24)} & Definirle solo se l'utente
del metodo \textbf{può ragionevolmente intraprendere un'azione di
recupero}. & \textbf{Preferite} per errori di validazione dello stato o
degli argomenti. \\
\end{longtable}

\begin{center}\rule{0.5\linewidth}{0.5pt}\end{center}

\subsection{``Exception Dirty Tricks'': Trasformare Checked in Unchecked
(Slide
25-26)}\label{chap16_exception-dirty-tricks-trasformare-checked-in-unchecked-slide-25-26}

Un pattern comune quando una checked exception sporcherebbe decine di
firme senza che il chiamante intermedio possa gestirla:

\begin{lstlisting}[language=Java]
public void bar() { // Nessuna dichiarazione throws richiesta!
    try {
        foo(); // foo() dichiara throws MyCheckedException
    } catch (MyCheckedException e) {
        throw new RuntimeException(e); // Incapsulata e rilanciata come Unchecked!
    }
}
\end{lstlisting}

L'eccezione originale non viene persa: è memorizzata come causa
(\emph{chained exception} visibile nel messaggio
\passthrough{\lstinline!Caused by:!} dello stacktrace).

\begin{center}\rule{0.5\linewidth}{0.5pt}\end{center}

\subsection{\texorpdfstring{Eccezioni nei Cicli (\emph{Exceptions and
Loops}, Slide
27-28)}{Eccezioni nei Cicli (Exceptions and Loops, Slide 27-28)}}\label{chap16_eccezioni-nei-cicli-exceptions-and-loops-slide-27-28}

\begin{itemize}
\tightlist
\item
  \textbf{Errore sulla singola iterazione (Slide 27)}: il blocco
  \passthrough{\lstinline!try-catch!} è \textbf{all'interno} del ciclo.
  Se l'iterazione fallisce, viene intercettata e il ciclo prosegue con
  la successiva (es. lettura da input utente fino a valore corretto).
\item
  \textbf{Errore fatale per l'intero ciclo (Slide 28)}: il blocco
  \passthrough{\lstinline!try-catch!} \textbf{avvolge} il ciclo. Se si
  verifica un'anomalia, il ciclo viene interrotto definitivamente e
  l'esecuzione salta direttamente al \passthrough{\lstinline!catch!}
  esterno.
\end{itemize}

\begin{center}\rule{0.5\linewidth}{0.5pt}\end{center}

\section{\texorpdfstring{Analisi Dettagliata del Codice
(\texttt{Lezione16/MavenDate})}{Analisi Dettagliata del Codice (Lezione16/MavenDate)}}\label{chap16_analisi-dettagliata-del-codice-lezione16mavendate}

In \passthrough{\lstinline!Lezione16/MavenDate!} il prof. Pasqua applica
i principi della Slide T17 rimuovendo i vecchi
\passthrough{\lstinline'System.out.println("Illegal date!")'} e
introducendo la gestione strutturata degli errori tramite il package
\passthrough{\lstinline!it.oop.exception!}.

\subsection{\texorpdfstring{Le Nuove Eccezioni:
\href{file:///home/wally/Documenti/GitHub/Programmazione-II/25-26/Teoria-e-Esercitazioni/Codice-20260902/Lezione16/MavenDate/src/main/java/it/oop/exception/IllegalDateException.java}{IllegalDateException.java}
e
\href{file:///home/wally/Documenti/GitHub/Programmazione-II/25-26/Teoria-e-Esercitazioni/Codice-20260902/Lezione16/MavenDate/src/main/java/it/oop/exception/OrderdPairException.java}{OrderdPairException.java}}{Le Nuove Eccezioni: IllegalDateException.java e OrderdPairException.java}}\label{chap16_le-nuove-eccezioni-illegaldateexception.java-e-orderdpairexception.java}

\subsubsection{\texorpdfstring{\href{file:///home/wally/Documenti/GitHub/Programmazione-II/25-26/Teoria-e-Esercitazioni/Codice-20260902/Lezione16/MavenDate/src/main/java/it/oop/exception/IllegalDateException.java}{IllegalDateException.java}
(Unchecked)}{IllegalDateException.java (Unchecked)}}\label{chap16_illegaldateexception.java-unchecked}

\begin{lstlisting}[language=Java]
package it.oop.exception;

public class IllegalDateException extends RuntimeException {
    public IllegalDateException(String message) {
        super(message);
    }
}
\end{lstlisting}

\begin{itemize}
\tightlist
\item
  Estende \passthrough{\lstinline!RuntimeException!}: è una
  \textbf{Unchecked Exception}.
\item
  Modella una violazione delle precondizioni sui valori numerici
  (giorno, mese, anno).
\item
  Non costringe tutti i client a dichiarare
  \passthrough{\lstinline!throws IllegalDateException!}.
\end{itemize}

\subsubsection{\texorpdfstring{\href{file:///home/wally/Documenti/GitHub/Programmazione-II/25-26/Teoria-e-Esercitazioni/Codice-20260902/Lezione16/MavenDate/src/main/java/it/oop/exception/OrderdPairException.java}{OrderdPairException.java}
(Checked)}{OrderdPairException.java (Checked)}}\label{chap16_orderdpairexception.java-checked}

\begin{lstlisting}[language=Java]
package it.oop.exception;

public class OrderdPairException extends Exception { // Notare il typo del docente "Orderd"
    public OrderdPairException(String message) {
        super(message);
    }
}
\end{lstlisting}

\begin{itemize}
\tightlist
\item
  Estende direttamente \passthrough{\lstinline!Exception!}: è una
  \textbf{Checked Exception}.
\item
  \textbf{Impatto immediato sulla firma}:

  \begin{itemize}
  \item
    In
    \href{file:///home/wally/Documenti/GitHub/Programmazione-II/25-26/Teoria-e-Esercitazioni/Codice-20260902/Lezione16/MavenDate/src/main/java/it/oop/core/OrderedPair.java\#L9}{OrderedPair.java}:

\begin{lstlisting}[language=Java]
public OrderedPair(T first, T second) throws OrderdPairException {
    this.first = first;
    this.second = second;
    if (first.compareTo(second) > 0)
        throw new OrderdPairException("Not orderd pair");
}
\end{lstlisting}
  \item
    In
    \href{file:///home/wally/Documenti/GitHub/Programmazione-II/25-26/Teoria-e-Esercitazioni/Codice-20260902/Lezione16/MavenDate/src/main/java/it/oop/core/DateInterval.java\#L6}{DateInterval.java}:

\begin{lstlisting}[language=Java]
public DateInterval(Date left, Date right) throws OrderdPairException {
    super(left, right); // Invoca super() che lancia la checked exception, quindi DateInterval DEVE dichiarare throws!
}
\end{lstlisting}
  \end{itemize}
\end{itemize}

\begin{center}\rule{0.5\linewidth}{0.5pt}\end{center}

\subsection{\texorpdfstring{Refactoring del Metodo \texttt{verify()} in
\href{file:///home/wally/Documenti/GitHub/Programmazione-II/25-26/Teoria-e-Esercitazioni/Codice-20260902/Lezione16/MavenDate/src/main/java/it/oop/core/Date.java\#L28-L35}{Date.java}}{Refactoring del Metodo verify() in Date.java}}\label{chap16_refactoring-del-metodo-verify-in-date.java}

\begin{lstlisting}[language=Java]
    void verify() {
        if (year < 0)
            throw new IllegalDateException("Illegal date: wrong year");
        if (month < 1 || month > 12)
            throw new IllegalDateException("Illegal date: wrong month");
        if (day < 1 || day > daysPerMonth(month))
            throw new IllegalDateException("Illegal date: wrong day");
    }
\end{lstlisting}

\begin{itemize}
\tightlist
\item
  \textbf{Miglioramento architetturale}:

  \begin{itemize}
  \tightlist
  \item
    Prima: un giorno errato (es. 35) stampava un messaggio sulla console
    ma creava comunque un oggetto \passthrough{\lstinline!Date!}
    incoerente nell'Heap.
  \item
    Ora: il costruttore viene interrotto da
    \passthrough{\lstinline!throw!} prima che l'assegnazione sia
    completata. \textbf{È impossibile creare un'istanza con stato non
    valido!}
  \end{itemize}
\end{itemize}

\begin{center}\rule{0.5\linewidth}{0.5pt}\end{center}

\subsection{\texorpdfstring{\href{file:///home/wally/Documenti/GitHub/Programmazione-II/25-26/Teoria-e-Esercitazioni/Codice-20260902/Lezione16/MavenDate/src/main/java/it/oop/ui/MainDate.java}{MainDate.java}
e l'Applicazione di Tutti i Pattern della Slide
T17}{MainDate.java e l'Applicazione di Tutti i Pattern della Slide T17}}\label{chap16_maindate.java-e-lapplicazione-di-tutti-i-pattern-della-slide-t17}

\begin{lstlisting}[language=Java]
package it.oop.ui;

import it.oop.core.*;
import it.oop.core.Date;
import it.oop.exception.IllegalDateException;
import it.oop.exception.OrderdPairException;

import java.util.InputMismatchException;
import java.util.Scanner;

public class MainDate {

    public static void main(String[] args) {
        System.out.println("Insert day, month, year as numbers");
        int d = -1, m = -1, y = -1;
        boolean correct = false;

        // 1. Pattern Slide 27: try-catch dentro il ciclo per recuperare l'errore di input dell'utente
        while (correct == false) {
            try {
                Scanner sc = new Scanner(System.in);
                d = sc.nextInt();
                m = sc.nextInt();
                y = sc.nextInt();
                correct = true; // Se uno dei nextInt() fallisce, questa riga non viene raggiunta
            } catch (InputMismatchException ime) {
                System.out.println("Input not valid, retry");
            }
        }

        // 2. Cattura della Unchecked Exception IllegalDateException per notifica pulita
        try {
            FormattedDate date = new ItalianDate(d, m, y);
            System.out.println(date.toString());
        } catch (IllegalDateException ide) {
            System.out.println("Date not valid: " + ide.getMessage());
        }

        // 3. Pattern Slide 26 ("Dirty Trick"): incapsulamento di Checked Exception in RuntimeException
        try {
            DateInterval di = new DateInterval(new Date(31, 1, 2025), new Date(1, 1, 2025));
        } catch (OrderdPairException ope) {
            ope.printStackTrace();
            throw new RuntimeException(ope); // Rilanciata come unchecked!
        }
    }
}
\end{lstlisting}

\begin{center}\rule{0.5\linewidth}{0.5pt}\end{center}

\section{Risultato di Esecuzione Reale da
Terminale}\label{chap16_risultato-di-esecuzione-reale-da-terminale}

\subsection{\texorpdfstring{Test 1: Input con data errata
(\texttt{35\ 1\ 2025})}{Test 1: Input con data errata (35 1 2025)}}\label{chap16_test-1-input-con-data-errata-35-1-2025}

Esecuzione:

\begin{lstlisting}[language=bash]
mvn exec:java -Dexec.mainClass="it.oop.ui.MainDate" <<< "35 1 2025"
\end{lstlisting}

Output:

\begin{lstlisting}
Insert day, month, year as numbers
Date not valid: Illegal date: wrong day
it.oop.exception.OrderdPairException: Not orderd pair
    at it.oop.core.OrderedPair.<init>(OrderedPair.java:13)
    at it.oop.core.DateInterval.<init>(DateInterval.java:7)
    at it.oop.ui.MainDate.main(MainDate.java:35)
...
Caused by: it.oop.exception.OrderdPairException: Not orderd pair
\end{lstlisting}

\subsection{Analisi dell'output:}\label{chap16_analisi-delloutput}

\begin{enumerate}
\def\labelenumi{\arabic{enumi}.}
\tightlist
\item
  \passthrough{\lstinline!Date not valid: Illegal date: wrong day!}: il
  costruttore di \passthrough{\lstinline!ItalianDate!} delega a
  \passthrough{\lstinline!Date.verify()!}, che lancia
  \passthrough{\lstinline!IllegalDateException("Illegal date: wrong day")!},
  catturata dal blocco
  \passthrough{\lstinline!catch (IllegalDateException ide)!} che stampa
  il messaggio di errore controllato.
\item
  \passthrough{\lstinline!DateInterval di = new DateInterval(31/1/2025, 1/1/2025)!}:
  poiché la data iniziale \passthrough{\lstinline!31/1/2025!} è
  cronologicamente successiva alla data finale
  \passthrough{\lstinline!1/1/2025!} (violazione dell'invariante
  \passthrough{\lstinline!start <= end!}),
  \passthrough{\lstinline!OrderedPair!} solleva la checked
  \passthrough{\lstinline!OrderdPairException!}. Il blocco catch stampa
  lo stack trace con \passthrough{\lstinline!ope.printStackTrace()!} e
  poi rilancia l'errore incapsulato con
  \passthrough{\lstinline!throw new RuntimeException(ope);!}, terminando
  l'esecuzione con codice di uscita d'errore (build failure) come
  previsto.
\end{enumerate}

\begin{center}\rule{0.5\linewidth}{0.5pt}\end{center}

Dimmi \textbf{``vai''} per procedere con il \textbf{Blocco 17} (Lezione
18 / Slide T18 --- \emph{Streams and Functional Programming}, Stream
API, operazioni intermedie vs terminali, \passthrough{\lstinline!map!},
\passthrough{\lstinline!filter!}, \passthrough{\lstinline!reduce!},
\passthrough{\lstinline!collect!} e il nuovo file
\passthrough{\lstinline!MainStream.java!} in
\passthrough{\lstinline!Lezione17/MavenDate!}).

\newpage
\chapter[T18 --- Interfacce Funzionali, Espressioni ]{T18 --- Interfacce Funzionali, Espressioni Lambda e Streams API}
\label{chap:t18}

{\small\sffamily\color{accentblue!80!black}\textit{Predicate, Function, Consumer, Supplier, Method Reference (::) e Pipeline Stream}}

\vspace{0.8em}

\section{Analisi Teorica Approfondita (Slide
T18)}\label{chap17_analisi-teorica-approfondita-slide-t18}

La lezione 18 introduce il paradigma di \textbf{programmazione
funzionale e dichiarativa} in Java (introdotto in Java 8), imperniato su
due pilastri: 1. \textbf{Interfacce Funzionali} e \textbf{Method
Reference} (\passthrough{\lstinline!::!}). 2. \textbf{Stream API}
(\passthrough{\lstinline!java.util.stream!}), per l'elaborazione fluente
e parallela di flussi di dati.

\begin{center}\rule{0.5\linewidth}{0.5pt}\end{center}

\subsection{Cos'è una Functional Interface? (Slide
3)}\label{chap17_cosuxe8-una-functional-interface-slide-3}

Un'interfaccia si dice \textbf{funzionale} se dichiara
\textbf{esattamente un unico metodo astratto} (pattern noto come
\textbf{SAM} --- \emph{Single Abstract Method}). * \textbf{Semantica}:
Puramente funzionale; l'esito della chiamata dipende esclusivamente
dagli argomenti ricevuti (nessun effetto collaterale occulto). *
\textbf{L'annotazione \passthrough{\lstinline!@FunctionalInterface!}}: *
Non è obbligatoria per far funzionare una lambda, ma è una best practice
fondamentale: istruisce il compilatore Java a verificare che
l'interfaccia contenga un solo metodo astratto. Se qualcuno tenta di
aggiungere un secondo metodo astratto, il compilatore rigetta il codice
con errore. * \textbf{Metodi esclusi dal conteggio SAM}: * Metodi con
implementazione predefinita (\passthrough{\lstinline!default!}). *
Metodi statici (\passthrough{\lstinline!static!}). * Metodi pubblici
astratti che fanno override di metodi di
\passthrough{\lstinline!java.lang.Object!} (es.
\passthrough{\lstinline!boolean equals(Object obj)!}).

\begin{center}\rule{0.5\linewidth}{0.5pt}\end{center}

\subsection{\texorpdfstring{Le 4 Interfacce Funzionali Fondamentali
(\texttt{java.util.function}) (Slide
4-12)}{Le 4 Interfacce Funzionali Fondamentali (java.util.function) (Slide 4-12)}}\label{chap17_le-4-interfacce-funzionali-fondamentali-java.util.function-slide-4-12}

Il package \passthrough{\lstinline!java.util.function!} standardizza i
contratti funzionali più frequenti:

\begin{longtable}[]{@{}
  >{\raggedright\arraybackslash}p{(\linewidth - 6\tabcolsep) * \real{0.2500}}
  >{\raggedright\arraybackslash}p{(\linewidth - 6\tabcolsep) * \real{0.2500}}
  >{\raggedright\arraybackslash}p{(\linewidth - 6\tabcolsep) * \real{0.2500}}
  >{\raggedright\arraybackslash}p{(\linewidth - 6\tabcolsep) * \real{0.2500}}@{}}
\toprule\noalign{}
\begin{minipage}[b]{\linewidth}\raggedright
Interfaccia
\end{minipage} & \begin{minipage}[b]{\linewidth}\raggedright
Firma Metodo SAM
\end{minipage} & \begin{minipage}[b]{\linewidth}\raggedright
Semantica Funzionale
\end{minipage} & \begin{minipage}[b]{\linewidth}\raggedright
Esempio d'Uso Tipico
\end{minipage} \\
\midrule\noalign{}
\endhead
\bottomrule\noalign{}
\endlastfoot
\passthrough{\lstinline!Predicate<T>!} &
\passthrough{\lstinline!boolean test(T t)!} & Filtro / Condizione
booleana & \passthrough{\lstinline!n -> n \% 2 == 0!} \\
\passthrough{\lstinline!Function<T, R>!} &
\passthrough{\lstinline!R apply(T t)!} & Trasformazione \(T \to R\) &
\passthrough{\lstinline!str -> str.length()!} \\
\passthrough{\lstinline!Consumer<T>!} &
\passthrough{\lstinline!void accept(T t)!} & Azione con effetto
collaterale &
\passthrough{\lstinline!str -> System.out.println(str)!} \\
\passthrough{\lstinline!Supplier<T>!} &
\passthrough{\lstinline!T get()!} & Generatore / Factory &
\passthrough{\lstinline!() -> Math.random()!} \\
\end{longtable}

\subsubsection{Varianti e Specializzazioni (Slide
6-12)}\label{chap17_varianti-e-specializzazioni-slide-6-12}

\begin{enumerate}
\def\labelenumi{\arabic{enumi}.}
\tightlist
\item
  \textbf{Versioni Primitive}: Per evitare l'overhead di
  boxing/unboxing, Java fornisce interfacce specializzate per
  \passthrough{\lstinline!int!}, \passthrough{\lstinline!long!},
  \passthrough{\lstinline!double!}:

  \begin{itemize}
  \tightlist
  \item
    \passthrough{\lstinline!IntPredicate!}
    (\passthrough{\lstinline!boolean test(int v)!}),
    \passthrough{\lstinline!LongPredicate!},
    \passthrough{\lstinline!DoublePredicate!}.
  \item
    \passthrough{\lstinline!IntFunction<R>!},
    \passthrough{\lstinline!ToIntFunction<T>!},
    \passthrough{\lstinline!IntConsumer!},
    \passthrough{\lstinline!IntSupplier!}.
  \end{itemize}
\item
  \textbf{Varianti a Due Argomenti}:

  \begin{itemize}
  \tightlist
  \item
    \passthrough{\lstinline!BiPredicate<T, U>!}:
    \passthrough{\lstinline!boolean test(T t, U u)!} (es.
    \passthrough{\lstinline!(a, b) -> a > b!}).
  \item
    \passthrough{\lstinline!BiFunction<T, U, R>!}:
    \passthrough{\lstinline!R apply(T t, U u)!} (es.
    \passthrough{\lstinline!(a, b) -> a + b!}).
  \item
    \passthrough{\lstinline!BiConsumer<T, U>!}:
    \passthrough{\lstinline!void accept(T t, U u)!}.
  \end{itemize}
\item
  \textbf{Operatori (Input e Output dello stesso tipo)}:

  \begin{itemize}
  \tightlist
  \item
    \passthrough{\lstinline!UnaryOperator<T> extends Function<T, T>!}:
    riceve un \(T\) e restituisce un \(T\).
  \item
    \passthrough{\lstinline!BinaryOperator<T> extends BiFunction<T, T, T>!}:
    riceve due \(T\) e restituisce un \(T\).
  \end{itemize}
\end{enumerate}

\begin{center}\rule{0.5\linewidth}{0.5pt}\end{center}

\subsection{\texorpdfstring{I Quattro Tipi di Method Reference
(\texttt{::}) (Slide
13-15)}{I Quattro Tipi di Method Reference (::) (Slide 13-15)}}\label{chap17_i-quattro-tipi-di-method-reference-slide-13-15}

Il \emph{Method Reference} è uno zucchero sintattico introdotto in Java
8 per rendere le espressioni lambda ancora più concise quando queste si
limitano a inoltrare direttamente i propri parametri a un metodo
esistente:

\begin{longtable}[]{@{}
  >{\raggedright\arraybackslash}p{(\linewidth - 4\tabcolsep) * \real{0.3333}}
  >{\raggedright\arraybackslash}p{(\linewidth - 4\tabcolsep) * \real{0.3333}}
  >{\raggedright\arraybackslash}p{(\linewidth - 4\tabcolsep) * \real{0.3333}}@{}}
\toprule\noalign{}
\begin{minipage}[b]{\linewidth}\raggedright
Categoria
\end{minipage} & \begin{minipage}[b]{\linewidth}\raggedright
Sintassi Method Reference
\end{minipage} & \begin{minipage}[b]{\linewidth}\raggedright
Equivalente Lambda Esplicita
\end{minipage} \\
\midrule\noalign{}
\endhead
\bottomrule\noalign{}
\endlastfoot
\textbf{1. Metodo Statico} &
\passthrough{\lstinline!Math::random!}\passthrough{\lstinline!Integer::parseInt!}
&
\passthrough{\lstinline!() -> Math.random()!}\passthrough{\lstinline!str -> Integer.parseInt(str)!} \\
\textbf{2. Metodo d'Istanza su Oggetto Esistente} (\emph{Bound}) &
\passthrough{\lstinline!System.out::println!}\passthrough{\lstinline!hexDigits::charAt!}
&
\passthrough{\lstinline!x -> System.out.println(x)!}\passthrough{\lstinline!i -> hexDigits.charAt(i)!} \\
\textbf{3. Metodo d'Istanza su Tipo Arbitrario} (\emph{Unbound}) &
\passthrough{\lstinline!String::toUpperCase!}\passthrough{\lstinline!String::length!}
&
\passthrough{\lstinline!str -> str.toUpperCase()!}\passthrough{\lstinline!str -> str.length()!} \\
\textbf{4. Costruttore} &
\passthrough{\lstinline!Person::new!}\passthrough{\lstinline!Integer[]::new!}
&
\passthrough{\lstinline!str -> new Person(str)!}\passthrough{\lstinline!size -> new Integer[size]!} \\
\end{longtable}

\begin{center}\rule{0.5\linewidth}{0.5pt}\end{center}

\subsection{\texorpdfstring{La Stream API (\texttt{java.util.stream})
(Slide
16-18)}{La Stream API (java.util.stream) (Slide 16-18)}}\label{chap17_la-stream-api-java.util.stream-slide-16-18}

Uno \passthrough{\lstinline!Stream<T>!} è una \textbf{sequenza di
elementi generata da una sorgente} che supporta operazioni aggregate di
calcolo.

\subsubsection{Le Tre Proprietà Chiave di uno
Stream:}\label{chap17_le-tre-proprietuxe0-chiave-di-uno-stream}

\begin{enumerate}
\def\labelenumi{\arabic{enumi}.}
\tightlist
\item
  \textbf{Pipelining}: Le operazioni intermedie restituiscono un nuovo
  stream, permettendo la concatenazione fluente (\emph{method
  chaining}).
\item
  \textbf{Internal Iteration}: A differenza delle collezioni dove lo
  sviluppatore controlla l'iterazione tramite
  \passthrough{\lstinline!for!}/\passthrough{\lstinline!while!}
  (\emph{iterazione esterna}), lo stream gestisce il ciclo internamente.
  Ciò consente alla JVM di ottimizzare l'esecuzione e abilitare il
  parallelismo trasparente (\passthrough{\lstinline!parallelStream()!}).
\item
  \textbf{Lazy Evaluation (Valutazione Pigra)}: Le operazioni intermedie
  non vengono eseguite quando vengono dichiarate; restano ``in attesa''.
  L'intera pipeline viene elaborata in un unico passaggio solo quando
  viene invocata un'operazione terminale!
\end{enumerate}

\begin{center}\rule{0.5\linewidth}{0.5pt}\end{center}

\subsection{Ciclo di Vita di una Pipeline Stream (Slide
18-24)}\label{chap17_ciclo-di-vita-di-una-pipeline-stream-slide-18-24}

Una pipeline si compone tassativamente di 3 fasi:

\[\text{Sorgente (Source)} \longrightarrow \text{Zero o più Operazioni Intermedie} \longrightarrow \text{Una Operazione Terminale}\]

\needspace{7cm}
\begin{asciibox}
┌───────────────┐     ┌─────────────┐     ┌───────────┐     ┌───────────────┐
│ Stream.of(...) ├──► │ .filter(...) ├──► │ .map(...) ├──►  │ .forEach(...) │
└───────────────┘     └─────────────┘     └───────────┘     └───────────────┘
   [Sorgente]          [Intermedia]        [Intermedia]        [Terminale]
\end{asciibox}

\subsubsection{Creazione della Sorgente (Slide 19,
21-22)}\label{chap17_creazione-della-sorgente-slide-19-21-22}

\begin{itemize}
\tightlist
\item
  Da array: \passthrough{\lstinline!Arrays.stream(arr)!} o
  \passthrough{\lstinline!Stream.of(v1, v2, v3)!}.
\item
  Da collezione: \passthrough{\lstinline!list.stream()!} o
  \passthrough{\lstinline!set.stream()!}.
\item
  \textbf{Generazione infinita}:

  \begin{itemize}
  \tightlist
  \item
    \passthrough{\lstinline!Stream.iterate(seed, UnaryOperator)!}:
    applica iterativamente l'operatore a partire dal seme. Es:
    \passthrough{\lstinline!Stream.iterate(0, i -> i + 1)!} genera
    \(0, 1, 2, 3, \dots\)
  \item
    \passthrough{\lstinline!Stream.generate(Supplier)!}: invoca
    continuamente il fornitore. Es:
    \passthrough{\lstinline!Stream.generate(Math::random)!}.
  \item
    Per evitare loop infiniti, si arresta la generazione con
    l'operazione intermedia \passthrough{\lstinline!.limit(n)!}.
  \end{itemize}
\end{itemize}

\subsubsection{Operazioni Intermedie (Slide 23, 25-29,
34-36)}\label{chap17_operazioni-intermedie-slide-23-25-29-34-36}

Restituiscono un nuovo \passthrough{\lstinline!Stream<R>!} e sono pigre
(\emph{lazy}): * \passthrough{\lstinline!filter(Predicate<T>)!}:
trattiene solo gli elementi per cui il predicato è
\passthrough{\lstinline!true!}. *
\passthrough{\lstinline!map(Function<T, R>)!}: trasforma ciascun
elemento \(T \to R\) (mappatura \(1 \to 1\)). *
\passthrough{\lstinline!flatMap(Function<T, Stream<R>>)!}: appiattisce
strutture annidate (es. matrici \passthrough{\lstinline!Integer[][]!} o
liste di liste) in un unico stream monodimensionale. *
\passthrough{\lstinline!distinct()!}: elimina i duplicati (in base a
\passthrough{\lstinline!equals!} e \passthrough{\lstinline!hashCode!}).
* \passthrough{\lstinline!sorted()!} /
\passthrough{\lstinline!sorted(Comparator)!}: ordina gli elementi. *
\passthrough{\lstinline!limit(n)!}: tronca lo stream ai primi \(n\)
elementi. * \passthrough{\lstinline!skip(n)!}: scarta i primi \(n\)
elementi.

\subsubsection{Operazioni Terminali (Slide 20, 23,
30-33)}\label{chap17_operazioni-terminali-slide-20-23-30-33}

Consumano lo stream e producono un risultato o un effetto collaterale: *
\textbf{Chiusura su Collezioni / Array (Slide 20)}: *
\passthrough{\lstinline!stream.toList()!}: raccoglie gli elementi in una
\passthrough{\lstinline!List!} immutabile (Java 16+). *
\passthrough{\lstinline!stream.collect(Collectors.toSet())!}: raccoglie
in un \passthrough{\lstinline!Set!}. *
\passthrough{\lstinline!stream.toArray(Integer[]::new)!}: esporta in
array tipizzato. * \textbf{Iterazione / Consumo}:
\passthrough{\lstinline!stream.forEach(Consumer<T>)!}. *
\textbf{Conteggio / Riduzione}:
\passthrough{\lstinline!stream.count()!},
\passthrough{\lstinline!stream.reduce(...)!}. * \textbf{Matching}:
\passthrough{\lstinline!anyMatch(Predicate)!},
\passthrough{\lstinline!allMatch(Predicate)!},
\passthrough{\lstinline!noneMatch(Predicate)!}. * \textbf{Ricerca \&
\passthrough{\lstinline!Optional<T>!} (Slide 31-32)}: *
\passthrough{\lstinline!findFirst()!},
\passthrough{\lstinline!findAny()!}: restituiscono un
\passthrough{\lstinline!Optional<T>!}, contenitore sicuro introdotto per
debellare le \passthrough{\lstinline!NullPointerException!}. * Metodi di
\passthrough{\lstinline!Optional!}:
\passthrough{\lstinline!isPresent()!}, \passthrough{\lstinline!get()!},
\passthrough{\lstinline!orElse(valoreDefault)!}.

\begin{studywarning}[Attenzione / Trappola d'Esame] \textbf{Regola Aurea degli Stream (Slide 24 --- Domanda
d'Esame):} \textbf{Uno Stream può essere consumato UNA SOLA VOLTA!} Se
si tenta di invocare una seconda operazione terminale sullo stesso
oggetto \passthrough{\lstinline!Stream!}, la JVM solleva a runtime:
\passthrough{\lstinline!java.lang.IllegalStateException: stream has already been operated upon or closed!}.
\end{studywarning}

\begin{center}\rule{0.5\linewidth}{0.5pt}\end{center}

\section{\texorpdfstring{Analisi Dettagliata del Codice
(\texttt{Lezione17/MavenDate})}{Analisi Dettagliata del Codice (Lezione17/MavenDate)}}\label{chap17_analisi-dettagliata-del-codice-lezione17mavendate}

In \passthrough{\lstinline!Lezione17/MavenDate!}, il prof. Pasqua
mantiene l'intera architettura pregressa e aggiunge il file
\href{file:///home/wally/Documenti/GitHub/Programmazione-II/25-26/Teoria-e-Esercitazioni/Codice-20260902/Lezione17/MavenDate/src/main/java/it/oop/ui/MainStream.java}{MainStream.java},
che esemplifica in modo magistrale tutte le funzionalità degli Stream
appena studiate.

\subsection{\texorpdfstring{Analisi Riga per Riga di
\href{file:///home/wally/Documenti/GitHub/Programmazione-II/25-26/Teoria-e-Esercitazioni/Codice-20260902/Lezione17/MavenDate/src/main/java/it/oop/ui/MainStream.java}{MainStream.java}}{Analisi Riga per Riga di MainStream.java}}\label{chap17_analisi-riga-per-riga-di-mainstream.java}

\begin{lstlisting}[language=Java]
package it.oop.ui;

import it.oop.core.AmericanDate;
import it.oop.core.FormattedDate;
import it.oop.core.ItalianDate;

import java.util.List;
import java.util.stream.Stream;
import it.oop.core.Date;
import static java.util.stream.Stream.generate;

public class MainStream {
    public static void main(String[] args) {
        // ESEMPIO 1: Generazione con iterate, filtro e limit
        // Sequenza generata: "a", "aa", "aaa", "aaaa", ...
        Stream<String> stream = Stream.iterate("a", s -> s + "a")
                .filter(s -> s.length() % 2 == 1) // Tiene solo stringhe di lunghezza DISPARI
                .limit(10);                       // Si arresta dopo le prime 10 stringhe dispari
        List<String> list = stream.toList();      // Operazione terminale: unwrap a List
        System.out.println(list.toString());

        // ESEMPIO 2: Generazione con generate, Supplier e test di primalità
        List<Integer> list2 = Stream.generate(() -> (int) (Math.random() * 20)) // Generatore infinito [0..19]
                .filter(i -> isPrime(i))                                        // Filtra numeri primi
                .limit(10)                                                      // Ne estrae esattamente 10
                .toList();                                                      // Operazione terminale
        System.out.println(list2.toString());

        // ESEMPIO 3: Pipeline su oggetti di dominio del progetto (Date)
        Stream.of(new ItalianDate(15,1,2025), new Date(16,1,2025), new ItalianDate(2,2,2025))
                .filter(d -> d instanceof ItalianDate) // Filtra escludendo le Date generiche
                .map(d -> new AmericanDate(d.getDay(), d.getMonth(), d.getYear())) // Converte in AmericanDate
                .forEach(System.out::println); // Method reference per stampa su stdout
    }

    // Algoritmo di primalità implementato interamente con Stream!
    private static boolean isPrime(int n) {
        return n == 0 ? false : Stream.iterate(1, i -> i + 1)
                .limit(n)              // Genera i numeri da 1 a n
                .map(i -> n % i)       // Calcola il resto della divisione (0 indica divisore esatto)
                .filter(i -> i == 0)   // Mantiene solo i resti nulli (i divisori)
                .count() <= 2;         // Se i divisori sono <= 2 (cioè 1 e se stesso), è primo!
    }
}
\end{lstlisting}

\begin{center}\rule{0.5\linewidth}{0.5pt}\end{center}

\section{Risultato di Compilazione ed Esecuzione
Reale}\label{chap17_risultato-di-compilazione-ed-esecuzione-reale}

Comando eseguito nel progetto
\href{file:///home/wally/Documenti/GitHub/Programmazione-II/25-26/Teoria-e-Esercitazioni/Codice-20260902/Lezione17/MavenDate}{Lezione17/MavenDate}:

\begin{lstlisting}[language=bash]
mvn compile exec:java -Dexec.mainClass="it.oop.ui.MainStream"
\end{lstlisting}

Output ottenuto:

\begin{lstlisting}
[a, aaa, aaaaa, aaaaaaa, aaaaaaaaa, aaaaaaaaaaa, aaaaaaaaaaaaa, aaaaaaaaaaaaaaa, aaaaaaaaaaaaaaaaa, aaaaaaaaaaaaaaaaaaa]
[7, 11, 13, 19, 3, 13, 17, 5, 19, 1]
1/15/2025
2/2/2025
\end{lstlisting}

\subsection{Decodifica dell'Output:}\label{chap17_decodifica-delloutput}

\begin{enumerate}
\def\labelenumi{\arabic{enumi}.}
\tightlist
\item
  \passthrough{\lstinline![a, aaa, aaaaa, ..., aaaaaaaaaaaaaaaaaaa]!}
  \(\rightarrow\) Lista delle prime 10 stringhe a lunghezza dispari (1,
  3, 5, 7, 9, 11, 13, 15, 17, 19 caratteri).
\item
  \passthrough{\lstinline![7, 11, 13, 19, 3, 13, 17, 5, 19, 1]!}
  \(\rightarrow\) Dieci numeri primi casuali estratti nell'intervallo
  \([0, 20)\) e validati dalla pipeline stream di
  \passthrough{\lstinline!isPrime(int n)!}.
\item
  \passthrough{\lstinline!1/15/2025!} e
  \passthrough{\lstinline!2/2/2025!} \(\rightarrow\) Le due istanze di
  \passthrough{\lstinline!ItalianDate!} (15/1/2025 e 2/2/2025) superano
  il \passthrough{\lstinline!filter(instanceof ItalianDate)!}, vengono
  convertite tramite \passthrough{\lstinline!map!} in
  \passthrough{\lstinline!AmericanDate!} (formato
  \passthrough{\lstinline!M/D/Y!}) e stampate a video tramite
  \passthrough{\lstinline!System.out::println!}. L'istanza generica
  \passthrough{\lstinline!new Date(16,1,2025)!} viene invece scartata
  dal filtro.
\end{enumerate}

\begin{center}\rule{0.5\linewidth}{0.5pt}\end{center}

Dimmi \textbf{``vai''} per procedere con il \textbf{Blocco 18} (Lezione
18 / Slide T19 --- \emph{Documentation and Unit Testing}, Javadoc, JUnit
3 vs JUnit 4 vs JUnit 5, asserzioni, ciclo di vita dei test
\passthrough{\lstinline!@BeforeEach!}/\passthrough{\lstinline!@AfterEach!},
e l'analisi di \passthrough{\lstinline!TestDate.java!} e
\passthrough{\lstinline!TestItalianDate.java!} in
\passthrough{\lstinline!Lezione18/MavenDate!}).

\newpage
\chapter[T19 --- Documentazione con Javadoc e Unit T]{T19 --- Documentazione con Javadoc e Unit Testing con JUnit 5}
\label{chap:t19}

{\small\sffamily\color{accentblue!80!black}\textit{Convenzioni Javadoc, Architettura JUnit 5, Annotazioni di Ciclo di Vita ed Asserzioni}}

\vspace{0.8em}

\section{Analisi Teorica Approfondita (Slide
T19)}\label{chap18_analisi-teorica-approfondita-slide-t19}

La lezione 19 approfondisce due pratiche ingegneristiche irrinunciabili
nello sviluppo professionale in Java: 1. \textbf{Documentazione
automatizzata del codice} tramite lo standard \textbf{Javadoc}. 2.
\textbf{Collaudo e Unit Testing automatizzato} tramite il framework
\textbf{JUnit 5 (Jupiter)} e il plugin Maven \textbf{Surefire}.

\begin{center}\rule{0.5\linewidth}{0.5pt}\end{center}

\subsection{Documentazione con Javadoc (Slide
3-11)}\label{chap18_documentazione-con-javadoc-slide-3-11}

La documentazione del codice sorgente è essenziale per manutenibilità,
collaborazione e pubblicazione di API riusabili. In Java la
documentazione vive a stretto contatto con il codice nei cosiddetti
\textbf{Doc Comments}: * \textbf{Sintassi}: Delimitati da
\passthrough{\lstinline!/**!} all'inizio e \passthrough{\lstinline!*/!}
alla fine. Ogni riga intermedia inizia convenzionalmente con
\passthrough{\lstinline!*!}:
\passthrough{\lstinline!java   /**    * Descrizione sintetica del componente.    * <p>Supporta l'uso di tag HTML come paragrafi, link e formattazione.</p>    */!}
* \textbf{I Tag Javadoc Standard (Slide 5-6)}: Iniziano con
\passthrough{\lstinline!@!}, sono \emph{case-sensitive} e devono
comparire a inizio riga: * \passthrough{\lstinline!@param <nome>!}:
descrive un parametro formale di un metodo o costruttore. *
\passthrough{\lstinline!@return!}: descrive il significato del valore
restituito da un metodo (omesso nei metodi
\passthrough{\lstinline!void!} e nei costruttori). *
\passthrough{\lstinline!@throws <Eccezione>!} (o
\passthrough{\lstinline!@exception!}): descrive le condizioni anomale
che causano il sollevamento di una determinata eccezione. *
\passthrough{\lstinline!\{@link <package.Classe\#metodo>\}!}: genera un
ipertesto cliccabile verso un'altra classe o metodo. *
\passthrough{\lstinline!@author!}: autore del modulo. *
\passthrough{\lstinline!@since <versione>!}: versione a partire dalla
quale la feature è disponibile. * \passthrough{\lstinline!@version!}:
versione corrente del sorgente. * \passthrough{\lstinline!@deprecated!}:
avvisa che l'elemento è obsoleto, illustrandone il motivo e
l'alternativa moderna da utilizzare.

\subsubsection{Regole di Visibilità e Generazione (Slide
9-11)}\label{chap18_regole-di-visibilituxe0-e-generazione-slide-9-11}

\begin{itemize}
\item
  \textbf{Default di Javadoc}: Per impostazione predefinita, Javadoc
  genera la documentazione per le sole entità
  \textbf{\passthrough{\lstinline!public!} e
  \passthrough{\lstinline!protected!}} (l'interfaccia pubblica esposta
  ai client). I campi e metodi \passthrough{\lstinline!private!} o
  \emph{package-private} vengono omessi.
\item
  \textbf{Inclusione del privato}: Si deve passare esplicitamente il
  flag \passthrough{\lstinline!-private!} da riga di comando:

\begin{lstlisting}[language=bash]
javadoc -private -d doc *
\end{lstlisting}
\item
  \textbf{Integrazione Maven
  (\passthrough{\lstinline!maven-javadoc-plugin!}, Slide 11)}:
  Aggiungendo il plugin nel \passthrough{\lstinline!pom.xml!}, la
  documentazione viene generata con un singolo comando standard:

\begin{lstlisting}[language=bash]
mvn javadoc:javadoc
\end{lstlisting}
\end{itemize}

\begin{center}\rule{0.5\linewidth}{0.5pt}\end{center}

\subsection{Unit Testing con JUnit (Slide
13-16)}\label{chap18_unit-testing-con-junit-slide-13-16}

Il testing programmatico garantisce la correttezza delle singole unità
software in isolamento (metodi o singole classi): * \textbf{Vantaggi
sistemici}: * \textbf{Test-Driven Development (TDD)}: scrivere i test
prima ancora di implementare il codice applicativo. * \textbf{Regression
Testing}: certezza che refactoring o nuove feature non rompano
comportamenti pregressi funzionanti. * Riduzione drastica del tempo
speso nel debugging manuale. * \textbf{Anatomia di un Test Case}: 1.
Preparazione dell'input (\emph{Arrange / Setup}). 2. Esecuzione del
metodo sotto test (\emph{Act}). 3. Confronto dell'output effettivo
(\emph{Actual}) con l'output atteso (\emph{Expected}) tramite
\textbf{asserzioni} (\emph{Assert}). * \textbf{Meccanismo di
Fallimento}: * Un metodo di test restituisce sempre
\passthrough{\lstinline!void!}. * Se tutte le asserzioni sono
soddisfatte, il metodo termina normalmente e JUnit lo contrassegna come
\textbf{PASSED} (verde). * Se un'asserzione fallisce, solleva un errore
di tipo \passthrough{\lstinline!AssertionError!} (o
\passthrough{\lstinline!AssertionFailedError!}). Il framework JUnit
cattura l'errore, contrassegna il test come \textbf{FAILED} (rosso) con
il report della discrepanza, e \textbf{prosegue regolarmente con i test
successivi} senza arrestare l'intera suite.

\begin{center}\rule{0.5\linewidth}{0.5pt}\end{center}

\subsection{\texorpdfstring{Asserzioni in JUnit 5
(\texttt{org.junit.jupiter.api.Assertions}, Slide
17-21)}{Asserzioni in JUnit 5 (org.junit.jupiter.api.Assertions, Slide 17-21)}}\label{chap18_asserzioni-in-junit-5-org.junit.jupiter.api.assertions-slide-17-21}

Il pacchetto Jupiter standardizza una vasta famiglia di metodi statici:

\begin{longtable}[]{@{}
  >{\raggedright\arraybackslash}p{(\linewidth - 2\tabcolsep) * \real{0.5000}}
  >{\raggedright\arraybackslash}p{(\linewidth - 2\tabcolsep) * \real{0.5000}}@{}}
\toprule\noalign{}
\begin{minipage}[b]{\linewidth}\raggedright
Metodo Asserzione
\end{minipage} & \begin{minipage}[b]{\linewidth}\raggedright
Comportamento
\end{minipage} \\
\midrule\noalign{}
\endhead
\bottomrule\noalign{}
\endlastfoot
\passthrough{\lstinline!assertEquals(expected, actual, [msg])!} &
Verifica che \passthrough{\lstinline!expected.equals(actual)!}. Valido
per primitivi e oggetti. \\
\passthrough{\lstinline!assertTrue(condition, [msg])!} & Verifica che la
condizione booleana sia \passthrough{\lstinline!true!}. \\
\passthrough{\lstinline!assertFalse(condition, [msg])!} & Verifica che
la condizione booleana sia \passthrough{\lstinline!false!}. \\
\passthrough{\lstinline!assertArrayEquals(expected, actual)!} &
Confronta il contenuto e l'ordinamento di due array elemento per
elemento. \\
\passthrough{\lstinline!assertNull(obj)!} /
\passthrough{\lstinline!assertNotNull(obj)!} & Verifica se il puntatore
è \passthrough{\lstinline!null!} o non nullo. \\
\passthrough{\lstinline!assertThrows(Exception.class, executable)!} &
Verifica che l'esecuzione di una lambda/blocco sollevi tassativamente
l'eccezione attesa! \\
\end{longtable}

\subsubsection{Best Practice: Static Imports (Slide
21)}\label{chap18_best-practice-static-imports-slide-21}

Per evitare di anteporre continuamente
\passthrough{\lstinline!Assertions.!} davanti a ogni chiamata, si usa
l'import statico:

\begin{lstlisting}[language=Java]
import static org.junit.jupiter.api.Assertions.assertEquals;
import static org.junit.jupiter.api.Assertions.assertThrows;
\end{lstlisting}

\begin{center}\rule{0.5\linewidth}{0.5pt}\end{center}

\subsection{Il Ciclo di Vita del Test e le Annotazioni JUnit 5 (Slide
22)}\label{chap18_il-ciclo-di-vita-del-test-e-le-annotazioni-junit-5-slide-22}

\needspace{7cm}
\begin{asciibox}
┌────────────────────────┐
│      @BeforeAll        │ (eseguito 1 sola volta prima di tutti, STATIC)
└───────────┬────────────┘
            │
┌───────────┴─────────────────┐
│   Per ogni metodo @Test:    │
│  ┌────────────────────────┐ │
│  │      @BeforeEach       │ │ (setup specifico per il test)
│  └───────────┬────────────┘ │
│              ▼              │
│  ┌────────────────────────┐ │
│  │         @Test          │ │ (esecuzione del caso di test)
│  └───────────┬────────────┘ │
│              ▼              │
│  ┌────────────────────────┐ │
│  │       @AfterEach       │ │ (teardown / pulizia risorse)
│  └────────────────────────┘ │
└───────────┬─────────────────┘
            │
┌───────────┴────────────┐
│       @AfterAll        │ (eseguito 1 sola volta alla fine, STATIC)
└────────────────────────┘
\end{asciibox}

\begin{itemize}
\tightlist
\item
  \passthrough{\lstinline!@BeforeAll!}: Inizializzazione ``pesante'' o
  condivisa tra tutti i test (es. apertura connessione, setup di un
  database o di costanti immutabili). \textbf{Deve essere
  \passthrough{\lstinline!static!}}.
\item
  \passthrough{\lstinline!@AfterAll!}: Chiusura finale delle risorse
  globali. \textbf{Deve essere \passthrough{\lstinline!static!}}.
\item
  \passthrough{\lstinline!@BeforeEach!}: Inizializzazione fresca prima
  di \emph{ciascun} test per garantire l'isolamento e l'indipendenza dei
  test (evitando che un test modifichi lo stato di un altro).
\item
  \passthrough{\lstinline!@AfterEach!}: Pulizia post-test (es.
  cancellazione di file temporanei creati durante il test).
\item
  \passthrough{\lstinline!@DisplayName("Descrizione chiara")!}:
  Personalizza il nome del test visualizzato nei report o nell'IDE.
\item
  \passthrough{\lstinline!@Disabled!}: Disabilita temporaneamente il
  test senza doverlo commentare o cancellare.
\end{itemize}

\begin{center}\rule{0.5\linewidth}{0.5pt}\end{center}

\subsection{Organizzazione del Progetto e Convenzioni Maven (Slide
23-24)}\label{chap18_organizzazione-del-progetto-e-convenzioni-maven-slide-23-24}

Maven e i moderni build tool impongono una struttura standard di
cartelle: * \textbf{\passthrough{\lstinline!src/main/java!}}: Contiene
il codice sorgente dell'applicazione (es.
\passthrough{\lstinline!it.oop.core.Date!}). *
\textbf{\passthrough{\lstinline!src/test/java!}}: Contiene
esclusivamente le classi di collaudo. * \textbf{Allineamento dei Package
(Regola Fondamentale)}: Una classe di test che verifica
\passthrough{\lstinline!it.oop.core.Date!} \textbf{deve risiedere nello
stesso identico package \passthrough{\lstinline!it.oop.core!}}
(all'interno di \passthrough{\lstinline!src/test/java!}). In questo modo
la classe di test ha accesso non solo ai membri
\passthrough{\lstinline!public!}, ma anche a tutti i metodi e campi con
visibilità di \textbf{package
(\passthrough{\lstinline!package-private!})}, facilitando il testing
interno senza forzare l'apertura a \passthrough{\lstinline!public!} di
dettagli architetturali riservati!

\begin{center}\rule{0.5\linewidth}{0.5pt}\end{center}

\section{\texorpdfstring{Analisi Dettagliata del Codice
(\texttt{Lezione18/MavenDate})}{Analisi Dettagliata del Codice (Lezione18/MavenDate)}}\label{chap18_analisi-dettagliata-del-codice-lezione18mavendate}

In \passthrough{\lstinline!Lezione18/MavenDate!} troviamo la
configurazione completa del plugin Maven Surefire e due classi di test
che collaudano la gerarchia di \passthrough{\lstinline!Date!}.

\subsection{\texorpdfstring{Configurazione in
\href{file:///home/wally/Documenti/GitHub/Programmazione-II/25-26/Teoria-e-Esercitazioni/Codice-20260902/Lezione18/MavenDate/pom.xml\#L44-L68}{pom.xml}}{Configurazione in pom.xml}}\label{chap18_configurazione-in-pom.xml}

\begin{lstlisting}[language=XML]
    <build>
        <plugins>
            <!-- Plugin Javadoc (Slide 11) -->
            <plugin>
                <groupId>org.apache.maven.plugins</groupId>
                <artifactId>maven-javadoc-plugin</artifactId>
                <version>3.6.2</version>
                <configuration>
                    <source>1.8</source>
                    <show>private</show> <!-- Include metodi e campi privati -->
                </configuration>
            </plugin>
        </plugins>
    </build>

    <dependencies>
        <!-- Motore di esecuzione JUnit 5 (Slide 24) -->
        <dependency>
            <groupId>org.junit.jupiter</groupId>
            <artifactId>junit-jupiter-engine</artifactId>
            <version>5.10.0</version>
            <scope>test</scope> <!-- Visibile solo durante la fase di test -->
        </dependency>
        <!-- Maven Surefire Plugin per l'esecuzione di 'mvn test' -->
        <dependency>
            <groupId>org.apache.maven.plugins</groupId>
            <artifactId>maven-surefire-plugin</artifactId>
            <version>3.5.4</version>
        </dependency>
    </dependencies>
\end{lstlisting}

\begin{center}\rule{0.5\linewidth}{0.5pt}\end{center}

\subsection{\texorpdfstring{\href{file:///home/wally/Documenti/GitHub/Programmazione-II/25-26/Teoria-e-Esercitazioni/Codice-20260902/Lezione18/MavenDate/src/test/java/it/oop/core/TestDate.java}{TestDate.java}}{TestDate.java}}\label{chap18_testdate.java}

Questa classe collauda sia il flusso nominale (costruzione valida) che
il flusso eccezionale (lancio di eccezioni):

\begin{lstlisting}[language=Java]
package it.oop.core;

import it.oop.exception.IllegalDateException;
import org.junit.jupiter.api.Assertions;
import org.junit.jupiter.api.Test;
import static org.junit.jupiter.api.Assertions.assertThrows;

public class TestDate {

    private Date date;

    // 1. Collaudo del caso nominale con assertEquals
    @Test
    public void testDateConstructor() {
        date = new Date(1, 12, 2025);
        Assertions.assertEquals(1, date.getDay());
        Assertions.assertEquals(12, date.getMonth());
        Assertions.assertEquals(2025, date.getYear());
    }

    // 2. Collaudo del caso eccezionale con assertThrows e Lambda
    @Test
    public void testDateConstructorException() {
        // Verifica che passando un anno negativo (-2), il costruttore sollevi tassativamente IllegalDateException
        assertThrows(IllegalDateException.class, () -> date = new Date(1, 12, -2));
    }
}
\end{lstlisting}

\begin{center}\rule{0.5\linewidth}{0.5pt}\end{center}

\subsection{\texorpdfstring{\href{file:///home/wally/Documenti/GitHub/Programmazione-II/25-26/Teoria-e-Esercitazioni/Codice-20260902/Lezione18/MavenDate/src/test/java/it/oop/core/TestItalianDate.java}{TestItalianDate.java}}{TestItalianDate.java}}\label{chap18_testitaliandate.java}

Esemplifica l'uso della fixture \passthrough{\lstinline!@BeforeAll!} e
mette in luce il superamento della vecchia parola chiave
\passthrough{\lstinline!assert!}:

\begin{lstlisting}[language=Java]
package it.oop.core;

import org.junit.jupiter.api.Assertions;
import org.junit.jupiter.api.BeforeAll;
import org.junit.jupiter.api.Test;

public class TestItalianDate {
    private static ItalianDate date;

    // Fixture statica eseguita una volta sola prima di tutti i test
    @BeforeAll
    public static void setup() {
        date = new ItalianDate(1, 1, 1970);
    }

    @Test
    public void printFormatTest() {
        // Verifica del formato con JUnit 5:
        Assertions.assertEquals("dd/mm/yyyy", date.printFormat());
        
        // Confronto con la vecchia sintassi 'assert' del linguaggio (commentata dal docente):
        // assert date.printFormat().equals("dd/mm/yyyy") : "wrong format";
    }
}
\end{lstlisting}

\begin{center}\rule{0.5\linewidth}{0.5pt}\end{center}

\section{Risultato di Esecuzione Reale da
Terminale}\label{chap18_risultato-di-esecuzione-reale-da-terminale}

Comando eseguito nel progetto
\href{file:///home/wally/Documenti/GitHub/Programmazione-II/25-26/Teoria-e-Esercitazioni/Codice-20260902/Lezione18/MavenDate}{Lezione18/MavenDate}:

\begin{lstlisting}[language=bash]
mvn test
\end{lstlisting}

Output ottenuto:

\begin{lstlisting}
[INFO] -------------------------------------------------------
[INFO]  T E S T S
[INFO] -------------------------------------------------------
[INFO] Running it.oop.core.TestDate
[INFO] Tests run: 2, Failures: 0, Errors: 0, Skipped: 0, Time elapsed: 0.106 s -- in it.oop.core.TestDate
[INFO] Running it.oop.core.TestItalianDate
[INFO] Tests run: 1, Failures: 0, Errors: 0, Skipped: 0, Time elapsed: 0.021 s -- in it.oop.core.TestItalianDate
[INFO] 
[INFO] Results:
[INFO] 
[INFO] Tests run: 3, Failures: 0, Errors: 0, Skipped: 0
[INFO] 
[INFO] ------------------------------------------------------------------------
[INFO] BUILD SUCCESS
[INFO] ------------------------------------------------------------------------
\end{lstlisting}

I 3 test sono stati eseguiti con successo, validando in modo formale la
suite del progetto e generando i relativi file di report XML e TXT nella
cartella \passthrough{\lstinline!target/surefire-reports/!}.

\begin{center}\rule{0.5\linewidth}{0.5pt}\end{center}

Dimmi \textbf{``vai''} per procedere con il \textbf{Blocco 19} (Lezione
19 / Slide T20 --- \emph{Java I/O}, stream di byte vs caratteri,
\passthrough{\lstinline!InputStream!},
\passthrough{\lstinline!OutputStream!},
\passthrough{\lstinline!Reader!}, \passthrough{\lstinline!Writer!},
\passthrough{\lstinline!BufferedReader!}, serializzazione, gestione
delle risorse con \passthrough{\lstinline!try-with-resources!} e analisi
di \passthrough{\lstinline!MainIO.java!} in
\passthrough{\lstinline!Lezione19/MavenDate!}).

\newpage
\chapter[T20 --- Java I/O, Streams, File, NIO.2 e Se]{T20 --- Java I/O, Streams, File, NIO.2 e Serializzazione}
\label{chap:t20}

{\small\sffamily\color{accentblue!80!black}\textit{Byte Streams vs Character Streams, Bufferizzazione, Decorator Pattern I/O e Serializzazione}}

\vspace{0.8em}

\section{Analisi Teorica Approfondita (Slide
T20)}\label{chap19_analisi-teorica-approfondita-slide-t20}

La lezione 20 conclude il percorso teorico del corso esplorando il
sottosistema di \textbf{Input/Output (I/O)} di Java
(\passthrough{\lstinline!java.io!} e
\passthrough{\lstinline!java.net!}), la persistenza su disco, la
comunicazione di rete, e la gestione di formati dati strutturati moderni
(\textbf{CSV} e \textbf{JSON}) tramite librerie terze gestite con Maven.

\begin{center}\rule{0.5\linewidth}{0.5pt}\end{center}

\subsection{L'Astrazione di I/O Stream (Slide
3-4)}\label{chap19_lastrazione-di-io-stream-slide-3-4}

In Java, qualsiasi trasferimento di dati da o verso l'esterno è
modellato tramite l'astrazione di \textbf{I/O Stream} (\emph{flusso
sequenziale unidirezionale di dati}): * \textbf{Differenza tassonomica
cruciale}: Gli \textbf{I/O Streams} (\passthrough{\lstinline!java.io!})
non vanno confusi con gli \textbf{Stream funzionali} di Java 8
(\passthrough{\lstinline!java.util.stream.Stream!}). Gli I/O Stream
trasportano byte o caratteri fisici verso periferiche o file; gli Stream
di Java 8 elaborano pipeline di trasformazione di oggetti in memoria. *
Un I/O Stream può essere agganciato a: * File su disco. * Flussi
standard di processo: \passthrough{\lstinline!System.in!} (standard
input), \passthrough{\lstinline!System.out!} (standard output),
\passthrough{\lstinline!System.err!} (standard error). * Connessioni di
rete (\emph{socket} TCP/IP o endpoint HTTP). * Buffer di memoria (array
di byte o stringhe).

\begin{center}\rule{0.5\linewidth}{0.5pt}\end{center}

\subsection{Il Dualismo dell'I/O in Java: Byte Streams vs Character
Streams (Slide
4-9)}\label{chap19_il-dualismo-dellio-in-java-byte-streams-vs-character-streams-slide-4-9}

L'architettura di \passthrough{\lstinline!java.io!} è rigorosamente
bipartita in due gerarchie parallele:

\needspace{7cm}
\begin{asciibox}
                      ┌─────────────────────────┐
                      │ java.io Stream Classes  │
                      └────────────┬────────────┘
                                   │
             ┌─────────────────────┴─────────────────────┐
             ▼                                           ▼
  ┌─────────────────────┐                     ┌─────────────────────┐
  │    Byte Streams     │                     │  Character Streams  │
  │  (8-bit dati grezzi)│                     │ (16-bit Unicode UTF)│
  └──────────┬──────────┘                     └──────────┬──────────┘
             │                                           │
  ┌──────────┴──────────┐                     ┌──────────┴──────────┐
  ▼                     ▼                     ▼                     ▼
┌─────────────┐ ┌──────────────┐       ┌─────────────┐ ┌─────────────┐
│ InputStream │ │ OutputStream │       │   Reader    │ │   Writer    │
└──────┬──────┘ └──────┬───────┘       └──────┬──────┘ └──────┬──────┘
       │               │                      │               │
┌──────┴──────┐ ┌──────┴───────┐       ┌──────┴──────┐ ┌──────┴──────┐
│FileInput-   │ │FileOutput-   │       │ FileReader  │ │ FileWriter  │
│Stream       │ │Stream        │       └──────┬──────┘ └──────┬──────┘
└──────┬──────┘ └──────┬───────┘              │               │
       │               │               ┌──────┴──────┐ ┌──────┴──────┐
┌──────┴──────┐ ┌──────┴───────┐       │Buffered-    │ │Buffered-    │
│Buffered-    │ │Buffered-     │       │Reader       │ │Writer       │
│InputStream  │ │OutputStream  │       └─────────────┘ └─────────────┘
└─────────────┘ └──────────────┘
\end{asciibox}

\subsubsection{\texorpdfstring{Byte Streams (\texttt{InputStream} e
\texttt{OutputStream}, Slide
5-7)}{Byte Streams (InputStream e OutputStream, Slide 5-7)}}\label{chap19_byte-streams-inputstream-e-outputstream-slide-5-7}

\begin{itemize}
\tightlist
\item
  \textbf{Unità di dato}: Singolo byte grezzo (8 bit, intervallo
  \([0, 255]\) restituito come \passthrough{\lstinline!int!}, dove il
  valore \(-1\) segnala la fine dello stream --- \emph{End Of File,
  EOF}).
\item
  \textbf{Destinazione d'uso}: Immagini, file multimediali, bytecode
  compilato (\passthrough{\lstinline!.class!}), file compressi
  (\passthrough{\lstinline!.zip!}), comunicazioni binarie di rete.
\item
  \textbf{Metodi cardine}:

  \begin{itemize}
  \tightlist
  \item
    \passthrough{\lstinline!int read()!}: legge il prossimo byte (o
    \(-1\) a fine stream).
  \item
    \passthrough{\lstinline!int read(byte[] b)!}: riempie il buffer di
    byte.
  \item
    \passthrough{\lstinline!void write(int b)!}: scrive un byte.
  \item
    \passthrough{\lstinline!void write(byte[] b)!}: scrive un intero
    array di byte.
  \item
    \passthrough{\lstinline!void close()!}: rilascia il descrittore del
    file del sistema operativo.
  \end{itemize}
\item
  \textbf{Implementazioni notevoli}:

  \begin{itemize}
  \tightlist
  \item
    \passthrough{\lstinline!FileInputStream!} /
    \passthrough{\lstinline!FileOutputStream!}: lettura/scrittura
    diretta su disco.
  \item
    \passthrough{\lstinline!BufferedInputStream!} /
    \passthrough{\lstinline!BufferedOutputStream!}: aggiunge un buffer
    in memoria per minimizzare le costose chiamate di sistema del kernel
    del SO.
  \item
    \passthrough{\lstinline!DataInputStream!} /
    \passthrough{\lstinline!DataOutputStream!}: serializzazione di
    primitivi Java (\passthrough{\lstinline!readInt()!},
    \passthrough{\lstinline!writeDouble()!}).
  \end{itemize}
\end{itemize}

\subsubsection{\texorpdfstring{Character Streams (\texttt{Reader} e
\texttt{Writer}, Slide
8-10)}{Character Streams (Reader e Writer, Slide 8-10)}}\label{chap19_character-streams-reader-e-writer-slide-8-10}

\begin{itemize}
\tightlist
\item
  \textbf{Unità di dato}: Caratteri Unicode (16 bit UTF-16).
\item
  \textbf{Destinazione d'uso}: Esclusivamente file di testo, file
  sorgente, documenti testuali.
\item
  \textbf{Le classi ponte (\emph{Bridge Adapters})}:

  \begin{itemize}
  \tightlist
  \item
    \passthrough{\lstinline!InputStreamReader!}: converte un
    \passthrough{\lstinline!InputStream!} di byte in un
    \passthrough{\lstinline!Reader!} di caratteri applicando una
    specifica codifica (es. UTF-8).
  \item
    \passthrough{\lstinline!OutputStreamWriter!}: converte caratteri in
    byte da inviare a un \passthrough{\lstinline!OutputStream!}.
  \end{itemize}
\item
  \textbf{Implementazioni notevoli}:

  \begin{itemize}
  \tightlist
  \item
    \passthrough{\lstinline!FileReader!} /
    \passthrough{\lstinline!FileWriter!}: lettura/scrittura di file di
    testo.
  \item
    \passthrough{\lstinline!BufferedReader!}: legge blocchi di testo
    bufferizzati e fornisce il comodissimo metodo
    \passthrough{\lstinline!String readLine()!} (restituisce un'intera
    riga o \passthrough{\lstinline!null!} a fine file).
  \item
    \passthrough{\lstinline!BufferedWriter!}: scrittura testuale con
    supporto al metodo \passthrough{\lstinline!newLine()!}.
  \end{itemize}
\end{itemize}

\begin{center}\rule{0.5\linewidth}{0.5pt}\end{center}

\subsection{\texorpdfstring{Gestione delle Risorse: Da
\texttt{try-finally} a \texttt{try-with-resources} (Slide 7,
10-11)}{Gestione delle Risorse: Da try-finally a try-with-resources (Slide 7, 10-11)}}\label{chap19_gestione-delle-risorse-da-try-finally-a-try-with-resources-slide-7-10-11}

\subsubsection{Il Vecchio Approccio (Pre-Java 7, Slide
10):}\label{chap19_il-vecchio-approccio-pre-java-7-slide-10}

\begin{lstlisting}[language=Java]
BufferedReader br = null;
try {
    br = new BufferedReader(new FileReader("i.txt"));
    // lettura...
} catch (IOException e) {
    // gestione errore...
} finally {
    if (br != null) {
        try {
            br.close(); // Ulteriore try-catch obbligatorio perché close() lancia IOException!
        } catch (IOException e) { ... }
    }
}
\end{lstlisting}

Questo pattern era verbose, faticoso da manutenere e fonte continua di
\emph{resource leaks} (descrittori di file aperti non chiusi se si
verificavano eccezioni nel \passthrough{\lstinline!finally!}).

\subsubsection{\texorpdfstring{L'Approccio Moderno:
\texttt{try-with-resources} (Java 7+, Slide
11)}{L'Approccio Moderno: try-with-resources (Java 7+, Slide 11)}}\label{chap19_lapproccio-moderno-try-with-resources-java-7-slide-11}

Tutte le classi che implementano l'interfaccia standard
\passthrough{\lstinline!java.lang.AutoCloseable!} possono essere
dichiarate all'interno delle parentesi tonde del blocco
\passthrough{\lstinline!try!}:

\begin{lstlisting}[language=Java]
try (BufferedReader br = new BufferedReader(new FileReader("i.txt"));
     BufferedWriter bw = new BufferedWriter(new FileWriter("o.txt"))) {
    String line;
    while ((line = br.readLine()) != null) {
        bw.write(line);
        bw.newLine();
    }
} catch (FileNotFoundException fnfe) {
    System.out.println("File not found...");
} catch (IOException ioe) {
    System.out.println("I/O Error...");
}
\end{lstlisting}

\begin{itemize}
\tightlist
\item
  \textbf{Garanzie della JVM}:

  \begin{enumerate}
  \def\labelenumi{\arabic{enumi}.}
  \tightlist
  \item
    Le risorse vengono chiuse \textbf{automaticamente} non appena il
    blocco \passthrough{\lstinline!try!} termina, sia in caso di
    completamento regolare sia in caso di eccezione o
    \passthrough{\lstinline!return!} anticipato.
  \item
    Vengono chiuse in \textbf{ordine inverso} rispetto alla loro
    dichiarazione (prima \passthrough{\lstinline!bw!}, poi
    \passthrough{\lstinline!br!}).
  \item
    Se sia il corpo del \passthrough{\lstinline!try!} che la chiamata a
    \passthrough{\lstinline!close()!} sollevano eccezioni, l'eccezione
    del corpo è quella principale propagata, mentre l'eccezione di
    chiusura viene salvata come \textbf{Suppressed Exception}
    (recuperabile con \passthrough{\lstinline!e.getSuppressed()!}).
  \end{enumerate}
\end{itemize}

\begin{center}\rule{0.5\linewidth}{0.5pt}\end{center}

\subsection{Gestione di File e Risorse di Rete (URL/URI) (Slide
13-17)}\label{chap19_gestione-di-file-e-risorse-di-rete-urluri-slide-13-17}

\begin{itemize}
\tightlist
\item
  \textbf{La classe \passthrough{\lstinline!java.io.File!} (Slide
  13-14)}: Rappresenta il percorso astratto di un file o di una
  directory. Non legge né scrive dati direttamente, ma permette di
  interrogare e modificare il filesystem:

  \begin{itemize}
  \tightlist
  \item
    \passthrough{\lstinline!exists()!},
    \passthrough{\lstinline!isFile()!},
    \passthrough{\lstinline!isDirectory()!},
    \passthrough{\lstinline!length()!},
    \passthrough{\lstinline!getAbsolutePath()!}.
  \item
    \passthrough{\lstinline!mkdir()!},
    \passthrough{\lstinline!delete()!},
    \passthrough{\lstinline!renameTo(File dest)!},
    \passthrough{\lstinline!listFiles()!}.
  \end{itemize}
\item
  \textbf{URL e URI (\passthrough{\lstinline!java.net!}, Slide 15-17)}:

  \begin{itemize}
  \item
    \passthrough{\lstinline!URI!} modella l'identificatore formale
    (\emph{RFC 2396}).
  \item
    \passthrough{\lstinline!URL!} modella la locazione fisica e il
    protocollo per accedere alla risorsa web.
  \item
    \emph{Avvertenza contemporanea (Slide 16)}: Il costruttore
    \passthrough{\lstinline!new URL("...")!} è \textbf{deprecato} da
    Java 20. La prassi moderna impone di creare un'istanza di
    \passthrough{\lstinline!URI!} e convertirla:

\begin{lstlisting}[language=Java]
URL url = new URI("https://info.cern.ch/index.html").toURL();
InputStream in = url.openStream(); // Apre una connessione HTTP e scarica i dati
\end{lstlisting}
  \end{itemize}
\end{itemize}

\begin{center}\rule{0.5\linewidth}{0.5pt}\end{center}

\subsection{Dati Strutturati: CSV e JSON (Slide
18-26)}\label{chap19_dati-strutturati-csv-e-json-slide-18-26}

Java non possiede parser nativi integrati nel runtime per file CSV e
JSON. Per questi formati si ricorre a librerie esterne integrate tramite
dipendenze Maven.

\subsubsection{File CSV con OpenCSV (Slide
19-22)}\label{chap19_file-csv-con-opencsv-slide-19-22}

\begin{itemize}
\item
  Un file CSV (\emph{Comma-Separated Values}) memorizza tabelle in
  formato testuale, separando le colonne con virgole o punti e virgola.
\item
  \textbf{Dipendenza Maven}:
  \passthrough{\lstinline!com.opencsv:opencsv:5.12.0!}.
\item
  \textbf{Scrittura}: \passthrough{\lstinline!CSVWriter!} scrive array
  di stringhe \passthrough{\lstinline!String[]!} gestendo
  automaticamente apici ed escape:

\begin{lstlisting}[language=Java]
CSVWriter csvw = new CSVWriter(new FileWriter("data.csv"));
csvw.writeNext(new String[]{ "id", "name", "address" });
\end{lstlisting}
\item
  \textbf{Lettura}: \passthrough{\lstinline!CSVReader!} con
  \passthrough{\lstinline!readNext()!} riga per riga o
  \passthrough{\lstinline!readAll()!}.
\end{itemize}

\subsubsection{File JSON con Google Gson (Slide
23-26)}\label{chap19_file-json-con-google-gson-slide-23-26}

\begin{itemize}
\item
  JSON (\emph{JavaScript Object Notation}) è lo standard dominante per
  lo scambio dati basato su mappe chiave-valore
  \passthrough{\lstinline!\{\}!} e liste ordinate
  \passthrough{\lstinline![]!}.
\item
  \textbf{Dipendenza Maven}:
  \passthrough{\lstinline!com.google.code.gson:gson:2.13.2!}.
\item
  \textbf{Serializzazione e Deserializzazione Automatica (Data
  Binding)}:

\begin{lstlisting}[language=Java]
Gson gson = new Gson();

// 1. Oggetto Java -> Stringa JSON (Serializzazione)
Date date = new Date(1, 1, 1970);
String json = gson.toJson(date); // Restituisce '{"day":1,"month":1,"year":1970}'

// 2. Stringa JSON -> Oggetto Java (Deserializzazione)
Date d = gson.fromJson(json, Date.class); // Ricostruisce l'istanza valorizzando i campi!
\end{lstlisting}

  Gson accede ai campi (anche \passthrough{\lstinline!private!} e
  \passthrough{\lstinline!final!}) tramite \textbf{Reflection}, senza
  richiedere getter/setter pubblici o costruttori senza argomenti!
\item
  \textbf{Pretty Printing}:
  \passthrough{\lstinline!new GsonBuilder().setPrettyPrinting().create()!}
  formatta il JSON con ritorni a capo e indentazione a 2 spazi.
\end{itemize}

\begin{center}\rule{0.5\linewidth}{0.5pt}\end{center}

\section{\texorpdfstring{Analisi Dettagliata del Codice
(\texttt{Lezione19/MavenDate})}{Analisi Dettagliata del Codice (Lezione19/MavenDate)}}\label{chap19_analisi-dettagliata-del-codice-lezione19mavendate}

In \passthrough{\lstinline!Lezione19/MavenDate!} il prof. Pasqua integra
tutti i concetti della Slide T20 all'interno della classe dimostrativa
\href{file:///home/wally/Documenti/GitHub/Programmazione-II/25-26/Teoria-e-Esercitazioni/Codice-20260902/Lezione19/MavenDate/src/main/java/it/oop/ui/MainIO.java}{MainIO.java}.

\subsection{\texorpdfstring{Configurazione delle Dipendenze in
\href{file:///home/wally/Documenti/GitHub/Programmazione-II/25-26/Teoria-e-Esercitazioni/Codice-20260902/Lezione19/MavenDate/pom.xml\#L68-L78}{pom.xml}}{Configurazione delle Dipendenze in pom.xml}}\label{chap19_configurazione-delle-dipendenze-in-pom.xml}

\begin{lstlisting}[language=XML]
        <!-- OpenCSV per parsing e generazione CSV (Slide 20) -->
        <dependency>
            <groupId>com.opencsv</groupId>
            <artifactId>opencsv</artifactId>
            <version>5.12.0</version>
        </dependency>
        <!-- Google Gson per serializzazione/deserializzazione JSON (Slide 24) -->
        <dependency>
            <groupId>com.google.code.gson</groupId>
            <artifactId>gson</artifactId>
            <version>2.13.2</version>
        </dependency>
\end{lstlisting}

\begin{center}\rule{0.5\linewidth}{0.5pt}\end{center}

\subsection{\texorpdfstring{Analisi Riga per Riga di
\href{file:///home/wally/Documenti/GitHub/Programmazione-II/25-26/Teoria-e-Esercitazioni/Codice-20260902/Lezione19/MavenDate/src/main/java/it/oop/ui/MainIO.java}{MainIO.java}}{Analisi Riga per Riga di MainIO.java}}\label{chap19_analisi-riga-per-riga-di-mainio.java}

\subsubsection{\texorpdfstring{Lettura Binaria di Bytecode con
\texttt{FileInputStream} e \texttt{FileChannel} (Righe
18-41)}{Lettura Binaria di Bytecode con FileInputStream e FileChannel (Righe 18-41)}}\label{chap19_lettura-binaria-di-bytecode-con-fileinputstream-e-filechannel-righe-18-41}

\begin{lstlisting}[language=Java]
        FileInputStream fis = null;
        try {
            // Apertura dello stream di byte su un file binario compilato (.class)
            fis = new FileInputStream("src/main/resources/Date.class");
            FileChannel fc = fis.getChannel();
            int b;
            // Lettura byte a byte
            while ((b = fis.read()) != -1) {
                System.out.println((char) b);
            }
            // Rewind dello stream riportando la posizione a 0 tramite FileChannel
            fc.position(0);
            StringBuilder sb = new StringBuilder();
            // Lettura dell'intero contenuto in un colpo solo con readAllBytes() (Java 9+)
            for (byte bb : fis.readAllBytes()) {
                sb.append((char) bb);
            }
            System.out.println("-----");
            System.out.println(sb.toString());
        } catch (IOException ioe) {
            System.out.println(ioe.getMessage());
        } finally {
            // Chiusura classica pre-Java 7 con try-catch protetto
            try {
                if (fis != null) fis.close();
            } catch (IOException ioe) {
                System.out.println(ioe.getMessage());
            }
        }
\end{lstlisting}

\subsubsection{\texorpdfstring{Scrittura di Testo con
\texttt{FileWriter} e \texttt{try-with-resources} (Righe
42-50)}{Scrittura di Testo con FileWriter e try-with-resources (Righe 42-50)}}\label{chap19_scrittura-di-testo-con-filewriter-e-try-with-resources-righe-42-50}

\begin{lstlisting}[language=Java]
        // Blocco try-with-resources: fw viene chiuso automaticamente
        try (FileWriter fw = new FileWriter("src/main/resources/file.txt")) {
            StringBuilder sb = new StringBuilder();
            // Genera stringhe di 'a' crescenti usando Stream di Java 8
            Stream.iterate("a", s -> s + "a")
                    .limit(15)
                    .forEach(s -> sb.append(s).append("\n"));
            fw.write(sb.toString());
        } catch (IOException ioe) {
            System.out.println(ioe.getMessage());
        }
\end{lstlisting}

\subsubsection{Generazione di Tabella CSV con OpenCSV (Righe
52-65)}\label{chap19_generazione-di-tabella-csv-con-opencsv-righe-52-65}

\begin{lstlisting}[language=Java]
        try (CSVWriter csvw = new CSVWriter(new FileWriter("src/main/resources/data.csv"))) {
            String[] row = { "id", "name", "address" };
            csvw.writeNext(row); // Scrive l'intestazione delle colonne

            row[0] = "3";
            row[1] = "Paul";
            row[2] = "Strada le Grazie, 15";
            csvw.writeNext(row);

            row[0] = "6";
            row[1] = "Sam";
            row[2] = "Strada le Grazie, 18";
            csvw.writeNext(row);
        } catch (IOException ioe) {
            System.out.println(ioe.getMessage());
        }
\end{lstlisting}

\subsubsection{Serializzazione e Deserializzazione con Gson (Righe
66-80)}\label{chap19_serializzazione-e-deserializzazione-con-gson-righe-66-80}

\begin{lstlisting}[language=Java]
        Gson gson = new Gson();
        Date date = new Date(1, 1, 1970);

        // Serializzazione dell'oggetto Date in formato JSON
        String json = gson.toJson(date);
        System.out.println(json); // Stampa: {"day":1,"month":1,"year":1970}

        // Deserializzazione da stringa JSON a nuova istanza Date
        Date date1 = gson.fromJson("{\"day\":1,\"month\":1,\"year\":1971}", Date.class);
        System.out.println(date1.toString()); // Stampa: y1971m1d1

        // Scrittura su file con Pretty Printing abilitato
        try (FileWriter fw = new FileWriter("src/main/resources/file.json")) {
            new GsonBuilder().setPrettyPrinting().create().toJson(date, fw);
        } catch (IOException ioe) {
            System.out.println(ioe.getMessage());
        }
\end{lstlisting}

\begin{center}\rule{0.5\linewidth}{0.5pt}\end{center}

\section{Risultato di Compilazione ed Esecuzione
Reale}\label{chap19_risultato-di-compilazione-ed-esecuzione-reale}

Comando eseguito nel progetto
\href{file:///home/wally/Documenti/GitHub/Programmazione-II/25-26/Teoria-e-Esercitazioni/Codice-20260902/Lezione19/MavenDate}{Lezione19/MavenDate}:

\begin{lstlisting}[language=bash]
mvn compile exec:java -Dexec.mainClass="it.oop.ui.MainIO"
\end{lstlisting}

Output ottenuto a video:

\begin{lstlisting}
[Dumping del bytecode del file Date.class con header CAFEBABE e constant pool]
-----
{"day":1,"month":1,"year":1970}
y1971m1d1
[INFO] BUILD SUCCESS
\end{lstlisting}

\subsection{Ispezione dei file generati su
disco:}\label{chap19_ispezione-dei-file-generati-su-disco}

\begin{enumerate}
\def\labelenumi{\arabic{enumi}.}
\item
  \passthrough{\lstinline!src/main/resources/file.txt!}:

\begin{lstlisting}
a
aa
aaa
...
aaaaaaaaaaaaaaa
\end{lstlisting}
\item
  \passthrough{\lstinline!src/main/resources/data.csv!}:

\begin{lstlisting}
"id","name","address"
"3","Paul","Strada le Grazie, 15"
"6","Sam","Strada le Grazie, 18"
\end{lstlisting}
\item
  \passthrough{\lstinline!src/main/resources/file.json!}:

\begin{lstlisting}
{
  "day": 1,
  "month": 1,
  "year": 1970
}
\end{lstlisting}
\end{enumerate}

\begin{center}\rule{0.5\linewidth}{0.5pt}\end{center}

\section{Traguardo Raggiunto: Revisione Integrale del Corso
Completata!}\label{chap19_traguardo-raggiunto-revisione-integrale-del-corso-completata}

Abbiamo completato con successo e senza tralasciare alcun dettaglio: *
Tutte le \textbf{20 slide teoriche} (da \passthrough{\lstinline!T01!} a
\passthrough{\lstinline!T20!}). * Tutti i \textbf{progetti e le lezioni
pratiche} (da \passthrough{\lstinline!Lezione04!} a
\passthrough{\lstinline!Lezione19!}). * Tutti i diagrammi, meme, pattern
architetturali, finezze di linguaggio, bug del docente, test Maven e
logiche di compilazione ed esecuzione.

Fammi sapere se desideri approfondire specifici temi per l'esame,
affrontare vecchi temi d'esame scritti/pratici, o dedicarti agli
esercizi del tuo workspace
(\passthrough{\lstinline!esercizi-wally-25-26!})!

\end{document}
